%*********************************************************************
% fduMaththesis: 复旦大学数学科学学院研究生毕业论文 LaTeX 模版
% 模版基于开源项目 by 曾祥东（https://github.com/stone-zeng/fduthesis）开发, 
% 部分内容摘自数学科学学院2023年毕业生邹森博士的毕业论文。
% 该模版旨在帮助复旦大学数学科学学院的研究生规范撰写毕业论文，符合学院的格式要求。
% 
% 注：
%   1. 请确保使用 UTF-8 编码保存
%   2. 请使用 XeLaTeX 或 LuaLaTeX 编译
%   3. 请阅读说明文档<复旦数学科学学院研究生毕业论文LaTeX模版使用说明>，细节设置可参考<源项目说明文档>以及源项目github地址。
%*********************************************************************
\newif\ifchinese
\chinesetrue
\chinesefalse

\ifchinese
  \documentclass[type=master]{fduthesis}
\else
  \documentclass[type=master]{fduthesis-en}
\fi

\newcommand{\bilingual}[2]{%
  \ifchinese
    #2\relax
  \else
    #1\relax
  \fi
}

% \documentclass[type=master]{fduthesis-en}%硕士
% \documentclass[type=doctor]{fduthesis}

% !TEX root = ./main.tex

\fdusetup{
	style = {font-size = -4,
			bib-backend =bibtex,
            bib-resource = {main.bib},
			bib-style = numerical
   },
	info = {
		title = {可无限分支的$N$-分支随机游走的渐进速率性质},
		title* = {Asymptotic Speed of the $N$-Branching Random Walk with Infinite Offspring},
		author = {王逸扬},
		author* = {Wang Yiyang},
		supervisor = {Élie Aïdékon 教~授},
		department = {数学科学学院},
		major = {概率论与数理统计},
		student-id = {23210180123},
		instructors = {{Élie E. F. Aïdékon\quad 教~授}},%,{某\quad 某\quad 教~授},{某某某\quad 教~授}},
		keywords = {$N$-分支随机游走, 小偏差, $\mathrm L^2$-极限, 渐进速率},
		keywords* = {$N$-BRW, Small Deviations, $\mathrm L^2$-Limit, Asymptotic Speed},
		clc = {O24},
		date = {2024年03月22日}
	}
}

% 序号
\usepackage{enumerate}

% 定义图片文件目录与扩展名
\graphicspath{{figs/}}
\DeclareGraphicsExtensions{.pdf,.eps,.png,.jpg,.jpeg}

% 子图宏包
\usepackage{subcaption}
\usepackage{bicaption}

% 确定浮动对象的位置，可以使用 [H]，强制将浮动对象放到这里（可能效果很差）
\usepackage{float}

% 表格排列
\usepackage{array}

% 固定宽度的表格
\usepackage{tabularx}

% 使用三线表：toprule，midrule，bottomrule。
\usepackage{booktabs}

% 表格中支持跨行
\usepackage{multirow}

% 使用长表格
\usepackage{longtable}

% 附带脚注的表格
\usepackage{threeparttable}

% 附带脚注的长表格
\usepackage{threeparttablex}

% 算法环境宏包
\usepackage[ruled,vlined,linesnumbered]{algorithm2e}
% \usepackage{algorithm, algorithmicx, algpseudocode}

% 代码环境宏包
\usepackage{listings}
\lstdefinestyle{lstStyleCode}{%
  aboveskip         = \medskipamount,
  belowskip         = \medskipamount,
  basicstyle        = \ttfamily\zihao{6},
  commentstyle      = \slshape\color{black!60},
  stringstyle       = \color{green!40!black!100},
  keywordstyle      = \bfseries\color{blue!50!black},
  extendedchars     = false,
  upquote           = true,
  tabsize           = 2,
  showstringspaces  = false,
  xleftmargin       = 1em,
  xrightmargin      = 1em,
  breaklines        = false,
  framexleftmargin  = 1em,
  framexrightmargin = 1em,
  columns           = flexible,
  keepspaces        = true,
  texcl             = true,
  mathescape        = true
}
\lstnewenvironment{codeblock}[1][]{%
  \lstset{style=lstStyleCode,#1}}{}

% 绘图宏包
\usepackage{tikz}
\usetikzlibrary{arrows.meta, shapes.geometric,cd}

% 超链接
\usepackage{xcolor}
\usepackage{hyperref}
\hypersetup{
    colorlinks=true, % 启用彩色链接
    linkcolor=black,  % 普通内部链接的颜色
    citecolor=black,  % 引用链接的颜色
    filecolor=black,  % 打开本地文件的链接颜色
    urlcolor=black    % 链接到URL的颜色
}

% 确保在导言区后期再重新设置一次
\AtEndPreamble{%
    \hypersetup{
        colorlinks=true,     % 启用彩色链接
        linkcolor=red,     % 普通内部链接的颜色
        citecolor=blue,     % 引用链接的颜色
        filecolor=black,     % 本地文件的链接颜色
        urlcolor=black,      % URL 链接的颜色
    }
}

% 页码显示
\usepackage{fancyhdr}
\pagestyle{fancy}
\fancyhf{} 
\fancyfoot[LE,RO]{\thepage} 
\fancypagestyle{plain}{
    \fancyhf{}
    \fancyfoot[LE,RO]{\thepage} 
}
% `L`, `C`, `R`分别对应左、中、右，`E`, `O`对应偶数页和奇数页


% 自定义引用样式，支持可选参数
\NewDocumentCommand{\mycite}{o m}{%
  \IfNoValueTF{#1}
    {\scalebox{1.2}[1.2]{\raisebox{-0.70ex}{\cite{#2}}}}
    { \scalebox{0.9}[0.9]{\raisebox{0.10ex}{\textnormal{[#1, {\scalebox{1.3}[1.3]{\raisebox{-0.80ex}{\citealp{#2}}}}]}}}}}


\setlength{\abovedisplayskip}{0pt}
\setlength{\belowdisplayskip}{0pt}
\usepackage{enumitem}
\setlist{nosep}
\bilingual{
    \linespread{1.2}
}{
    \linespread{1.3}
}


\usepackage{mathrsfs}
\usepackage{anyfontsize}
% \def\QED{\square}
\allowdisplaybreaks[4]

\newtheorem[
  header-font=\normalfont\bfseries\upshape, 
  body-font=\normalfont\mdseries\itshape,
  head-format={\NAME\,\NUMBER\NOTE},
  note-brace-left={\mdseries\upshape(},
  note-brace-right={\mdseries\upshape)},
]{assumption}{\bilingual{Assumption}{假设}}

\setcounter{tocdepth}{1}

\begin{document}

% 前置部分包含目录、中英文摘要以及符号表等
\frontmatter

% 目录
\tableofcontents

% 插图索引
%\listoffigures*
% 表格索引
%\listoftables*
% 算法索引
%\listofalgorithms*

% 插图目录
\listoffigures

% !TEX root = ../main.tex

\bilingual{

}{
  \begin{abstract}

    本文研究有限方差 ($\alpha=2$) 普适性类下,~ 具有一般分支律的 $N$-分支随机游走 ($N$-BRW) 前沿演化的渐近性质. 本文所谓$N$-分支随机游走,~ 是指从$\lfloor N^a\rfloor$ ($a\in[0,1]$) 个粒子出发,~ 每一代根据一般分支律 $\Xi$ (可数无限的点过程) 繁殖,~ 并仅保留所有子代粒子中位置最靠前的 $N$ 个粒子的演化过程.

    主要工作围绕 Brunet--Derrida 关于渐进速度修正项$(\log N)^{-2}$ 标度律的猜想展开：首先,~ 在脊柱分解框架下,~ 利用 Mogul'ski\u{i} 小偏差定理建立水平和倾斜管道的概率估计. ~通过引入最左侧位移的正指数矩假设与分支律交叉项的负指数约束,~ 本文将 A\"id\'ekon 与 Hu 针对特定 Poisson 分支律模型建立的 $L^2$ 轨迹极限定理推广至一般分布. ~证明了在三次对数时间与一次对数空间缩放尺度下,~ 前沿位置收敛于线性轨迹 $-a + \frac{\pi^2}{2}\sigma^2 t$,~ 确立了该轨迹作为系统本征属性的普适性. ~其次,~ 针对单粒子起始情形 ($a=0$) 改进了证明策略,~ 利用超临界性质通过位移截断构造至多二分叉过程,~ 并引入``重启机制''克服截断后的灭绝风险,~ 保障了极限轨迹在主尺度下的稳定性.

    最后,~ 利用上述极限定理,~ 本文严格验证了一般分支律下速度修正项的 $(\log N)^{-2}$ 标度律,~ 并给出了包含显式常数的渐近速度表达式. ~本文结果精细刻画了前沿演化的动态过程,~ 为理解有限种群选择机制下的微观涨落与宏观速度偏移之间的关联提供了完备的数学基础.
  \end{abstract}
}

\ifchinese
  \begin{abstract*}
\else
  \begin{abstract}
\fi
  This thesis investigates the asymptotic front evolution of the $N$-branching random walk ($N$-BRW) within the finite variance ($\alpha=2$) universality class with general branching laws. The $N$-branching random walk considered here starts with $\lfloor N^a\rfloor$ ($a\in[0,1]$) particles, where each generation reproduces according to a general branching law $\Xi$ (a countably infinite point process), and only the $N$ leftmost particles among all offspring are retained for the subsequent evolution.

  This thesis mainly focuses on the Brunet--Derrida conjecture regarding the speed correction: First, within the framework of spine decomposition, the Mogul'ski\u{i} small deviation theorem is utilized to establish small deviation estimates in horizontal and tilted tubes. By imposing a positive exponential moment assumption on the leftmost displacement and negative exponential constraints on the branching cross-terms, the $L^2$ trajectory limit theorem of A\"id\'ekon and Hu has been generalized to general branching laws. It has been proven that under triple-logarithmic time and single-logarithmic space scaling, the front position converges to a linear trajectory $-a + \frac{\pi^2}{2}\sigma^2 t$, confirming its universality as an intrinsic property of the system. Second, for the case of a single-particle start ($a=0$), the proof strategy is enhanced by constructing an at-most binary branching process via displacement truncation and implementing a "re-booting mechanism" to mitigate extinction risks, ensuring the stability of the limit trajectory under the primary scale.

  Finally, based on the aforementioned limit theorems, it has been rigorously verified that the $(\log N)^{-2}$ scaling law of the speed correction under general laws holds and the explicit constant is derived. These results provide a refined characterization of front dynamics and offer a comprehensive mathematical foundation for understanding selection mechanisms in finite populations.
\ifchinese
  \end{abstract*}
\else
  \end{abstract}
\fi

% 符号表
% 语法与 LaTeX 表格一致：列用 & 区分，行用 \\ 区分
% 如需修改格式，可以使用可选参数：
%   \begin{notation}[ll]
%     $x$ & 坐标 \\
%     $p$ & 动量
%   \end{notation}
% 可选参数与 LaTeX 标准表格的列格式说明语法一致
% 这里的 “ll” 表示两列均为自动宽度，并且左对齐
\begin{notation}[rl]
    $\mathscr F$         & $\sigma$-Algebra \\
    $\mathbb P(\cdot)$   & Probability measure \\
    $\mathbb E[\cdot]$ or $\mathbb E_{\mathbb P}[\cdot]$
                         & Expectation w.r.t. $\mathbb P$ \\
    $\mathbb E(\cdot|\cdot)$ or $\mathbb E_{\mathbb P}(\cdot|\cdot)$
                         & Conditional expectation w.r.t. $\mathbb P$ \\
    ${\mathrm{Var}}[\cdot]$ or ${\mathrm{Var}}_{\mathbb P}[\cdot]$
                         & Variance w.r.t. $\mathbb P$ \\
    $\left.f(a)\right|_{a=\xi}
    $                    & Composition of measurable function $f$ and random variable $\xi$ \\
    ${\bf 1}_{\{\cdot\}}$
                         & Indicator function \\
    $\mathbb V$          & Space of tree vertices of Ulam--Harris--Neveu terminology \\
    $\mathbb T$ and $\mathbb T_0$
                         & Space of trees and locally finite trees \\
    $T$ and $\boldsymbol T$   
                         & Tree (in $\mathbb T$) and trees (in $\mathbb T^{\mathbb I}$ with index set $\mathbb I$) \\
    $\mathcal T$ and $\boldsymbol{\mathcal T}$
                         & Random tree and random trees \\
    $u\in T$             & The vertex $u\in\mathbb V$ is on tree $T$ \\
    $V$ and $\mathcal V$ 
                         & Mark function $\boldsymbol T\to\mathbb R$ w.r.t. $\boldsymbol T$, and random mark function\\
    $(\boldsymbol T,V)$ and $(\boldsymbol{\mathcal T},\mathcal V)$ 
                         & Marked tree(s) and random marked tree(s) \\
    $\mathcal V^{N}$ or $\mathcal V^{N|a}$   
                         & $\mathcal V$ for $N$-BRW starting with $\lfloor N^a\rfloor$ particles at $0$ \\
    $V^{(i)}$ or $V^{\cdot|(i)}$
                         & Mark function of $V$ for index $i$ \\
    $\Xi$                & Point process of offspring in BRW \\
    $\Xi_i$              & $i$-th smallest point in $\Xi$ \\
    $|u|$                & Generation of a vertex $u\in\mathbb V$\\
    $\Theta_{|u|=n}$ e.g. $\sum_{|u|=n}$
                         & Operation $\Theta$ over [all $u\in T$ or $\mathcal T$ with $|u|=n$] or\\
                         & over [all $i\in\mathbb I$ and $u\in\boldsymbol T(i)$ or $\boldsymbol{\mathcal T}(i)$ with $|u|=n$]\\
    $M_n$ and $\mathcal M_n$
                         & Infimum of $V$ and $\mathcal V$ at generation $n$ \\
    $[\![\,\cdot\,,\,\cdot\,]\!]$
                         &  Closed integer interval\\
    $(\cdot)_+$ and $(\cdot)_-$
                         & Positive and negative part \\
    $O_N(\cdot)$ and $o_N(\cdot)$
                         & Big O and little o notation for sequences indexed by $N$
%   $u$                  & 标量值函数        \\
%   $\boldsymbol{u}$        & 向量值函数        \\
%   $\mathcal{K}$            & 算子 \\
%   $[A,B]:=AB-BA$      & 交换子 \\
%   $\boldsymbol{A}$        & 系数矩阵     \\
%   $\mathbb{A}$        & 系数张量     \\
\end{notation}

\mainmatter
% 建议采用多文件编译的方式
% 比较好的做法是把每一章放进一个单独的 tex 文件里，并在这里用 \include 导入，例如
%   \include{chapter1}
%   \include{chapter2}
%   \include{chapter3}

% \chapter{简介}
复旦大学数学科学学院研究生毕业论文 LaTeX 模版基于开源模版fduthesis开发by 曾祥东\footnote{https://github.com/stone-zeng/fduthesis}，我们对他表示感谢。本模板的第一版由复旦大学数学科学学院2021级罗心悦整理完成，目前为迭代的第二版模板，由2019级冯典在第一版的基础上进行了更新。关于模板更新的内容，我们将在下文详细指出。部分内容摘自复旦大学数学科学学院2023年毕业生邹森博士的毕业论文。说明性文字参考了上海交通大学学位论文模版\footnote{https://github.com/sjtug/SJTUThesis}。 

本模板依然在持续的迭代更新中，学院每学年将在官网\footnote{https://math.fudan.edu.cn/c1/39/c30396a639289/page.htm}提供一次大的更新，随时迭代更新的版本请跳转\href{https://github.com/VeMath/fduMath_thesis/tree/v1.0.0}{{\color{red} github仓库}}\footnote{https://github.com/VeMath/fduMath\_thesis/tree/v1.0.0}。
欢迎大家指出模板的不足之处，有任何意见与建议可以在github仓库中提交issue或者联系维护管理员罗心悦：xinyueluo21@m.fudan.edu.cn和张晨阳：24110180062@m.fudan.edu.cn。


该模版旨在帮助复旦大学数学科学学院的研究生规范撰写毕业论文，符合学院的格式要求，更多细节可以参考《复旦数学科学学院研究生毕业论文\LaTeX{}模版使用说明》。

\section{更新说明}
这一节将给出与上一节的差别，其中对于模板本身的更改仅需了解。对于论文书写比较关键的更改，我们将会用{\color{blue} 蓝色}进行标注。
\subsection{关于引用}

事实上，如果使用\textbackslash cite命令，我们将发现引用的上标会出现在文字的右上角，比如文献\cite{ChenPDE}。{\color{blue} 我们在setup.tex文件中对引用命令进行了重写，如有需求，请使用\textbackslash mycite命令，比如文献\mycite{ChenPDE}。}

\textbf{如果有更好的解决办法，欢迎大家与我们联系。}

\subsection{关于页码位置}

第一版中的页码出现在每一页的正下方，我们对此进行了调整，现在的页码位置分奇偶，奇数页出现在右下角，而偶数页出现在左下角。

目前模板默认的方式是奇数页在右下角，偶数页在左下角。如果确实需要修改回居中格式，可在文件setup.tex中进行修改。
只需将setup.tex中对应代码片段的LE,RO更改为C即可。

奇数页右下角，偶数页左下角：
\begin{codeblock}
\fancyfoot [LE,RO] { \thepage }
\fancypagestyle{plain}{
    \fancyhf{}
    \fancyfoot[LE,RO]{\thepage} 
}
\end{codeblock}

页码居中：
\begin{codeblock}
\fancyfoot [C] {  \thepage }
\fancypagestyle{plain}{
    \fancyhf{}
    \fancyfoot[C]{\thepage} 
}
\end{codeblock}

\subsection{关于目录索引颜色}
我们新增了一种超链接的颜色，这使得原先红色的目录及索引，以黑色的形式呈现。


% % !TEX root = ../main.tex

\chapter{数学语言}

\section{数学符号和公式}

\begin{itemize}
    \item 公式应另起一行居中排版。公式后应注明编号，按章顺序编排，编号右端对齐。
    \item 公式末尾是需要添加标点符号的，至于用逗号还是句号，取决于公式下面一句是接着公式说的，还是另起一句。
    \begin{equation}
  \frac{2h}{\pi}\int_{0}^{\infty}\frac{\sin\left( \omega\delta \right)}{\omega}
  \cos\left( \omega x \right) \mathrm{d}\omega = 
  \begin{cases}
    h, & \left| x \right| < \delta, \\
    \frac{h}{2}, & x = \pm \delta, \\
    0, & \left| x \right| > \delta.
  \end{cases}
\end{equation}
    \item 公式较长时最好在等号 ``$=$'' 处转行。
    \begin{align}
    & I (X_3; X_4) - I (X_3; X_4 \mid X_1) - I (X_3; X_4 \mid X_2) \nonumber \\
  = & [I (X_3; X_4) - I (X_3; X_4 \mid X_1)] - I (X_3; X_4 \mid \tilde{X}_2) \\
  = & I (X_1; X_3; X_4) - I (X_3; X_4 \mid \tilde{X}_2).
\end{align}
\item 公式较长时最好在等号 ``$=$'' 处转行。
如果在等号处转行难以实现，也可在 $+$、$-$、$\times$、$\div$ 运算符号处转行，转行
时运算符号仅书写于转行式前，不重复书写。
\begin{multline}
  \frac{1}{2} \Delta (f_{ij} f^{ij}) =
    2 \left(\sum_{i<j} \chi_{ij}(\sigma_{i} - \sigma_{j})^{2}
    + f^{ij} \nabla_{j} \nabla_{i} (\Delta f) \right. \\
  \left. + \nabla_{k} f_{ij} \nabla^{k} f^{ij} +
    f^{ij} f^{k} \left[2\nabla_{i}R_{jk}
    - \nabla_{k} R_{ij} \right] \vphantom{\sum_{i<j}} \right).
\end{multline}
\end{itemize}



\section{定理环境}

这里举一个 ``定义'' 和 ``定理'' 的例子。

\begin{definition}[非负整指数索伯列夫空间\mycite{ChenPDE}]
设$\Omega\subset\mathbb{R}^n$是一个给定区域，令$m$为非负整数，$1\leq p\leq \infty$.定义索伯列夫空间$W^{m,p}$为
$$
W^{m,p}(\Omega)=\left\{u\in L^p(\Omega): \text{对任意}\alpha\in \mathbb{N}_{0}^n\text{满足 }0\leq |\alpha|\leq m, D^{\alpha}u\in L^p(\Omega) \right\},
$$
并装备以范数
\begin{equation*}
\begin{aligned}
\|u\|_{W^{m,p}(\Omega)} &= \left( \sum_{|\alpha|\leq m} \|D^{\alpha}u\|_{L^p(\Omega)}^p \right)^{1/p},\quad 1\leq p <\infty, \label{eq:sobolev-norm}\\
\|u\|_{W^{m,\infty}(\Omega)} &= \max_{|\alpha|\leq m} \|D^{\alpha}u\|_{L^{\infty}(\Omega)}.
\end{aligned}
\end{equation*}
\end{definition}


\begin{theorem}[索伯列夫嵌入定理{\mycite{AF2003}}]\label{thm:sobolev-embedding}
令$\Omega$为$\mathbb{R}^n$上一个有$C^{\infty}$光滑边界的有界区域.令$j\geq 0,m\geq 1$为整数，$1\leq p<\infty$，如下嵌入关系成立：
\begin{enumerate}
\item 如果$mp>n>(m-1)p$，则
$$
W^{m,p}(\Omega) \hookrightarrow C^{0,s}(\overline{\Omega}), \quad \text{对于 }0<s\leq m-\frac{n}{p}.
$$
如果$n=(m-1)p$，则
$$
W^{m.p}(\Omega) \hookrightarrow C^{0,s}(\overline{\Omega}), \quad \text{对于 }0<s<1.
$$
\item 如果$1\leq k\leq n$，以及$mp=n$，则
$$
W^{j+m,p}(\Omega)\hookrightarrow W^{j,q}(\Omega), \quad \text{对于 }p\leq q < \infty.
$$
\item 如果$mp<n$，则
$$
W^{j+m,p}(\Omega) \hookrightarrow W^{j,q}(\Omega), \quad \text{对于 }p\leq q \leq p^{*}:=\frac{np}{n-mp}.
$$
\end{enumerate}
\end{theorem}



\section{线性方程的适定性理论}\label{sec:linear-wellposedness}
本节摘自邹森博士的毕业论文，包含一些数学符号的使用、定理的叙述与证明、数学式的引用。

在这一节中，我们讨论线性薛定谔方程
\begin{equation}\label{eq:schrodinger-0}
\left\{\begin{aligned}
-\Delta u - k^2 u +c(x)u &=F, && x\in\Omega,\\
u &=f, && x\in\partial\Omega
\end{aligned}\right.
\end{equation}
弱解的适定性，从而定义反问题的观测数据DtN算子.此外，为研究非线性方程的适定性，我们使用了线性亥姆霍兹方程
\begin{equation}\label{eqn:Helmholtz}
\left\{\begin{aligned}
-\Delta u - k^2 u &=F, && x\in\Omega,\\
u &=f, && x\in\partial\Omega
\end{aligned}\right.
\end{equation}
的强解理论，我们也在本节中加以介绍.这些结论都可以从一般形式椭圆方程的弱解和强解理论导出.为简洁起见，我们首先给出如下定义：我们定义散度型微分算子为
\begin{equation}\label{eq:divergence-form}
L:=-\sum_{i=1}^n D_i(a^{ij}(x)D_{j})+c(x).
\end{equation}
称散度型微分算子$L$在区域$\Omega$上满足椭圆性条件，如果对于任意$x\in\Omega$，$\xi\in\mathbb{R}^n$，存在常数$\lambda,\Lambda>0$，使得
\begin{equation}\label{eq:strict-elliptic}
\lambda |\xi|^2 \leq \sum_{i,j=1}^n a^{ij}(x)\xi_{i}\xi_{j} \leq \Lambda|\xi|^2 .
\end{equation}
另外我们定义非散度型微分算子
\begin{equation}\label{eq:nondivergence-form}
L=-\sum_{i,j=1}^{n}a^{ij}(x)D_{ij} + c(x).
\end{equation}
类似地，对于非散度型微分算子，我们也可以定义其椭圆性条件：对于任意$x\in\Omega$，$\xi\in\mathbb{R}^n$，存在常数$\lambda,\Lambda>0$，使得
\begin{equation*}
\lambda |\xi|^2 \leq \sum_{i,j=1}^n a^{ij}(x)\xi_{i}\xi_{j} \leq \Lambda|\xi|^2 .
\end{equation*}
对于以上两种形式的微分算子，在后文的讨论中，除非特殊说明，我们假定对于$1\leq i,j\leq n$，有$a^{ij}(x)\in C^{\infty}(\overline{\Omega})$以及$c\in L^{\infty}(\Omega)$.

\subsection{弱解的存在性和连续依赖性}\label{sec:weak-solution}

我们首先叙述一般形式椭圆方程的Dirichlet问题弱解的存在性：
\begin{theorem}[Dirichlet问题弱解的存在性{\mycite{Chen2elliptic}}]\label{thm:weak-solution}
令$L$为\eqref{eq:divergence-form}中定义的散度型椭圆微分算子，令$\Omega\subset\mathbb{R}^n$为有界开区域，使得Sobolev嵌入定理在其上成立.对于$\mu\in\mathbb{R}$，Dirichlet问题
\begin{equation}\label{eq:dirichlet-equation}
\left\{\begin{aligned}
& Lu+\mu u = F,\\
& u-g\in H_0^1(\Omega)\end{aligned}\right.
\end{equation}
的解有两种情况：
\begin{enumerate}[(1)]
\item 对任意$F\in H^{-1}(\Omega)$，$g\in H^{1}(\Omega)$，\eqref{eq:dirichlet-equation}存在唯一弱解$u\in H^1(\Omega)$；
\item 存在一个非平凡的$u\in H^1_0(\Omega)$使得$Lu+\mu u=0$.
\end{enumerate}
更进一步，满足情形(2)的全体$\mu$构成的集合为可数离散的，$\infty$是唯一可能的极限点，我们称这样的$\mu$为$L$算子在$\Omega$上的Dirichlet特征值，记为$\operatorname{Spec}_{\Omega}(L)$，或者简写为$\operatorname{Spec}(L)$.对于每一Dirichlet特征值$\mu$，其对应的特征函数空间是有限维的.
\end{theorem}

如果增加$c\geq 0$几乎处处成立的假设条件，有如下结论：
\begin{theorem}[Dirichlet问题的$H^1$弱解{\mycite{Chen2elliptic}}]\label{thm:regularity}
令$L$为\eqref{eq:divergence-form}中定义的散度型椭圆微分算子，令$\Omega\subset\mathbb{R}^n$为有界开区域，使得其上成立索伯列夫嵌入定理.进一步假设$c\geq 0$几乎处处成立.对于任意$F\in H^{-1}(\Omega)$和$g\in H^1(\Omega)$，Dirichlet问题
\begin{equation}
\left\{\begin{aligned}
&Lu = F, \\
& u-g\in H^{1}_{0}(\Omega)
\end{aligned}\right.
\end{equation}
存在唯一弱解，而且满足估计式
\begin{equation}
\|u\|_{H^1_0(\Omega)} \leq C\left(\|F\|_{H^{-1}(\Omega)} +\|g\|_{H^1(\Omega)} \right),
\end{equation}
其中常数$C$与$F$和$g$无关.
\end{theorem}

在\eqref{eq:schrodinger-0}中我们仅假设$c\in L^{\infty}(\Omega)$，不满足定理\ref{thm:regularity}中$c\geq 0$几乎处处成立的条件，但是通过避开特征值的假设，我们依然可以利用定理\ref{thm:weak-solution}和定理\ref{thm:regularity}得到\eqref{eq:schrodinger-0}弱解的适定性结论：
\begin{corollary}[线性薛定谔方程\eqref{eq:schrodinger-0}的适定性]\label{cor:Schrodinger-weak-solution}
令$\Omega\subset\mathbb{R}^n$为带$C^1$边界的有界开区域.如果$k^2$不是$-\Delta+c$在$\Omega$上的Dirichlet特征值，则对于任意$F\in H^{-1}(\Omega)$，$f\in H^{1/2}(\partial\Omega)$，Dirichlet问题\eqref{eq:schrodinger-0}存在唯一解$u\in H^1(\Omega)$，且满足估计式
\begin{equation}\label{eq:regularity-weak-solution}
\|u\|_{H^1(\Omega)} \leq C\left(\|F\|_{H^{-1}(\Omega)} +\|f\|_{H^{1/2}(\partial\Omega)} \right),
\end{equation}
其中常数$C$仅仅依赖于$k,c$和$\Omega$.
\end{corollary}
注意到当$c(x)\equiv 0$时，上述推论依然成立，因此我们也可以得到\eqref{eqn:Helmholtz}的适定性.

\begin{proof}
根据定理\ref{thm:weak-solution}，由于$k^2$不是$-\Delta+c$在$\Omega$上的Dirichlet特征值，故我们可以得到$u$的存在唯一性.现在我们证明弱解$u$满足\eqref{eq:regularity-weak-solution}.不失一般性，我们只考虑齐次边界条件$f=0$的情形.

由于$c\in L^{\infty}(\Omega)$，故存在常数$s\in\mathbb{R}$，使得$c(x)+s\leq 0$对于$x\in\Omega$几乎处处成立.我们记$L:=-\Delta +c$，$L_{s}=-\Delta +s +c$.令定理\ref{thm:regularity}中$g\equiv 0$，我们可以得到$L_{s}$为$H_{0}^1(\Omega)\rightarrow H^{-1}(\Omega)$的有界线性算子，且存在连续的逆算子$L_{s}^{-1}:H^{-1}(\Omega)\rightarrow H^{1}_0(\Omega)$.对于$s$使得$k^2+s\ne 0$,我们记$\sigma=(k^2+s)^{-1}$，则对于$u\in H^{1}_0(\Omega)$有如下等价形式
\begin{equation}\label{eq:eu0}
Lu=k^2 u +F \quad \Leftrightarrow \quad L_{s} u = (k^2+s)j\circ i(u) + F
\quad \Leftrightarrow \quad (\sigma I-K)(i(u)) = \sigma i L_s^{-1}F
\end{equation}
其中$i: H^{1}_0(\Omega)\rightarrow L^2(\Omega)$，$j:L^2(\Omega)\rightarrow H^{-1}(\Omega)$为自然的包含映射，$K:=i\circ L_s^{-1} \circ j$.

由于我们假设$k^2\notin\operatorname{Spec}(-\Delta+c)$，根据定理\ref{thm:regularity}的证明过程，$K$为紧自伴算子{\mycite{Chen2elliptic}}.故我们可以取得$s$使得$\sigma\notin\operatorname{Spec}(K)$，此时$\sigma I-K$为$L^2(\Omega)\rightarrow L^2(\Omega)$的可逆线性算子.也就是说，对于任意$F\in H^{-1}(\Omega)$，存在$\tilde{u}\in L^2(\Omega)$，使得
\begin{equation}\label{eq:eu1}
(\sigma I-K)\tilde{u} = \sigma i L_s^{-1}F.
\end{equation}
结合\eqref{eq:eu0}和\eqref{eq:eu1}我们可以得出存在$u\in H^{1}_0(\Omega)$使得$\tilde{u}=i(u)$.根据定理\ref{thm:regularity}以及$\sigma I-K$和$L_{s}$的可逆性，我们有
\begin{equation*}
\begin{aligned}
\|u\|_{H^1(\Omega)} &\leq C\left( \left\|(k^2+s) j(\tilde{u})\right\|_{H^{-1}(\Omega)} + \|F\|_{H^{-1}(\Omega)} \right)\\
&\leq C\left( \|\tilde{u}\|_{L^2(\Omega)} + \|F\|_{H^{-1}(\Omega)} \right)\\
& \leq C\left( \left\|\sigma(\sigma I-K)^{-1}\circ i\circ L^{-1}_s F\right\|_{L^2(\Omega)} + \|F\|_{H^{-1}(\Omega)} \right)\\
& \leq C\|F\|_{H^{-1}(\Omega)}.
\end{aligned}
\end{equation*}

对于非齐次边界条件的情况，即$f\ne 0$，我们可以对$w=u-Ef$应用以上过程得到\eqref{eq:regularity-weak-solution}，其中$E$为$\operatorname{Tr}:H^{1}(\Omega)\rightarrow H^{1/2}(\partial\Omega)$的右逆算子.
\end{proof}


% % !TeX root = ../main.tex

\chapter{图表}
\section{插图}
插图功能是利用TeX的特定编译程序提供的机制实现的，不同的编译程序支持不同的图形方式。XeTeX支持插入 EPS、PDF、PNG、JPEG 格式的图片。一般图形都是处在浮动环境中。之所以称为浮动是指最终排版效果图形的位置不一定与源文
件中的位置对应。如果要强制固定浮动图形的位置，请使用float宏包，它提供了 \texttt{[H]} 参数。

\subsection{单个图形}
\begin{figure}[H]
    \begin{center}
        \includegraphics[width=0.5\textwidth]{./figs/cost.png}
\caption[Alessandrini-PIE等式与组合解$u_{S}^{(0)}$]{在非线性次数$m=3$和波数$k=15$时，组合解$u_{S}^{(0)}$的生成.\textbf{红色矩形（虚线）}:$|\xi| = 45.6$时的组合解$u_{S}^{(0)}$.
\textbf{蓝色矩形（实线）:} Dirichlet边值 $u_{S}^{(0)} \big|_{\partial\Omega}$ （\textcolor{red}{红色曲线}）和线性化Neumann数据的近似$\big( \partial_{\nu} u_{S} - \partial_{\nu} u_{S}^{(0)} \big) \big|_{\partial\Omega}$ （\textcolor{blue}{蓝色曲线}）.
\textbf{蓝色圆圈}：根据Alessandrini-PIE等式做组合.}
\label{fig:superposition}
\end{center}
\end{figure}

\subsection{多个图形}
\begin{figure}[H]
\centering
\textbf{非线性次数} $m = 5$ \\[1ex]
\,\hfill \textbf{(i)} $k = 5$ \hfill\,\\
\includegraphics[width=0.45\textwidth,trim=20 60 20 35,clip]{./figs/pwc_5_Ic_0050.png}
\includegraphics[width=0.54\textwidth,trim=10 60 30 90,clip]{./figs/pwc_5_Fc_0050.png}\\
\,\hfill \textbf{(i)} $k = 10$ \hfill\,\\
\includegraphics[width=0.45\textwidth,trim=20 60 20 35,clip]{./figs/pwc_5_Ic_0100.png}
\includegraphics[width=0.54\textwidth,trim=10 60 30 90,clip]{./figs/pwc_5_Fc_0100.png}\\
\,\hfill \textbf{(i)} $k = 15$ \hfill\,\\
\includegraphics[width=0.45\textwidth,trim=20 60 20 35,clip]{./figs/pwc_5_Ic_0150.png}
\includegraphics[width=0.54\textwidth,trim=10 60 30 90,clip]{./figs/pwc_5_Fc_0150.png}\\
\caption[算法2：不连续位势的重构结果]{当非线性次数$m=5$，波数\textbf{(i)} $k = 5$,\textbf{(ii)} $k = 10$和\textbf{(iii)} $k = 15$时的重构结果.
第一列：利用频域截断$|\xi| \leq (m+1)k$得到的重构位势函数$c_{\textrm{inv}}(x)$.
第二列：真实位势函数和重构位势函数的角度误差$|c(x) - c_{\textrm{inv}}(x)|$.
第三列：重构位势函数的傅立叶系数的绝对值$\mathcal{F}[c_{\textrm{inv}}](\xi)|$.
}
\label{fig:5_pwc}
\end{figure}

如果多个图形相互独立，并不共用一个图形计数器，那么用 \texttt{minipage} 或者
\texttt{parbox} 就可以，如图\ref{fig:parallel1} 与图\ref{fig:parallel2}。

\begin{figure}[H]
  \centering
  \begin{minipage}{0.48\textwidth}
    \centering
    \includegraphics[height=1.5cm]{./figs/fudan-emblem.pdf}
    \caption{并排第一个图}
    \label{fig:parallel1}
  \end{minipage}\hfill
  \begin{minipage}{0.48\textwidth}
    \centering
    \includegraphics[height=1.5cm]{./figs/fudan-emblem.pdf}
    \caption{并排第二个图}
    \label{fig:parallel2}
  \end{minipage}
\end{figure}

如果要为共用一个计数器的多个子图添加子图题，建议使用较新的 subcaption 宏包，不建议使用 subfigure 或 subfig 等宏包。

推荐使用 subcaption宏包的 subcaptionbox并排子图，子图题置于子图之下，子图号用 a)、b) 等表示。也可以使用subcaption宏包的 subcaption（放在 minipage中，用法同 caption）。

subcaption 宏包也提供了subfigure和 subtable 环境，如
图\ref{fig:subfigure}。

\begin{figure}[!htp]
  \centering
  \begin{subfigure}{0.3\textwidth}
    \centering
    \includegraphics[height=2cm]{./figs/fudan-emblem.pdf}
    \caption{校徽}
  \end{subfigure}
  \hspace{1cm}
  \begin{subfigure}{0.4\textwidth}
    \centering
    \includegraphics[height=1.5cm]{./figs/fudan-emblem.pdf}
    \caption{校徽。这个图略矮些，subfigure中同一行的子图在顶端对齐。}
  \end{subfigure}
  \caption{包含子图题的范例（使用 subfigure）}
  \label{fig:subfigure}
\end{figure}

\subsection{Tik作图}
以下摘自Hansimov的知乎文章[LaTeX 绘图指南 - 001] TikZ 的简介、资源以及学习方法\footnote{https://zhuanlan.zhihu.com/p/48300815}。

TikZ 是 LaTeX 下的一个（著名的）绘图宏包。TikZ 的德文原文是 TikZ ist kein Zeichenprogramm， 这是一个 ``GNU's Not Unix!'' 式的递归缩写。翻译成英文就是 TikZ is not a drawing program，中文意思是 ``TikZ 不是一个绘图程序''。（程序员式冷幽默）

PGF/TikZ 相关的学习资源很多，可以参考这个项目：xiaohanyu/awesome-tikz\footnote{https://github.com/xiaohanyu/awesome-tikz}。基本列出了常见的高质资源，语言大多为英文。此外还有一些重要的资源:

\begin{itemize}
    \item 英文文档：pgfmanual\footnote{http://mirrors.ctan.org/graphics/pgf/base/doc/pgfmanual.pdf}
    \item PGF/TikZ 中文手册(在翻)：pgfmanual-zh\footnote{https://github.com/Hansimov/pgfmanual-zh}
    \item 命令大全：VisualTikZ\footnote{http://mirrors.ctan.org/info/visualtikz/VisualTikZ.pdf}
     \item 各种样例：TikZ and PGF examples\footnote{http://www.texample.net/tikz/examples/all/}
     \item TeX 社区：Questions tagged [tikz-pgf]\footnote{https://tex.stackexchange.com/questions/tagged/tikz-pgf}
\end{itemize}

下面是来自于邹森博士论文中的一个样例。
\begin{figure}[H]
\centering
%----------------------%
% original, STYLE 1
\begin{tikzpicture}[>=latex,scale=1.0]
\draw[rounded corners=8,dashed,fill=gray!15!white]
 (0.6,-0.8) rectangle (2.6,0.8) node at ++(-0.3,-0.3) {$\Omega$};
\node  (f) at (0,0) {$f$};
\node  (u) at (2,0) {$u$};
\node (du) at (4,0) {$\partial_{\nu} u$};
\draw[->,thick,blue] (f) -- node[above] {$c$}  (u) node[pos=0.5,below] {\footnotesize$(I)$};
\draw[->,thick] (u) -- (du) node[pos=0.5,below] {\footnotesize$\partial_{\nu}$};
\draw[rounded corners=15,->,thick,blue] (f) |- (2,1.2) node[above] {$\Lambda_{c}$} -| (du);
\node at (2,-1.3) {\textbf{(i)} 原始问题};
\end{tikzpicture}%\\
%----------------------%
% linearized
\hspace{4ex}
\begin{tikzpicture}[>=latex,scale=1.0]
\draw[rounded corners=8,dashed,fill=gray!15!white]
 (0.5,-0.8) rectangle (4.6,0.8) node at ++(-0.3,-0.3) {$\Omega$};
\node   (f) at (0,0) {$f$};
\node  (u0) at (2,0) {$u^{(0)}$};
\node  (u1) at (4,0) {$u^{(1)}$};
\node (du1) at (6,0) {$\partial_{\nu} u^{(1)}$};
\draw[->,thick]  (f) --  (u0) node[pos=0.5,below] {\footnotesize$(I_{0})$};
\draw[->,thick,blue] (u0) -- node[above] {$c$}  (u1) node[pos=0.5,below] {\footnotesize$(I_{1})$};
\draw[->,thick] (u1) -- (du1) node[pos=0.5,below] {\footnotesize$\partial_{\nu}$};
\draw[rounded corners=15,->,thick,blue] (f) |- (3,1.2) node[above] {$\Lambda'_{c}$} -| (du1);
\node at (3,-1.3) {\textbf{(ii)} 线性化系统};
\end{tikzpicture}\\
%----------------------%
% PIE type
\vspace{1ex}
\begin{tikzpicture}[>=latex,scale=1.0]
\draw[rounded corners=8,dashed,fill=gray!15!white]
 (0.5,-1) rectangle (6.6,1) node at ++(-0.3,-0.3) {$\Omega$};
\node   (f) at (0,0) {$f_{j}$};
\draw[rounded corners=2,red] (-0.5,0.7) rectangle (2.45,-0.7) node at ++(-0.4,0.2) {\tiny$j \in S$};
\node  (u0) at (2,0) {$u_{j}^{(0)}$};
\node  (v0) at (4,0) {$u_{S}^{(0)}$};
\node  (v1) at (6,0) {$u_{S}^{(1)}$};
\node (dv1) at (8,0) {$\partial_{\nu} u_{S}^{(1)}$};
\draw[->,thick]  (f) --  (u0) node[pos=0.5,below] {\footnotesize$(I_{0})$};
\draw[->,thick,red] (u0) --  (v0) node[pos=0.5,below] {\footnotesize$\sum\limits_{j \in S}$};
\draw[->,thick,blue] (v0) -- node[above] {$c$}  (v1) node[pos=0.5,below] {\footnotesize$(I_{S}^{(1)})$};
\draw[->,thick] (v1) -- (dv1) node[pos=0.5,below] {\footnotesize$\partial_{\nu}$};
\draw[rounded corners=2] (-0.7,1.3) rectangle (8.7,-1.3) node at ++(-0.4,0.2) {\tiny$S \subseteq U$};
\draw[rounded corners=15,->,thick,blue] (f)++(0,0.7) |- (4,1.6) node[above] {$\Lambda'_{c}$} -| (dv1);
\node at (4,-1.8) {\textbf{(iii)} PIE型线性化系统};
\end{tikzpicture}
%----------------------%
\caption[Alessandrini等式示意图]{三种Alessandrini型等式的对比图示}
\label{fig:identity}
\end{figure}


\section{表格}

\subsection{基本表格}

编排表格应简单明了，表达一致，明晰易懂，表文呼应、内容一致。表题置于表上。

表格的编排建议采用国际通行的三线表\footnote{三线表，以其形式简洁、功能分明、阅读
方便而在科技论文中被推荐使用。三线表通常只有 3 条线，即顶线、底线和栏目线，没有竖线。}。三线表可以使用 booktabs 提供的 toprule、midrule 和bottomrule。它们与 longtable能很好的配合使用。

\begin{table}[!hpt]
  \caption[一个颇为标准的三线表]{一个颇为标准的三线表\footnotemark}
  \label{tab:firstone}
  \centering
  \begin{tabular}{@{}llr@{}} \toprule
    \multicolumn{2}{c}{Item} \\ \cmidrule(r){1-2}
    Animal & Description & Price (\$)\\ \midrule
    Gnat  & per gram  & 13.65 \\
          & each      & 0.01 \\
    Gnu   & stuffed   & 92.50 \\
    Emu   & stuffed   & 33.33 \\
    Armadillo & frozen & 8.99 \\ \bottomrule
  \end{tabular}
\end{table}
\subsection{复杂表格}

我们经常会在表格下方标注数据来源，或者对表格里面的条目进行解释。可以用threeparttable实现带有脚注的表格，如表\ref{tab:footnote}。

\begin{table}[!htpb]
  \bicaption{一个带有脚注的表格的例子}{A Table with footnotes}
  \label{tab:footnote}
  \centering
  \begin{threeparttable}[b]
     \begin{tabular}{ccccccc}
      \toprule
      \multirow{2}*{total} & \multicolumn{2}{c}{20\tnote{a}} & \multicolumn{2}{c}{40} & \multicolumn{2}{c}{60} \\
      \cmidrule(lr){2-3}\cmidrule(lr){4-5}\cmidrule(lr){6-7}
      & www & \multicolumn{1}{c}{k} & www & k & www & k \\ % 使用说明符 d 的列会自动进入数学模式，使用 \multicolumn 对文字表头做特殊处理
      \midrule
      & $\underset{(2.12)}{4.22}$ & 120.0140\tnote{b} & 333.15 & 0.0411 & 444.99 & 0.1387 \\
      & 168.6123 & 10.86 & 255.37 & 0.0353 & 376.14 & 0.1058 \\
      & 6.761    & 0.007 & 235.37 & 0.0267 & 348.66 & 0.1010 \\
      \bottomrule
    \end{tabular}
    \begin{tablenotes}
    \item [a] the first note.
    \item [b] the second note.
    \end{tablenotes}
  \end{threeparttable}
\end{table}

如某个表需要转页接排，可以用longtable实现。接排时表题省略，表头应重复书
写，并在右上方写 ``续表 xx''，如表\ref{tab:performance}。

\begin{ThreePartTable}
  \begin{TableNotes}
    \item[a] 一个脚注
    \item[b] 另一个脚注
  \end{TableNotes}
  \begin{longtable}[c]{c*{6}{r}}
    \bicaption{实验数据}{Experimental data}
    \label{tab:performance} \\
    \toprule
    测试程序 & \multicolumn{1}{c}{正常运行} & \multicolumn{1}{c}{同步}
      & \multicolumn{1}{c}{检查点} & \multicolumn{1}{c}{卷回恢复}
      & \multicolumn{1}{c}{进程迁移} & \multicolumn{1}{c}{检查点} \\
    & \multicolumn{1}{c}{时间 (s)} & \multicolumn{1}{c}{时间 (s)}
      & \multicolumn{1}{c}{时间 (s)} & \multicolumn{1}{c}{时间 (s)}
      & \multicolumn{1}{c}{时间 (s)} &  文件（KB）\\
    \midrule
    \endfirsthead
    \multicolumn{7}{l}{\textbf{续表~\thetable}} \\
    % 英语论文：\multicolumn{7}{r}{\textbf{Table~\thetable~(continued)}} \\
    \toprule
    测试程序 & \multicolumn{1}{c}{正常运行} & \multicolumn{1}{c}{同步}
      & \multicolumn{1}{c}{检查点} & \multicolumn{1}{c}{卷回恢复}
      & \multicolumn{1}{c}{进程迁移} & \multicolumn{1}{c}{检查点} \\
    & \multicolumn{1}{c}{时间 (s)} & \multicolumn{1}{c}{时间 (s)}
      & \multicolumn{1}{c}{时间 (s)} & \multicolumn{1}{c}{时间 (s)}
      & \multicolumn{1}{c}{时间 (s)}&  文件（KB）\\
    \midrule
    \endhead
    \hline
    \multicolumn{7}{r}{续下页}
    \endfoot
    \insertTableNotes
    \endlastfoot
    CG.A.2 & 23.05 & 0.002 & 0.116 & 0.035 & 0.589 & 32491 \\
    CG.A.4 & 15.06 & 0.003 & 0.067 & 0.021 & 0.351 & 18211 \\
    CG.A.8 & 13.38 & 0.004 & 0.072 & 0.023 & 0.210 & 9890 \\
    CG.B.2 & 867.45 & 0.002 & 0.864 & 0.232 & 3.256 & 228562 \\
    CG.B.4 & 501.61 & 0.003 & 0.438 & 0.136 & 2.075 & 123862 \\
    CG.B.8 & 384.65 & 0.004 & 0.457 & 0.108 & 1.235 & 63777 \\
    MG.A.2 & 112.27 & 0.002 & 0.846 & 0.237 & 3.930 & 236473 \\
    MG.A.4 & 59.84 & 0.003 & 0.442 & 0.128 & 2.070 & 123875 \\
    MG.A.8 & 31.38 & 0.003 & 0.476 & 0.114 & 1.041 & 60627 \\
    MG.B.2 & 526.28 & 0.002 & 0.821 & 0.238 & 4.176 & 236635 \\
    MG.B.4 & 280.11 & 0.003 & 0.432 & 0.130 & 1.706 & 123793 \\
    MG.B.8 & 148.29 & 0.003 & 0.442 & 0.116 & 0.893 & 60600 \\
    LU.A.2 & 2116.54 & 0.002 & 0.110 & 0.030 & 0.532 & 28754 \\
    LU.A.4 & 1102.50 & 0.002 & 0.069 & 0.017 & 0.255 & 14915 \\
    LU.A.8 & 574.47 & 0.003 & 0.067 & 0.016 & 0.192 & 8655 \\
    LU.B.2 & 9712.87 & 0.002 & 0.357 & 0.104 & 1.734 & 101975 \\
    LU.B.4 & 4757.80 & 0.003 & 0.190 & 0.056 & 0.808 & 53522 \\
    LU.B.8 & 2444.05 & 0.004 & 0.222 & 0.057 & 0.548 & 30134 \\
    EP.A.2 & 123.81 & 0.002 & 0.010 & 0.003 & 0.074 & 1834 \\
    EP.A.4 & 61.92 & 0.003 & 0.011 & 0.004 & 0.073 & 1743 \\
    EP.A.8 & 31.06 & 0.004 & 0.017 & 0.005 & 0.073 & 1661 \\
    EP.B.2 & 495.49 & 0.001 & 0.009 & 0.003 & 0.196 & 2011 \\
    EP.B.4 & 247.69 & 0.002 & 0.012 & 0.004 & 0.122 & 1663 \\
    EP.B.8 & 126.74 & 0.003 & 0.017 & 0.005 & 0.083 & 1656 \\
    SP.A.2 & 123.81 & 0.002 & 0.010 & 0.003 & 0.074 & 1854 \\
    SP.A.4 & 51.92 & 0.003 & 0.011 & 0.004 & 0.073 & 1543 \\
    SP.A.8 & 31.06 & 0.004 & 0.017 & 0.005 & 0.073 & 1671 \\
    SP.B.2 & 495.49 & 0.001 & 0.009 & 0.003 & 0.196 & 2411 \\
    SP.B.4 \tnote{a} & 247.69 & 0.002 & 0.014 & 0.006 & 0.152 & 2653 \\
    SP.B.8 \tnote{b} & 126.74 & 0.003 & 0.017 & 0.005 & 0.082 & 1755 \\
    \bottomrule
  \end{longtable}
\end{ThreePartTable}

\section{算法环境}

算法环境可以使用 agorithms宏包或者较新的algorithm2e实现。
算法\ref{algo:algorithm} 是一个使用algorithm2e的例子。关于排版算法环境
的具体方法，请阅读相关宏包的官方文档。

\begin{algorithm}[htb]
  \caption{算法示例}
  \label{algo:algorithm}
  \small
  \SetAlgoLined
  \KwData{this text}
  \KwResult{how to write algorithm with \LaTeXe }

  initialization\;
  \While{not at end of this document}{
    read current\;
    \eIf{understand}{
      go to next section\;
      current section becomes this one\;
    }{
      go back to the beginning of current section\;
    }
  }
\end{algorithm}

同样可以考虑以如下表格的形式实现。
\begin{table}[H] 
\centering
\begin{tabular}{p{\textwidth}}
\toprule
\textbf{算法1：3维情况下部分边界观测反问题$\lambda_c\rightarrow c$的神经网络结构}  \\%
\midrule
\textbf{输入：} %
$\lambda\in\mathbb{R}^{N_{m_1}\times N_{m_2}\times N_{h_1}\times N_{h_2}}$，参数$channel,N_{z},w,n_{cnn}$；
\textbf{输出：} %
$c\in \mathbb{R}^{N_x\times N_y\times N_y}$. \\[-15pt]%
\begin{enumerate}[1:]
  \item 将$\lambda_c$的后两个维度向量化，形成一个三维张量，记为$\lambda$;
  \item $\tilde{\lambda}_{\hat{h}}\leftarrow\texttt{Encoding3d}[channel](\lambda)$; %
  \item  $\tilde{c}_{\hat{z}_3}\leftarrow\texttt{BCR-Net2d}(\tilde{\lambda}_{\hat{h}}) $; %
  \item $\bar{c}\leftarrow\texttt{Decoding3d}[N_z](\tilde{c}_{\hat{z_3}})$;%
  \item $c\leftarrow\texttt{CNN3d}[w,n_{cnn}](\bar{c})$;%
  \item 返回 $c$ %
\end{enumerate} \\%
\bottomrule
\end{tabular}
\caption*{算法 1: 线性薛定谔位势反问题的部分边界观测问题的神经网络结构\mycite{FY2020}}
\end{table}





\chapter{Introduction}
\section{Background and Motivation}
\section{Assumptions}
\begin{assumption}[Supercritical Regime]
    \label{assumption:psi-zero}
    $\psi(0)>0$.
\end{assumption}
\begin{assumption}[At Least One Child]
    \label{assumption:at-least-one-child}
    Under \Cref{assumption:psi-zero}, we further assume that
    $\#\Xi\ge 1$ a.s.
\end{assumption}
\begin{assumption}[Exponential Moment]
    \label{assumption:psi-one}
    $\psi(1)=0$.
\end{assumption}
\begin{assumption}[Zero Mean]
    \label{assumption:psi-prime-one}
    $\psi'(1)=0$.
\end{assumption}
\begin{assumption}[Finite Variance]
  \label{assumption:finite-variance}
  \begin{equation*}
        \mathbb E[\xi^2]
        =
        \mathbb E\!\left[
            \sum_{x\in\Xi}
            x^2e^{-x}
        \right]
        =\sigma^2<\infty.
  \end{equation*}
\end{assumption}
\begin{assumption}[Finite Cross Log-Laplacian]
    \label{assumption:log-expt-cross-sum-exp-neg-sum}
    \begin{equation*}
        \log\mathbb E\!\left[
            \sum_{
                i,j\in\mathbb N,\,
                i\neq j
            }
            e^{
                -(\Xi_i+\Xi_j)
            }
        \right]
        <
        \infty.
    \end{equation*}
\end{assumption}
\begin{assumption}[First Moment of Positive Part of Leftmost Step]
    \label{assumption:moment-1}
    Under \Cref{assumption:at-least-one-child}, we further assume that
    $\mathbb E[(\Xi_1)_+]<\infty$.
\end{assumption}
\begin{assumption}[Fourth Moment of Positive Part of Leftmost Step]
    \label{assumption:moment-4}
    Under \Cref{assumption:moment-1}, we further assume that
    $\mathbb E[(\Xi_1)_+^4]<\infty$.
\end{assumption}
\begin{assumption}[At Least Two Children]
    \label{assumption:at-least-two-children}
    Under \Cref{assumption:at-least-one-child}, we further assume that
    $\#\Xi\ge 2$ a.s.
\end{assumption}
\begin{assumption}[Exponential Moment of Second Leftmost Step]
    \label{assumption:moment-exp-2}
    Under \Cref{assumption:at-least-two-children}, we further assume that there exists some $\lambda>0$ such that
    \begin{equation*}
        \mathbb E\left[e^{\lambda\Xi_2}\right]<\infty.
    \end{equation*}
\end{assumption}
\section{Main Results}
\section{Methodologies}

\chapter{Assumptions and Known Results}
% !TeX root = ../main.tex
% !TeX root = ../../main.tex
\section{\bilingual{Formalization and Branching Properties}{形式化与分支性质}}
\label{section:formalization-and-branching-properties}

\bilingual{
    Consider an implicit probability space $(\Omega,\mathcal F,\mathbb P)$, on which all random objects in the following are defined. For the value space of vertices of trees, we adopt the terminology of Ulam--Harris--Neveu \mycite{Neveu1986}, which is defined as follows
    \begin{equation*}
        \mathbb V:=\bigcup_{n=0}^\infty\mathbb N^n,
    \end{equation*}
    where $\mathbb N^0:=\{\emptyset\}$, $\emptyset$ is the root of the tree, $\mathbb N^n$ is the $n$-dimensional natural number space. The length function $|\cdot|:\mathbb V\to\mathbb N$ defined on it is called the generation.  Denote $u_i$ as the ancestor of $u\in\mathbb V$ at generation $i$. A {\itshape tree} (tree) or {\itshape rooted labeled tree} is a subset $T$ of $\mathbb V$ that satisfies the following conditions:
    \begin{itemize}
        \item $\emptyset\in T$;
        \item If $u\in T$ and $u=vw$ where $v,w\in\mathbb V$, $vw$ denotes the concatenation of the two, then $v\in T$;
        \item If $uj\in T$, $j\in\mathbb N$, then $uj'\in T$ for any $j'\in[\![1,j]\!]$.
    \end{itemize}
    The set of all trees is denoted by $\mathbb T$. A {\itshape finite tree} is a tree with a finite number of vertices. A {\itshape locally finite tree} is a tree whose truncation at each generation is finite, and the set of all locally finite trees is denoted by $\mathbb T_0$. The space $\mathbb T$ naturally defines a metric
    \begin{equation*}
        d(T_1,T_2):=e^{-\inf\{n\ge 0:T_1\cap\mathbb N^n\neq T_2\cap\mathbb N^n\}},
    \end{equation*}
    where $T_1,T_2\in\mathbb T$, but it is only complete on the space of locally finite trees $\mathbb T_0$, whose countable dense set is the set of all finite trees. For the $\sigma$-algebra on $\mathbb T$, we use the coarser $\sigma$-algebra $\mathcal G_{\mathbb T}$ generated by the cylinder sets
    \begin{equation*}
        \left\{
            T\in\mathbb T: u\in T
        \right\},\quad u\in\mathbb V
    \end{equation*}
    Denote the random tree taking values in $(\mathbb T,\mathcal G_{\mathbb T})$ as $\mathcal T$.
}{
    考虑一个隐含的概率空间$(\Omega,\mathcal F,\mathbb P)$, 后文的所有随机对象都定义在这个概率空间上. 对于树上节点的取值空间, 采用 Ulam--Harris--Neveu \mycite{Neveu1986} 的术语, 其定义如下
    \begin{equation*}
        \mathbb V:=\bigcup_{n=0}^\infty\mathbb N^n,
    \end{equation*}
    其中 $\mathbb N^0:=\{\emptyset\}$, $\emptyset$ 是树的根节点, $\mathbb N^n$ 是 $n$ 维的自然数空间. 其上的长度函数$|\cdot|:\mathbb V\to\mathbb N$称为代数. 记$u_i$为$u\in\mathbb V$的第$i$代的祖先. 所谓{\kaishu 树} (tree) 或{\kaishu 有根标签树} (rooted labeled tree), 是指 $\mathbb V$ 的一个子集 $T$ 满足如下条件:
    \begin{itemize}
        \item $\emptyset\in T$;
        \item 若 $u\in T$ 且 $u=vw$ 其中 $v,w\in\mathbb V$, $vw$指两者的连接, 则 $v\in T$;
        \item 若 $uj\in T$, $j\in\mathbb N$, 则 $uj'\in T$ 对任意 $j'\in[\![1,j]\!]$ 成立.
    \end{itemize}
    树的全体记作$\mathbb T$. 所谓{\kaishu 有限树} (finite tree), 是指节点数有限的树. 所谓{\kaishu 局部有限树} (locally finite tree), 是指每一层的截断都有限的树, 全体记作$\mathbb T_0$. 一般采用$\mathbb T_0$空间作为考察对象, 但本文需要考虑无限分支的分支律, 所得到的树可能不满足局部有限的条件, 因此也需要考虑$\mathbb T$空间. 树空间$\mathbb T$上自然定义有一个度量
    \begin{equation*}
        d(T_1,T_2):=e^{-\inf\{n\ge 0:T_1\cap\mathbb N^n\neq T_2\cap\mathbb N^n\}},
    \end{equation*}
    其中 $T_1,T_2\in\mathbb T$, 但它仅在局部有限树空间$\mathbb T_0$上完备, 其可数稠密集为有限树全体. 对于$\mathbb T$上的$\sigma$-代数, 转而采用柱集
    \begin{equation*}
        \left\{
            T\in\mathbb T: u\in T
        \right\},\quad u\in\mathbb V
    \end{equation*}
    生成的更粗的$\sigma$-代数$\mathcal G_{\mathbb T}$. 记取值于$(\mathbb T,\mathcal G_{\mathbb T})$的随机树为$\mathcal T$.
}

\bilingual{
    Sometimes, we also consider a family of trees indexed by $\mathbb I$ (with a common root $\emptyset$), whose whole is $\mathbb T^{\mathbb I}$, and its $\sigma$-algebra is $\mathcal G_{\mathbb T}^{\otimes\mathbb I}$, which is also generated by cylinder sets. Denote the random tree family taking values in $(\mathbb T^{\mathbb I},\mathcal G_{\mathbb T}^{\otimes\mathbb I})$ as $\boldsymbol{\mathcal T}$.  Moreover, for some $(u,i)\in\boldsymbol T$, denote $\boldsymbol T_{u,i}$ as the subtree rooted at $u$, which is defined as
    \begin{equation*}
        \boldsymbol T_{u,i}:=\{v\in\mathbb V:uv\in\boldsymbol T(i)\}.
    \end{equation*}
}{
    有时也会考虑由$\mathbb I$索引的一族 (有共同根$\emptyset$的) 树$\boldsymbol T$, 其全体为$\mathbb T^{\mathbb I}$, 其$\sigma$-代数为$\mathcal G_{\mathbb T}^{\otimes\mathbb I}$, 事实上也是由柱集生成的. 记取值于$(\mathbb T^{\mathbb I},\mathcal G_{\mathbb T}^{\otimes\mathbb I})$的随机树族为$\boldsymbol{\mathcal T}$. 另外, 对于某个$(u,i)\in\boldsymbol T$, 记$\boldsymbol T_{u,i}$为以$u$为根节点的子树, 其定义为
    \begin{equation*}
        \boldsymbol T_{u,i}:=\{v\in\mathbb V:uv\in\boldsymbol T(i)\}.
    \end{equation*}
}

\bilingual{
    The so-called {\itshape marked tree} is defined as follows: it is a tree $T\in\mathbb T$ equipped with a marking function $V:T\to\mathbb R$; denoted as $(T,V)$ or simply $V$. We require that $V(\emptyset)=0$. The set of all marked trees is denoted by
    \begin{equation*}
        \mathbb M:=\{(T,V):T\in\mathbb T,\,V:T\to\mathbb R\}.
    \end{equation*}
    The $\sigma$-algebra $\mathcal G_{\mathbb M}$ on it is generated by the cylinder sets
    \begin{equation*}
        \left\{(T,V)\in\mathbb M: u\in T,\,V(u)\in A\right\},\quad u\in\mathbb V,\,A\in\mathcal B(\mathbb R).
    \end{equation*}
    Denote the random marked tree taking values in $(\mathbb M,\mathcal G_{\mathbb M})$ as $(\mathcal T,\mathcal V)$, sometimes simply denoted as $\mathcal V$, and called a {\itshape branching random walk} (BRW).
}{
    所谓{\kaishu 标记树} (marked tree), 是指在某棵树$T\in\mathbb T$的基础上, 还定义了一个标记函数 $V:T\to\mathbb R$; 记作$(T,V)$或直接记为$V$. 要求$V(\varnothing)=0$. 标记树全体记为
    \begin{equation*}
        \mathbb M:=\{(T,V):T\in\mathbb T,\,V:T\to\mathbb R\}.
    \end{equation*}
    其上的$\sigma$-代数$\mathcal G_{\mathbb M}$由柱集
    \begin{equation*}
        \left\{(T,V)\in\mathbb M: u\in T,\,V(u)\in A\right\},\quad u\in\mathbb V,\,A\in\mathcal B(\mathbb R)
    \end{equation*}
    生成, 记取值于$(\mathbb M,\mathcal G_{\mathbb M})$的随机标记树为$(\mathcal T,\mathcal V)$, 有时也直接记为$\mathcal V$, 称为{\kaishu 分支随机游走} (branching random walk, BRW).
}

\bilingual{
    Similarly, we also consider a family of marked trees indexed by $\mathbb I$. For a family of trees $\boldsymbol T\in\mathbb T^{\mathbb I}$, the marking function is still denoted as $V$, but its domain is considered as
    $
        \left\{
            (u,i):u\in\boldsymbol T(i),\,i\in\mathbb I
        \right\}
    $. We require that $V(\varnothing,i)=0$. Such marked trees are denoted as $(\boldsymbol T,V)$ or simply $V$, and the set of all such marked trees is denoted by
    \begin{equation*}
        \mathbb M^{\mathbb I}:=\{(\boldsymbol T,V):\boldsymbol T\in\mathbb T^{\mathbb I},\,V:\{(u,i):u\in\boldsymbol T(i),\,i\in\mathbb I\}\to\mathbb R\}.
    \end{equation*}
    Its $\sigma$-algebra $\mathcal G_{\mathbb M}^{\otimes\mathbb I}$ is generated by the cylinder sets
    \begin{equation*}
        \left\{(\boldsymbol T,V)\in\mathbb M^{\mathbb I}: u\in\boldsymbol T(i),\,V(u,i)\in A\right\},\quad u\in\mathbb V,\,i\in\mathbb I,\,A\in\mathcal B(\mathbb R)
    \end{equation*}
    Denote the random marked tree family taking values in $(\mathbb M^{\mathbb I},\mathcal G_{\mathbb M}^{\otimes\mathbb I})$ as $(\boldsymbol{\mathcal T},\mathcal V)$, sometimes simply denoted as $\mathcal V$.
}{
    同样地, 也会考虑由$\mathbb I$索引的一族标记树. 对于一族树$\boldsymbol T\in\mathbb T^{\mathbb I}$, 标记函数仍记作$V$, 但其定义域视作
    $
        \left\{
            (u,i):u\in\boldsymbol T(i),\,i\in\mathbb I
        \right\}
    $.
    要求$V(\varnothing,i)=0$. 这样的标记树族记为$(\boldsymbol T,V)$或直接记为$V$, 其全体记为
    \begin{equation*}
        \mathbb M^{\mathbb I}:=\{(\boldsymbol T,V):\boldsymbol T\in\mathbb T^{\mathbb I},\,V:\{(u,i):u\in\boldsymbol T(i),\,i\in\mathbb I\}\to\mathbb R\}.
    \end{equation*}
    其$\sigma$-代数$\mathcal G_{\mathbb M}^{\otimes\mathbb I}$由柱集
    \begin{equation*}
        \left\{(\boldsymbol T,V)\in\mathbb M^{\mathbb I}: u\in\boldsymbol T(i),\,V(u,i)\in A\right\},\quad u\in\mathbb V,\,i\in\mathbb I,\,A\in\mathcal B(\mathbb R)
    \end{equation*}
    生成, 记取值于$(\mathbb M^{\mathbb I},\mathcal G_{\mathbb M}^{\otimes\mathbb I})$的随机标记树族为$(\boldsymbol{\mathcal T},\mathcal V)$, 或直接记为$\mathcal V$.
}

\bilingual{
    Now consider the countable-length {\itshape real vector space} $\mathbb R_{\mathrm V}$, which is defined as
    \begin{equation*}
        \mathbb R_{\mathrm V}:=\mathbb R^{\mathbb N}\cup\bigcup_{n=0}^\infty\mathbb R^n,
    \end{equation*}
    where $\mathbb R^0:=\{\emptyset\}$, $\emptyset$ is the zero vector of zero dimension. The length function $\#:\mathbb R_{\mathrm V}\to\mathbb N\cup\{\infty\}$ and the component function $(\cdot)^i:\mathbb R_{\mathrm V}\to\mathbb R$ are defined on it. The $\sigma$-algebra $\mathcal G_{\mathbb R_{\mathrm V}}$ on it is generated by the family of sets
    \begin{equation*}
        \left\{
            x\in\mathbb R_{\mathrm V}: \#x\ge i,\,x^i\in A
        \right\},\quad i\in\mathbb N,\,A\in\mathcal B(\mathbb R)
    \end{equation*}
    The so-called branching law $\Xi$ is a random vector taking values in $(\mathbb R_{\mathrm V},\mathcal G_{\mathbb R_{\mathrm V}})$, whose component $\Xi^i$ represents the displacement of the $i$-th offspring, and its length $\#\Xi$ represents the number of offspring. In the following, under certain assumptions, if there exists a leftmost component, we will denote $\Xi_i$ as the $i$-th leftmost (smallest) component of $\Xi$.
}{
    现在考虑包含可数长度向量的{\kaishu 实向量空间} $\mathbb R_{\mathrm V}$, 即
    \begin{equation*}
        \mathbb R_{\mathrm V}:=\mathbb R^{\mathbb N}\cup\bigcup_{n=0}^\infty\mathbb R^n,
    \end{equation*}
    其中 $\mathbb R^0:=\{\emptyset\}$, $\emptyset$ 是零维的零向量. 其上定义有长度函数$\#:\mathbb R_{\mathrm V}\to\mathbb N\cup\{\infty\}$ 与分量函数$(\cdot)^i:\mathbb R_{\mathrm V}\to\mathbb R$. 其上的$\sigma$-代数$\mathcal G_{\mathbb R_{\mathrm V}}$由如下集族
    \begin{equation*}
        \left\{
            x\in\mathbb R_{\mathrm V}: \#x\ge i,\,x^i\in A
        \right\},\quad i\in\mathbb N,\,A\in\mathcal B(\mathbb R)
    \end{equation*}
    生成. 所谓分支律$\Xi$, 就是一个取值于$(\mathbb R_{\mathrm V},\mathcal G_{\mathbb R_{\mathrm V}})$的随机向量, 其分量$\Xi^i$表示第$i$个子代的位移, 其长度$\#\Xi$表示子代的数量. 下文中, 在一定的假设下, 若存在最左侧分量, 将用$\Xi_i$表示$\Xi$的最左侧 (最小的) 第$i$个分量.
}

\bilingual{
    Given a branching law $\Xi$, consider an index set $\mathbb I$, then a family of i.i.d. branching laws $\{\Xi_{u,i}\}_{u\in\mathbb V,\,i\in\mathbb I}$ determines a random marked tree family $(\boldsymbol{\mathcal T},\mathcal V)$, where $\boldsymbol{\mathcal T}(i)$ is the Galton--Watson tree generated by this family of branching laws in sequence, $\mathcal V(u,i)$ is the displacement of vertex $u$ on tree $\boldsymbol{\mathcal T}(i)$, satisfying the following relation
    \begin{equation*}
        \Xi_{u,i}=\mathcal V(u,i)-\mathcal V(u_{|u|-1},i),
        \quad i\in\mathbb I,\,u\in\boldsymbol{\mathcal T}(i)-\{\emptyset\}.
    \end{equation*}
    In the absence of ambiguity, the index $i$ is usually omitted. Moreover, denote the set of all particles in the $n$-th generation as
    \begin{equation*}
        \mathcal Z_n:=\{(u,i):u\in\boldsymbol{\mathcal T}(i),\,i\in\mathbb I,\,|u|=n\},
    \end{equation*}
    and its quantity as $\#\mathcal Z_n$. Also denote the infimum of the particle positions in the $n$-th generation as
    \begin{equation*}
        \mathcal M_n:=\inf\{\mathcal V(u,i):(u,i)\in\mathcal Z_n\}.
    \end{equation*}
    In the following, under certain assumptions, this infimum exists and is a minimum.
}{
    给定分支律 $\Xi$ 与索引集 $\mathbb I$. 对每个 $(u,i)\in\mathbb V\times\mathbb I$ 取一个独立同分布的分支律 $\Xi_{u,i}$, 由此逐代生成随机标记树族 $(\boldsymbol{\mathcal T},\mathcal V)$. 每棵树 $\boldsymbol{\mathcal T}(i)$ 都是 Galton--Watson 树, 节点位移满足
    \begin{equation*}
        \Xi_{u,i}=\mathcal V(u,i)-\mathcal V(u_{|u|-1},i),
        \quad i\in\mathbb I,\,u\in\boldsymbol{\mathcal T}(i)-\{\emptyset\}.
    \end{equation*}
    在不造成歧义的情况下, 通常省略索引$i$. 另外记第$n$代的粒子全体为
    \begin{equation*}
        \mathcal Z_n:=\{(u,i):u\in\boldsymbol{\mathcal T}(i),\,i\in\mathbb I,\,|u|=n\},
    \end{equation*}
    其数量记为$\#\mathcal Z_n$. 还记第$n$代的粒子位置的下确界为
    \begin{equation*}
        \mathcal M_n:=\inf\{\mathcal V(u,i):(u,i)\in\mathcal Z_n\}.
    \end{equation*}
    下文中, 在一定的假设下, 此下确界存在且为最小值.
}

\bilingual{
    Now consider the generation filtration generated by the family of branching laws $(\Xi_{u,i})$ on the implicit probability space. Define
    \begin{equation*}
        \mathcal F_n:=\begin{cases}
            \sigma\left(
                \Xi_{u,i}:u\in\mathbb V,\,i\in\mathbb I,\,|u|\le n-1
            \right),& n\in\mathbb Z_+,\\
            \{\emptyset,\Omega\},& n=0,
        \end{cases}
    \end{equation*}
    which is also equal to
    $
        \sigma\left(
            \mathcal V(u,i):u\in\mathbb V,\,i\in\mathbb I,\,|u|\le n
        \right)
    $.
}{
    现在考虑分支律族$(\Xi_{u,i})$在隐含概率空间上生成的逐代域流. 定义
    \begin{equation*}
        \mathcal F_n:=\begin{cases}
            \sigma\left(
                \Xi_{u,i}:u\in\mathbb V,\,i\in\mathbb I,\,|u|\le n-1
            \right),& n\in\mathbb Z_+,\\
            \{\emptyset,\Omega\},& n=0,
        \end{cases}
    \end{equation*}
    它也等于
    $
        \sigma\left(
            \mathcal V(u,i):u\in\mathbb V,\,i\in\mathbb I,\,|u|\le n
        \right)
    $.
}

\bilingual{
    The following branching properties are important tools in this thesis.
}{
    下面的分支性质是本文的重要工具.
}

\begin{lemma}[\bilingual{Branching Property}{分支性质}]\it\label{lemma:branching-property}
    \bilingual{
        Given a branching law $\Xi$. With the same notation as above. For any $n\in\mathbb Z_+$, conditional on $\mathcal F_n$ and $\#\mathcal Z_n\ge 1$, the family of sub-marked trees rooted at $\{(u,i):u\in\mathcal Z_n\}$, i.e., $\{(\boldsymbol{\mathcal T}\!\!{}_{u,i},\,v\mapsto\mathcal V(uv,i)-\mathcal V(u,i)):(u,i)\in\mathcal Z_n\}$, is a family of i.i.d. random marked trees, whose distribution is the same as some $(\boldsymbol{\mathcal T}(i_0),\mathcal V(\cdot,i_0))$, $i_0\in\mathbb I$.
    }{
        给定分支律$\Xi$, 以及由它决定的随机标记树族$(\boldsymbol{\mathcal T},\mathcal V)$, 则对于任意 $n\in\mathbb Z_+$, 条件于 $\mathcal F_n$ 以及 $\#\mathcal Z_n\ge 1$, 以 $\{(u,i):u\in\mathcal Z_n\}$ 为根节点的子标记树族$\{(\boldsymbol{\mathcal T}\!\!{}_{u,i},\,v\mapsto\mathcal V(uv,i)-\mathcal V(u,i)):(u,i)\in\mathcal Z_n\}$ 是一族独立同分布的随机标记树, 其分布与某个$(\boldsymbol{\mathcal T}(i_0),\mathcal V(\cdot,i_0))$, $i_0\in\mathbb I$相同.
    }
\end{lemma}

\begin{lemma}[\bilingual{Branching Property at Stopping Times}{停时处的分支性质}]\it
    \label{lemma:branching-property-stopping-time}
    \bilingual{
        Given a branching law $\Xi$. With the same notation as above. For an $(\mathcal F_n)$-stopping time $\tau$, conditional on $\mathcal F_\tau$ and $\#\mathcal Z_\tau\ge 1$, the family of sub-marked trees rooted at $\{(u,i):u\in\mathcal Z_\tau\}$, i.e., $\{(\boldsymbol{\mathcal T}\!\!{}_{u,i},\,v\mapsto\mathcal V(uv,i)-\mathcal V(u,i)):(u,i)\in\mathcal Z_\tau\}$, is a family of i.i.d. random marked trees, whose distribution is the same as some $(\boldsymbol{\mathcal T}(i_0),\mathcal V(\cdot,i_0))$, $i_0\in\mathbb I$.
    }{
        给定分支律 $\Xi$. 记号同上. 对于$(\mathcal F_n)$-停时$\tau$, 条件于 $\mathcal F_\tau$ 以及 $\#\mathcal Z_\tau\ge 1$, 以 $\{(u,i):u\in\mathcal Z_\tau\}$ 为根节点的子标记树族$\{(\boldsymbol{\mathcal T}\!\!{}_{u,i},\,v\mapsto\mathcal V(uv,i)-\mathcal V(u,i)):(u,i)\in\mathcal Z_\tau\}$ 是一族独立同分布的随机标记树, 其分布与某个$(\boldsymbol{\mathcal T}(i_0),\mathcal V(\cdot,i_0))$, $i_0\in\mathbb I$相同.
    }
\end{lemma}
\bilingual{
    The above branching properties are special cases of stronger branching properties about {\itshape stopping lines} \mycite{chauvin1986propriete,JAGERS1989183}. A {\itshape line} $\ell$ of a tree is a set of vertices in $\mathbb V$ that are pairwise non-ancestral, the corresponding $\sigma$-algebra $\mathcal G_\ell$ is generated by the branching r.v. at points not in $\ell$ and their descendants, a stopping line is a random line such that any event of being before any deterministic line is measurable.
}{
    \noindent 上述分支性质是更强版本 \mycite{chauvin1986propriete,JAGERS1989183} 的关于{\kaishu 停止线} (stopping lines) 的分支性质的特例. 所谓{\kaishu 线} (line)是指$\mathbb V$中两两无继承关系的点集 $\ell$, 对应的$\sigma$-代数$\mathcal G_\ell$由非$\ell$及其后代处的分支生成, 所谓停止线是指使得其位置在任何确定线前的事件均可测的随机线.
}

\raggedbottom

% !TeX root = ../../main.tex
\section{\bilingual{Many-to-One Formula and Spine Decomposition}{多对一公式与脊柱分解}}
    \bilingual{
        Define the log-Laplace transform of a branching law $\Xi$ as
        \begin{equation*}
            \psi(\lambda):=\log\mathbb E\left[\sum_{x\in\Xi}e^{-\lambda x}\right],\quad \lambda\in\mathbb R,
        \end{equation*}
        whose basic properties are given in Appendix \ref{section:lemmas-on-assumptions}. Given a branching law $\Xi$ and a random marked tree $\mathcal V$ determined by it (we do not consider random marked tree families in this section). If there exists $\lambda>0$ such that $\psi(\lambda)<\infty$ (guaranteed by \Cref{assumption:psi-one} in our assumptions, i.e., $\psi(1)=0<\infty$), define the {\itshape additive martingale} $W_n=W_n^\lambda$ for $\Xi$ as
        \begin{equation*}
            W_n^\lambda:=\sum_{|u|=n}e^{-\lambda\mathcal V(u)-\psi(\lambda)n},\quad n\in\mathbb Z_+,
        \end{equation*}
        where the summation is taken over $\{u\in\mathcal T:|u|=n\}$. This additive martingale is adapted to the natural filtration.
    }{
        定义分支律$\Xi$的log-Laplace变换为
        \begin{equation*}
            \psi(\lambda):=\log\mathbb E\left[\sum_{x\in\Xi}e^{-\lambda x}\right],\quad \lambda\in\mathbb R,
        \end{equation*}
        其基本性质见 附录 \ref{section:lemmas-on-assumptions}. 给定分支律$\Xi$与它决定的一个随机标记树$\mathcal V$ (本节不考虑随机标记树族). 若存在$\lambda>0$使得$\psi(\lambda)<\infty$ (在本文的假设中已由\Cref{assumption:psi-one} 保证, i.e., $\psi(1)=0<\infty$), 为$\Xi$定义{\kaishu 加性鞅} (additive martingale) $W_n=W_n^\lambda$为
        \begin{equation*}
            W_n^\lambda:=\sum_{|u|=n}e^{-\lambda\mathcal V(u)-\psi(\lambda)n},\quad n\in\mathbb Z_+,
        \end{equation*}
        其中求和范围是$\{u\in\mathcal T:|u|=n\}$. 此加性鞅适应于自然流.
    }

    \bilingual{
        Let $\xi$ be a random variable determined by $\Xi$ with the following distribution, i.e., for any measurable function $h:\mathbb R\to\mathbb R_+$, we have
        \begin{equation*}
            \mathbb E[h(\xi)]
            =
            \frac{
                \mathbb E\left[
                    \sum_{x\in\Xi}
                    e^{-\lambda x}h(x)
                \right]
            }{
                \mathbb E\left[
                    \sum_{x\in\Xi}
                    e^{-\lambda x}
                \right]
            }.
        \end{equation*}
        Consider the random walk $Y_n:=\sum_{i=1}^n \xi_i$, where $\xi_i$ are i.i.d. copies of $\xi$. Then we have the following important tool, which transforms certain expectations about branching random walks into certain expectations about the above random walk, which is the so-called {\itshape many-to-one formula}, see Shi \mycite{shi2015branching}.
    }{
        令$\xi$为由$\Xi$决定的具有如下分布的随机变量, 即对任意可测函数 $h:\mathbb R\to\mathbb R_+$, 有
        \begin{equation*}
            \mathbb E[h(\xi)]
            =
            \frac{
                \mathbb E\left[
                    \sum_{x\in\Xi}
                    e^{-\lambda x}h(x)
                \right]
            }{
                \mathbb E\left[
                    \sum_{x\in\Xi}
                    e^{-\lambda x}
                \right]
            }.
        \end{equation*}
        考虑由$\xi$决定的随机游走$Y_n:=\sum_{i=1}^n \xi_i$, 其中 $\xi_i$是$\xi$的独立同分布的拷贝. 则有如下重要的工具, 它将某些关于分支随机游走的期望转化为某些关于上述随机游走的期望, 这就是所谓{\kaishu 多对一公式} (many-to-one formula), 见 Shi \mycite{shi2015branching}.
    }
    \begin{lemma}[\bilingual{Many-to-One Formula}{多对一公式}]\it\label{lemma:many-to-one}
        \bilingual{
            For any $n\in\mathbb Z_+$ and any measurable function $g:\mathbb R^n\to\mathbb R_+$,~ we have
        }{
            对任意 $n\in\mathbb Z_+$ 和任意可测函数 $g:\mathbb R^n\to\mathbb R_+$, 有
        }
        \begin{equation*}
            \mathbb E\left[
                \sum_{|u|=n}
                g(\mathcal V(u_1),\ldots,\mathcal V(u_n))
            \right]
            =
            \mathbb E\left[
                e^{\lambda Y_n+\psi(\lambda)n}g(Y_1,\ldots,Y_n)
            \right],
        \end{equation*}
        或者
        \begin{equation*}
            \mathbb E\left[
                \sum_{|u|=n}
                e^{-\lambda \mathcal V(u)-\psi(\lambda)n}
                g(\mathcal V(u_1),\ldots,\mathcal V(u_n))
            \right]
            =
            \mathbb E\left[
                g(Y_1,\ldots,Y_n)
            \right],
        \end{equation*}
        \bilingual{
            where $\lambda$ is any positive number such that $\psi(\lambda)<\infty$,~ and $Y_n$ is the random walk determined by $\xi$.
        }{
            其中$\lambda$是满足$\psi(\lambda)<\infty$的任意正数, $Y_n$是由$\xi$决定的随机游走.
        }
    \end{lemma}

    \bilingual{
        The above many-to-one formula is a consequence of the following spine decomposition theory. Consider an additional {\itshape infinite ray} space $\partial\mathbb V=\mathbb N^{\mathbb N}$. Given a branching law $\Xi$, we call a random variable $(\hat{\mathcal T},\hat{\mathcal V},w)$ taking values in $\mathbb M\times\partial\mathbb V$ a {\itshape size-biased BRW} determined by $\Xi$, if:
        \begin{itemize}
            \item $w_0=\emptyset$ is the root, the vertex $w_{i+1}=(w(0),\cdots,w(i+1))$, $i\in\mathbb N$ is an offspring of $w_i$ according to the branching law $\hat\Xi$, where the distribution of $\hat\Xi$ is given by, for any bounded measurable function $F$ on the space of branching laws, we have
            \begin{equation*}
                \mathbb E[F(\hat\Xi)]
                =
                \mathbb E\left[
                    F(\Xi)
                    \sum_{x\in\Xi}
                    e^{-\lambda x-\psi(\lambda)}
                \right]
                =
                \mathbb E\left[
                    F(\Xi)
                    W_1^\lambda
                \right],
            \end{equation*}
            and the distribution of the natural number-valued random variable $w(i+1)$ is given by
            \begin{equation*}
                \mathbb P(w(i+1)=k)
                =
                \frac{
                    e^{-\lambda\hat\Xi^k}
                }{
                    \sum_{k\in\mathbb N}
                    e^{-\lambda\hat\Xi^k}
                },
                \quad k\in\mathbb N.
            \end{equation*}
            \item The branching law of the branching random walk $(\hat{\mathcal T},\hat{\mathcal V})$ at vertex $w_i$, $i\in\mathbb N$ is $\hat\Xi$, while at other vertices it is $\Xi$.
        \end{itemize}
        Then, the random walk $\hat{\mathcal V}(w_n)$ here has the same distribution as $Y_n$ defined above (thus the above \nameref{lemma:many-to-one} holds for it); and if we perform a measure change with Radon--Nikodym derivative $W_n^\lambda$ on the implicit probability space at each generation $n\in\mathbb Z_+$, then for the branching random walk $(\mathcal T,\mathcal V)$ determined by $\Xi$, its truncated distribution up to generation $n$ under the new measure is the same as the distribution of $(\hat{\mathcal T},\hat{\mathcal V})$ under the old measure $\mathbb P$. This means that, by Kolmogorov's extension theorem, there always exists a probability measure $\mathbb Q$ such that the branching random walk almost surely does not become extinct, and there exists a ``spine'', such that the original branching random walk can be decomposed into a spine and the results obtained by further branching at each vertex on this spine, and the overall behavior of the original branching random walk can be analyzed through the behavior of the spine using \nameref{lemma:many-to-one}.
    }{
        上述随机游走也称为原分支随机游走的{\kaishu 脊柱}(spine). 此因我们有如下脊柱分解理论. 考虑额外的一个{\kaishu 无限射线} (infinite ray) 空间 $\partial\mathbb V=\mathbb N^{\mathbb N}$. 给定分支律$\Xi$, 称取值于$\mathbb M\times\partial\mathbb V$的随机变量$(\hat{\mathcal T},\hat{\mathcal V},w)$为$\Xi$决定的{\kaishu 有偏分支随机游走} (size-biased BRW), 若:
        \begin{itemize}
            \item $w_0=\emptyset$ 是根节点, 节点$w_{i+1}=(w(0),\cdots,w(i+1))$, $i\in\mathbb N$是$w_i$根据分支律$\hat\Xi$得到的一个子代, 其中$\hat\Xi$的分布为, 对任意实向量空间上的有界可测函数$F$成立
            \begin{equation*}
                \mathbb E[F(\hat\Xi)]
                =
                \mathbb E\left[
                    F(\Xi)
                    \sum_{x\in\Xi}
                    e^{-\lambda x-\psi(\lambda)}
                \right]
                =
                \mathbb E\left[
                    F(\Xi)
                    W_1^\lambda
                \right],
            \end{equation*}
            而自然数取值的随机变量$w(i+1)$的分布为
            \begin{equation*}
                \mathbb P(w(i+1)=k)
                =
                \frac{
                    e^{-\lambda\hat\Xi^k}
                }{
                    \sum_{k\in\mathbb N}
                    e^{-\lambda\hat\Xi^k}
                },
                \quad k\in\mathbb N.
            \end{equation*}
            \item 分支随机游走$(\hat{\mathcal T},\hat{\mathcal V})$在节点$w_i$, $i\in\mathbb N$处的分支律为$\hat\Xi$, 在其余节点处的分支律为$\Xi$.
        \end{itemize}
        那么, 这里的随机游走$\hat{\mathcal V}(w_n)$与上文定义的$Y_n$同分布 (从而上述\nameref{lemma:many-to-one}就对它成立); 且在每一代$n\in\mathbb Z_+$若对隐含概率空间进行Radon--Nikodym导数为$W_n^\lambda$的测度变换, 则对于由$\Xi$决定的分支随机游走$(\mathcal T,\mathcal V)$而言, 它的在前$n$代的截断在新测度下的分布与$(\hat{\mathcal T},\hat{\mathcal V})$的在旧测度$\mathbb P$下的分布相同. 这意味着, 通过 Kolmogorov 扩张定理, 总是存在一个概率测度$\mathbb Q$使得分支随机游走几乎必然不灭绝, 且存在一根``脊柱'', 使得原分支随机游走可被分解为一根脊柱与这根脊柱上各节点处再行分支所得到的结果, 并且原分支随机游走的整体行为可以通过\nameref{lemma:many-to-one}用脊柱的行为来分析.
    }

\raggedbottom

% !TeX root = ../../main.tex
\section{\bilingual{Explanations on Assumptions}{有关本文所作假设的说明}}\label{section:explanations-on-assumptions}
    \bilingual{
        The following assumptions are widely used in the literature on $N$-branching random walks, see Chapter 5 of Shi \mycite{shi2015branching} and Mallein~\mycite{Mallein_2017}.
    }{
        以下假设在$N$-分支随机游走的文献中被广泛使用,参见 Shi \mycite{shi2015branching} 第 5 章以及 Mallein~\mycite{Mallein_2017}.
    }

    \bilingual{
        (\Cref{assumption:psi-zero}). i.e., $\psi(0)=\log \mathbb E[\#\Xi]>0$. In this case, the Galton--Watson process corresponding to the branching law $\Xi$ is supercritical and has a positive probability of survival, but still has $\mathbb P(\#\Xi=0)\ge 0$, so the GW process can still become extinct with positive probability;~ in particular, when $\mathbb P(\#\Xi=0)>0$, the $N$-branching random walk must become extinct, because we can consider the population process as a Markov chain on the state space $[\![0,N]\!]$: since there is a positive lower bound for the extinction probability of any generation, it will almost surely be absorbed into the extinction state.
    }{
        (\Cref{assumption:psi-zero}).即$\psi(0)=\log\mathbb E[\#\Xi]>0$. 此时,该分支律$\Xi$对应的Galton-Watson过程是超临界的,具有正的生存概率,但仍有$\mathbb P(\#\Xi=0)\ge 0$,~ 该GW过程仍可以正概率灭绝;特别地, 当$\mathbb P(\#\Xi=0)>0$时,该$N$-分支随机游走必然灭绝,这是因为,若考虑粒子总数过程作为$[\![0,N]\!]$状态空间上的Markov链:因为整代灭绝概率有正下界,所以它几乎必然被吸入灭绝状态.
    }

    \bilingual{
        (\Cref{assumption:at-least-one-child}). i.e., $\#\Xi\ge 1$ almost surely. Or equivalently, $\mathbb P(\#\Xi=0)=0$, which excludes the above situation that leads to the extinction of the $N$-branching random walk. In fact, in this case, the process almost surely survives. To extend to the general supercritical case, we can use the reduced tree, which retains only the nodes with infinite descendants to form a ``skeleton'' subtree; moreover, we can also discuss under the $\mathbb Q$-measure that ensures the permanent survival of the spine through an additive martingale measure change.
    }{
        (\Cref{assumption:at-least-one-child}).即$\#\Xi\ge 1$几乎必然成立.或者等价地,~ $\mathbb P(\#\Xi=0)=0$,~ 这排除了上述导致$N$-分支随机游走灭绝的情形. 事实上此时该过程几乎必然存活.若需推广至一般的超临界情形,可通过{\kaishu 化简树} (reduced tree)实现,即仅保留具有无限后代的节点构成的``骨干''子树;此外,亦可通过加性鞅的测度变换，在确保脊柱永久存活的 $\mathbb Q$-测度下进行讨论.
    }

    \bilingual{
        (\Cref{assumption:psi-one,assumption:psi-prime-one}). i.e., $\psi(1)=0=\psi'(1)$. According to \nameref{lemma:many-to-one},
    }{
        (\Cref{assumption:psi-one,assumption:psi-prime-one}). 即$\psi(1)=0=\psi'(1)$. ~根据\nameref{lemma:many-to-one},
    }
    \begin{equation*}
        0=\psi'(1)=
        \frac{
            \mathbb E\left[
                \sum_{x\in\Xi}
                e^{-x}(-x)
            \right]
        }{
            \mathbb E\left[
                \sum_{x\in\Xi}
                e^{-x}
            \right]
        }
        =
        \frac{
            \mathbb E\left[
                -\xi
            \right]
        }{
            e^{\psi(1)}
        }
        =
        -\mathbb E[\xi]
        ,
    \end{equation*}
    \bilingual{
        thus the single-step increment $\xi$ of the spine random walk has zero mean. Moreover, for the general case, i.e., there exists $\lambda>0$ such that $\psi(\lambda)<\infty$, in most cases it can be transformed to satisfy the conditions of this assumption. In fact, as long as there exists $0<\lambda^*<\inf\{\lambda>0:\psi(\lambda)<\infty\}$ such that $\psi'(\lambda^*)=\frac{\psi(\lambda^*)}{\lambda^*}$, then we can use the transformation $\tilde{\mathcal V}(x):=\lambda^*\mathcal V(x)+\psi(\lambda^*)|x|$ to make the new branching law satisfy this assumption. As for how to transform a general branching random walk to the boundary case as this assumption, see Jaffuel \mycite{jaffuel2009criticalrandombarriersurvival}.
    }{
        因此脊柱随机游走的单步步长$\xi$具有零均值.另外,对于一般的情形,即存在$\lambda>0$使得$\psi(\lambda)<\infty$的情形,在绝大多数情况下可以通过变换满足本假设的条件.事实上,只要存在$0<\lambda^*<\inf\{\lambda>0:\psi(\lambda)<\infty\}$使得$\psi'(\lambda^*)=\frac{\psi(\lambda^*)}{\lambda^*}$,~ 就可以使用变换$\tilde{\mathcal V}(x):=\lambda^*\mathcal V(x)+\psi(\lambda^*)|x|$使得新的分支律满足本假设. 对于如何将一般的分支随机游走转化为此假设中的边界情形,参见 Jaffuel \mycite{jaffuel2009criticalrandombarriersurvival}.
    }

    \bilingual{
        In particular, when there exists $\lambda>0$ such that $\psi(\lambda)<\infty$, the leftmost position $\Xi_1$ of the distribution $\Xi$ is well-defined, and the infimum position $\mathcal M_n$ defined above for the non-selective branching random walk is almost surely attainable, thus it is a minimum. In fact, for any real number $A$, the total number of particles $N(A)$ produced by the distribution $\Xi$ that fall into $(-\infty,A]$ is almost surely finite ({\itshape locally finite on the left}), because
    }{
        特别地,当存在$\lambda>0$使得$\psi(\lambda)<\infty$时,上文提到的分支律$\Xi$的最左侧位置$\Xi_1$是良定义的,另外对于无选择的分支随机游走,上文定义的下确界位置 $\mathcal M_n$ 几乎必然可达到,从而是最小值.事实上,对于任意实数$A$,~ 分布律$\Xi$产生的子代落于$(-\infty,A]$的粒子总数$N(A)$几乎必然有限 ({\kaishu 左侧局部有限}),此因
    }
    \begin{align*}
        \mathbb E[N(A)]
        &
        =
        \mathbb E\left[
            \sum_{x\in\Xi}
            \mathbf 1_{\{x\le A\}}
        \right]
        =
        e^{\lambda A}
        \mathbb E\left[
            \sum_{x\in\Xi}
            e^{-\lambda A}
            \mathbf 1_{\{x\le A\}}
        \right]
        \\ &
        \le
        e^{\lambda A}
        \mathbb E\left[
            \sum_{x\in\Xi}
            e^{-\lambda x}
        \right]
        =
        e^{\lambda A+\psi(\lambda)}
        <\infty.
    \end{align*}

    \bilingual{
        (\Cref{assumption:finite-variance}). i.e., $\psi''(1)=\sigma^2<\infty$. According to \nameref{lemma:many-to-one} and the above assumption,
    }{
        (\Cref{assumption:finite-variance}).即$\psi''(1)=\sigma^2<\infty$. ~根据\nameref{lemma:many-to-one}和上条假设,
    }
    \begin{equation*}
        \psi''(1)
        =
        \left.\left(
            \frac{
                \mathbb E\left[
                    \sum_{x\in\Xi}
                    e^{-\lambda x}(-\lambda x)
                \right]
            }{
                e^{\psi(\lambda)}
            }
        \right)'\right|_{\lambda=1}
        =
        \mathbb E\left[
            \sum_{x\in\Xi}
            e^{-x}x^2
        \right]
        =
        \mathbb E[\xi^2]
        <\infty,
    \end{equation*}
    \bilingual{
        Thus the single-step increment $\xi$ of the spine random walk has finite variance. This makes this assumption a special case of the single-step increment with $\alpha$-stable distribution when $\alpha=2$~\mycite{Mallein_2017}.
    }{
        因此脊柱随机游走单步步长$\xi$具有有限方差.这使得本假设成为具有$\alpha$-稳定分布的单步步长情形在$\alpha=2$时的特例\mycite{Mallein_2017}.
    }

    \bilingual{
        (\Cref{assumption:moment-1,assumption:moment-4,assumption:moment-exp-1}). i.e., $\mathbb E[(\Xi_1)_+]<\infty$, $\mathbb E[(\Xi_1)_+^4]<\infty$ and even there exists $\lambda>0$ such that $\mathbb E[e^{\lambda \Xi_1}]<\infty$. These assumptions restrict the right tail heaviness of the leftmost offspring, preventing its single extreme displacement from interfering with the asymptotic stability of the frontier. It should be noted that even if the spine random walk has zero mean (\Cref{assumption:psi-prime-one}), due to the negative exponential weighting in \nameref{lemma:many-to-one}, any positive displacement will be exponentially compressed and ``ineffective'' in the calculation of its contribution to the mean of the spine, thus allowing for extreme values on the right side of the leftmost displacement $\Xi_1$, which is why this assumption is needed. As a comparison, \Cref{assumption:psi-one} (i.e., $\psi(1)=0$) ensures that the left tail of the branching law has at least an exponential decay rate. This double-sided moment constraint together ensures that when branching at each generation, the spread of the particle cloud is effectively controlled, thus making the leftmost position $\mathcal M_n$ have good stochastic stability.
    }{
        (\Cref{assumption:moment-1,assumption:moment-4,assumption:moment-exp-1}).即$\mathbb E[(\Xi_1)_+]<\infty$, $\mathbb E[(\Xi_1)_+^4]<\infty$ 乃至存在$\lambda>0$ 使得 $\mathbb E[e^{\lambda \Xi_1}]<\infty$. ~这些假设限制了最左侧子代{\kaishu 向右跳跃的尾部厚度} (right tail heaviness),防止其单次极端位移干扰前沿的渐近稳定性.需要注意,即使脊柱随机游走是零均值的 (\Cref{assumption:psi-prime-one}),根据\nameref{lemma:many-to-one}中的负指数权重,任何正向位移对脊柱均值的贡献都会被指数级压缩而在计算中``失效'',于是最左侧位移$\Xi_1$也被允许有右侧极端值,故而需要本条的假设.作为比较，\Cref{assumption:psi-one} (i.e. $\psi(1)=0$) 则保证了分支律的{\kaishu 左尾} (left tail) 具有至少指数级的衰减速率.这种双侧的矩约束共同确保了随机标记树在每一代分支时,粒子云团的{\kaishu 跨度} (spread) 受到有效控制,从而使得最左侧位置 $\mathcal M_n$ 具有良好的随机稳定性.
    }

    \bilingual{
        Under these assumptions, some properties of the leftmost offspring distribution $\Xi_1$ and the leftmost position $\mathcal M^N$ of the $N$-branching random walk can be derived, see \Cref{lemma:moment-1:neg,lemma:M:exp-neg-moment,lemma:M:moment-1:neg,lemma:M:moment-4:neg,lemma:M:moment-1:pos,lemma:M:moment-4:pos}.
    }{
        在这些假设下能推得最左侧子代分布$\Xi_1$、以及$N$-分支随机游走最左侧位置$\mathcal M^N$的一些性质,详见\Cref{lemma:moment-1:neg,lemma:M:exp-neg-moment,lemma:M:moment-1:neg,lemma:M:moment-4:neg,lemma:M:moment-1:pos,lemma:M:moment-4:pos}.
    }

    \bilingual{
        The last assumption, \Cref{assumption:log-expt-cross-sum-exp-neg-sum}, i.e.,
    }{
        余下的\Cref{assumption:log-expt-cross-sum-exp-neg-sum},即
    }
    \begin{equation*}
        \log\mathbb E\!\left[
            \sum_{
                i,j\in\mathbb N,\,
                i\neq j
            }
            e^{
                -(\Xi_i+\Xi_j)
            }
        \right]
        <
        \infty,
    \end{equation*}
    \bilingual{
        is a technical condition that is used in the proof of the second-order result in this thesis to ensure the finiteness of certain second moments, thus allowing the use of the Paley--Zygmund inequality, see \Cref{lemma:4.5}. On the contrary, if this expectation is infinite, then in the same branching, there may be situations where two or more offspring simultaneously make extreme left jumps, which can lead to an explosive leftward spread of the population in a certain generation, potentially breaking the asymptotic stability of the frontier. In fact, using truncation methods, this condition can be further weakened.
    }{
        在本文的二阶方法证明中被用来保证某些二阶矩的有限性,从而得以使用Paley--Zygmund不等式,详见\Cref{lemma:4.5}.相反,如果这个期望值达到无穷,在同一次分支中,可能会出现两个或以上的子代同时向左极端跳跃的情形,这会导致种群在某一代爆发式地向左极端扩散,可能使得前沿的渐近稳定性被破坏.事实上,使用截断方法,这一条件可以被进一步减弱.
    }

\raggedbottom

% !TeX root = ../../main.tex
\section{\bilingual{\texorpdfstring{Mogul'ski\u i}{Mogul'skii} Theorem on Small Deviation of Random Walks with \texorpdfstring{$\alpha=2$}{alpha=2}}{\texorpdfstring{$\alpha=2$}{alpha=2}步长随机游走的小偏差\texorpdfstring{Mogul'ski\u i}{Mogul'skii}定理}}
    \bilingual{
        Consider some independent identically distributed random variables $(\xi_i)_{i\ge 1}$ and some sequence $(x(n))$ indexed by $n$. Let
    }{

        Mogul'ski\u i \mycite{doi:10.1137/1119081} 给出了一个关于具$\alpha$-稳定步长随机游走的小偏差结果. ~考虑某个独立同分布的随机变量列 $(\xi_i)_{i\ge 1}$ 和某个以 $n$ 为索引的序列 $(x(n))$. ~令
    }
    \begin{equation}
        s_n(t):=\frac{1}{x(n)}\sum_{i=0}^{\lfloor nt\rfloor}\xi_i,\quad
        t\in[0,1],\,\xi_0=0.
        \label{eq:mogulskii-s-n-t}
    \end{equation}
    \bilingual{
        $s_n$ is thus defined as a random real function taking values in $D(0,1)$ (i.e., real functions on $[0,1]$ such that càdlàg on $[0,1)$ and left-continuous at $1$).
    }{
        因此,~ $s_n$ 被定义为一个随机实函数,~ 取值于 $D(0,1)$ (即在 $[0,1)$ 上是{\kaishu 左极右连} (càdlàg) 的、在 $1$ 上是左连续的实函数).
    }

    \bilingual{
        Mogul'ski\u i defined some functional $H_\alpha$ for $\alpha\in(0,2]$ on some set $\mathfrak N$ of $D(0,1)$. We state here only the definition for $G\in\mathfrak N$ of the form
    }{
        Mogul'ski\u i 为 $\alpha\in(0,2]$ 和 $D(0,1)$ 某子集 $\mathfrak N$ 中函数定义了一个泛函 $H_\alpha$. ~这里仅陈述 $G\in\mathfrak N$ 的定义,~ 其形式如下
    }
    \begin{equation}
        G=\{f\in D(0,1):f(0)=0,\,S_2(t)<f(t)<S_1(t),\,\forall t\in[0,1]\},
        \label{eq:mogulskii-G}
    \end{equation}
    \bilingual{
        where $S_1,S_2$ are right-continuous and piecewise constant such that
    }{
        其中 $S_1,S_2$ 是右连续且分段常数的函数,~ 满足
    }
    \begin{equation}
        S_2(t-0)\vee S_2(t+0)
        <
        S_1(t-0)\wedge S_1(t+0),\quad
        t\in[0,1].
        \label{eq:mogulskii-S-1-S-2}
    \end{equation}
    \bilingual{
        and we have the definition
    }{
        并且有定义
    }
    \begin{equation}
        H_\alpha(G):=\int_0^1(S_1(t)-S_2(t))^{-\alpha}\,\mathrm dt.
        \label{eq:mogulskii-h-alpha}
    \end{equation}

    \bilingual{
        Recall the Mogul'ski\u i's Theorem for $\alpha=2$.
    }{
        这里引用 $\alpha=2$ 的 Mogul'ski\u i 定理.
    }

    \begin{theorem}[\bilingual{\texorpdfstring{Mogul'ski\u i}{Mogul'skii}'s Theorem}{\texorpdfstring{Mogul'ski\u i}{Mogul'skii}定理}]\it\label{theorem:mogulskii}
        \bilingual{
            Suppose that $\mathbb E\xi_i=0$, $\mathbb E\xi_i^2=1$ and that the sequence $(x(n))$ is such that
        }{
            假设 $\mathbb E\xi_i=0$,~ $\mathbb E\xi_i^2=1$ 并且序列 $(x(n))$ 满足
        }
        \begin{equation*}
            \lim_{n\to\infty}x(n)=\infty,\quad
            \lim_{n\to\infty}x(n)n^{-1/2}=0.
        \end{equation*}
        \bilingual{
            Then the relation
        }{
            则关系式
        }
        \begin{equation*}
            \log\mathbb P(s_n\in G)\sim-\frac{\pi^2}{2}H_2(G)nx^{-2}(n)
        \end{equation*}
        \bilingual{
            holds for set $G\in\mathfrak N$.
        }{
            对于集合 $G\in\mathfrak N$ 成立.
        }
    \end{theorem}
    \bilingual{
        Its meaning is that, for a random walk with zero mean and unit variance, to obtain small deviation results, it is necessary to scale the space-time as the number of steps of the walk grows, as long as the growth rate of the time scale is faster than the square of the growth rate of the space scale, then the probability that the trajectory of this random walk falls into some small deviation space (the function set $G$) in the scaled space-time will decay exponentially, with a decay rate being the growth rate of the time scale divided by the square of the growth rate of the space scale, and the constant in the decay term being $-\frac{\pi^2}{2}H_2(G)$, which only depends on the shape of $G$. In this thesis, this theorem will be used to analyze the survival probability of $N$-branching random walks in a specified small deviation scaled space-time $G$ (which will usually be a tube), which can first be transformed into the corresponding survival probability of spine random walks in that space-time $G$ using \nameref{lemma:many-to-one}, and then obtain exact decay rates and constant $-\frac{\pi^2}2$.
    }{
        \noindent 它的含义是说,~ 对于具有零均值与单位方差的随机游走,~ 为了得到小偏差的结果,~ 需要随游走的步数增长而对时空进行缩放,~ 只要时间尺度的增长速率快于空间尺度的增长速率的平方,~ 那么该随机游走在缩放后的时空中的轨迹落于某个小偏差时空区域 (函数集合 $G$) 的概率,~ 将以指数级衰减,~ 衰减的速率是时间尺度的增长速率除去空间尺度增长速率的平方,~ 衰减项的常数为 $-\frac{\pi^2}{2}H_2(G)$,~ 只与$G$的形状有关. ~在本文中,~ 该定理将被用来分析 $N$-分支随机游走在指定的小偏差缩放时空区域$G$ (通常将是一个管道) 中的存活概率,~ 这可以先利用\nameref{lemma:many-to-one}转化为对应的脊柱随机游走在该时空区域$G$中的存活概率,~ 进而得到确切的衰减速率与常数$-\frac{\pi^2}2$.
    }

\raggedbottom

% !TeX root = ../../main.tex
\section{\bilingual{Asymptotic Speeds of the Extremal Positions in \texorpdfstring{$N$}{N}-BRW of I.I.D. Bi-Branching Law}{同分布二分支\texorpdfstring{$N$}{N}-分支随机游走最值的渐近速率}}
    \bilingual{
        B\'erard and Gou\'er\'e \mycite{bérard2010brunetderridabehaviorbranchingselectionparticle} used Kingman's subadditive ergodic theorem to show that, for the binary branching law $\Xi=(p,p)$, if the branching law satisfies the condition that there exists $\lambda>0$ such that $\psi(\lambda)>0$ and $\psi(-\lambda)>0$, then the asymptotic speeds of the leftmost and rightmost positions of the $N$-branching random walk starting with $N$ particles exist in the sense of $\mathrm L^1$. This paper will follow this method, but under a more general branching law that can be infinitely branching, it will be shown that the asymptotic speed of the single-sided leftmost position of the $N$-branching random walk starting with at most $N$ particles exists in the sense of $\mathrm L^1$, see \Cref{corollary:leftmost-speed:N0}.

        Moreover, they also showed that this asymptotic speed has an order of $v_\infty(p)-\chi(p)\log^{-2}N+o(\log^{-2}N)$, where
        \begin{equation*}
            \chi(p)=\frac{\pi^2}{2}\lambda^*\psi''(\lambda^*),
        \end{equation*}
        and $\lambda^*>0$ satisfies $\lambda^*\psi'(-\lambda^*)-\psi(-\lambda^*)=\log 2$. The final result about the order of the asymptotic speed of the single-sided leftmost position of $N$-branching random walk will also be given in this thesis, see \Cref{theorem:speed}, but mainly using the conclusion of \nameref{theorem:mogulskii} and the method of A\"id\'ekon and Hu \mycite{aïdékon2025averageweightpercolationcompletegraph}.
    }{
        B\'erard 与 Gou\'er\'e \mycite{bérard2010brunetderridabehaviorbranchingselectionparticle} 利用 Kingman 次可加遍历定理得出, 在二分支的分支律$\Xi=(p,p)$下, 若分支律满足存在$\lambda>0$使得$\psi(\lambda)>0$与$\psi(-\lambda)>0$的条件, 则从$N$个粒子出发的$N$-分支随机游走最左侧位置与最右侧位置渐近速率在$\mathrm L^1$的意义下存在. 本文将遵照这一方法, 但在更一般的可无限分支的分支律下, 证明从少于等于$N$个粒子出发的$N$-分支随机游走单侧最左侧位置渐近速率在$\mathrm L^1$的意义下存在, 详见\Cref{corollary:leftmost-speed:N0}.

        另外, 他们还证得该渐近速率具有阶数$v_\infty(p)-\chi(p)\log^{-2}N+o(\log^{-2}N)$, 其中
        \begin{equation*}
            \chi(p)=\frac{\pi^2}{2}\lambda^*\psi''(\lambda^*),
        \end{equation*}
        其中$\lambda^*>0$满足$\lambda^*\psi'(-\lambda^*)-\psi(-\lambda^*)=\log 2$. 本文最后也将给出一个关于从少于等于$N$个粒子出发的$N$-分支随机游走单侧最左侧位置渐近速率阶数的结论, 详见\Cref{theorem:speed}, 但主要用到\nameref{theorem:mogulskii}的结论和A\"id\'ekon 与 Hu \mycite{aïdékon2025averageweightpercolationcompletegraph}的方法.
    }

\raggedbottom

% !TeX root = ../../main.tex
\section{\bilingual{Limit Theorem for Log-Scaled Space-Time \texorpdfstring{$N$}{N}-BRW}{对数缩放后一类 \texorpdfstring{\!$N$}{N}-分支随机游走的极限定理}}

    \bilingual{
        A\"id\'ekon and Hu \mycite{aïdékon2025averageweightpercolationcompletegraph} proved a limit theorem for a special branching law $\Xi$ under which the $N$-branching random walk is defined. This branching law is given by
        \begin{equation*}
            \Xi=\sum_{i=1}^N\delta_{eT_i-1},
        \end{equation*}
        where $0<T_1<T_2<\cdots$ is a Poisson point process on $\mathbb R_+$ with parameter $1$. Their limit theorem states that, for the $N$-branching random walk starting with $\lfloor N^a\rfloor$ particles at position $0$ at time $0$ ($a\in[0,1]$), under the space-time scaling of log-cubed in time and single logarithm in space, the position of the leftmost particle at each time slice $t>0$ converges in $\mathrm L^2$ to a position on a straight line $-a+\frac{\pi^2}2t$ in this log-scaled space-time as $N\to\infty$, where $\frac{\pi^2}2$ comes from the application of \nameref{theorem:mogulskii}. In this thesis, following their method, we will generalize their result to more general assumptions in \Cref{proposition:3.1} in Section \ref{section:main-results}. Moreover, in the final conclusion about the order of the asymptotic speed of the leftmost particle of $N$-branching random walk (\Cref{theorem:speed}), this generalized result will be applied.
    }{
        A\"id\'ekon 与 Hu \mycite{aïdékon2025averageweightpercolationcompletegraph} 为一种特殊分支律 $\Xi$ 下的 $N$-分支随机游走证得了一个极限定理. 该分支律为
        \begin{equation*}
            \Xi=\sum_{i=1}^N\delta_{eT_i-1},
        \end{equation*}
        其中 $0<T_1<T_2<\cdots$ 是一个 $\mathbb R_+$ 上以 $1$ 为参数的 Poisson 点过程. 他们的极限定理表明,~ 对于零时刻从位置 $0$ 以 $\lfloor N^a\rfloor$ 个粒子出发的 $N$-分支随机游走 ($a\in[0,1]$),~ 在三次对数的时间与一次对数的空间下,~ 每一个时间切片 $t>0$ 处的最左侧粒子的位置在 $N\to\infty$ 下的 $\mathrm L^2$-极限位置落于该对数缩放时空下的一条直线$-a+\frac{\pi^2}2t$上该时间切片处,~ 其中 $\frac{\pi^2}2$ 来自于对\nameref{theorem:mogulskii}的应用. ~本文将在第 \ref{section:main-results} 节中的\Cref{proposition:3.1},~ 遵照他们的方法,~ 推广他们的结果至更一般的假设条件. ~并且,~ 在最后的关于 $N$-分支随机游走最左侧粒子的渐近速率阶数的结论 (\Cref{theorem:speed}) 中,~ 这一推广的结果将被应用.
    }



\chapter{Generalization of Known Results and Their Corollaries}
\section{Lemmas on Applying Mogul'skii's Theorem}

Now consider a sequence $(L_\Delta)_{\Delta\ge 0}$ indexed by $\Delta$ with $\lim_{\Delta\to\infty}L_\Delta=\infty$. Basically, when applying Mogul'skii's Theorem, we are letting $n$ there to be $L_\Delta$ and $x(n)$ there to be $\Delta$, to form an inverse sequence. We are in this section to make rigorous this statement.

Let
\begin{align*}
    x(n)&:=\inf\{\Delta\ge 0: L_\Delta+1\ge n\},&n\ge 0,
    \\ y(n)&:= x(n+1)-1,&n\ge 0.
\end{align*}

\begin{lemma}\it\label{lemma:mogulskii-x-y}
The two sequences $x$ and $y$ satisfy
\begin{equation*}
    \lim_{n\to\infty}x(n)=\lim_{n\to\infty}y(n)=\infty.
\end{equation*}
If we further require that $\lim_{\Delta\to\infty}\frac{L_\Delta}{\Delta^2}=\infty$, we have
\begin{equation*}
    \lim_{n\to\infty}x(n)n^{-1/2}
    =\lim_{n\to\infty}y(n)n^{-1/2}
    =0.
\end{equation*}
\end{lemma}
The proof for Lemma \ref{lemma:mogulskii-x-y} is omitted. We shall use the fact $L_{x(n)}\ge n-1$ and thus $0<\frac{x^2(n)}{n}<\frac{x^2(n)}{L_{x(n)}}$.

Consider the random walk $(Y_n)_{n\ge 0}$ with zero mean and finite variance $0<\sigma^2<\infty$ obtained from the many-to-one lemma. For any $a\in(0,1)$, we denote in this section
\begin{align*}
    q(\Delta,n)&:=\log\mathbb P\left(
        -a\Delta\le Y_j\le (1-a)\Delta,\,
        \forall j\in[\![1,n]\!]
    \right),\\
    p(\Delta,n)&:=\frac{\Delta^2}nq(\Delta,n).
\end{align*}
For any $\Delta\ge 0$, there exists $n=n(\Delta)=L_\Delta+1$ such that $x(n)\le\Delta<x(n+1)$,  and we can get by monotonicity of probability that
\begin{lemma} \it\label{lemma:mogulskii-mono}
    For $\Delta_1\le\Delta_2$ such that $L:=L_{\Delta_1}=L_{\Delta_2}$, we have $q(\Delta_1,L)\le q(\Delta_2,L)$.
\end{lemma}

We are now able to get to the main lemma with an additional condition for $(L_\Delta)$.
\begin{lemma}
    \it
    \label{lemma:mogulskii-cal}
    \quad
    \begin{enumerate}[label=(\roman*)]
        \item By Mogul'skii's Theorem, we have, indexed by $n$,
            \begin{equation*}
                \lim_{n\to\infty}p(x(n),n-1)
                =\lim_{n\to\infty}p(y(n),n-1)
                =-\frac{\pi^2}2\sigma^2.
            \end{equation*}
        \item For $n=n(\Delta)$, and indexed by $\Delta$,
            \begin{equation*}
                \lim_{\Delta\to\infty}p(x(n),n-1)
                =\lim_{\Delta\to\infty}p(y(n),n-1)
                =-\frac{\pi^2}2\sigma^2.
            \end{equation*}
        \item Under the additional condition for $(L_\Delta)$ that, when $\Delta\to\infty$,
            \begin{equation*}
                N_\Delta:=\#\left\{\hat\Delta\ge 0: L_\Delta=L_{\hat\Delta}\right\}
                =O_\Delta(1),
            \end{equation*}
            we have
            \begin{equation*}
                \lim_{\Delta\to\infty}p(\Delta,L_\Delta)=-\frac{\pi^2}2\sigma^2.
            \end{equation*}
        \item Under the same condition as (iii), for $v>0$,
            \begin{align*}
                &\quad
                \lim_{\Delta\to\infty}
                \frac{\Delta^2}{L_\Delta}
                \log\mathbb P\left(
                    -a\Delta\le Y_j\le (1-a)\Delta,\,
                    \forall j\in[\![1,L_\Delta]\!],\,
                    Y_{L_\Delta}\le(v-a)\Delta
                \right)
                \nonumber \\ &
                =
                -\frac{\pi^2}{2}\sigma^2.
            \end{align*}
    \end{enumerate}
\end{lemma}
\begin{proof}[Proof of Lemma~\ref{lemma:mogulskii-cal}]
    For (i) and (ii), we prove here only the result for $x$. The result for $y$ is the same since the conditions for Mogul'skii's Theorem (Theorem \ref{theorem:mogulskii}) are satisfied by both $x$ and $y$ as stated in Lemma \ref{lemma:mogulskii-x-y}.
    \begin{enumerate}[label=(\roman*)]
        \item  We denote here the step random variables of $(Y_i)$ as $(\sigma\xi_i)$. Define $s_n$ as in \eqref{eq:mogulskii-s-n-t}. We can express $q(x(n),n-1)$ in $s_n$ as
        \begin{align*}
            q(x(n),n-1)
            &=
            \log\mathbb P
            \left(
                \frac 1\sigma(-a)\le s_n(t)\le\frac 1\sigma(1-a),\,
                t\in[0,1]
            \right) \nonumber
            \\ &=:\log\mathbb P(s_n\in G_0),
            \quad n\ge 1.
        \end{align*}
        For each $m\ge 1$, we apply Mogul'skii's Theorem for $G=G(m)$ of the form \eqref{eq:mogulskii-G} with
        \begin{equation*}
            S_1=\frac{1}\sigma\left(1-a-\frac 1{2m}\right),\,
            S_2=\frac{1}\sigma\left(-a+\frac 1{2m}\right),
        \end{equation*}
        and we can use \eqref{eq:mogulskii-h-alpha} to calculate the corresponding $H_2(G)$ as
        \begin{equation*}
            H_2(G)=\sigma^2\left(1-\frac 1m\right)^{-2}.
        \end{equation*}
        Hence, by Mogul'skii's Theorem, we have
        \begin{align*}
            \lim_{n\to\infty}
            \frac{x^2(n)}{n}
            \log\mathbb P(s_n\in G(m))
            =
            -\frac{\pi^2}2\sigma^2\left(1-\frac 1m\right)^{-2}.
        \end{align*}
        Since $G(m)\subset G_0$, by monotonicity of probability, 
        \begin{align*}
            \liminf_{n\to\infty}p(x(n),n-1)
            &\ge\liminf_{n\to\infty}\frac{x^2(n)}n\log\mathbb P(s_n\in G(m))
            \nonumber\\ &
            =-\frac{\pi^2}2\sigma^2\left(1-\frac 1m\right)^{-2}.
        \end{align*}
        Let $m\to\infty$ and we get
        \begin{equation*}
            \liminf_{n\to\infty}p(x(n),n-1)\ge-\frac{\pi^2}{2}\sigma^2.
        \end{equation*}
        Similarly, we apply Mogul'skii's Theorem for $G=\widetilde G(m)$ of the form \eqref{eq:mogulskii-G} with
        \begin{equation*}
            S_1=\frac{1}\sigma\left(1-a+\frac 1{2m}\right),\,
            S_2=\frac{1}\sigma\left(-a-\frac 1{2m}\right),
        \end{equation*}
        and we can get
        \begin{equation*}
            \limsup_{n\to\infty}p(x(n),n-1)\le-\frac{\pi^2}2\sigma^2\left(1-\frac 1m\right)^{-2}.
        \end{equation*}
        Let $m\to\infty$ and the result would follow.
        \item When $\Delta\to\infty$, $n=n(\Delta)$ goes to infity because of the original condition on $(L_\Delta)$, and the result follows from (i) for subsequences.
        \item By Lemma \ref{lemma:mogulskii-mono} and since $x(n)\le\Delta\le x(n+1)-1=y(n)$, we have the lower and upper bound for $p(\Delta,L_\Delta)$ as
        \begin{align*}
            p(\Delta,L_\Delta)
            &=\frac{\Delta^2}{L_\Delta}q(\Delta,L_\Delta)
            \ge\frac{\Delta^2}{L_\Delta}q(x(n),L_\Delta)
            =\frac{\Delta^2}{x^2(n)}p(x(n),n-1)
            \\
            p(\Delta,L_\Delta)
            &=\frac{\Delta^2}{L_\Delta}q(\Delta,L_\Delta)
            \le\frac{\Delta^2}{L_\Delta}q(y(n),L_\Delta)
            =\frac{\Delta^2}{y^2(n)}p(y(n),n-1)
        \end{align*}
        Let $\Delta\to\infty$ in both equation*s and it suffices to show that
        \begin{equation*}
            \lim_{\Delta\to\infty}\frac{\Delta}{x(n)}
            =\lim_{\Delta\to\infty}\frac{\Delta}{y(n)}
            =1.
        \end{equation*}
        Indeed, we have $x(n)\le\Delta\le x(n)+N_\Delta$ and $y(n)-N_\Delta\le\Delta\le y(n)$, and the result follows from the additional condition.
        \item The proof is analagous to (i) to (iii). We only need to calculate the value of $H_2(G)$ for $G$ of the form \eqref{eq:mogulskii-G} with
        \begin{equation*}
            S_1(t)=\begin{cases}
                \frac 1\sigma\left(1-a\mp\frac{v\wedge 1}{2m}\right),
                &t\in[0,\frac{n-1}n),\\
                \frac 1\sigma\left(v-a\mp\frac{v\wedge 1}{2m}\right),
                &t\in[\frac{n-1}n,1],
            \end{cases}
            \quad
            S_2=\frac 1\sigma\left(-a\pm\frac{v\wedge 1}{2m}\right).
        \end{equation*}
        We use \eqref{eq:mogulskii-h-alpha} to calculate
        \begin{align*}
            H_2(G)
            &=
            \sigma^2\left(1-\frac{v\wedge 1}m\right)^{-2}\frac{n-1}n
            +
            \sigma^2\left(v-\frac{v\wedge 1}m\right)^{-2}\frac 1n
            \\ &\sim
            \sigma^2\left(1-\frac{v\wedge 1}m\right)^{-2},
        \end{align*}
        when $n\to\infty$.
    \end{enumerate}
\end{proof}

\begin{example}
    \label{example:mogulskii-L-1}
    For $v>0$, $u>0$ and $L_\Delta=\lfloor u\Delta^{2+v}\rfloor$, we have all the conditions satisfied, i.e.,
    \begin{equation*}
        \lim_{\Delta\to\infty}\frac{L_\Delta}{\Delta^2}=\infty,\,
        N_\Delta = O_\Delta(1).
    \end{equation*}
    Indeed, $L_\Delta>u\Delta^{2+v}-1$, and $N_{\Delta}=1$ for $\Delta$ large enough.
\end{example}

\begin{example}
    \label{example:mogulskii-L-2}
    For $v>0$, and $L_\Delta=\lfloor\Delta^{2+v}\rfloor-\lfloor\Delta^2\rfloor$, we have all the conditions satisfied. Indeed, $L_\Delta>\Delta^{2+v}-\Delta^2-2$, and $N_\Delta=1$ for $\Delta$ large enough.
\end{example}
% !TEX root = ../main.tex

\section{\bilingual{Lemmas on Spine Random Walk}{有关脊柱随机游走的引理}}\label{section:spine-random-walk-lemmas}

\bilingual{
    The results from this section onward are based on the works of A\"id\'ekon and Hu \mycite{aïdékon2025averageweightpercolationcompletegraph}. Consider the random walk $(Y_n)$ obtained from \nameref{lemma:many-to-one}.
}{
    从本节开始的结论基于并推广了A\"id\'ekon与Hu \mycite{aïdékon2025averageweightpercolationcompletegraph} 的工作.考虑由\nameref{lemma:many-to-one}得到的随机游走$(Y_n)$.
}

\begin{lemma}[\bilingual{Small Deviation}{水平管道意义下的小偏差}, A\"id\'ekon--Hu \mycite{aïdékon2025averageweightpercolationcompletegraph}]\label{lemma:4.1}
    \bilingual{
        Let $v>0$ and $0<d<\frac{\pi^2}2\sigma^2$. There exists $\Delta_0>0$ such that for all $\Delta\ge\Delta_0$, $r_0\le 0\le r_0+\Delta$ and $L\ge 1$,
    }{
        设 $v>0$ 且 $0<d<\frac{\pi^2}2\sigma^2$. ~存在 $\Delta_0>0$ 使得对所有 $\Delta\ge\Delta_0$,~ $r_0\le 0\le r_0+\Delta$ 且 $L\ge 1$,
    }
    {
        \setlength{\belowdisplayskip}{0pt}
        \begin{equation}
            \mathbb P(r_0\le Y_j\le r_0+\Delta,\,\forall j\in[\![1,L]\!])\le\exp\left( \Delta^v-\frac{d}{\Delta^2}L\right).
            \label{eq:4.1}
        \end{equation}
    }
\end{lemma}

% \vspace{-15pt}
\begin{figure}[H]
    \centering
    \begin{tikzpicture}[
        scale=1.5,
        >=Stealth,
        % 随机游走路径样式
        rw-path/.style={thick, blue!70!black, mark=*, mark size=1.5pt},
        % 管道填充样式
        tube-fill/.style={gray!10},
        % 标注样式
        label-font/.style={font=\small}
    ]

        % 1. 绘制背景管道填充 (r_0 到 r_0 + delta)
        \fill[tube-fill] (0, -1.5) rectangle (4, 0.5);

        % 2. 绘制辅助虚线
        \draw[dashed, black!60] (0, -2) -- (4.5, -2) node[right] {$-\Delta$};
        % \draw[dashed, black!60] (0, 0) -- (4.5, 0) node[right] {$y = 0$};
        \draw[dashed, black!60] (4, -2.5) -- (4, 1) node[above] at (4, 1) {$L$};

        % 3. 绘制管道实线 (r_0 到 r_0 + delta)
        \draw[thick, black] (0, -1.5) -- (4.5, -1.5) node[right] {$r_0$};
        \draw[thick, black] (0, 0.5) -- (4.5, 0.5) node[right] {$r_0 + \Delta$};

        % 4. 绘制坐标轴 (Y轴实线)
        \draw[->, thick] (0, -2.5) -- (0, 1.2) node[above] {$Y_n$};
        \draw[->, thick, dashed] (0, 0) -- (5, 0) node[right] {$n$};

        % 5. 绘制随机游走 (起点在 0,0，始终保持在 [-1.5, 0.5] 之间)
        % 选取了一些看起来随机的点
        \draw[rw-path] plot coordinates {
            (0, 0) 
            (0.8, -0.6) 
            (1.5, 0.2) 
            (2.2, -1.1) 
            (3.1, -0.4) 
            (4, -0.8)
        };

        % 6. 标注 Delta (使用大括号)
        \draw[decorate, decoration={brace, amplitude=5pt}] 
            (-0.1, -1.5) -- (-0.1, 0.5) node[midway, left=6pt] {$\Delta$};

        % 标注特定的点
        % \node[below left] at (0,0) {$0$};

    \end{tikzpicture}
    \caption{\bilingual{Illustration of the small deviation of the random walk \texorpdfstring{$Y_n$}{Yn} within the horizontal tube \texorpdfstring{$[r_0, r_0+\Delta]$}{[r_0, r_0+Delta]}}{随机游走 \texorpdfstring{$Y_n$}{Yn} 在水平管道 \texorpdfstring{$[r_0, r_0+\Delta]$}{[r_0, r_0+Delta]} 内的小偏差示意图}}
\end{figure}

\begin{proof}\label{proof:lemma:4.1}
    \bilingual{
        By applying \nameref{theorem:mogulskii} (see \Cref{lemma:mogulskii-cal} (ii))% and \Cref{example:mogulskii-L-1})
        , for all $a\in(0,1)$ and $\ell=\ell(\Delta)=\lfloor\frac 1d\Delta^{2+v}\rfloor$,
    }{
        通过应用 \nameref{theorem:mogulskii} (参见\Cref{lemma:mogulskii-cal} (ii))% 和\Cref{example:mogulskii-L-1})
        ,对所有 $a\in(0,1)$ 和 $\ell=\ell(\Delta)=\lfloor\frac 1d\Delta^{2+v}\rfloor$,~
    }
    \begin{equation*}
        \lim_{\Delta\to\infty}
        \frac{\Delta^2}{\ell}
        \log\mathbb P\left(
                -a\Delta\le Y_j\le(1-a)\Delta,\,
                \forall j\in[\![1,\ell]\!]
        \right)
        =
        -\frac{\pi^2}2\sigma^2.
        \label{eq:lemma:4.1:mogulskii-applied}
    \end{equation*}

    \bilingual{
        Fix $m\ge 1$ so that $d{(1+\frac 1m)}^2<\frac{\pi^2}2\sigma^2$. Apply the above equation to $(1+\frac 1m)\Delta$ in lieu of $\Delta$ and
    }{
        固定 $m\ge 1$ 使得 $d{(1+\frac 1m)}^2<\frac{\pi^2}2\sigma^2$. ~应用上式,应用时将上式中的 $\Delta$ 替换为 $(1+\frac 1m)\Delta$,~ 并且令
    }
    \begin{equation*}
        a=\frac{1-\frac km}{1+\frac 1m},\quad
        k\in[\![0,m-1]\!],
    \end{equation*}
    \bilingual{
        and we can get, for all $k\in[\![0,m-1]\!]$,
    }{
        从而得,对所有 $k\in[\![0,m-1]\!]$,~
    }
    \begin{align*}
        &\quad\lim_{\Delta\to\infty}
        \frac 
        {
            {\left(1+\frac 1m\right)}^2
            \Delta^2
        } \ell
        \log\mathbb P\left(
                \left(-1+\frac km\right)\Delta
                \le Y_j
                \le\frac{k+1}{m}\Delta,\,
                \forall j\in[\![1,\ell]\!]
        \right)
        \\ &=
        -\frac{\pi^2}{2}\sigma^2
        <-d{\left(1+\frac 1m\right)}^2.
    \end{align*}
    \bilingual{
        Then we find $\Delta_0=\Delta(d)>0$ such that for all $\Delta\ge\Delta_0$ and $k\in[\![0,m-1]\!]$,
    }{
        于是可得 $\Delta_0=\Delta(d)>0$ 使得对所有 $\Delta\ge\Delta_0$ 和 $k\in[\![0,m-1]\!]$,~ 
    }
    \begin{equation*}
        \mathbb P\left(
            \left(-1+\frac km\right)\Delta
            \le Y_j
            \le\frac{k+1}{m}\Delta,\,
            \forall j\in[\![1,\ell]\!]
        \right)
        \le
        \exp\left(
            -\frac{d}{\Delta^2}\ell
        \right).
    \end{equation*}
    \bilingual{
        Now for $r_0\le 0\le r_0+\Delta$, we can find some $k\in[\![0,m-1]\!]$ such that $\frac{r_0}\Delta+1\in[\frac km,\frac{k+1}m]$, and hence, for all $\Delta\ge\Delta_0$,
    }{
        现在对 $r_0\le 0\le r_0+\Delta$,~ 存在 $k\in[\![0,m-1]\!]$ 使得 $\frac{r_0}\Delta+1\in[\frac km,\frac{k+1}m]$,~ 因此对所有 $\Delta\ge\Delta_0$,~ 
    }
    \begin{align*}
        &\quad\mathbb P\left(
                r_0\le Y_j\le r_0+\Delta,\,
                \forall j\in[\![1,\ell]\!]
        \right)
        \\ &\le\mathbb P\left(
            \left(-1+\frac km\right)\Delta
            \le Y_j
            \le\frac{k+1}{m}\Delta,\,
            \forall j\in[\![1,\ell]\!]
        \right)
        \le
        \exp\left(
            -\frac{d}{\Delta^2}\ell
        \right).
    \end{align*}
    
    \bilingual{
        Now for any $L\ge 1$, split the interval $[\![1,L]\!]$ into $q$ possible intervals of length $\ell$ and one possible interval of length $r$, i.e., $L=q\ell+r$ with $q\ge 0$ and $0\le r<\ell$. Denote in this proof
    }{
        现在对任何 $L\ge 1$,~ 将区间 $[\![1,L]\!]$ 分成 $q$ 个长度为 $\ell$ 的区间和一个长度为 $r$ 的区间,即 $L=q\ell+r$,~ 其中 $q\ge 0$ 且 $0\le r<\ell$. ~在本证明中记
    }
    \begin{equation*}
        A_k
        :=
        \{
            r_0\le Y_j\le r_0+\Delta,\,
            \forall j\in[\![(k-1)\ell+1,k\ell]\!]
        \},
        \quad
        k\in[\![1,q]\!].
    \end{equation*}
    \bilingual{
        We have
    }{
        从而
    }
    \begin{align*}
        &\quad
        \mathbb P(
            r_0\le Y_j\le r_0+\Delta,\,
            \forall j\in[\![1,L]\!]
        )
        \\ &
        \le
        \mathbb P(
            \cap_{k=1}^qA_k
        )
        =
        \prod_{k=1}^q\mathbb P\left(
            A_k
            \middle|
            \cap_{l=1}^{k-1}A_l
        \right)
        \\ &
        =
        \prod_{k=1}^q\left(
            \frac 1{
                \mathbb P\left(
                    \cap_{l=1}^{k-1}A_l
                \right)
            }
            \mathbb E\left[
                \mathbb E\left(
                    {\bf 1}_{A_k}
                    \middle|
                    \mathscr F_{(k-1)\ell}
                \right)
                {\bf 1}_{
                    \cap_{l=1}^{k-1}A_l
                }
            \right]
        \right)
        \\ &
        \le
        \prod_{k=1}^q
        \sup_{x\in r_0+[0,\Delta]}
        \mathbb P(
            r_0\le Y_j+x\le r_0+\Delta,\,
            \forall j\in[\![1,\ell]\!]
        )
        \\ &
        \le
        \left[
            \exp\left(
                -\frac d{\Delta^2}\ell
            \right)
        \right]^q
        =
        \exp\left(\frac{d}{\Delta^2}r-\frac{d}{\Delta^2}L\right)
        % \\ &
        \le
        \exp\left(\Delta^v-\frac{d}{\Delta^2}L\right),
    \end{align*}
    \bilingual{
        where $r_0-x\in[-\Delta,0]$ for $x\in r_0+[0,\Delta]$, and the last inequality follows from $r<\ell\le\frac{1}{d}\Delta^{2+v}$.
    }{
        其中 $r_0-x\in[-\Delta,0]$ 对 $x\in r_0+[0,\Delta]$ 成立,且最后一式是因 $r<\ell\le\frac{1}{d}\Delta^{2+v}$.
    }
\end{proof}

\begin{remark}
    \bilingual{
        For $u>0$, if we take $\ell=\ell(\Delta)=\lfloor\frac ud\Delta^{2+v}\rfloor$ in lieu of $\ell$ in the previous proof, we can get an upper bound of
    }{
        对 $u>0$,~ 如果在前面的证明中将 $\ell$ 替换为 $\lfloor\frac ud\Delta^{2+v}\rfloor$,~ 可得上界
    }
    $
        e^{
            u\Delta^v-\frac{d}{\Delta^2}L
        }.
    $
\end{remark}

\bilingual{
    Based on the previous lemma, we have the following lemma.
}{
    基于前一个引理,我们有下面的引理.
}

\begin{lemma}[\bilingual{Small Deviation under No Large Drop}{无大回撤意义下的小偏差}, A\"id\'ekon--Hu \mycite{aïdékon2025averageweightpercolationcompletegraph}]
    \label{lemma:2.3}
    \bilingual{
        Let $c>0$. For any $d\in(0,\frac{\pi^2}2\sigma^2)$, there exists $\Delta_0=\Delta_0(c,d)\ge 1$ such that for all $\Delta\ge\Delta_0$ and $L\ge\Delta^{2+c}$, we have
    }{
        设 $c>0$. ~对任何 $d\in(0,\frac{\pi^2}2\sigma^2)$,~ 存在 $\Delta_0=\Delta_0(c,d)\ge 1$ 使得对所有 $\Delta\ge\Delta_0$ 和 $L\ge\Delta^{2+c}$,~ 成立
    }
    {
        \setlength\belowdisplayskip{-20pt}
        \begin{equation*}
            \mathbb P\left(
                Y_L
                \le
                \frac\Delta c,\,
                \min_{0\le i\le j\le L}
                (Y_j-Y_i)
                \ge
                -\Delta
            \right)
            \le
            e^{-\frac{dL}{\Delta^2}}.
        \end{equation*}
    }
\end{lemma}

\begin{figure}[htbp]
    \centering
    \begin{tikzpicture}[
        xscale=1.7, % 修改这里控制横向（代数 n）的间距
        yscale=1.1, % 修改这里控制纵向（空间位置 V）的伸展程度
        >=Stealth,
        % 密集随机游走样式
        rw-path/.style={thick, blue!70!black},
        % 动态屏障样式
        barrier-line/.style={thick, red!70!black, dashed},
        % 参考线样式
        ref-line/.style={thin, gray!60, dashed}
    ]

        % 0. 绘制填充阴影 (放在最底层)
        \fill[blue!5] (0,0) 
            -- (0.2, 0.6) -- (0.4, -0.3) -- (0.6, 0.2) -- (0.8, 1.1) 
            -- (1.1, 0.3) -- (1.3, 0.7) -- (1.6, -0.5) -- (1.9, 0.2) 
            -- (2.2, 1.5) -- (2.5, 0.4) -- (2.8, 0.9) -- (3.1, -0.45) 
            -- (3.4, 0.1) -- (3.7, -0.4) -- (4.0, 0.6) -- (4.3, -0.2) -- (4.5, 0.3) % 蓝色路径结束
            -- (4.5, -0.5) -- (2.2, -0.5) -- (2.1, -0.9) -- (0.8, -0.9) -- (0.7, -1.4) -- (0.2, -1.4) -- (0, -2.0) % 红色路径逆序
            -- cycle;

        % 1. 坐标轴
        \draw[->, thick] (0, -2.5) -- (0, 2.2) node[above] {$Y_n$};
        \draw[->, thick] (0, 0) -- (5.5, 0) node[right] {$n$};

        % 2. 终点限制线 y = Delta/c (假设 Delta/c = 0.8)
        \draw[thick, black!80] (5, 0.8) -- (5.4, 0.8) node[right] {$\frac{\Delta}{c}$};
        \draw[ref-line] (0, 0.8) -- (5, 0.8);

        % 3. 时间 L 的垂线
        \draw[dashed, black!80] (4.5, -2.5) -- (4.5, 2.0) node[above] {$L \ge \Delta^{2+c}$};

        % 4. 密集随机游走路径 (零均值特征)
        % 定义一组波动点
        \draw[rw-path] (0,0) 
            -- (0.2, 0.6) -- (0.4, -0.3) -- (0.6, 0.2) -- (0.8, 1.1) 
            -- (1.1, 0.3) -- (1.3, 0.7) -- (1.6, -0.5) -- (1.9, 0.2) 
            -- (2.2, 1.5) -- (2.5, 0.4) -- (2.8, 0.9) -- (3.1, -0.45) 
            -- (3.4, 0.1) -- (3.7, -0.4) -- (4.0, 0.6) -- (4.3, -0.2) -- (4.5, 0.3);

        % 5. 动态屏障 Max Y_i - Delta (假设 Delta = 2.0)
        % 初始 Max=0 -> 屏障 -2.0
        % 第一高峰 0.6 -> 屏障 -1.4
        % 第二高峰 1.1 -> 屏障 -0.9
        % 第三高峰 1.5 (最高) -> 屏障 -0.5
        \draw[barrier-line] 
            (0, -2.0) -- (0.2, -1.4) -- (0.7, -1.4) -- (0.8, -0.9) 
            -- (2.1, -0.9) -- (2.2, -0.5) -- (4.5, -0.5) 
            node[right] {$\max_{i\le n} Y_i - \Delta$};

        % 6. 标注 Delta 间距 (在最高点处)
        \draw[ref-line] (0, 1.5) -- (2.2, 1.5);
        \draw[<->] (2.3, 1.5) -- (2.3, -0.5) node[midway, left] {$\Delta$};

        % 7. 关键点标注
        \node[below left] at (0,0) {$0$};
        \fill[blue!70!black] (4.5, 0.3) circle (1.5pt) node[right, black] {$Y_L$};

    \end{tikzpicture}
    \caption{\bilingual{Illustration of the historical drawdown of the random walk \texorpdfstring{$Y_n$}{Yn} always constrained by \texorpdfstring{$\Delta$}{Delta} and the endpoint below \texorpdfstring{$\frac{\Delta}{c}$}{Delta/c}}{随机游走 \texorpdfstring{$Y_n$}{Yn} 的历史回撤始终受限于 \texorpdfstring{$\Delta$}{Delta} 且终点位于 \texorpdfstring{$\frac{\Delta}{c}$}{Delta/c} 之下}}
\end{figure}

\begin{proof}
    \label{proof:lemma:2.3}\quad

    \bilingual{
        We denote
    }{
        记事件
    }
    \begin{equation*}
        B = B(c,\Delta,L) := \left\{ Y_L \le \frac\Delta c, \, \min_{0\le i\le j\le L} (Y_j-Y_i) \ge -\Delta \right\}.
    \end{equation*}
    \bilingual{
        Note that $B(c_1,\Delta,L) \subset B(c_2,\Delta,L)$ for all $c_1 \ge c_2$ and fixed $\Delta, L$.
    }{
        注意到对于固定的 $\Delta$ 和 $L$,~ 若 $c_1 \ge c_2$,~ 则有 $B(c_1,\Delta,L) \subset B(c_2,\Delta,L)$.
    }

    \bilingual{
        For the case $c \ge \frac{1}{2}$, it suffices to prove the result for $c' \in (0, \frac{1}{2})$ and $L \ge \Delta^{\frac{2}{1-2c'}}$. Indeed, for any $c \ge \frac{1}{2}$, there exists $c'$ such that
    }{
        对于 $c \ge \frac{1}{2}$ 的情形,只需证明原结论对 $c' \in (0, \frac{1}{2})$ 且 $L \ge \Delta^{\frac{2}{1-2c'}}$ 成立.事实上,对任何 $c \ge \frac{1}{2}$,~ 总能选取 $c'$ 满足
    }
    \begin{equation*}
        0 < \frac{1}{2} - \frac{1}{2+c} \le c' < \frac{1}{2} \le c.
    \end{equation*}
    \bilingual{
        Applying the result for $c'$ and $L \ge \Delta^{\frac{2}{1-2c'}}$, there exists $\Delta_1 = \Delta_1(c',d) = \Delta_1(c,d) \ge 1$ such that for all $\Delta \ge \Delta_1$ and $L \ge \Delta^{\frac{2}{1-2c'}}$
    }{
        对此 $c'$ 和 $L \ge \Delta^{\frac{2}{1-2c'}}$ 应用结论,存在 $\Delta_1 = \Delta_1(c',d) = \Delta_1(c,d) \ge 1$ 使得对所有 $\Delta \ge \Delta_1$ 且 $L \ge \Delta^{\frac{2}{1-2c'}}$ 成立
    }
    \begin{equation*}
        \mathbb{P}(B(c,\Delta,L)) \le \mathbb{P}(B(c',\Delta,L)) \le e^{-\frac{dL}{\Delta^2}}.
    \end{equation*}
    \bilingual{
        Since $\Delta^{\frac{2}{1-2c'}} \le \Delta^{2+c}$, the above inequality holds for all $L \ge \Delta^{2+c}$. In this regime, for $0 < d < \frac{\pi^2}{2}\sigma^2$, we choose $\eta \in (0,1)$ sufficiently small such that
    }{
        由于 $\Delta^{\frac{2}{1-2c'}} \le \Delta^{2+c}$,~ 上式对所有 $L \ge \Delta^{2+c}$ 均成立.此时,对于 $0 < d < \frac{\pi^2}{2}\sigma^2$,~ 选取足够小的 $\eta \in (0,1)$ 使得
    }
    \begin{equation*}
        (2+\eta)\left(\frac{1}{2}-c\right) < 1, \quad d(1+2\eta)^4 < \frac{\pi^2}{2}\sigma^2.
    \end{equation*}

    \bilingual{
        For $0 < c < \frac{1}{2}$, we directly prove the result for $L \ge \Delta^{2+c}$ using the same method, only requiring $\eta \in (0,1)$ to be small enough to satisfy $\eta < c$ and the second inequality above.
    }{
        对于 $0 < c < \frac{1}{2}$ 的情形,直接用相同方法证明 $L \ge \Delta^{2+c}$ 时的结论,区别仅在于需令 $\eta \in (0,1)$ 足够小以满足 $\eta < c$ 且成立上述第二个不等式.
    }

    \bilingual{
        Note that in both cases, we have
    }{
        注意到,无论上述哪种情况,均满足
    }
    \begin{equation*}
        \frac{1}{c} \ge 2, \quad L > \lfloor\Delta^{2+\eta}\rfloor, \quad \lim_{\Delta\to\infty} \frac{L}{\lfloor\Delta^{2+\eta}\rfloor} = \infty.
    \end{equation*}
    \bilingual{
        We divide the interval $[\![0,L]\!]$ into smaller intervals of length $\lfloor\Delta^{2+\eta}\rfloor$, and the space into intervals of length $\lceil\eta\Delta\rceil$. Let $q := \lfloor\frac{L}{\lfloor\Delta^{2+\eta}\rfloor}\rfloor$, $\ell_0 := 0$, and
    }{
        将时间区间 $[\![0,L]\!]$ 分成长度为 $\lfloor\Delta^{2+\eta}\rfloor$ 的小区间,并将空间划分为长度为 $\lceil\eta\Delta\rceil$ 的层级.令 $q := \lfloor\frac{L}{\lfloor\Delta^{2+\eta}\rfloor}\rfloor$,~ $\ell_0 := 0$,~ 且定义
    }
    \begin{align*}
        \ell_j &:= L - q\lfloor\Delta^{2+\eta}\rfloor + (j-1)\lfloor\Delta^{2+\eta}\rfloor, \quad j \in [\![1, q+1]\!], \\
        I_j &:= [\![\ell_j, \ell_{j+1}]\!], \quad j \in [\![0, q]\!].
    \end{align*}
    \bilingual{
        On the event $B$, we have
    }{
        在事件 $B$ 上,成立
    }
    \begin{equation*}
        \min_{0\le j\le L} Y_j \ge -\Delta, \quad \max_{0\le i\le L} Y_i \le Y_L + \Delta \le \left(\frac{1}{c}+1\right)\Delta.
    \end{equation*}
    \bilingual{
        Let $C := \frac{1}{c}+2 \ge 4$ so that the total range is within $C\Delta$. For any $j \in [\![0,q]\!]$, there exists some random $k_j \in \mathbb{N}$ such that
    }{
        令 $C := \frac{1}{c}+2 \ge 4$,~ 此时路径取值的跨度不超过 $C\Delta$. ~对于任何 $j \in [\![0,q]\!]$,~ 存在随机变量 $k_j \in \mathbb{N}$ 满足
    }
    \begin{equation*}
        \left(\frac{1}{c}+1\right)\Delta - (k_j+1)\lceil\eta\Delta\rceil < \max_{0 \le \ell \le \ell_{j+1}} Y_\ell \le \left(\frac{1}{c}+1\right)\Delta - k_j\lceil\eta\Delta\rceil := M_j.
    \end{equation*}
    \bilingual{
        Note that $M_j$ can take at most $\lceil \frac C\eta \rceil$ possible values. Let $\mathcal{M}$ be the deterministic set of these possible values. Define the index set
    }{
        注意 $M_j$ 的可能取值个数不超过 $\lceil \frac C\eta \rceil$ 个.令 $\mathcal{M}$ 为 $M_j$ 可能取值的确定性集合.定义下标集
    }
    \begin{equation*}
        J := \left\{ j \in [\![0,q-1]\!] : M_j \neq M_{j+1} \right\}.
    \end{equation*}
    \bilingual{
        Since $M_j$ is non-decreasing in $j$, we have $\#J \le \lceil \frac C\eta \rceil$.
    }{
        由于 $M_j$ 随 $j$ 非递减,故有 $\#J \le \lceil \frac C\eta \rceil$.
    }

    \bilingual{
        For $j \notin J$, we have $M_j = M_{j+1}$, which implies that on $B$, $Y_i \in [M_j - \lceil\eta\Delta\rceil - \Delta, M_j]$ for all $i \in I_{j+1}$. (See \Cref{figure:tube-constraint}.)
    }{
        若 $j \notin J$,~ 则 $M_j = M_{j+1}$,~ 由此推得,在事件 $B$ 上,对所有 $i \in I_{j+1}$,~ 都有 $Y_i \in [M_j - \lceil\eta\Delta\rceil - \Delta, M_j]$. (见\Cref{figure:tube-constraint}.)
    }

    \begin{figure}[H]
        \centering
        \begin{tikzpicture}[
            xscale=1.2, 
            yscale=1.2, 
            >=Stealth,
            rw-path/.style={thick, blue!70!black},
            running-max/.style={thin, gray!80, dashed},
            barrier-red/.style={thick, red!70!black, dashed},
            barrier-green/.style={thick, green!60!black, dashed},
            ref-line/.style={thin, black!80, dashed}
        ]

            % 1. 参数设定
            % Delta = 2.2
            % Mj = 1.8, Mj - etaDelta = 1.2
            % 绿线 = 1.2 - 2.2 = -1.0
            
            % 2. 绘制阴影区域 (Y_n 与 红色动态屏障之间)
            \fill[blue!5] 
                (0,0) -- (0.5, 1.1) -- (1.5, -0.8) -- (2.5, 1.6) -- (3.5, 0.4) -- (4, 1.5) % I_j 结束
                -- (5, -0.5) -- (6, 1.7) -- (7, -0.2) -- (8, 0.5) % I_{j+1} 结束
                -- (8, -0.5) -- (6, -0.5) -- (5.95, -0.6)-- (2.5, -0.6) -- (2.3, -1.1) -- (0.5, -1.1) -- (0, -2.2) % 红色路径逆序
                -- cycle;

            % 3. 坐标轴
            % \draw[->, thick] (-0.2, -2.5) -- (8.5, -2.5) node[right] {$n$};
            \draw[->, thick] (-0.2, -2.5) -- (-0.2, 2.5) node[above] {$Y_n$};

            % 4. 离散层级标注 (仅标注这两根线)
            \draw[barrier-green] (-0.3, 1.8) -- (8.2, 1.8) node[right] {$M_j$};
            \draw[black!80, dashed] (-0.3, 1.2) -- (8.2, 1.2) node[right] {$M_j - \lceil\eta\Delta\rceil$};

            % 5. 时间区间划分
            \draw[ref-line] (0, -2.6) -- (0, 2.2) node[below=5.6cm] {$\ell_j$};
            \draw[ref-line] (4, -2.6) -- (4, 2.2) node[below=5.6cm] {$\ell_{j+1}$};
            \draw[ref-line] (8, -2.6) -- (8, 2.2) node[below=5.6cm] {$\ell_{j+2}$};
            \node at (2, -2.8) {$I_j$};
            \node at (6, -2.8) {$I_{j+1}$};

            % 6. 静态绿线 (区间起始限制)
            \draw[barrier-green] (4, -1.0) -- (8.2, -1.0) node[right] {$(M_j - \lceil\eta\Delta\rceil) - \Delta$};

            % 7. 零均值随机游走路径
            \draw[rw-path] (0,0) 
                -- (0.5, 1.1) -- (1.5, -0.8) -- (2.5, 1.6) -- (3.5, 0.4) -- (4, 1.5) 
                -- (5, -0.5) -- (6, 1.7) -- (7, -0.2) -- (8, 0.5);

            % 8. 区间最大值(Running Max 阶梯线)
            \draw[running-max] (0,0) -- (0.5, 1.1) -- (2.3, 1.1) -- (2.5, 1.6) -- (5.95, 1.6) -- (6, 1.7) -- (8, 1.7);

            % 9. 动态红线 (max Y_i - Delta)
            % 随创新高而阶梯式上升
            \draw[barrier-red] (0, -2.2) -- (0.5, -1.1) -- (2.3, -1.1) -- (2.5, -0.6) -- (5.95, -0.6) -- (6, -0.5) -- (8, -0.5)
                node[right=0.25cm] {$\max_{i\le n} Y_i - \Delta$};

            % 10. 标注 Delta 间距 (在 I_{j+1} 内部展示)
            \draw[<->] (6.2, 1.7) -- (6.2, -0.5);
            \node[right] at (6.2, 0.6) {$\Delta$};

            % 11. 标注起点 0
            % \node[left] at (0,0) {$0$};

        \end{tikzpicture}
        \caption{\bilingual{Illustration of the tube constraints for the random walk when \texorpdfstring{$j\notin J$}{j not in J} within intervals \texorpdfstring{$I_j$}{Ij} and \texorpdfstring{$I_{j+1}$}{Ij+1}}{\texorpdfstring{$j\notin J$}{j not in J}时区间 \texorpdfstring{$I_j$}{Ij} 与 \texorpdfstring{$I_{j+1}$}{Ij+1} 内随机游走的管道约束}}
        \label{figure:tube-constraint}
    \end{figure}

    \bilingual{
        Now, let $\mathscr{J}$ be the family of all possible indices where the discretized maximum can jump
    }{
        令 $\mathscr{J}$ 为各区间最大值所有可能跳跃点的下标集族
    }
    \begin{equation*}
        \mathscr{J} := \left\{ J_0 \subseteq [\![0, q-1]\!] : \#J_0 \le \left\lceil \frac{C}{\eta} \right\rceil \right\}.
    \end{equation*}
    \bilingual{
        For a fixed $J_0 \in \mathscr{J}$, let its complement be $J_0^{\mathrm{c}} = \{ j_1 > j_2 > \cdots > j_{\#J_0^{\mathrm{c}}} \} := [\![0, q-1]\!] \setminus J_0$. Note that $\#J_0^{\mathrm{c}} \ge q - \lceil \frac C\eta \rceil$.
    }{
        对于固定的 $J_0 \in \mathscr{J}$,~ 令其补集为 $J_0^{\mathrm{c}} = \{ j_1 > j_2 > \cdots > j_{\#J_0^{\mathrm{c}}} \} := [\![0, q-1]\!] \setminus J_0$. ~注意 $\#J_0^{\mathrm{c}} \ge q - \lceil \frac C\eta \rceil$.
    }
    
    \bilingual{
        We estimate $\mathbb{P}(B, J=J_0)$ by applying the Markov property at times $\ell_{j_k+1}$ in decreasing order of $k$. Since $B \cap \{J=J_0\}$ implies that for all $j \in J_0^{\mathrm{c}}$, the maximum $M_j$ remains constant over the interval $I_{j+1}$, the path is constrained within a tube of width $\lceil\eta\Delta\rceil + \Delta$. We obtain
    }{
        通过按 $k$ 降序对时刻 $\ell_{j_k+1}$ 应用马尔可夫性质来估计 $\mathbb{P}(B, J=J_0)$. ~由于事件 $B \cap \{J=J_0\}$ 蕴含了:对所有 $j \in J_0^{\mathrm{c}}$,~ 区间最大值 $M_j$ 在区间 $I_{j+1}$ 上保持不变,此时路径被限制在宽度为 $\lceil\eta\Delta\rceil + \Delta$ 的管道内.由此得
    }
    \begin{align*}
        \mathbb{P}(B, J=J_0) 
        &\le \mathbb{E} \left[ \prod_{k=1}^{\#J_0^{\mathrm{c}}} \mathbf{1}_{ \left\{ \forall i \in I_{j_k+1}, \, M_{j_k}-\lceil\eta\Delta\rceil-\Delta \le Y_i \le M_{j_k} \right\} } \right] \\
        &= \mathbb{E} \left[ \prod_{k=2}^{\#J_0^{\mathrm{c}}} \mathbf{1}_{\{\dots\}} \mathbb{E} \left( \mathbf{1}_{ \left\{ \forall i \in I_{j_1+1}, \, M_{j_1}-\lceil\eta\Delta\rceil-\Delta \le Y_i \le M_{j_1} \right\} } \middle| \mathcal{F}_{\ell_{j_1+1}} \right) \right].
    \end{align*}
    \bilingual{
        Since $M_{j_1}$ is $\mathcal{F}_{\ell_{j_1+1}}$-measurable, we can evaluate the inner conditional expectation by considering the starting position $a = Y_{\ell_{j_1+1}}$. On the event $B$, $a$ is constrained within $[M_{j_1}-\lceil\eta\Delta\rceil-\Delta, M_{j_1}]$. Thus, for any $m \in \mathcal{M}$, we have
    }{
        由于 $M_{j_1}$ 是 $\mathcal{F}_{\ell_{j_1+1}}$-可测的,我们可以通过考察起始位置 $a = Y_{\ell_{j_1+1}}$ 来计算内部条件期望.在事件 $B$ 上,~ $a$ 被限制在 $[M_{j_1}-\lceil\eta\Delta\rceil-\Delta, M_{j_1}]$ 内.因此,对于任何 $m \in \mathcal{M}$,~ 成立
    }
    \begin{align*}
        &\quad \mathbb{E} \left( \mathbf{1}_{\{\dots\}} \middle| \mathcal{F}_{\ell_{j_1+1}} \right) \\
        &=\sum_{m\in\mathcal M} \mathbf{1}_{\left\{M_{j_1}=m\right\}} \mathbb{P} \left.\left( \forall i \in [\![1, \lfloor\Delta^{2+\eta}\rfloor]\!], \, m-\lceil\eta\Delta\rceil-\Delta \le Y_i + a \le m \right)\right|_{a=Y_{\ell_{j_1+1}}} \\
        &\le\sum_{m\in\mathcal M} \mathbf{1}_{\left\{M_{j_1}=m\right\}} \sup_{a \in [m-\lceil\eta\Delta\rceil-\Delta, m]} \mathbb{P} \left( \forall i \in [\![1, \lfloor\Delta^{2+\eta}\rfloor]\!], \, m-\lceil\eta\Delta\rceil-\Delta \le Y_i + a \le m \right) \\
        &\le \max_{m \in \mathcal{M}} \sup_{a \in [m-\lceil\eta\Delta\rceil-\Delta, m]} \mathbb{P} \left( \forall i \in [\![1, \lfloor\Delta^{2+\eta}\rfloor]\!], \, m-\lceil\eta\Delta\rceil-\Delta \le Y_i + a \le m \right) \\
        &\le e^{
            (\lceil\eta\Delta\rceil+\Delta)^{\eta/2}
            -
            \frac{
                d(1+2\eta)^4
            }{
                (\lceil\eta\Delta\rceil+\Delta)^2
            }
            \lfloor\Delta^{2+\eta}\rfloor
        }
        \le
        e^{
            (2\eta\Delta+\Delta)^{\eta/2}
            -
            \frac{
                d(1+2\eta)^4
            }{
                (2\eta\Delta+\Delta)^2
            }
            (\Delta^{2+\eta}-\Delta^{2+\eta/2})
        }
        \\ &
        \le
        e^{
            2d\eta(1+2\eta)\Delta^\eta
            -
            d(1+2\eta)^2\Delta^\eta
        }
        =
        e^{
            -d(1+2\eta)\Delta^\eta
        },
    \end{align*}
    \bilingual{
        where the fourth inequality follows from \Cref{lemma:4.1} applied to a tube of width $\lceil\eta\Delta\rceil+\Delta$. This holds for some $\Delta_1=\Delta_1(c,d)>0$ and all $\lceil\eta\Delta\rceil+\Delta\ge\Delta_1$, with parameters $v=\frac\eta 2$, $d(1+2\eta)^4$ in place of $d$, and $\lfloor\Delta^{2+\eta}\rfloor$ in place of $L$. The final inequality holds for some $\Delta_2=\Delta_2(c,d)>\frac 1{1+\eta}\Delta_1$ and all $\Delta\ge\Delta_2$. Iterating this argument for all $k$ yields
    }{
        其中第四个不等式对宽度为 $\lceil\eta\Delta\rceil+\Delta$ 的管道应用\Cref{lemma:4.1} 得到,对于某个 $\Delta_1=\Delta_1(c,d)>0$ 和所有 $\lceil\eta\Delta\rceil+\Delta\ge\Delta_1$ 均成立.参数设置为 $v=\frac\eta 2$,~ 用 $d(1+2\eta)^4$ 替代 $d$,~ 且用 $\lfloor\Delta^{2+\eta}\rfloor$ 替代 $L$. ~而最后一个不等式对某个 $\Delta_2=\Delta_2(c,d)>\frac 1{1+\eta}\Delta_1$ 和所有 $\Delta\ge\Delta_2$ 成立.对所有 $k$ 迭代此过程可得
    }
    \begin{align*}
        \mathbb P(B,J=J_0)
        &
        \le
        \left( e^{ -d(1+2\eta)\Delta^\eta } \right)^{\#J_0^{\mathrm c}}
        \le
        e^{
            -d(1+2\eta)\Delta^\eta
            \left(
                \lfloor \frac L{\lfloor\Delta^{2+\eta}\rfloor} \rfloor
                -
                \lceil \frac C\eta \rceil
            \right)
        }
        \\ &
        \le
        e^{
            -d(1+2\eta)\Delta^\eta
            \left(
                \frac L{\Delta^{2+\eta}}
                -1
                -
                2\frac C\eta
            \right)
        }
        =
        e^{
            -d
            (1+2\eta)
            \frac L{\Delta^2}
            + d
            (1+\eta)(1+2\frac C\eta)
            \Delta^\eta
        }
        \\ &
        \le
        e^{
            -d
            (1+2\eta)
            \frac L{\Delta^2}
            + d
            \eta
            \frac L{\Delta^2}
        }
        =
        e^{
            -d
            \left(1+\eta\right)
            \frac L{\Delta^2}
        },
    \end{align*}
    \bilingual{
        where the last inequality holds for some $\Delta_3=\Delta_3(c,d)>\Delta_2$ and all $\Delta\ge\Delta_3$. Finally, summing over all $J_0 \in \mathscr{J}$, we have
    }{
        其中最后一个不等式对某个 $\Delta_3=\Delta_3(c,d)>\Delta_2$ 和所有 $\Delta\ge\Delta_3$ 成立.于是
    }
    \begin{align*}
        \mathbb P(B)
        &
        =
        \sum_{J_0\in\mathscr J}
        \mathbb P(B,J=J_0)
        \le
        \binom{q}{\lceil C/\eta \rceil}
        e^{
            -d
            \left(1+\eta\right)
            \frac L{\Delta^2}
        }
        \le
        2^{q}
        e^{
            -d
            \left(1+\eta\right)
            \frac L{\Delta^2}
        }
        \\ &
        \le
        e^{
            \frac L{\Delta^{2+\eta/2}}\log 2
            -d
            (1+\eta)
            \frac L{\Delta^2}
        }
        \le
        e^{
            d
            \eta
            \frac L{\Delta^2}
            -d
            (1+\eta)
            \frac L{\Delta^2}
        }
        =
        e^{
            -d
            \frac L{\Delta^2}
        },
    \end{align*}
    \bilingual{
        where the inequalities hold for some $\Delta_0=\Delta_0(c,d)>\Delta_3$ and all $\Delta\ge\Delta_0$.
    }{
        其中不等式对某个 $\Delta_0=\Delta_0(c,d)>\Delta_3$ 和所有 $\Delta\ge\Delta_0$ 成立,证毕.
    }
\end{proof}

\bilingual{
    Based on \Cref{lemma:2.3}, \Cref{lemma:4.1} can be extended to include the case of a small drift.
}{
    基于\Cref{lemma:2.3},~ \Cref{lemma:4.1} 可以推广到包含小漂移的情形.
}

\begin{lemma}[\bilingual{Small Deviation with Small Drift}{小漂移倾斜管道意义下的小偏差}]\label{lemma:4.2}
    \bilingual{
        Let $v>0$, $0<d<\frac{\pi^2}2\sigma^2$ and $c\ge 0$. There exists $\Delta_0>0$ such that for all $\Delta\ge\Delta_0$, $r_0\le 0\le r_0+\Delta$ and $L\ge 1$,
    }{
        设 $v>0$,~ $0<d<\frac{\pi^2}2\sigma^2$ 且 $c\ge 0$. ~存在 $\Delta_0>0$ 使得对所有 $\Delta\ge\Delta_0$,~ $r_0\le 0\le r_0+\Delta$ 且 $L\ge 1$,
    }
    \begin{equation}
        \mathbb P\left(
            r_j\le Y_j\le s_j,\,
            \forall j\in[\![1,L]\!]
        \right)
        \le
        \exp\left(
            \Delta^v-\frac d{\Delta^2}L
        \right),
        \label{eq:4.2}
    \end{equation}
    \bilingual{
        where $r_j:=r_0+\frac c{\Delta^2}j$ and $s_j:=r_j+\Delta$.
    }{
        其中 $r_j:=r_0+\frac c{\Delta^2}j$ 且 $s_j:=r_j+\Delta$.
    }
\end{lemma}

\begin{figure}[htbp]
    \centering
    \begin{tikzpicture}[
        xscale=1.5,
        yscale=1.0,
        >=Stealth,
        rw-path/.style={thick, blue!70!black, mark=*, mark size=1.5pt},
        tube-fill/.style={gray!10},
        label-font/.style={font=\small}
    ]

        % 1. 计算坐标参数 (斜率为 0.2)
        % x=0: r0 = -1.5, s0 = 0.5
        % x=4: r4 = -1.5 + 0.2*4 = -0.7, s4 = 0.5 + 0.2*4 = 1.3
        % x=4.5: r4.5 = -1.5 + 0.2*4.5 = -0.6, s4.5 = 0.5 + 0.2*4.5 = 1.4

        % 2. 背景倾斜填充 (平行四边形区域)
        \fill[tube-fill] (0, -1.5) -- (4, -0.7) -- (4, 1.3) -- (0, 0.5) -- cycle;

        % 3. 参考虚线
        \draw[dashed, black!60] (0, -2) -- (4.5, -2) node[right] {$-\Delta$};
        \draw[dashed, black!60] (4, -2.5) -- (4, 1.8) node[above] {$L$};

        % 4. 倾斜管道边界 (实线)
        % 下边界: y = 0.2x - 1.5
        \draw[thick, black] (0, -1.5) -- (4.5, -0.6) node[right] {$r_n$};
        % 上边界: y = 0.2x + 0.5
        \draw[thick, black] (0, 0.5) -- (4.5, 1.4) node[right] {$s_n$};

        % 5. 坐标轴 (n轴虚线化，突出零基准)
        \draw[->, thick] (0, -2.5) -- (0, 2) node[above] {$Y_n$};
        \draw[->, thick, dashed] (0, 0) -- (5, 0) node[right] {$n$};

        % 6. 随机游走路径 (根据斜率向上微调)
        \draw[rw-path] plot coordinates {
            (0, 0) 
            (0.8, -0.3) 
            (1.6, 0.6) 
            (2.4, -0.1) 
            (3.2, 0.8) 
            (4, 0.4)
        };

        % 7. 左侧跨度标注 (x=0)
        \draw[decorate, decoration={brace, amplitude=5pt}] 
            (-0.1, -1.5) -- (-0.1, 0.5) node[midway, left=6pt] {$\Delta$};

        % 8. 右侧跨度标注 (x=l 的右边)
        \draw[decorate, decoration={brace, amplitude=5pt, mirror}] 
            (4.1, -0.7) -- (4.1, 1.3) node[midway, right=6pt] {$\Delta$};

        % 标注起始点 r0
        \node[left] at (0, -1.5) {$r_0$};

    \end{tikzpicture}
    \caption{\bilingual{Illustration of the small deviation of the random walk \texorpdfstring{$Y_n$}{Yn} within the tilted tube \texorpdfstring{$[r_n, s_n]$}{[r_n, s_n]} with small drift}{随机游走 \texorpdfstring{$Y_n$}{Yn} 在具有小漂移的倾斜管道 \texorpdfstring{$[r_n, s_n]$}{[r_n, s_n]} 内的小偏差示意图}}
\end{figure}

\begin{proof}\label{proof:lemma:4.2}
    \quad

    \bilingual{
        Fix $0<\eta<v\wedge 1$. Let $\ell=\ell(\Delta):=\lfloor\Delta^{2+\eta}\rfloor$. We first note that if $L<\ell$, then for some $\Delta_1=\Delta_1(c,d,v)>0$ and all $\Delta\ge\Delta_1$, we have $d\Delta^\eta<\Delta^v$ since $\eta<v$. Hence, the right-hand side of \eqref{eq:4.2} is greater than $1$ and the inequality trivially holds.
    }{
        固定 $0<\eta<v\wedge 1$. ~令 $\ell=\ell(\Delta):=\lfloor\Delta^{2+\eta}\rfloor$. ~首先注意到如果 $L<\ell$,~ 则对于某个 $\Delta_1=\Delta_1(c,d,v)>0$ 和所有 $\Delta\ge\Delta_1$,~ 成立 $d\Delta^\eta<\Delta^v$, 此因 $\eta<v$. ~因此,~ \eqref{eq:4.2} 的右侧大于 $1$ 并且不等式显然成立.
    }

    \bilingual{
        Now assume $L \ge \ell$. Divide the interval $[\![1, L]\!]$ into $q = \lfloor\frac L\ell \rfloor \ge 1$ intervals of length $\ell$, and a remainder interval of length $r = L - q\ell \in [0, \ell)$. 
        Define the event $B := \{ r_j \le Y_j \le s_j, \, \forall j \in [\![1, L]\!] \}$. 
        For each $k \in [\![1, q]\!]$, let $A_k$ be the event that the random walk stays in the tilted tube during the $k$-th interval, i.e.,
    }{
        现假设 $L \ge \ell$. ~将区间 $[\![1, L]\!]$ 分成 $q = \lfloor \frac L\ell \rfloor \ge 1$ 个长度为 $\ell$ 的区间,以及一个余项区间,长度为 $r = L - q\ell \in [0, \ell)$.
        定义事件 $B := \{ r_j \le Y_j \le s_j, \, \forall j \in [\![1, L]\!] \}$. 
        对于每个 $k \in [\![1, q]\!]$,~ 令 $A_k$ 为随机游走在第 $k$ 个区间内保持在倾斜管道中的事件,即
    }
    $
        A_k := \left\{ r_j \le Y_j \le s_j, \, \forall j \in [\![(k-1)\ell+1, k\ell]\!] \right\}.
    $
    \bilingual{
        Clearly $B \subseteq \bigcap_{k=1}^q A_k$. By the Markov property,
    }{
        显然 $B \subseteq \bigcap_{k=1}^q A_k$. ~利用Markov 性质有
    }
    \begin{equation*}
        \mathbb{P}(B) \le \mathbb{E} \left[ \mathbb{E} \left( \mathbf{1}_{A_q} \middle| \mathcal{F}_{(q-1)\ell} \right) \mathbf{1}_{\bigcap_{k=1}^{q-1} A_k} \right].
    \end{equation*}
    \bilingual{
        For any state $y$ such that $Y_{(q-1)\ell} = y$ on  $\bigcap_{k=1}^{q-1} A_k$, we have $y \in [r_{(q-1)\ell}, s_{(q-1)\ell}]$. The conditional probability is
    }{
        对于事件 $\bigcap_{k=1}^{q-1} A_k$ 上的任何状态 $y = Y_{(q-1)\ell}$,~ 满足 $y \in [r_{(q-1)\ell}, s_{(q-1)\ell}]$. ~其条件概率为
    }
    \begin{align*}
        \mathbb{P}_y(A_q) 
        &= \mathbb{P} \left( r_{(q-1)\ell+j} \le y + Y_j \le s_{(q-1)\ell+j}, \, \forall j \in [\![1, \ell]\!] \right) \\
        &= \mathbb{P} \left( r_{(q-1)\ell} + \textstyle\frac c{\Delta^2}j \le y + Y_j \le r_{(q-1)\ell} + \frac c{\Delta^2}j + \Delta, \, \forall j \in [\![1, \ell]\!] \right) \\
        &= \mathbb{P} \left( r_0' + \textstyle\frac c{\Delta^2}j \le Y_j \le r_0' + \frac c{\Delta^2}j + \Delta, \, \forall j \in [\![1, \ell]\!] \right),
    \end{align*}
    \bilingual{
        where $r_0' := r_{(q-1)\ell} - y$. Since $y \in [r_{(q-1)\ell}, r_{(q-1)\ell} + \Delta]$, we have $r_0' \in [-\Delta, 0]$. 
        On this event, for any $0 \le i < j \le \ell$, the drawdown satisfies $Y_j - Y_i \ge \frac c{\Delta^2}(j-i) - \Delta \ge -\Delta$. 
        Also, the terminal value satisfies $Y_\ell \le r_0' + \frac c{\Delta^2}\ell + \Delta \le \frac c{\Delta^2}\ell + \Delta$. 
        Since $\frac c{\Delta^2}\ell \le \frac{c}{\Delta^2} \Delta^{2+\eta} = c \Delta^\eta$ and $\eta < 1$, we have $Y_\ell \le \Delta (1 + c \Delta^{\eta-1}) \le \frac {2\Delta} \eta$ for some $\Delta_2=\Delta_2(c,d,v)>\Delta_1$ and all $\Delta \ge \Delta_2$. 
        Applying \Cref{lemma:2.3} with $\frac \eta 2$ in lieu of $c$, we have, for some $\Delta_0=\Delta_0(c,d,v)>\Delta_2$ and all $\Delta \ge \Delta_0$,
    }{
        其中 $r_0' := r_{(q-1)\ell} - y$. ~由于 $y \in [r_{(q-1)\ell}, r_{(q-1)\ell} + \Delta]$,~ 成立 $r_0' \in [-\Delta, 0]$. 
        在此事件上,对于任何 $0 \le i < j \le \ell$,~ 回撤满足 $Y_j - Y_i \ge \frac c{\Delta^2}(j-i) - \Delta \ge -\Delta$.
        同时,终点值满足 $Y_\ell \le r_0' + \frac c{\Delta^2}\ell + \Delta \le \frac c{\Delta^2}\ell + \Delta$.
        由于 $\frac c{\Delta^2}\ell \le \frac{c}{\Delta^2} \Delta^{2+\eta} = c \Delta^\eta$ 且 $\eta < 1$,~ 对于某个 $\Delta_2=\Delta_2(c,d,v)>\Delta_1$ 和所有 $\Delta \ge \Delta_2$,~ 成立 $Y_\ell \le \Delta (1 + c \Delta^{\eta-1}) \le \frac {2\Delta} \eta$. 
        应用\Cref{lemma:2.3},参数设为$\frac\eta 2$代替$c$,~ 可得对于某个$\Delta_0=\Delta_0(c,d,v)>\Delta_2$和所有$\Delta\ge\Delta_0$,~ 有
    }
    \begin{equation*}
        \mathbb{E} \left( \mathbf{1}_{A_q} \middle| \mathcal{F}_{(q-1)\ell} \right) \le \sup_{r_0' \in [-\Delta, 0]} \mathbb{P}_{y} (A_q) \le \exp\left( -d \frac{\ell}{\Delta^2} \right).
    \end{equation*}
    \bilingual{
        By induction, we obtain
    }{
        通过归纳迭代,得到
    }
    \begin{equation*}
        \mathbb{P}(B) \le \prod_{k=1}^q \exp\left( -d \frac{\ell}{\Delta^2} \right) = \exp\left( - \frac{d (q\ell)}{\Delta^2} \right).
    \end{equation*}
    \bilingual{
        Substituting $q\ell = L - r$ and noting that $r<\ell$ and $\Delta\ge\Delta_1$, we have
    }{
        代入 $q\ell = L - r$,~ 注意到$r<\ell$与$\Delta\ge\Delta_1$,~ 有
    }
    \begin{equation*}
        \mathbb{P}(B) \le \exp\left( - \frac{dL}{\Delta^2} + \frac{dr}{\Delta^2} \right) \le \exp\left( - \frac{dL}{\Delta^2} + d\Delta^\eta \right) \le \exp\left( \Delta^v - \frac{dL}{\Delta^2} \right),
    \end{equation*}
    \bilingual{
        which completes the proof.
    }{
        证毕.
    }
\end{proof}

\begin{remark}
    \bilingual{
        For $u>0$, if we take $d\Delta^\eta<u\Delta^v$ in lieu of $d\Delta^\eta<\Delta^v$ in the previous proof, we can get an upper bound of
    }{
        对 $u>0$,~ 在前面的证明中用 $d\Delta^\eta<u\Delta^v$ 代替 $d\Delta^\eta<\Delta^v$,~ 可得上界
    }
    $
        e^{
            u\Delta^v-\frac{d}{\Delta^2}L
        }.
    $
\end{remark}

\bilingual{
    The following lemma, using \nameref{theorem:mogulskii}, is similar to \Cref{lemma:4.1} and provides a lower bound for the small deviation probability under the tilted tube end constraint.
}{
    以下引理类似\Cref{lemma:4.1},利用\nameref{theorem:mogulskii},提供了倾斜管道末端限制下的小偏差下界.
}

\begin{lemma}[\bilingual{Lower Bound for Small Deviation under Tilted Tube End Constraint}{倾斜管道末端限制下的小偏差下界}]\label{lemma:4.3}
    \bilingual{
        Let $v>0$, $d'>\frac{\pi^2}{2}\sigma^2$, $c\ge 0$ and $\varepsilon>0$. There exists $\Delta_0>0$ such that for all $\Delta\ge\Delta_0$, $r_0\le 0\le r_0+\Delta$ and $L\ge\lfloor\Delta^{2+\epsilon}\rfloor$,
    }{
        设 $v>0$,~ $d'>\frac{\pi^2}{2}\sigma^2$,~ $c\ge 0$ 且 $\varepsilon>0$. ~存在 $\Delta_0>0$ 使得对所有 $\Delta\ge\Delta_0$,~ $r_0\le 0\le r_0+\Delta$ 且 $L\ge\lfloor\Delta^{2+\epsilon}\rfloor$,
    }
    \begin{align}
        \mathbb P\left(
            r_j\le Y_j\le s_j,\,
            \forall j\in[\![1,L]\!],\,
            Y_L\in[r_L,r_L+v\Delta]
        \right)
        &\ge
        \exp\left(-\frac{d'}{\Delta^2}L\right),
        \label{eq:4.3}
        \\
        \mathbb P\left(
            r_j\le Y_j\le s_j,\,
            \forall j\in[\![1,L]\!],\,
            Y_L\in[s_L-v\Delta,s_L]
        \right)
        &\ge
        \exp\left(-\frac{d'}{\Delta^2}L\right),
        \label{eq:4.3:2}
    \end{align}
    \bilingual{
        where $r_j=r_0+\frac{c}{\Delta^2}j$ and $s_j=r_j+\Delta$.
    }{
        其中 $r_j=r_0+\frac{c}{\Delta^2}j$ 且 $s_j=r_j+\Delta$.
    }
\end{lemma}

\begin{figure}[htbp]
    \centering
    \begin{tikzpicture}[
        xscale=1.7, % 修改这里控制横向（代数 n）的间距
        yscale=1.1, % 修改这里控制纵向（空间位置 V）的伸展程度
        >=Stealth,
        rw-path/.style={thick, blue!70!black, mark=*, mark size=1.5pt},
        tube-fill/.style={gray!10},
        label-font/.style={font=\small}
    ]

        % 1. 坐标参数定义 (斜率 0.2, v=0.4)
        % x=0: r0 = -1.5, s0 = 0.5
        % x=4: rl = -0.7, sl = 1.3, rl+vDelta = -0.7 + 0.8 = 0.1

        % 2. 背景倾斜填充
        \fill[tube-fill] (0, -1.5) -- (4, -0.7) -- (4, 1.3) -- (0, 0.5) -- cycle;

        % 3. 参考虚线与时间尺度标注
        \draw[dashed, black!60] (0, -2) -- (4.5, -2) node[right] {$-\Delta$};
        \draw[dashed, black!60] (4, -2.5) -- (4, 1.8) node[above, black] {$L \ge \Delta^{2+\varepsilon}$};

        % 4. 倾斜管道边界 (实线)
        \draw[thick, black] (0, -1.5) -- (4.5, -0.6) node[right] {$r_n$};
        \draw[thick, black] (0, 0.5) -- (4.5, 1.4) node[right] {$s_n$};

        % 5. 末端限制平台 (x=l处的小平台)
        \draw[dashed, black] (4, -0.7) -- (6, -0.7);
        \draw[dashed, black] (4, 0.1) -- (6, 0.1);
        % 标注 v*Delta 间距
        \draw[<->, black] (6, -0.7) -- (6, 0.1) node[midway, right] {$v\Delta$};
        \node at (6, -0.7) [below] {$r_L$};

        % 6. 坐标轴
        \draw[->, thick] (0, -2.5) -- (0, 2) node[above] {$Y_n$};
        \draw[->, thick, dashed] (0, 0) -- (5, 0) node[right] {$n$};

        % 7. 密集的随机游走路径 (根据斜率向上微调并落在目标区间 [-0.7, 0.1] 内)
        \draw[rw-path] (0, 0) 
            -- (0.2, 0.3) -- (0.4, -0.2) -- (0.6, 0.1) -- (0.8, -0.4) 
            -- (1.0, 0.2) -- (1.2, -0.1) -- (1.4, 0.4) -- (1.6, 0.0) 
            -- (1.8, 0.6) -- (2.0, 0.1) -- (2.2, 0.4) -- (2.4, -0.2) 
            -- (2.6, 0.2) -- (2.8, -0.3) -- (3.0, 0.3) -- (3.2, -0.1) 
            -- (3.4, 0.4) -- (3.6, 0.0) -- (3.8, -0.4) -- (4.0, -0.2);
        \fill[blue!70!black] (4, -0.2) circle (1pt);

        % 8. 跨度标注
        \draw[decorate, decoration={brace, amplitude=5pt}] 
            (-0.1, -1.5) -- (-0.1, 0.5) node[midway, left=6pt] {$\Delta$};

        \draw[decorate, decoration={brace, amplitude=5pt, mirror}] 
            (4.1, -0.7) -- (4.1, 1.3) node[midway, right=6pt] {$\Delta$};

        % 标注起始点 r0
        \node[left] at (0, -1.5) {$r_0$};

    \end{tikzpicture}
    \caption{\bilingual{Illustration of the small deviation lower bound within the tilted tube \texorpdfstring{$[r_L, r_L + v\Delta]$}{[r_l, r_l + v Delta]} with end constraints}{带有末端限制 \texorpdfstring{$[r_L, r_L + v\Delta]$}{[r_l, r_l + v Delta]} 的倾斜管道小偏差下界示意图}}
\end{figure}

\bilingual{
    We prove the first inequality \eqref{eq:4.3}, where the event is denoted by $E$. The proof for \eqref{eq:4.3:2} follows from a symmetric construction. To avoid the use of exponential moments, we divide the evolution into two stages, namely a transition stage, whose duration is of order $\Delta^2$, so that the random walk $Y$ can enter the interior of the tube between $\frac{1}{3}$ and $\frac{2}{3}$, and a drifting stage using Mogul'ski\u i's theorem on a sequence of horizontal sub-tubes like a staircase, with each segment lasting of order $\Delta^{2+\eta}$, where $\eta<\varepsilon\wedge 1$, and the last segment is extended to the final time $L$, with duration of order $2\Delta^{2+\eta}$.
}{
    这里证明第一个不等式 \eqref{eq:4.3},记其中的事件为$E$. ~\eqref{eq:4.3:2} 的证明可通过对称构造得出.为了避开指数矩,下面将演化分为两个阶段,即一个过渡阶段,时长的数量级为$\Delta^2$,~ 使得随机游走$Y$能进入到管道内部$\frac 13$至$\frac 23$处,以及在随后阶梯状的一系列水平子管道上分别应用 Mogul'ski\u i 定理的漂移阶段,每一段时长的数量级为$\Delta^{2+\eta}$,~ 其中$\eta<\varepsilon\wedge 1$,~ 最后一段拉长到终点时刻$L$,~ 其数量级为$2\Delta^{2+\eta}$.
}

\begin{proof}\label{proof:lemma:4.3}
    \bilingual{
        Fix $d''\in(\frac{\pi^2}{2}\sigma^2, d')$.
    }{
        固定 $d''\in(\frac{\pi^2}{2}\sigma^2, d')$.
    }

    \bilingual{
        First, we consider the transition from the starting point $0$ to the interior of the tube. Let $t_1 := \lfloor \Delta^2 \rfloor$. We define the transition event as $E_0$, i.e.,
    }{
        首先,考虑从起点 $0$ 进入管道内部的过渡.令 $t_1 := \lfloor \Delta^2 \rfloor$. ~将该过渡事件定义为 $E_0$,~ 即
    }
    \begin{equation*}
        E_0 := \left\{ r_j \le Y_j \le s_j, \, \forall j \in [\![1, t_1]\!],\, Y_{t_1} \in W_1 \right\},
    \end{equation*}
    \bilingual{
        where $W_1 := [ r_{t_1} + \frac{1}{3}\Delta, r_{t_1} + \frac{2}{3}\Delta ]$ is the central window at time $t_1$. By the local limit theorem for random walks in a strip, the probability that the walk stays inside the moving boundaries $r_j\le Y_j\le s_j$ for $j\le t_1$ and ends in the central window $W_1$ is of order $O(1/\sqrt{t_1})=O(\frac 1\Delta)$. Indeed, after scaling, the process converges weakly to a Brownian motion conditioned to stay within $[0,1]$ (starting from $0$), whose terminal distribution has a strictly positive density on $(\frac 13, \frac 23)$. Consequently, 
    }{
        其中 $W_1 := [ r_{t_1} + \frac{1}{3}\Delta, r_{t_1} + \frac{2}{3}\Delta ]$ 是时刻 $t_1$ 的中心窗口.根据条带内随机游走的局部极限定理,游走在 $j\le t_1$ 的移动边界 $r_j\le Y_j\le s_j$ 内保持并且以 $W_1$ 结尾的概率的数量级为 $O(1/\sqrt{t_1})=O(\frac 1\Delta)$. ~事实上,经过缩放后,该过程弱收敛到一个条件布朗运动,该布朗运动被限制在 $[0,1]$ 内 (从 $0$ 开始),其终点分布在 $(\frac 13, \frac 23)$ 上具有严格正密度.因此,
    }
    \begin{equation*}
        \mathbb{P}(E_0) \ge \frac{C}{\Delta} \ge \exp\left(-\frac{d'-d''}{\Delta^2}L\right),
    \end{equation*}
    \bilingual{
        where $C > 0$, and it holds for some $\Delta_1=\Delta_1(d')>0$ and all $\Delta\ge\Delta_1$.
    }{
        其中$C > 0$,~ 对某个$\Delta_1=\Delta_1(d')>0$和所有$\Delta\ge\Delta_1$成立.
    }

    \bilingual{
        Next, we partition the remaining time $L - t_1$. Fix $\eta \in (0, \varepsilon \wedge 1)$ and let the block length be $\ell = \ell(\Delta) := \lfloor \Delta^{2+\eta} \rfloor$, so that we can get $q := \lfloor \frac{L - t_1}{\ell} \rfloor$ blocks. Since $\eta < \varepsilon$, for some $\Delta_2=\Delta_2(\varepsilon,d')>\Delta_1$ and all $\Delta\ge\Delta_2$, we have $q \ge 1$. Let $t_k := t_1 + (k-1)\ell$, $k \in [\![2, q]\!]$, and $t_{q+1} := L$, and let the $k$-th block be $I_k := [\![t_k, t_{k+1}]\!]$, $k\in[\![1, q]\!]$. Note that the last block has length $\ell \le |I_q| = L - t_q < 2\ell$. For each block $I_k$, define the horizontal sub-tube $H_k := [r_{t_{k+1}}, s_{t_k}]$. This tube is well-defined, since $\eta < 1$, so that for some $\Delta_3=\Delta_3(c,\varepsilon)>\Delta_2$ and all $\Delta \ge \Delta_3$, we have $8\frac{c}{\Delta^2}\ell < \Delta$. Hence the tube width is
    }{
        其次,划分剩余的 $L - t_1$. ~固定 $\eta \in (0, \varepsilon \wedge 1)$ 并令分块长度为 $\ell = \ell(\Delta) := \lfloor \Delta^{2+\eta} \rfloor$,~ 从而得到 $q := \lfloor \frac{L - t_1}{\ell} \rfloor$ 个区间. 由于$\eta<\varepsilon$, ~对某个$\Delta_2=\Delta_2(\varepsilon,d')>\Delta_1$和所有$\Delta\ge\Delta_2$成立$q\ge 1$. ~记时刻 $t_k := t_1 + (k-1)\ell$,~ $k \in [\![2, q]\!]$,~ 且 $t_{q+1} := L$,~ 并记第 $k$ 个区间为 $I_k := [\![t_k, t_{k+1}]\!]$,~ $k\in[\![1, q]\!]$. ~注意最后一个区间长度有$\ell\le|I_q|=L-t_q< 2\ell$. ~在每个区间 $I_k$ 内,定义水平子管道 $H_k := [r_{t_{k+1}}, s_{t_k}]$. ~该管道是良定义的,此因$\eta<1$,~ 从而对某个$\Delta_3=\Delta_3(c,\varepsilon)>\Delta_2$和所有 $\Delta \ge \Delta_3$ 成立 $8\frac{c}{\Delta^2}\ell < \Delta$. ~因而管道宽度为
    }
    \begin{equation*}
        \Delta(k) := s_{t_k}-r_{t_{k+1}}
        =-\frac{c}{\Delta^2}(t_{k+1}-t_k) + \Delta
        \ge -\frac{c}{\Delta^2} 2\ell + \Delta > 0,
        \quad
        k\in[\![1, q]\!].
    \end{equation*}
    \bilingual{
        And $H_k\cap H_{k-1}\neq\emptyset$ for all $k\in[\![2, q]\!]$, which is because
    }{
        并且$H_k\cap H_{k-1}\neq\emptyset$对所有$k\in[\![2, q]\!]$成立,此因
    }
    \begin{equation*}
        s_{t_{k-1}}-r_{t_{k+1}}
        =-\frac{c}{\Delta^2}(t_{k+1}-t_{k-1}) + \Delta
        \ge -\frac{c}{\Delta^2} 3\ell + \Delta > 0.
    \end{equation*}
    \bilingual{
        Additionally, at time $t_k$, $k\in[\![2, q]\!]$, define the target window for the interval $I_{k-1}$ as $W_k := [r_{t_{k+1}} + \frac{1}{3}\Delta(k), r_{t_{k+1}} + \frac{2}{3}\Delta(k)]$ and define the final window as $W_{q+1} := [r_L, r_L + v\Delta]$. Note that for $\Delta \ge \Delta_3$, we have $W_1 \subseteq H_1$ and $W_k \subseteq H_k \cap H_{k-1}$ for all $k \in [\![2, q]\!]$, which is because
    }{
        另外在 $t_k$,~ $k\in[\![2, q]\!]$ 处定义区间$I_{k-1}$的目标窗口为 $W_k:=[r_{t_{k+1}} + \frac{1}{3}\Delta(k), r_{t_{k+1}} + \frac{2}{3}\Delta(k)]$且定义最终窗口为 $W_{q+1} := [r_L, r_L + v\Delta]$. ~注意到, 在$\Delta\ge\Delta_3$时,~ $W_1\subseteq H_1$,~ 并且$W_k\subseteq H_k\cap H_{k-1}$ 对所有 $k\in[\![2, q]\!]$ 成立,此因
    }
    \begin{equation*}
        8\frac{c}{\Delta^2}\ell < \Delta
        \implies
        r_{t_{k+1}}
        <
        r_{t_k} + \frac 13\Delta(k)
        <
        r_{t_k} + \frac 23\Delta(k)
        <
        s_{t_{k-1}},
        \quad
        k\in[\![1, q]\!].
    \end{equation*}

    \bilingual{
        For each $k \in [\![1, q]\!]$, let $E_k$ be the event that the walk stays in $H_k$ during $I_k$ and ends in $W_{k+1}$, i.e.,
    }{
        对于每个 $k \in [\![1, q]\!]$,~ 记 $E_k$ 为随机游走 $Y$ 在 $I_k$ 期间保持在 $H_k$ 内且结束于 $W_{k+1}$ 的事件,即
    }
    \begin{equation*}
        E_k := \left\{ r_{t_{k+1}} \le Y_j \le s_{t_k}, \, \forall j \in I_k,\, Y_{t_{k+1}} \in W_{k+1} \right\}.
    \end{equation*}
    \bilingual{
        According to the Markov property, we have
    }{
        根据 Markov 性质,成立
    }
    \begin{align*}
        \mathbb P\left(E_k \middle| \mathcal F_{t_k}\right)
        &=
        \left.\mathbb P\left( \begin{aligned}  - y + r_{t_{k+1}} \le Y_j - y \le  - y + r_{t_{k+1}} + \Delta(k), \, \forall j \in I_k, \\ Y_{t_{k+1}} - y\in -y+W_{k+1} \end{aligned} \right)\right|_{y=Y_{t_k}}
        \\ &=:
        \left.\mathbb P\left(\; \begin{aligned} r'_0 \le \tilde Y_j \le r'_0 + \Delta(k), \, \forall j \in [\![1, |I_k|]\!], \\ \tilde Y_{|I_k|}\in [r'_0 + v_1\Delta(k), r'_0 + v_2\Delta(k)] \end{aligned} \;\right)\right|_{r'_0=-Y_{t_k}+r_{t_{k+1}}},
    \end{align*}
    \bilingual{
        where for any starting position $Y_{t_k} = y \in W_k$, let $\tilde Y_j := Y_{t_k+j} - y$, $j\in[\![1, |I_k|]\!]$ be the random walk shifted to start from $0$ on the interval $I_k$, and let $r'_0 = r'_0(k) := -y + r_{t_{k+1}}$ be the lower bound of the initial position of this random walk, which satisfies $r'_0 \in [-\frac{2}{3}\Delta(k), -\frac{1}{3}\Delta(k)]$. Additionally, let
    }{
        其中,对任意起始位置 $Y_{t_k} = y \in W_k$,~ 记 $\tilde Y_j:=Y_{t_k+j}-y$,~ $j\in[\![1, |I_k|]\!]$ 是区间 $I_k$ 上平移到从 $0$ 开始的随机游走,并记 $r'_0=r'_0(k):=-y+r_{t_{k+1}}$ 是该随机游走的初始位置下界,成立$r'_0\in[-\frac 23\Delta(k),-\frac 13\Delta(k)]$. ~另外,记
    }
    \begin{align*}
        v_1=v_1(\Delta(k)) &:= \begin{cases}
            \frac 1{\Delta(k)}\left(\frac 23\frac c{\Delta^2}|I_{k+1}|+\frac 13\Delta\right), & k\in[\![1, q-1]\!], \\
            0, & k=q
        \end{cases} 
        =O_{\Delta(k)}(1),\\
        v_2=v_2(\Delta(k)) &:= \begin{cases}
            \frac 1{\Delta(k)}\left(\frac 13\frac c{\Delta^2}|I_{k+1}|+\frac 23\Delta\right), & k\in[\![1, q-1]\!], \\
            \frac{v\Delta}{\Delta(k)}, & k=q
        \end{cases}
        =O_{\Delta(k)}(1).
    \end{align*}
    \bilingual{
        Note that, when $\Delta \ge \Delta_3$,
    }{
        注意到,在$\Delta\ge\Delta_3$时,
    }
    $
        8\frac{c}{\Delta^2}\ell < \Delta
        \implies
        v_1(\Delta(k))
        <1.
    $
    \bilingual{
        Apply \nameref{theorem:mogulskii} (see \Cref{lemma:mogulskii-cal} (iii), with $\Delta(k)$ in lieu of $\Delta$, $|I_k|$ in lieu of $L_\Delta$ and $\tilde Y$ in lieu of $Y$), since
    }{
        应用\nameref{theorem:mogulskii}(参见\Cref{lemma:mogulskii-cal} (iii),参数设置为$\Delta(k)$代替$\Delta$,~ $|I_k|$代替$L_\Delta$, ~$\tilde Y$代替$Y$),由于
    }
    \begin{equation*}
        \lim_{\Delta(k)\to\infty} \frac\Delta{\Delta(k)} = 1,
        \quad
        \lim_{\Delta(k)\to\infty} |I_k| = \infty,
        \quad
        \lim_{\Delta(k)\to\infty} \frac{|I_k|}{\Delta(k)^2} = \infty,
    \end{equation*}
    \bilingual{
        and we get that for all $a\in(0,1)$,
    }{
        得到对所有$a\in(0,1)$, 
    }
    \begin{equation}
        % \quad&
        % \lim_{\Delta\to\infty}\frac{\Delta^2}{|I_k|}
        % \log\mathbb P\left(
        %     \begin{aligned}
        %         -a\Delta(k)\le \tilde Y_j\le(1-a)\Delta(k),\,
        %         \forall j\in[\![1,|I_k|]\!],\\
        %         (v_1-a)\Delta(k)\le \tilde Y_{|I_k|}\le(v_2-a)\Delta(k)
        %     \end{aligned}
        % \right)
        % \nonumber \\
        % &
        % =
        \lim_{\Delta(k)\to\infty}
        % \frac{\Delta^2}{\Delta(k)^2}
        \frac{\Delta(k)^2}{|I_k|}
        \log\mathbb P\left(
            \begin{aligned}
                -a\Delta(k)\le \tilde Y_j\le(1-a)\Delta(k),\,
                \forall j\in[\![1,|I_k|]\!],\\
                (v_1-a)\Delta(k)\le \tilde Y_{|I_k|}\le(v_2-a)\Delta(k)
            \end{aligned}
        \right)
        % \nonumber \\ &
        =
        \frac{\pi^2}2\sigma^2.
        \label{eq:lemma:4.3:1}
    \end{equation}
    \bilingual{
        Following the reasoning as in the \nameref{proof:lemma:4.1}, we can deduce that there exists $\Delta_0=\Delta_0(v,d',\varepsilon,c)>\Delta_3$ such that for all $\Delta\ge\Delta_0$, all $k\in[\![1,q]\!]$ and all $r'_0\in[-\frac 23\Delta(k),-\frac 13\Delta(k)]$,
    }{
        参照\nameref{proof:lemma:4.1}的推理,可推得存在 $\Delta_0=\Delta_0(v,d',\varepsilon,c)>\Delta_3$,~ 使得对所有 $\Delta\ge\Delta_0$,~ $k\in[\![1,q]\!]$ 且 $r'_0\in[-\frac 23\Delta(k),-\frac 13\Delta(k)]$, 成立
    }
    \begin{equation*}
        \mathbb P\left(\;
            \begin{aligned}
                r'_0\le \tilde Y_j\le r'_0+\Delta(k),\,
                \forall j\in[\![1,|I_k|]\!],\\
                \tilde Y_{|I_k|}\in[r'_0+v_1\Delta(k),r'_0+v_2\Delta(k)]
            \end{aligned}
        \;\right)
        \ge
        \exp\left(
            -\frac{d''}{\Delta^2}|I_k|
        \right).
    \end{equation*}
    \bilingual{
        Indeed, fix $m\ge 1$ so that $d''{(1-\frac 1{3m})}^2>\frac{\pi^2}2\sigma^2$. Apply~\eqref{eq:lemma:4.3:1} to $(1-\frac 1{3m})\Delta(k)$ in lieu of $\Delta(k)$ and
    }{
        事实上,固定 $m\ge 1$ 使得 $d''{(1-\frac 1{3m})}^2>\frac{\pi^2}2\sigma^2$. ~将\eqref{eq:lemma:4.3:1} 应用于 $(1-\frac 1{3m})\Delta(k)$ 替代 $\Delta(k)$ 和
    }
    \begin{equation*}
        a=\frac{\frac 23-\frac{l+1}{3m}}{1-\frac 1{3m}}\in(0,1),\quad
        l\in[\![0,m-1]\!],
    \end{equation*}
    \bilingual{
        and we can get, for all $l\in[\![0,m-1]\!]$,
    }{
        可得到,对所有 $l\in[\![0,m-1]\!]$,~
    }
    \begin{align*}
        &\quad
        \lim_{\Delta(k)\to\infty}
        \textstyle\frac 
        {
            {\left(1-\frac 1{3m}\right)}^2
            \Delta(k)^2
        }{|I_k|}
        \log\mathbb P\left(
            \begin{aligned}
                \textstyle-1+\frac{\frac{l+1}m+1}3
                \le\frac{\tilde Y_j}{\Delta(k)}
                \le\frac{\frac lm+1}3,\,
                \forall j,\\
                \textstyle\frac{\frac lm+1}3-1+v_1\le
                \frac{\tilde Y_{|I_k|}}{\Delta(k)},\\
                \textstyle\frac{\tilde Y_{|I_k|}}{\Delta(k)}
                \le\frac{\frac lm+1}3-1+v_2
            \end{aligned}
        \right)
        \\
        &=
        \lim_{\Delta\to\infty}
        \textstyle\frac{\Delta(k)^2}{\Delta^2}
        \frac 
        {
            {\left(1-\frac 1{3m}\right)}^2
            \Delta^2
        }{|I_k|}
        \log\mathbb P\left(
            \begin{aligned}
                \textstyle-1+\frac{\frac{l+1}m+1}3
                \le\frac{\tilde Y_j}{\Delta(k)}
                \le\frac{\frac lm+1}3,\,
                \forall j,\\
                \textstyle\frac{\frac lm+1}3-1+v_1\le
                \frac{\tilde Y_{|I_k|}}{\Delta(k)},\\
                \textstyle\frac{\tilde Y_{|I_k|}}{\Delta(k)}
                \le\frac{\frac lm+1}3-1+v_2
            \end{aligned}
        \right)
        \\ &
        =
        -\frac{\pi^2}{2}\sigma^2
        >
        -d''{\left(1-\frac 1{3m}\right)}^2.
    \end{align*}
     \bilingual{
        Then we find $\Delta_0>0$ such that for all $\Delta\ge\Delta_0$ and $l\in[\![0,m-1]\!]$,
    }{
        于是可得 $\Delta_0>\Delta_3$,~ 使得对所有 $\Delta\ge\Delta_0$ 且 $l\in[\![0,m-1]\!]$,
    }
    \begin{equation*}
        \mathbb P\left(
            \begin{aligned}
                \textstyle-1+\frac{\frac{l+1}m+1}3
                \le\frac{\tilde Y_j}{\Delta(k)}
                \le\frac{\frac lm+1}3,\,
                \forall j,\\
                \textstyle\frac{\frac lm+1}3-1+v_1\le
                \frac{\tilde Y_{|I_k|}}{\Delta(k)},\\
                \textstyle\frac{\tilde Y_{|I_k|}}{\Delta(k)}
                \le\frac{\frac lm+1}3-1+v_2
            \end{aligned}
        \right)
        \ge
        \exp\left(
            -\frac{d''}{\Delta^2}|I_k|
        \right).
    \end{equation*}
    \bilingual{
        Now for $r'_0\in[-\frac 23\Delta(k),-\frac 13\Delta(k)]$, we can find some $l\in[\![0,m-1]\!]$ such that $3(\frac{r'_0}{\Delta(k)}+1)-1\in[\frac lm,\frac{l+1}m]$, and we get by monotonicity of probability that
    }{
        现在对 $r'_0\in[-\frac 23\Delta(k),-\frac 13\Delta(k)]$,~ 存在 $l\in[\![0,m-1]\!]$ 使得 $3(\frac{r'_0}{\Delta(k)}+1)-1\in[\frac lm,\frac{l+1}m]$,~ 根据概率的单调性,
    }
    \begin{equation*}
        \mathbb P\left(\;
            \begin{aligned}
                r'_0\le \tilde Y_j\le r'_0+\Delta(k),\,
                \forall j,\\
                r'_0+v_1\Delta(k)\le \tilde Y_{|I_k|},\\
                \tilde Y_{|I_k|}\le r'_0+v_2\Delta(k)
            \end{aligned}
        \;\right)
        \ge
        \mathbb P\left(
            \begin{aligned}
                \textstyle-1+\frac{\frac{l+1}m+1}3
                \le\frac{\tilde Y_j}{\Delta(k)}
                \le\frac{\frac lm+1}3,\,
                \forall j,\\
                \textstyle\frac{\frac lm+1}3-1+v_1\le
                \frac{\tilde Y_{|I_k|}}{\Delta(k)},\\
                \textstyle\frac{\tilde Y_{|I_k|}}{\Delta(k)}
                \le\frac{\frac lm+1}3-1+v_2
            \end{aligned}
        \right)
        \ge
        e^{
            -\frac{d''}{\Delta^2}|I_k|
        },
    \end{equation*}
    \bilingual{
        Finally, using the Markov property in the reverse order, the probability of the event $E$ in \eqref{eq:4.3} is bounded below by the product of probabilities, which yields, for $\Delta\ge\Delta_0$,
    }{
        最后,倒序利用 Markov 性质,~ \eqref{eq:4.3} 中事件 $E$ 的概率被各分块概率之积所下界控制,从而得到,在 $\Delta\ge\Delta_0$ 时，
    }
    \begin{align*}
        \mathbb{P}(E) 
        &\ge
        \mathbb E\left[
            \mathbf 1_{E_0} \prod_{k=1}^{q} \mathbf 1_{E_k}
        \right]
        =
        \mathbb E\left[
            \mathbf 1_{E_0} \prod_{k=1}^{q-1} \mathbf 1_{E_k}
            \mathbb P\left(E_q \middle| \mathcal F_{t_q}\right)
        \right]
        \\
        &
        \ge \mathbb{P}(E_0) \prod_{k=1}^{q} e^{-\frac{d''}{\Delta^2}|I_k|}
        \ge e^{-\frac{d' - d''}{\Delta^2}L} e^{-\frac{d''}{\Delta^2}(L-t_1)}
        \ge e^{-\frac{d'}{\Delta^2}L},
    \end{align*}
    \bilingual{
        which completes the proof of \eqref{eq:4.3}.
    }{
        此即证得 \eqref{eq:4.3}.
    }
    \qedhere
\end{proof}
% !TEX root = ../main.tex

\section{\bilingual{Lemmas on Branching Random Walk}{有关分支随机游走的引理}}

\begin{lemma}
    \label{lemma:4.5}
    \bilingual{
        Let $v>0$, $0<d<\frac{\pi^2}2\sigma^2$, and let $\Delta$, $r_j\le s_j$ as in \Cref{lemma:4.2}. For any $\ell\ge 1$ and $\delta>r_\ell$, we have
    }{
        设 $v>0$, $0<d<\frac{\pi^2}2\sigma^2$, $\Delta$ 和 $r_j\le s_j$ 如\Cref{lemma:4.2} 所示. 对于任何 $\ell\ge 1$ 和 $\delta>r_\ell$，我们有
    }
    \begin{equation}
        \mathbb E\!\left[
            \left(
                \sum_{|x|=\ell}
                \!\!{\bf 1}_{
                    \{
                        r_j\le\mathcal V(x_j)\le s_j,\,
                        \forall j\in[\![1,\ell]\!],\,
                        \mathcal V(x)\le\delta
                    \}
                }
            \right)^2
        \right]
        \le
        e^{4\Delta^v+2\delta-2d\frac{\ell-1}{\Delta^2}}
        \sum_{m=0}^{\ell-1}
        e^{d\frac m{\Delta^2}-r_m}
        +
        e^{\delta+\Delta^v-\frac d{\Delta^2}\ell}.
        \label{eq:4.6}
    \end{equation}
\end{lemma}
\begin{proof}[Proof of \Cref{lemma:4.5}]
    For any $|x|=\ell\ge 1$, we let
    \begin{equation*}
        F_x
        :=
        {\bf 1}_{
            \{
                r_j\le\mathcal V(x_j)\le s_j,\,
                \forall j\in[\![1,\ell]\!],\,
                \mathcal V(x)\le\delta
            \}
        }.
    \end{equation*}
    For any $|x|=|y|=\ell$ with $x\neq y$, there is an ancester $z$ of $|z|=m\in[\![0,\ell-1]\!]$ such that $z$ is the most recent common ancestor of $x$ and $y$, and then there exists $i,j\in\mathbb N$, $i\neq j$ such that $x_{m+1}=z^i$ and $y_{m+1}=z^j$ where $z^i$ is the $i$-th child of $z$ and that $F_xF_y$ can be only found in $\Sigma(z^i)\Sigma(z^j)$, where
    \begin{equation*}
        \Sigma(u)
        :=
        \sum_{|x|=\ell,\,x_{m+1}=u}
        F_x,
        \quad
        u\in\mathcal T,\,
        |u|=m+1.
    \end{equation*}
    We can then calculate
    \begin{equation*}
        \left(
            \sum_{|x|=\ell}
            F_x
        \right)^2
        =
        \sum_{|x|=|y|=\ell}
        F_xF_y
        =
        \sum_{m=0}^{\ell-1}
        \sum_{|z|=m}
        \sum_{i,j\in\mathbb N,\,i\neq j}
        \!\!\!
        \Sigma(z^i)\Sigma(z^j)
        +
        \sum_{|x|=\ell}
        F_x^2.
    \end{equation*}
    Rewrite $\Sigma(u)$ as
    \begin{equation*}
        \Sigma(u)
        =
        \sum_{|x|=\ell,\,x_{m+1}=u}
        \!\!\!
        F_x
        =
        {\bf 1}_{
            \{
                r_j\le\mathcal V(u_j)\le s_j,\,
                \forall j\in[\![1,m]\!]
            \}
        }
        \widetilde\Sigma(u)
    \end{equation*}
    with
    \begin{equation*}
        \widetilde\Sigma(u)
        :=
        \sum_{|x|=\ell,\,x_{m+1}=u}
        \!\!\!
        {\bf 1}_{
            \{
                r_j\le\mathcal V(x_j)\le s_j,\,
                \forall j\in[\![m+1,\ell]\!],\,
                \mathcal V(x)\le\delta
            \}
        }.
    \end{equation*}
    By the branching property and the many-to-one formula, we have
    \begin{align*}
        \mathbb E\!\left(\!
            \widetilde\Sigma(u)
            \middle|
            \mathcal F_{m+1}
        \!\right)
        &
        =
        \left.
            \mathbb E\!\left[
                \sum_{|x|=\ell-m-1}
                \!\!\!\!{\bf 1}_{
                    \{
                        r_{j+m+1}\le\mathcal V(x_j)+a\le s_{j+m+1},\,
                        \forall j\in[\![0,\ell-m-1]\!],\,
                        \mathcal V(x)+a\le\delta
                    \}
                }
            \right]
        \right|_{a=\mathcal V(u)}
        \\ &
        =
        \left.
            \mathbb E\!\left[
                e^{Y_{\ell-m-1}}
                {\bf 1}_{
                    \{
                        r_{j+m+1}\le Y_j+a\le s_{j+m+1},\,
                        \forall j\in[\![0,\ell-m-1]\!],\,
                        Y_{\ell-m-1}+a\le\delta
                    \}
                }
            \right]
        \right|_{a=\mathcal V(u)}
        \\ &
        \le\left.
            e^{\delta-a}
            \mathbb E\!\left[
                {\bf 1}_{
                    \{
                        r_{j+m+1}\le Y_j+a\le s_{j+m+1},\,
                        \forall j\in[\![0,\ell-m-1]\!]
                    \}
                }
            \right]
        \right|_{a=\mathcal V(u)}
        \\ &
        =\left.
            e^{\delta-a}
            {\bf 1}_{
                r_{m+1}\le a\le s_{m+1}
            }
            \mathbb E\!\left[
                {\bf 1}_{
                    \{
                        r_{m+1}+qj
                        \le Y_j+a
                        \le s_{m+1}+qj,\,
                        \forall j\in[\![1,\ell-m-1]\!]
                    \}
                }
            \right]
        \right|_{a=\mathcal V(u)}
        \\ &
        \le\left.
            e^{\delta-a}
            e^{
                \Delta^v
                -
                d\frac{\ell-m-1}{\Delta^2}
            }
        \right|_{a=\mathcal V(u)}
        =
        e^{\delta-\mathcal V(u)}
        e^{
            \Delta^v
            -
            d\frac{\ell-m-1}{\Delta^2}
        },
    \end{align*}
    where the last inequality holds for some $\Delta_1=\Delta_1(v,d,c)>0$, all $\Delta\ge\Delta_1$, and all $a\not\in[r_{m+1},s_{m+1}]$ as well as $a\in[r_{m+1},s_{m+1}]$ since in the latter case $r_{m+1}-a\in[-\Delta,0]$ and we can invoke \Cref{lemma:4.2} for $r_{m+1}-a$ in lieu of $r_0$. Then we have
    \begin{align*}
        \mathbb E\!\left(\!
            \sum_{i,j\in\mathbb N,\,i\neq j}
            \!\!\!
            \widetilde\Sigma(z^i)
            \widetilde\Sigma(z^j)
            \middle|
            \mathcal F_m
        \!\right)
        &
        =
        \mathbb E\!\left(
            \mathbb E\!\left(\!
                \sum_{i,j\in\mathbb N,\,i\neq j}
                \!\!\!
                \widetilde\Sigma(z^i)
                \widetilde\Sigma(z^j)
                \middle|
                \mathcal F_{m+1}
            \!\right)
            \middle|
            \mathcal F_m
        \!\right)
        \\ &
        =
        \mathbb E\!\left(
            \sum_{i,j\in\mathbb N,\,i\neq j}
            \!\!\!
            \mathbb E\!\left(\!
                \widetilde\Sigma(z^i)
                \widetilde\Sigma(z^j)
                \middle|
                \mathcal F_{m+1}
            \!\right)
            \middle|
            \mathcal F_m
        \!\right)
        \\ &
        =
        \mathbb E\!\left(
            \sum_{i,j\in\mathbb N,\,i\neq j}
            \!\!\!
            \mathbb E\!\left(\!
                \widetilde\Sigma(z^i)
                \middle|
                \mathcal F_{m+1}
            \!\right)
            \mathbb E\!\left(\!
                \widetilde\Sigma(z^j)
                \middle|
                \mathcal F_{m+1}
            \!\right)
            \middle|
            \mathcal F_m
        \!\right)
        \\ &
        \le
        e^{
            2\delta
            +2\Delta^v
            -2d\frac{\ell-m-1}{\Delta^2}
        }
        \mathbb E\!\left(
            \sum_{i,j\in\mathbb N,\,i\neq j}
            \!\!\!
            e^{
                -(
                    \mathcal V(z^i)
                    +\mathcal V(z^j)
                )
            }
            \middle|
            \mathcal F_m
        \!\right)
        \\ &
        =
        e^{
            2\delta
            +2\Delta^v
            -2d\frac{\ell-m-1}{\Delta^2}
        }
        \left.
            \mathbb E\!\left[
                \sum_{i,j\in\mathbb N,\,i\neq j}
                \!\!\!
                e^{
                    -(
                        (\mathcal V(\varnothing^i)+a)
                        +(\mathcal V(\varnothing^j)+a)
                    )
                }
            \!\right]
        \right|_{a=\mathcal V(z)}
        \\ &
        =
        e^{
            2\delta
            +2\Delta^v
            -2d\frac{\ell-m-1}{\Delta^2}
            -2\mathcal V(z)
        }
        \mathbb E\!\left[
            \sum_{i,j\in\mathbb N,\,i\neq j}
            \!\!\!
            e^{-(\Xi_i+\Xi_j)}
        \!\right].
    \end{align*}
    Under \Cref{assumption:log-expt-cross-sum-exp-neg-sum} that
    \begin{equation*}
        \log
        \mathbb E\!\left[
            \sum_{i,j\in\mathbb N,\,i\neq j}
            \!\!\!
            e^{-(\Xi_i+\Xi_j)}
        \!\right]
        =
        C
        <
        \infty,
    \end{equation*}
    we obtain that
    \begin{equation*}
        \mathbb E\!\left(
            \sum_{i,j\in\mathbb N,\,i\neq j}
            \!\!\!
            \widetilde\Sigma(z^i)
            \widetilde\Sigma(z^j)
            \middle|
            \mathcal F_m
        \!\right)
        \le
        e^{
            2\delta
            +3\Delta^v
            -2d\frac{\ell-m-1}{\Delta^2}
            -2\mathcal V(z)
        }
    \end{equation*}
    holds for some $\Delta_0=\Delta_0(v,d,c)>\Delta_1$ and all $\Delta\ge\Delta_0$. Then we have
    \begin{align*}
        &\quad
        \mathbb E\!\left(
            \sum_{m=0}^{\ell-1}
            \sum_{|z|=m}
            \sum_{i,j\in\mathbb N,\,i\neq j}
            \!\!\!
            \Sigma(z^i)
            \Sigma(z^j)
        \right)
        \\
        &
        =
        \sum_{m=0}^{\ell-1}
        \mathbb E\!\left[
            \sum_{|z|=m}
            {\bf 1}_{
                \{
                    r_j\le\mathcal V(z_j)\le s_j,\,
                    \forall j\in[\![1,m]\!]
                \}
            }
            \mathbb E\!\left(
                \sum_{i,j\in\mathbb N,\,i\neq j}
                \!\!\!
                \widetilde\Sigma(z^i)
                \widetilde\Sigma(z^j)
                \middle|
                \mathcal F_m
            \right)
        \right]
        \\ &
        \le
        e^{2\delta+3\Delta^v}
        \sum_{m=0}^{\ell-1}
        e^{-2d\frac{\ell-m-1}{\Delta^2}}
        \mathbb E\!\left[
            \sum_{|z|=m}
            {\bf 1}_{
                \{
                    r_j\le\mathcal V(z_j)\le s_j,\,
                    \forall j\in[\![1,m]\!]
                \}
            }
            e^{-2\mathcal V(z)}
        \right]
        \\ &
        \le
        e^{2\delta+3\Delta^v}
        \sum_{m=0}^{\ell-1}
        e^{-2d\frac{\ell-m-1}{\Delta^2}}
        e^{-r_m}
        \mathbb E\!\left[
            \sum_{|z|=m}
            {\bf 1}_{
                \{
                    r_j\le\mathcal V(z_j)\le s_j,\,
                    \forall j\in[\![1,m]\!]
                \}
            }
            e^{-\mathcal V(z)}
        \right]
        \\ &
        =
        e^{2\delta+3\Delta^v}
        \sum_{m=0}^{\ell-1}
        e^{-2d\frac{\ell-m-1}{\Delta^2}}
        e^{-r_m}
        \mathbb P(
            r_j\le Y_j\le s_j,\,
            \forall j\in[\![1,m]\!]
        )
        \\ &
        \le
        e^{2\delta+3\Delta^v}
        \sum_{m=0}^{\ell-1}
        e^{-2d\frac{\ell-m-1}{\Delta^2}}
        e^{-r_m}
        e^{\Delta^v-d\frac m{\Delta^2}}
        =
        e^{4\Delta^v+2\delta-2d\frac{\ell-1}{\Delta^2}}
        \sum_{m=0}^{\ell-1}
        e^{d\frac m{\Delta^2}-r_m}.
    \end{align*}
    This gives the first term on the right-hand side of \eqref{eq:4.6}. For the second term, we have
    \begin{equation*}
        \mathbb E\!\left(
            \sum_{|x|=\ell}
            F_x^2
        \right)
        =
        \mathbb E\!\left(
            \sum_{|x|=\ell}
            F_x
        \right)
        =
        \mathbb E\!\left[
            e^{Y_\ell}
            {\bf 1}_{
                \{
                    r_j\le Y_j\le s_j,\,
                    \forall j\in[\![1,\ell]\!],\,
                    Y_\ell\le\delta
                \}
            }
        \right]
        \le
        e^{\delta+\Delta^v-\frac d{\Delta^2}\ell},
    \end{equation*}
    which completes the proof.
\end{proof}

\begin{lemma}
    \label{lemma:4.6}
    Let $c>0$. For $L\ge 1$, let $r_j:=\frac{\pi^2\sigma^2}{2c^2}L^{-2/3}j$ and $s_j:=r_j+cL^{1/3}$ for $j\ge 0$. For any $v>0$, we have
    \begin{equation}
        \lim_{L\to\infty}
        L^{-\frac 13}
        \log\mathbb P(
            \exists |x|=L:\,
            r_j\le\mathcal V(x_j)\le s_j,\,
            \forall j\in[\![1,L]\!],\,
            \mathcal V(x)\le r_L+vL^{\frac 13}
        )
        =
        0.
        \label{eq:4.7}
    \end{equation}
    Moreover, 
    \begin{equation}
        \lim_{L\to\infty}
        L^{-\frac 13}
        \log\mathbb E\!\left[
            \sum_{1\le |x|\le L}
            {\bf 1}_{
                \{
                    r_j\le\mathcal V(x_j)\le s_j,\,
                    \forall j\in[\![1,|x|]\!]
                \}
            }
        \right]
        =
        c.
        \label{eq:4.8}
    \end{equation}
\end{lemma}

We first prove \eqref{eq:4.7}. It suffices to prove the lower bound, since the upper bound is always $0$. For the lower bound, we invoke the Paley--Zygmund inequality to turn the problem into one lower bound and another upper bound for a second moment term. The lower bound can be obtained by \Cref{lemma:4.3} and the upper bound can be obtained by \Cref{lemma:4.5,lemma:4.2}. Then we proceed to prove \eqref{eq:4.8}. The upper bound can be obtained by \Cref{lemma:4.2} and the lower bound can be obtained by \Cref{lemma:4.3}.

\begin{proof}[Proof of \Cref{lemma:4.6}]
    Let for $|x|=L$,
    \begin{equation*}
        F_x
        :=
        {\bf 1}_{
            \{
                r_j\le\mathcal V(x_j)\le s_j,\,
                \forall j\in[\![1,L]\!],\,
                \mathcal V(x)\le r_L+vL^{1/3}
            \}
        }.
    \end{equation*}
    Apply \eqref{eq:4.3} in \Cref{lemma:4.3} with $\Delta=cL^{1/3}$, $\frac vc$ in lieu of $v$, $0$ in lieu of $r_0$, and $\frac{\pi^2}2\sigma^2$ in lieu of $c$, we obtain that for any $d'>\frac{\pi^2}2\sigma^2$, some $L_1=L_1(c,v,d')>0$ and all $L\ge L_1$,
    \begin{align*}
        \mathbb E\!\left[
            \sum_{|x|=L}
            F_x
        \right]
        &
        =
        \mathbb E\!\left[
            e^{Y_L}
            {\bf 1}_{
                \{
                    r_j\le Y_j\le s_j,\,
                    \forall j\in[\![1,L]\!],\,
                    Y_L\le r_L+vL^{1/3}
                \}
            }
        \right]
        \\ &
        \ge
        e^{r_L}
        \mathbb P(
            r_j\le Y_j\le s_j,\,
            \forall j\in[\![1,L]\!],\,
            Y_L\le r_L+vL^{1/3}
        )
        \\ &
        \ge
        e^{r_L}
        e^{-\frac{d'}{(cL^{1/3})^2}L}
        =
        e^{\varepsilon'L^{1/3}},
    \end{align*}
    where we denote $\varepsilon'=\varepsilon'(c,d'):=(d'-\frac{\pi^2}{2}\sigma^2)c^{-2}$. Similarly, apply \eqref{eq:4.2} in \Cref{lemma:4.2} with $\Delta=cL^{1/3}$, $\frac 12$ in lieu of $v$, $0$ in lieu of $r_0$, and $\frac{\pi^2}2\sigma^2$ in lieu of $c$, we obtain that for any $d<\frac{\pi^2}2\sigma^2$, some $L_2=L_2(c,v,d)>0$ and all $L\ge L_2$, we have
    \begin{align*}
        \mathbb E\!\left[
            \sum_{|x|=L}
            F_x
        \right]
        &
        =
        \mathbb E\!\left[
            e^{Y_L}
            {\bf 1}_{
                \{
                    r_j\le Y_j\le s_j,\,
                    \forall j\in[\![1,L]\!],\,
                    Y_L\le r_L+vL^{1/3}
                \}
            }
        \right]
        \\ &
        \le
        e^{r_L+vL^{1/3}}
        \mathbb P(
            r_j\le Y_j\le s_j,\,
            \forall j\in[\![1,L]\!]
        )
        \\ &
        \le
        e^{r_L+vL^{1/3}}
        e^{(cL^{1/3})^{1/2}-\frac{d}{(cL^{1/3})^2}L}
        \\ &
        =
        e^{(v+(\frac{\pi^2}2\sigma^2-d)c^{-2})L^{1/3}}
        e^{c^{1/2}L^{1/6}}
        \le
        e^{(v+2\varepsilon)L^{1/3}},
    \end{align*}
    where $\varepsilon=\varepsilon(c,d):=(\frac{\pi^2}{2}\sigma^2-d)c^{-2}$, and the last inequality holds since
    $%\begin{equation*}
        c^{1/2}L^{1/6}
        \le
        \varepsilon L^{1/3}
    $%\end{equation*}
    for some $L_2>L_3=L_3(c,d)>0$ and all $L\ge L_3$.

    We compute now the secomd moment of $\sum_{|x|=L}F_x$ using \eqref{eq:4.6} in \Cref{lemma:4.5} with $\Delta=cL^{1/3}$, $\ell=L$, $\frac 12$ in lieu of $v$, $r_j$ and $s_j$ as defined in the statement of this lemma (i.e. $0$ in lieu of $r_0$ and $\frac{\pi^2}2\sigma^2$ in lieu of $c$ for \Cref{lemma:4.2}), and $\delta=r_L+vL^{1/3}$. We check that for $m\le L$, the terms in \eqref{eq:4.6} satisfy
    \begin{align*}
        d\frac{m}{\Delta^2}-r_m
        &
        =
        d\frac{m}{(cL^{1/3})^2}-\frac{\pi^2\sigma^2}{2c^2}L^{-2/3}m
        =
        -\varepsilon L^{-2/3}m,
        \\ 
        \delta-\frac d{\Delta^2}\ell
        &
        =
        \frac{\pi^2\sigma^2}{2c^2}L^{-2/3}L+vL^{1/3}-d\frac{L}{(cL^{1/3})^2}
        =
        \left(v+\varepsilon\right)L^{1/3}.
    \end{align*}
    Then for some $L_4=L_4(c,v,d)>L_2$ and all $L\ge L_4$,
    \begin{align*}
        \mathbb E\!\left[
            \left(
                \sum_{|x|=L}
                F_x
            \right)^2
        \right]
        &
        \le
        e^{4c^{1/2}L^{1/6}+2(v+\varepsilon)(L-1)^{1/3}}
        \sum_{m=0}^{L-1}
        e^{-\varepsilon L^{-2/3}m}
        +
        e^{c^{1/2}L^{1/6}+(v+\varepsilon)L^{1/3}}
        \\ &
        \le
        e^{(2v+6\varepsilon)L^{1/3}}
        L
        +e^{(v+2\varepsilon)L^{1/3}}
        \le
        (L+1)e^{(2v+6\varepsilon)L^{1/3}}
        \\ &
        \le
        e^{(2v+7\varepsilon)L^{1/3}},
    \end{align*}
    where the last inequality follows from the fact that $\log(L+1)\le\varepsilon L^{1/3}$ for some $L_4>L_5=L_5(c,d)>0$ and all $L\ge L_5$. We deduce from the Paley--Zygmund inequality that for some $L_0=L_0(c,v,d,d')=L_1+L_4$ and all $L\ge L_0$,
    \begin{align*}
        \mathbb P\left(
            \sum_{|x|=L}
            F_x
            \ge
            1
        \right)
        &
        \ge
        \frac{
            \mathbb E\!\left[
                \sum_{|x|=L}
                F_x
            \right]^2
        }{
            \mathbb E\!
            \left[
                \left(
                    \sum_{|x|=L}
                    F_x
                \right)^2
            \right]
        }
        \ge
        \frac{
            e^{-2\varepsilon'L^{1/3}}
        }{
            e^{(2v+7\varepsilon)L^{1/3}}
        }
        =
        e^{-(2v+7\varepsilon+2\varepsilon')L^{1/3}}.
    \end{align*}
    Hence for any $\eta>0$, we can find a $d=d(c,\eta)$ and a $d'=d'(c,\eta)$ such that $7\varepsilon+2\varepsilon'\le\eta$ and then
    \begin{equation*}
        \liminf_{L\to\infty}
        L^{-\frac 13}
        \log
        \mathbb P\left(
            \sum_{|x|=L}
            F_x
            \ge
            1
        \right)
        \ge
        -(2v+7\varepsilon+2\varepsilon')
        \ge
        -(2v+\eta).
    \end{equation*}
    The left-hand side depend on $v$ but not $\eta$. Let $\eta\to 0$. Moreover, the probability on the left-hand side of the above inequality is exactly the probability in \eqref{eq:4.7} and is decreasing when $v$ is decreasing. Hence, let $v$ goes to $0$ and
    \begin{equation*}
        \liminf_{L\to\infty}
        L^{-\frac 13}
        \log\mathbb P\left(
            \sum_{|x|=L}
            F_x
            \ge
            1
        \right)
        \ge
        \sup_{v_0\in(0,v)}
        (-2v_0)
        =
        0,
    \end{equation*}
    which gives \eqref{eq:4.7}.

    Now we prove \eqref{eq:4.8}. Let $\ell\le L$. Apply \eqref{eq:4.2} in \Cref{lemma:4.2} of the version \eqref{eq:4.1:remark} with $\Delta=cL^{1/3}$, any $v\in(0,1)$, $u=c^{-v}$, $0$ in lieu of $r_0$, and $\frac{\pi^2}2\sigma^2$ in lieu of $c$, we have that for any $d<\frac{\pi^2}2\sigma^2$, any $v\in(0,1)$, some $L'_0=L'_0(c,v,d)$ and all $L\ge L'_0$,
    \begin{align*}
        \mathbb E\!\left[
            \sum_{|x|=\ell}
            {\bf 1}_{
                \{
                    r_j\le\mathcal V(x_j)\le s_j,\,
                    \forall j\in[\![1,\ell]\!]
                \}
            }
        \right]
        &
        =
        \mathbb E\!\left[
            e^{Y_\ell}
            {\bf 1}_{
                \{
                    r_j\le Y_j\le s_j,\,
                    \forall j\in[\![1,\ell]\!]
                \}
            }
        \right]
        \\ &
        \le
        e^{s_\ell}
        \mathbb P(
            r_j\le Y_j\le s_j,\,
            \forall j\in[\![1,\ell]\!]
        )
        \\ &
        \le
        e^{s_\ell}
        e^{c^{-v}(cL^{1/3})^v-\frac{d}{(cL^{1/3})^2}\ell}
        =
        e^{\varepsilon L^{-2/3}\ell+L^{v/3}+cL^{1/3}},
    \end{align*}
    where $\varepsilon=(\frac{\pi^2}{2}\sigma^2-d)c^{-2}$ as above. Then take the sum over all $\ell$ to obtain
    \begin{align*}
        \mathbb E\!\left[
            \sum_{1\le |x|\le L}
            {\bf 1}_{
                \{
                    r_j\le\mathcal V(x_j)\le s_j,\,
                    \forall j\in[\![1,|x|]\!]
                \}
            }
        \right]
        &
        =
        \sum_{\ell=1}^L
        \mathbb E\!\left[
            \sum_{|x|=\ell}
            {\bf 1}_{
                \{
                    r_j\le\mathcal V(x_j)\le s_j,\,
                    \forall j\in[\![1,\ell]\!]
                \}
            }
        \right]
        \\ &
        \le
        \sum_{\ell=1}^L
        e^{\varepsilon L^{-2/3}\ell+L^{v/3}+cL^{1/3}}
        \le
        Le^{\varepsilon L^{1/3}+L^{v/3}+cL^{1/3}},
    \end{align*}
    and hence
    \begin{align*}
        &\quad
        \limsup_{L\to\infty}
        L^{-\frac 13}
        \log\mathbb E\!\left[
            \sum_{1\le |x|\le L}
            {\bf 1}_{
                \{
                    r_j\le\mathcal V(x_j)\le s_j,\,
                    \forall j\in[\![1,|x|]\!]
                \}
            }
        \right]
        \\
        &
        \le
        \varepsilon+c+
        \limsup_{L\to\infty}
        \left(
            \frac{\log L}{L^{1/3}}
            +
            \frac{1}{L^{(1-v)/3}}
        \right)
        % \\ &
        =
        \varepsilon+c
        \le
        \eta+c,
    \end{align*}
    which gives the upper bound in \eqref{eq:4.8} since $\eta$ can be arbitrarily small. For the lower bound, we apply \eqref{eq:4.3:2} in \Cref{lemma:4.3} with $\Delta=cL^{1/3}$, some $v>0$, $\frac vc$ in lieu of $v$, $0$ in lieu of $r_0$, and $\frac{\pi^2}2\sigma^2$ in lieu of $c$, we have that for any $d'>\frac{\pi^2}2\sigma^2$, some $L''_0=L''_0(c,v,d')>0$ and all $L\ge L''_0$,
    \begin{align*}
        &\quad
        \mathbb E\!\left[
            \sum_{1\le |x|\le L}
            {\bf 1}_{
                \{
                    r_j\le\mathcal V(x_j)\le s_j,\,
                    \forall j\in[\![1,|x|]\!]
                \}
            }
        \right]
        \\
        &
        \ge
        \mathbb E\!\left[
            \sum_{|x|=L}
            {\bf 1}_{
                \{
                    r_j\le\mathcal V(x_j)\le s_j,\,
                    \forall j\in[\![1,L]\!]
                \}
            }
        \right]
        \\ &
        \ge
        \mathbb E\!\left[
            e^{Y_L}
            {\bf 1}_{
                \{
                    r_j\le Y_j\le s_j,\,
                    \forall j\in[\![1,L]\!],\,
                    Y_L\ge s_L-vL^{1/3}
                \}
            }
        \right]
        \\ &
        \ge
        e^{s_L-vL^{1/3}}
        \mathbb P(
            r_j\le Y_j\le s_j,\,
            \forall j\in[\![1,L]\!],\,
            Y_L\ge s_L-vL^{1/3}
        )
        \\ &
        \ge
        e^{s_L-vL^{1/3}}
        e^{-\frac{d'}{(cL^{1/3})^2}L}
        =
        e^{(c-v-\varepsilon')L^{1/3}},
    \end{align*}
    and hence
    \begin{equation*}
        \liminf_{L\to\infty}
        L^{-\frac 13}
        \log\mathbb E\!\left[
            \sum_{1\le |x|\le L}
            {\bf 1}_{
                \{
                    r_j\le\mathcal V(x_j)\le s_j,\,
                    \forall j\in[\![1,|x|\!]
                \}
            }
        \right]
        \ge
        c-v-\varepsilon'
        \ge
        c-v-\eta,
    \end{equation*}
    which completes the proof for \eqref{eq:4.8} since $v$ and $\eta$ can be arbitrarily small.
\end{proof}

\begin{lemma}
    \label{lemma:4.7}
    Let $b>v>0$, $r:=-b\log n$. We have
    \begin{equation}
        \lim_{n\to\infty}
        \frac{1}{\log n}
        \log\mathbb P\!\left(
                \exists |x|
                =
                \lfloor\log^{5/2}n\rfloor:
                \begin{aligned}
                    r\le\mathcal V(x_j)\le 0,\,
                    \forall j\in[\![1,|x|]\!],\\
                    \mathcal V(x)\le r+v\log n
                \end{aligned}
        \right)
        =
        v-b.
        \label{eq:4.9}
    \end{equation}
    Moreover,
    \begin{equation}
        \mathbb E\!\left[
            \sum_{\ell=1}^\infty
            \sum_{|x|=\ell}
            {\bf 1}_{
                \{
                    r\le\mathcal V(x_j)\le 0,\,
                    \forall j\in[\![1,\ell]\!]
                \}
            }
        \right]
        =
        O_n(\log n)
        \label{eq:4.10}
    \end{equation}
    and
    \begin{equation}
        \mathbb E\!\left[
            \sum_{
                |x|=\lfloor\log^{5/2}n\rfloor
            }
            \!\!\!
            {\bf 1}_{
                \{
                    r\le\mathcal V(x_j)\le 0,\,
                    \forall j\in[\![1,|x|]\!],\,
                    \mathcal V(x)\le r+v\log n
                \}
            }
        \right]
        =
        n^{-b+v+o_n(1)}.
        \label{eq:4.11}
    \end{equation}
\end{lemma}

We first prove \eqref{eq:4.11}. \Cref{lemma:4.1,lemma:4.3} are used for upper and lower bounds respectively. Note that the tighter lower boundary $r+v'\log n$ for some $v'\in(0,v)$ is used in the lower bound. Then we prove \eqref{eq:4.9}. For the upper bound, the results follows from the proven \eqref{eq:4.11}. For the lower bound, the Paley--Zygmund inequality and the proven \eqref{eq:4.11} are invoked to turn the problem into an upper bound for a second moment term with respect to the tighter lower boundary $r+v'\log n$ mentioned above. This upper bound can be obtained by \Cref{lemma:4.6}. The remaining proof for \eqref{eq:4.10} follows from standard properties of random walks.

\begin{proof}[Proof of Lemma \ref{lemma:4.7}]
    \quad

    Let for $|x|=\lfloor\log^{5/2}n\rfloor$,
    \begin{equation*}
        F_x
        =
        F_x(b,v,n)
        :=
        {\bf 1}_{
            \{
                r\le\mathcal V(x_j)\le 0,\,
                \forall j\in[\![1,|x|]\!],\,
                \mathcal V(x)\le r+v\log n
            \}
        }.
    \end{equation*}
    We first prove \eqref{eq:4.11}, or equivalently
    \begin{equation*}
        \lim_{n\to\infty}
        \frac{1}{\log n}
        \log\mathbb E\!\left[
            \sum_{
                |x|=\lfloor\log^{5/2}n\rfloor
            }
            \!\!\!
            F_x
        \right]
        =
        -b+v.
    \end{equation*}
    For the upper bound, apply \eqref{eq:4.1} in \Cref{lemma:4.1} with $\Delta=-r$, $r_0=r$, $\ell=\lfloor\log^{5/2}n\rfloor$ and $v=\frac 12$, and we have that for any fixed $0<d<\frac{\pi^2}2\sigma^2$, some $n_1=n_1(d)>0$ and all $n\ge n_1$,
    \begin{align*}
        % &\quad
        \mathbb E\!\left[
            \sum_{
                |x|=\lfloor\log^{5/2}n\rfloor
            }
            \!\!\!
            F_x
        \right]
        % \\
        &
        =
        \mathbb E\!\left[
            e^{
                Y_{\lfloor\log^{5/2}n\rfloor}
            }
            {\bf 1}_{
                \{
                    r\le Y_j\le 0,\,
                    \forall j\in[\![1,\lfloor\log^{5/2}n\rfloor]\!],\,
                    Y_{\lfloor\log^{5/2}n\rfloor}\le r+v\log n
                \}
            }
        \right]
        \nonumber
        \\ &
        \le
        e^{r+v\log n}
        \mathbb E\!\left[
            {\bf 1}_{
                \{
                    r\le Y_j\le 0,\,
                    \forall j\in[\![1,\lfloor\log^{5/2}n\rfloor]\!]
                \}
            }
        \right]
        \nonumber
        \\ &
        \le
        n^{-b+v}
        e^{
            (b\log n)^{1/2}
            -
            \frac{d}{(b\log n)^2}\lfloor\log^{5/2}n\rfloor
        }
        \nonumber
        \\ &
        \le
        n^{-b+v+
            (b^{1/2}-db^{-2})\log^{-1/2}n
        }
        =
        n^{-b+v+o_n(1)},
        % \label{eq:lemma:4.7:upper}
    \end{align*}
    and hence
    \begin{equation*}
        \limsup_{n\to\infty}
        \frac{1}{\log n}
        \log\mathbb E\!\left[
            \sum_{
                |x|=\lfloor\log^{5/2}n\rfloor
            }
            \!\!\!
            F_x
        \right]
        \le
        -b+v.
    \end{equation*}
    For the lower bound, let for $|x|=\lfloor\log^{5/2}\rfloor$,
    \begin{equation*}
        \widetilde F_x
        =
        \widetilde F_x(
            b,v,v',n
        )
        :=
        {\bf 1}_{
            \{
                r+v'\log n\le\mathcal V(x_j)\le 0,\,
                \forall j\in[\![1,|x|]\!],\,
                \mathcal V(x)\le r+v\log n
            \}
        },
    \end{equation*}
    where $v'\in(0,v)$. Apply \eqref{eq:4.3} in \Cref{lemma:4.3} with $\Delta=-r-v'\log n=(b-v')\log n$, $r_0=r+v'\log n$, $c=0$, $\ell=\lfloor\log^{5/2}n\rfloor$ and $\frac{v-v'}{b-v'}$ in lieu of $v$ with any $v'\in(0,v)$, and we have that for any fixed $d'>\frac{\pi^2}2\sigma^2$, any $v'\in(0,v)$, some $n_2=n_2(b,v,v',d')>0$ and all $n\ge n_2$,
    \begin{align*}
        % &\quad
        \mathbb E\!\left[
            \sum_{
                |x|=\lfloor\log^{5/2}n\rfloor
            }
            \!\!\!
            F_x
        \right]
        % \\
        &
        \ge
        \mathbb E\!\left[
            \sum_{
                |x|=\lfloor\log^{5/2}n\rfloor
            }
            \!\!\!
            \widetilde F_x
        \right]
        \\ &
        =
        \mathbb E\!\left[
            e^{
                Y_{\lfloor\log^{5/2}n\rfloor}
            }
            {\bf 1}_{
                \{
                    r+v'\log n\le Y_j\le 0,\,
                    \forall j\in[\![1,\lfloor\log^{5/2}n\rfloor]\!],\,
                    Y_{\lfloor\log^{5/2}n\rfloor}\le r+v\log n
                \}
            }
        \right]
        \\ &
        \ge
        e^{r+v'\log n}
        \mathbb E\!\left[
            {\bf 1}_{
                \{
                    r+v'\log n\le Y_j\le 0,\,
                    \forall j\in[\![1,\lfloor\log^{5/2}n\rfloor]\!],\,
                    Y_{\lfloor\log^{5/2}n\rfloor}\le r+v\log n
                \}
            }
        \right]
        \\ &
        \ge
        n^{-b+v'}e^{
            -\frac{d'}{((b-v')\log n)^2}\lfloor\log^{5/2}n\rfloor
        }
        \ge
        n^{-b+v'}e^{
            -\frac{d'}{((b-v')\log n)^2}\log^{11/4}n
        }
        \\ &
        =
        n^{-b+v'-d'(b-v')^{-2}\log^{-1/4}n}
        =
        n^{-b+v'+o_n(1)},
    \end{align*}
    where the last equality $\log^{5/2}n\le\log^{11/4}n$ holds for some $0<n_3<n_2$ and all $n\ge n_3$. Hence,
    \begin{equation*}
        \liminf_{n\to\infty}
        \frac 1{\log n}
        \log\mathbb E\!\left[
            \sum_{
                |x|=\lfloor\log^{5/2}n\rfloor
            }
            \!\!\!
            F_x
        \right]
        \ge
        \sup_{v'\in(0,v)}
        (-b+v')
        =
        -b+v,
    \end{equation*}
    which completes the proof for \eqref{eq:4.11}. Moreover, substituting $\widetilde F_x$ for $F_x$ in the above proof yields
    \begin{equation}
        \mathbb E\!\left[
            \sum_{
                |x|=\lfloor\log^{5/2}n\rfloor
            }
            \!\!\!
            \widetilde F_x
        \right]
        =
        n^{-b+v+o_n(1)}.
        \label{eq:4.11:tilde}
    \end{equation}

    Then we proceed to prove \eqref{eq:4.9}, equivalently
    \begin{equation*}
        \lim_{n\to\infty}
        \frac{1}{\log n}
        \log\mathbb P\left(
            \sum_{
                |x|=\lfloor\log^{5/2}n\rfloor
            }
            \!\!\!
            F_x
            \ge
            1
        \right)
        =
        -b+v.
    \end{equation*}
    For the upper bound, we simply invoke the proved \eqref{eq:4.11}. Indeed,
    \begin{align*}
        &\quad
        \mathbb P\left(
            \sum_{
                |x|=\lfloor\log^{5/2}n\rfloor
            }
            \!\!\!
            F_x
            \ge
            1
        \right)
        \\
        &
        \le
        \mathbb E\!\left[
            \left(
                \sum_{
                        |x|=\lfloor\log^{5/2}n\rfloor
                    }
                    \!\!\!
                    F_x
            \right)
            {\bf 1}_{
                \{
                    \sum_{
                        |x|=\lfloor\log^{5/2}n\rfloor
                    }
                    F_x
                    \ge
                    1
                \}
            }
            +
            \left(
                \sum_{
                        |x|=\lfloor\log^{5/2}n\rfloor
                    }
                    \!\!\!
                    F_x
            \right)
            {\bf 1}_{
                \{
                    \sum_{
                        |x|=\lfloor\log^{5/2}n\rfloor
                    }
                    F_x
                    =
                    0
                \}
            }
        \right]
        \\ &
        =
        \mathbb E\!\left[
            \sum_{
                |x|=\lfloor\log^{5/2}n\rfloor
            }
            \!\!\!
            F_x
        \right].
    \end{align*}
    For the lower bound, we again utilize $\widetilde F_x$ for better bound for $F_x$ and proceed to apply the Paley--Zygmund inequality. For the second moment term in the inequality, apply \eqref{eq:4.6} in \Cref{lemma:4.5} with $\Delta=(b-v')\log n$, $r_0=r+v'\log n$, $c=0$, $\ell=\lfloor\log^{5/2}n\rfloor$, $\frac 12$ in lieu of $v$, and $\delta=r+v\log n$, and we have that for any fixed $0<d<\frac{\pi^2}2\sigma^2$, any $0<v'<v<b$, some $n_4=n_4(b,v,v',d)>0$ and all $n\ge n_4$, we calculate
    \begin{align*}
        e^{d\frac{m}{\Delta^2}-r_m}
        &
        \le
        e^{-r_m}
        =
        n^{b-v'},
        \\
        e^{\Delta^{1/2}}
        &
        =
        e^{
            ((b-v')\log n)^{1/2}
        }
        =
        n^{(b-v')^{1/2}\log^{-1/2}n}
        =
        n^{o_n(1)},
        \\
        \ell
        &
        =
        n^{\frac 1{\log n}\log\lfloor\log^{5/2}n\rfloor}
        \le
        n^{\frac 1{\log n}\log\log^{11/4}n}
        =
        n^{o_n(1)},
        \\
        e^{\delta-\frac{d}{\Delta^2}\ell}
        &
        =
        e^{
            \frac 1{\log n}\left(
                -b\log n+v\log n
                -\frac{d}{((b-v')\log n)^2}
                \lfloor\log^{5/2}n\rfloor
            \right)
        }
        =
        n^{-b+v+o_n(1)},
    \end{align*}
    and hence
    \begin{align*}
        \mathbb E\!\left[
            \left(
                \sum_{
                    |x|=\lfloor\log^{5/2}n\rfloor
                }
                \!\!\!
                \widetilde F_x
            \right)^2
        \right]
        &
        \le
        e^{4\Delta^{1/2}+2(\delta-\frac{d}{\Delta^2}(\ell-1))}
        \sum_{m=0}^{\ell-1}
        e^{d\frac{m}{\Delta^2}-r_m}
        +
        e^{\Delta^{1/2}+\delta-\frac{d}{\Delta^2}\ell}
        \\ &
        \le
        n^{4o_n(1)+2(-b+v+o_n(1))}
        n^{o_n(1)+b-v'}
        +
        n^{o_n(1)+(-b+v+o_n(1))}
        \\ &
        =
        n^{-b+2v-v'+o_n(1)}.
    \end{align*}
    Finally, we deduce from the Paley--Zygmund inequality, \eqref{eq:4.11:tilde} and the above inequality that for all $n\ge n_4$,
    \begin{align*}
        \mathbb P\left(
            \sum_{
                |x|=\lfloor\log^{5/2}n\rfloor
            }
            \!\!\!
            F_x
            \ge
            1
        \right)
        &
        \ge
        \mathbb P\left(
            \sum_{
                |x|=\lfloor\log^{5/2}n\rfloor
            }
            \!\!\!
            \widetilde F_x
            \ge
            1
        \right)
        \\ &
        \ge
        \frac{
            \mathbb E\!\left[
                \sum_{
                    |x|=\lfloor\log^{5/2}n\rfloor
                }
                \widetilde F_x
            \right]^2
        }{
            \mathbb E\!\left[
                \left(
                    \sum_{
                        |x|=\lfloor\log^{5/2}n\rfloor
                    }
                    \widetilde F_x
                \right)^2
            \right]
        }
        \ge
        \frac{
            (n^{-b+v+o_n(1)})^2
        }{
            n^{-b+2v-v'+o_n(1)}
        }
        =
        n^{-b+v'+o_n(1)}
        ,
    \end{align*}
    and hence
    \begin{equation*}
        \liminf_{n\to\infty}
        \frac{1}{\log n}
        \log\mathbb P\left(
            \sum_{
                |x|=\lfloor\log^{5/2}n\rfloor
            }
            \!\!\!
            F_x
            \ge
            1
        \right)
        \ge
        \sup_{v'\in(0,v)}
        (-b+v')
        =
        -b+v
        ,
    \end{equation*}
    which completes the proof for \eqref{eq:4.9}.

    Finally, we prove \eqref{eq:4.10}. Denote the first exit time of $[r,0]$ by $Y$ as
    \begin{equation*}
        \tau
        :=
        \inf\left\{
            j\ge 1: Y_j\notin[r,0]
        \right\}.
    \end{equation*}
    By the many-to-one formula, we have that for any $\ell\ge 1$,
    \begin{align*}
        \mathbb E\!\left[
            \sum_{|x|=\ell}
            {\bf 1}_{
                \{
                    r\le\mathcal V(x_j)\le 0,\,
                    \forall j\in[\![1,\ell]\!]
                \}
            }
        \right]
        &
        =
        \mathbb E\!\left[
            e^{Y_\ell}
            {\bf 1}_{
                \{
                    r\le Y_j\le 0,\,
                    \forall j\in[\![1,\ell]\!]
                \}
            }
        \right]
        \\ &
        \le
        e^0
        \mathbb E\!\left[
            {\bf 1}_{
                \{
                    r\le Y_j\le 0,\,
                    \forall j\in[\![1,\ell]\!]
                \}
            }
        \right]
        % \\ &
        =
        \mathbb P(\tau>\ell).
    \end{align*}
    Hence,
    \begin{align*}
        \mathbb E\!\left[
            \sum_{\ell=1}^\infty
            \sum_{|x|=\ell}
            {\bf 1}_{
                \{
                    r\le\mathcal V(x_j)\le 0,\,
                    \forall j\in[\![1,\ell]\!]
                \}
            }
        \right]
        &
        =
        \sum_{\ell=1}^\infty
        \mathbb E\!\left[
            \sum_{|x|=\ell}
            {\bf 1}_{
                \{
                    r\le\mathcal V(x_j)\le 0,\,
                    \forall j\in[\![1,\ell]\!]
                \}
            }
        \right]
        \\ &
        \le
        \sum_{\ell=1}^\infty
        \mathbb P(\tau>\ell)
        =
        \mathbb E[\tau]
        -
        \mathbb P(\tau>0)
        % \\ &
        \le
        \mathbb E[\tau],
    \end{align*}
    and the result follows from the {\color{blue}standard property}
    $
        \mathbb E[\tau]
        =
        {\color{blue}
            O_r(|r|)
        }
        =
        O_n(\log n)
    $.
\end{proof}
% !TeX root = ../main.tex
% !TeX root = ../../main.tex
% !TEX root = ../main.tex

\section{\bilingual{Lemmas on \texorpdfstring{$N$}{N}-BRW}{有关\texorpdfstring{$N$}{N}-分支随机游走的引理}}\label{section:lemmas-on-n-brw}

\bilingual{
    We use the notation $\mathcal C$ to represent the set of all finite counting measures on $\mathbb R$. For $\mu\in\mathcal C$, we write $M(\mu)$ for the total mass of $\mu$, i.e., the number of particles.
}{
    记 $\mathcal C$ 为实数轴上所有有限计数测度的集合.~ 对于 $\mu\in\mathcal C$,~ 记 $M(\mu)$ 为 $\mu$ 的总质量,~ 即其粒子数.
}

\bilingual{
    We also write $\mu_1\succcurlyeq\mu_2$ if $M(\mu_1)\le M(\mu_2)$ and the atoms of $\mu_1$ are greater than or equal to those of $\mu_2$ in the sense that if $\mu_1=\sum_{i=1}^{M(\mu_1)}\delta_{x_i}$ and $\mu_2=\sum_{i=1}^{M(\mu_2)}\delta_{y_i}$ with $x_1\le x_2\le\cdots\le x_{M(\mu_1)}$ and $y_1\le y_2\le\cdots\le y_{M(\mu_2)}$, then $x_i\ge y_i$ for all $i\in[\![1,M(\mu_1)]\!]$. In the degenerate case where $M(\mu_1)=0$, we have $\mu_1\succcurlyeq\mu_2$ for any $\mu_2$.
}{
    称$\mu_1$占优于$\mu_2$,~ 记作$\mu_1\succcurlyeq\mu_2$,~ 若 $\mu_1$的粒子数不多于$\mu_2$,~ 即$M(\mu_1)\le M(\mu_2)$,~ 且 $\mu_1$ 的各原子位置依次大于等于 $\mu_2$ 靠左侧的各原子位置,~ 即,~ 若记 $$\mu_1=\sum_{i=1}^{M(\mu_1)}\delta_{x_i},\quad \mu_2=\sum_{i=1}^{M(\mu_2)}\delta_{y_i},$$ 其中 $x_1\le x_2\le\cdots\le x_{M(\mu_1)}$ 且 $y_1\le y_2\le\cdots\le y_{M(\mu_2)}$,~ 则 $x_i\ge y_i$ 对所有 $i\in[\![1,M(\mu_1)]\!]$ 成立.~ 在退化的情形 $M(\mu_1)=0$,~ 对任何 $\mu_2$,~ 都有 $\mu_1\succcurlyeq\mu_2$.
}

\bilingual{
    For all $N\ge 1$, let $\mathcal C(N)$ be the set of finite counting measures on $\mathbb R$ with total mass equal to $N$. We identify $\mathcal C(N)$ with the family of sets of real numbers with $N$ elements, and we endow $\mathcal C(N)$ with the partial order $\succcurlyeq$ defined above. For example, the population of particles $\mathcal Z^N_n$ at time $n$ in the $N$-BRW $\mathcal V^N$ can be identified with an element of $\mathcal C(N)$, and we write $\mathcal Z^N_n\in\mathcal C(N)$ for all $n\ge 0$. Recall that we may also write $\mathcal V^N_n$ in lieu of $\mathcal Z^N_n$ to represent the population of particles at time $n$ in $\mathcal V^N$.
}{
    对任一 $N\ge 1$,~ 记 $\mathcal C(N)$ 为实数轴上总质量为 $N$ 的有限计数测度全体组成的集合.~ 并且,~ 对于$\mathcal C(N)$中的某一测度,~ 它可被视作为元素个数为 $N$ 的实数集合.~ 例如,~ 在 $N$-分支随机游走 $\mathcal V^N$ 的术语中,~ 时刻 $n$ 的粒子全体 $\mathcal Z^N_n$ 可视作为 $\mathcal C(N)$ 中的一个测度,~ 此时可记 $\mathcal Z^N_n\in\mathcal C(N)$.~ 回顾上文,~ 也可以用 $\mathcal V^N_n$ 来表示时刻 $n$ 的粒子全体 $\mathcal Z^N_n$.~ 另外,~ 在 $\mathcal C(N)$ 上也可类似定义上述占优关系 $\succcurlyeq$.
}

\bilingual{
    We also denote by $\max\mu$ and $\min\mu$ the maximum and minimum of the support of $\mu\in\mathcal C$, respectively. For example, the leftmost position $\mathcal M^N_n$ at time $n$ in the $N$-BRW $\mathcal V^N$ can be identified with $\min\mathcal Z^N_n$.
}{
    对 $\mu\in\mathcal C$,~ 记 $\max\mu$ 和 $\min\mu$ 分别为 $\mu$ 的支撑集的最大值和最小值.~ 例如,~ 在 $N$-分支随机游走 $\mathcal V^N$ 的术语中,~ 时刻 $n$ 的最左侧位置 $\mathcal M^N_n$ 可视作为 $\min\mathcal Z^N_n$.
}

% !TeX root = ../../main.tex
\subsection{\bilingual{Coupling of two \texorpdfstring{$N$}{N}-BRWs}{两个\texorpdfstring{$N$}{N}-分支随机游走的耦合}}

\bilingual{
    This subsection extends the coupling result of B\'erard and Gour\'er\'e~\mycite{bérard2010brunetderridabehaviorbranchingselectionparticle} for two binary $N$-BRWs to the case of two (or more) infinite-branching $N$-BRWs.
}{
    本小节把B\'erard与Gour\'er\'e~\mycite{bérard2010brunetderridabehaviorbranchingselectionparticle}关于两个二分支$N$-分支随机游走的耦合的结论推广到两个 (乃至多个) 无限分支$N$-分支随机游走的情形.
}

\begin{lemma}[\bilingual{Monotonicity}{单调性}]
    \label{lemma:1}
    \bilingual{
        For all $1\le N_1\le N_2$, and $\mu_1\in\mathcal C(N_1)$, $\mu_2\in\mathcal C(N_2)$ such that $\mu_1\succcurlyeq\mu_2$ (holds almost surely if $\mu_1,\mu_2$ are random), there exists a pair of random variables $(Z^1,Z^2)$ taking values in $\mathcal C(N_1)\times\mathcal C(N_2)$ such that
    }{
        对任一 $1\le N_1\le N_2$, 以及 $\mu_1\in\mathcal C(N_1)$, $\mu_2\in\mathcal C(N_2)$ 满足 $\mu_1\succcurlyeq\mu_2$ (若 $\mu_1,\mu_2$ 是随机的, 则此处指该关系几乎必然成立), 存在一对随机变量 $(Z^1,Z^2)$ 取值于 $\mathcal C(N_1)\times\mathcal C(N_2)$, 使得以下条件成立:
    }
    \begin{itemize}
        \item \bilingual{
            the distribution of $Z^i$ for $i=1,2$ is that of the population of particles obtained by performing one branching-selection step with branching distribution $\Xi$ and from the population $\mu_i$ of $N_i$ initial particles;
        }{
            $Z^i$的分布为从初始粒子群 $\mu_i$ 出发, 按照分支律 $\Xi$ 进行一次分支-选择步骤后得到的粒子群的分布, 其中 $i=1,2$.
        }
        \item \bilingual{
                $Z^1\succcurlyeq Z^2$ almost surely.
            }{
                $Z^1\succcurlyeq Z^2$ 几乎必然成立.
        }
    \end{itemize}
\end{lemma}
\vspace{-1.8\baselineskip}
\begin{proof}
    \quad

    \bilingual{
        Consider an i.i.d.~family on the same probability space of random variables $(\Xi_{i,j})$ for $i\in[\![1,N_2]\!]$ and $j\ge 1$ with common distribution $\Xi$. For $k=1,2$, write $\mu_k=\sum_{i=1}^{N_k}\delta_{x_i^k}$ with $x_1^k\le x_2^k\le\cdots\le x_{N_k}^k$. Let
    }{
        考虑在同一概率空间上具有共同分布 $\Xi$ 的独立同分布随机变量族 $(\Xi_{i,j})_i$ 其中 $i\in[\![1,N_2]\!]$ 且 $j\ge 1$. 对 $k=1,2$, 记 $\mu_k=\sum_{i=1}^{N_k}\delta_{x_i^k}$ 其中 $x_1^k\le x_2^k\le\cdots\le x_{N_k}^k$. 定义
    }
    $$
        T^k
        :=
        \sum_{i=1}^{N_k}
        \sum_{j=1}^{\infty}
        \delta_{x_i^k+\Xi_{i,j}}
    ,$$
    \bilingual{
        and define $Z^k$ as being performed by the $N_k$ leftmost particles in $T^k$.
    }{
        并且定义 $Z^k$ 为从 $T^k$ 中选取最左侧的 $N_k$ 个粒子组成的测度.
    }

    \bilingual{
        From the assumption $\mu_1\succcurlyeq\mu_2$, we have $x_i^1\ge x_i^2$ for all $i\in[\![1,N_1]\!]$, whence
    }{
        根据假设 $\mu_1\succcurlyeq\mu_2$, 对所有 $i\in[\![1,N_1]\!]$, 都有 $x_i^1\ge x_i^2$, 从而
    }
    \begin{equation*}
        x_i^1+\Xi_{i,j}\ge x_i^2+\Xi_{i,j},\quad i\in[\![1,N_1]\!],\ j\ge 1.
    \end{equation*}
    \bilingual{
        It follows that $T^1\succcurlyeq T^2$, and thus $Z^1\succcurlyeq Z^2$ by definition.
    }{
        从而 $T^1\succcurlyeq T^2$, 因此根据定义 $Z^1\succcurlyeq Z^2$.
    }
\end{proof}

\begin{figure}[H]
    \centering
    \begin{tikzpicture}[
        xscale=1,
        yscale=1.5,
        >=Stealth,
        % 样式定义
        node-both/.style={circle, fill=blue!70!black, draw=green!60!black, line width=1.2pt, inner sep=1.2pt},
        node-N/.style={circle, fill=blue!70!black, inner sep=1.8pt},
        node-k/.style={circle, fill=green!60!black, inner sep=1.8pt},
        node-N-dead/.style={circle, draw=blue!70!black, fill=white, inner sep=1.4pt, cross out, draw=blue!80!black, line width=1.2pt},
        node-k-dead/.style={circle, draw=green!60!black, fill=white, inner sep=1.4pt, cross out, draw=green!60!black, line width=1.2pt},
        node-dead/.style={circle, draw=black!40, fill=white, inner sep=1.4pt},
        % 连线样式
        edge-N/.style={thick, blue!70!black},
        edge-k/.style={thick, green!60!black},
        edge-N-dashed/.style={dashed, blue!70!black},
        edge-k-dashed/.style={dashed, green!60!black},
        edge-both/.style={
            thick,
            blue!70!black, % 底色：蓝色
            postaction={
                draw,
                green!60!black,
                % on 2pt 代表绿色段长度，off 2pt 代表空隙（露出底下的蓝色）
                % 若要更频繁，将 2pt 改为 1pt
                dash pattern=on 5pt off 5pt
            }
        },
        edge-dashed/.style={dashed, black!30},
        % 灭杀边界
        kill-line/.style={red!80!black, dotted, ultra thick}
    ]

        % --- 坐标轴与背景 ---
        \foreach \y in {0,1} \draw[-, black!20, dashed] (-4.5, \y) -- (2.5, \y);
        \node[right=10pt] at (2.7, 0) {$n=0$};
        \node[right=10pt] at (2.7, 1) {$n=1$};

        % ================= n=0 =================
        \node[node-k] (n0-k-1) at (-1.6, 0) {};
        \node[below] at (n0-k-1) {{\color{green!60!black}$\Xi_1$}};
        \node[node-k] (n0-k-2) at (2.2, 0) {};
        \node[below] at (n0-k-2) {{\color{green!60!black}$\Xi_2$}};
        \node[node-N] (n0-N-1) at (-2.2, 0) {};
        \node[below] at (n0-N-1) {{\color{blue!70!black}$\Xi_1$}};
        \node[node-N] (n0-N-2) at (-0.8, 0) {};
        \node[below] at (n0-N-2) {{\color{blue!70!black}$\Xi_2$}};
        \node[node-N] (n0-N-3) at (-0.2, 0) {};
        \node[below] at (n0-N-3) {{\color{blue!70!black}$\Xi_3$}};

        % ================= n=1 =================

        \node[node-k] (n1-k-11) at (-3.0, 1) {};
        \draw[edge-k] (n0-k-1) -- (n1-k-11);
        \node[node-k] (n1-k-12) at (-2.1, 1) {};
        \draw[edge-k] (n0-k-1) -- (n1-k-12);
        \node[node-k-dead] (n1-k-13) at (-1.3, 1) {};
        \draw[edge-k-dashed] (n0-k-1) -- (n1-k-13);

        \node[node-N] (n1-N-11) at (-3.6, 1) {};
        \draw[edge-N] (n0-N-1) -- (n1-N-11);
        \node[node-N] (n1-N-12) at (-2.7, 1) {};
        \draw[edge-N] (n0-N-1) -- (n1-N-12);
        \node[node-N-dead] (n1-N-13) at (-1.9, 1) {};
        \draw[edge-N-dashed] (n0-N-1) -- (n1-N-13);

        \node[node-N] (n1-N-21) at (-4.0, 1) {};
        \draw[edge-N] (n0-N-2) -- (n1-N-21);
        \node[node-k-dead] (n1-k-21) at (-1.0, 1) {};
        \draw[edge-k-dashed] (n0-k-2) -- (n1-k-21);

        \node[node-N-dead] (n1-N-31) at (-0.4, 1) {};
        \draw[edge-N-dashed] (n0-N-3) -- (n1-N-31);
        \node[node-N-dead] (n1-N-32) at (0.8, 1) {};
        \draw[edge-N-dashed] (n0-N-3) -- (n1-N-32);
        \node[node-N-dead] (n1-N-33) at (1.9, 1) {};
        \draw[edge-N-dashed] (n0-N-3) -- (n1-N-33);

    \end{tikzpicture}
    \caption{\bilingual{Illustration of a \texorpdfstring{$3$}{3}-BRW (blue) and a dominating \texorpdfstring{$2$}{2}-BRW (green) coupling}{一个\texorpdfstring{$3$-BRW}{3-BRW} (蓝色) 与一个占优的\texorpdfstring{$2$-BRW}{2-BRW} (绿色) 的耦合示例图}}
\end{figure}

\begin{corollary}
    \label{corollary:2}
    \bilingual{
        For all $1\le N_1\le N_2$, and $\mu_1\in\mathcal C(N_1)$, $\mu_2\in\mathcal C(N_2)$ such that $\mu_1\succcurlyeq\mu_2$ (holds almost surely if $\mu_1,\mu_2$ are random), there exists a coupling $(Z^1_n,Z^2_n)_{n\ge 0}$ of $\mathcal V^{N_1}$ and $\mathcal V^{N_2}$ (on possibly different probability space) such that $Z^1_0=\mu_1$, $Z^2_0=\mu_2$, and $Z^1_n\succcurlyeq Z^2_n$ almost surely for all $n\ge 0$. In particular, we have $\min Z^1_n\ge \min Z^2_n$ almost surely for all $n\ge 0$.
    }{
        对任一 $1\le N_1\le N_2$, 以及 $\mu_1\in\mathcal C(N_1)$, $\mu_2\in\mathcal C(N_2)$ 满足 $\mu_1\succcurlyeq\mu_2$ (若 $\mu_1,\mu_2$ 是随机的, 则此处指该关系几乎必然成立), 存在 (可能定义于不同的概率空间的) $\mathcal V^{N_1}$ 和 $\mathcal V^{N_2}$ 的一个耦合 $(Z^1_n,Z^2_n)_{n\ge 0}$ 成立 $Z^1_0=\mu_1$, $Z^2_0=\mu_2$, 且 $Z^1_n\succcurlyeq Z^2_n$ 几乎对所有 $n\ge 0$必然成立. 特别地, 对所有 $n\ge 0$, 都有 $\min Z^1_n\ge \min Z^2_n$ 几乎必然成立.
    }
\end{corollary}

\bilingual{
    Note that the above coupling can be extended to more than two $N$-BRWs, since we can use the same family of random variables $(\Xi_{i,j})$ for all the $N$-BRWs involved, and thus can be defined on the same probability space. In particular, we have the following translation invariance property of $N$-BRW.
}{
    注意到上述耦合可以推广到两个以上的 $N$-分支随机游走, 因为可以对所有涉及的 $N$-分支随机游走使用同一族随机变量 $(\Xi_{i,j})$, 从而定义于同一概率空间上. 特别地, 我们有以下 $N$-分支随机游走的平移不变性.
}

\begin{corollary}[\bilingual{Translation Invariance}{平移不变性}]
    \label{corollary:translation}
    \bilingual{
        For all $N\ge 1$, $a\in\mathbb R$ and $\mu=N\delta_a\in\mathcal C(N)$, there exists a coupling $(Z^1_n,Z^2_n)_{n\ge 0}$ of two $N$-BRW with initial population $N\delta_0$ and $\mu$ respectively such that $Z^1_0=N\delta_0$, $Z^2_0=\mu$, and $Z^2_n=a+Z^1_n$ almost surely for all $n\ge 0$, where $a+Z^1_n$ denotes the measure obtained by shifting all atoms of $Z^1_n$ by $a$.
    }{
        对任一 $N\ge 1$, $a\in\mathbb R$ 和 $\mu=N\delta_a\in\mathcal C(N)$, 对分别以$N\delta_0$ 与 $\mu$ 为初始粒子分布的两个 $N$-分支随机游走, 存在它们的一个耦合 $(Z^1_n,Z^2_n)_{n\ge 0}$ 使得 $Z^1_0=N\delta_0$, $Z^2_0=\mu$, 且 $Z^2_n=a+Z^1_n$ 对所有 $n\ge 0$几乎必然成立, 其中 $a+Z^1_n$ 表示通过将 $Z^1_n$ 的所有原子位置平移 $a$ 得到的测度.
    }
\end{corollary}

\bilingual{
    Similarly, the couplings in the above results can be extended to more than two $N$-BRWs with combined relations of both translation and ordering.
}{
    类似地, 上述结果中的耦合可以推广到两个以上的 $N$-分支随机游走, 结论可以包括多个平移与占优关系的组合.
}

% !TeX root = ../../main.tex
\subsection{\bilingual{Existence for Leftmost Asymptotic Speed in \texorpdfstring{$N$}{N}-BRW}{\texorpdfstring{$N$}{N}-分支随机游走最左侧渐近速率的存在性}}

\bilingual{
    This subsection extends the result of B\'erard and Gour\'er\'e~\mycite{bérard2010brunetderridabehaviorbranchingselectionparticle} on the existence of the leftmost asymptotic speed for binary $N$-BRW to the case of infinite-branching $N$-BRW.
}{
    本小节把B\'erard与Gour\'er\'e~\mycite{bérard2010brunetderridabehaviorbranchingselectionparticle}关于二分支$N$-分支随机游走的渐近速率存在性的结论推广到无限分支$N$-分支随机游走的情形.
}

\begin{lemma}
    \label{proposition:2}
    \bilingual{
        Under \Cref{assumption:psi-one,assumption:moment-1} (i.e. $\psi(1)=0$ and $\mathbb E[(\Xi_1)_+]<\infty$), there exists $v_N=v_N(\Xi)\in\mathbb R$ such that, with probability one, and in $\mathrm L^1(\mathbb P)$, for the $N$-BRW $\mathcal V^N$ with initial population $\mathcal Z^N_0=N\delta_0$,
    }{
        在\Cref{assumption:psi-one,assumption:moment-1} (即 $\psi(1)=0$ 且 $\mathbb E[(\Xi_1)_+]<\infty$) 的条件下, 存在 $v_N=v_N(\Xi)\in\mathbb R$ 使得对于以 $\mathcal Z^N_0=N\delta_0$ 为初始粒子分布的 $N$-分支随机游走 $\mathcal V^N$, 几乎必然成立, 且在 $\mathrm L^1(\mathbb P)$ 意义下成立,
    }
    {
        \setlength{\belowdisplayskip}{-10pt}
        \begin{equation*}
            \lim_{n\to\infty}
            \frac 1n
            \mathcal M^N_n
            =
            v_N.
        \end{equation*}
    }
\end{lemma}

\begin{proof}
    \quad

    \bilingual{
        We are to apply Kingman's subadditive ergodic theorem. Consider an i.i.d.~family on the same probability space of random variables $(\Xi_{\ell,i,j})$ for $\ell\ge 0$, $i\in[\![1,N]\!]$ and $j\ge 1$ so that $\{(\Xi_{\ell,i,j})_{j\ge 1}\}_{\ell,i}$ have common distribution $\Xi$. The index $\ell$ will be used to shift the first branch from $\Xi_{0,\cdot}$ to $\Xi_{\ell,\cdot}$. Define $(W_{\ell,k})_{k\ge 0}$ for $\ell\ge 0$ as follows.
    }{
        Kingman次可加遍历定理将被应用以证明本引理. 考虑在同一概率空间上具有共同分布 $\Xi$ 的独立同分布随机变量族 $(\Xi_{\ell,i,j})_{\ell,i}$ 其中 $\ell\ge 0$, $i\in[\![1,N]\!]$ 且 $j\ge 1$. 索引 $\ell$ 将被用来把第$\ell$个对象的第一次分支的分支律从 $\Xi_{0,\cdot}$ 沿时间轴正方向移动到 $\Xi_{\ell,\cdot}$. 对 $\ell\ge 0$ 定义 $(W_{\ell,k})_{k\ge 0}$ 如下.
    }
    \begin{itemize}
        \item
        \bilingual{
            For $k=0$, let $W_{\ell,0}=N\delta_0$ for all $\ell\ge 0$.
        }{
            对 $k=0$, 对所有 $\ell\ge 0$ 记 $W_{\ell,0}=N\delta_0$.
        }
        \item
        \bilingual{
            For $k\ge 1$, suppose $W_{\ell,k-1}$ has been defined for all $\ell\ge 0$. We write
        }{
            对 $k\ge 1$, 假设 $W_{\ell,k-1}$ 已对所有 $\ell\ge 0$ 定义. 记
        }
        \begin{equation*}
            W_{\ell,k-1}=\sum_{i=1}^N\delta_{x_i},
        \end{equation*}
        \bilingual{
            with $x_1\le x_2\le\cdots\le x_N$. The population $T_{\ell,k}$ derived from $W_{\ell,k-1}$ at time $k$ by branching is
        }{
            其中 $x_1\le x_2\le\cdots\le x_N$. 由 $W_{\ell,k-1}$ 在时刻 $k$ 通过分支得到的粒子总数 $T_{\ell,k}$ 为
        }
        \begin{equation*}
            T_{\ell,k}
            :=
            \sum_{i=1}^N
            \sum_{j=1}^\infty
            \delta_{x_i+\Xi_{\ell+k,i,j}}.
        \end{equation*}
        \bilingual{
            Then $W_{\ell,k}$ is defined as being performed by the $N$ leftmost particles in $T_{\ell,k}$.
        }{
            然后 $W_{\ell,k}$ 定义为从 $T_{\ell,k}$ 中选取最左侧的 $N$ 个粒子组成的测度.
        }
    \end{itemize}
    \bilingual{
        We apply an improved version (Liggett \mycite{liggett1985improved}) of Kingman's subadditive ergodic theorem to the sequence of random variables
    }{
        Kingman次可加遍历定理的一个改进版本 (Liggett \mycite{liggett1985improved}) 将被应用于以下随机变量序列:
    }
    \begin{equation*}
        X_{m,n}:=-\min W_{m,n-m},\quad 0\le m<n.
    \end{equation*}
    \bilingual{
        We verify the assumptions of the theorem as follows.
    }{
        以下是对定理假设的验证.
    }
    \begin{itemize}
        \item (\bilingual{Subadditivity}{次可加性}).
        \bilingual{
            $X_{0,n}\le X_{0,m}+X_{m,n}$ for all $0\le m<n$ by the definition of the branching-selection procedure. Indeed,
        }{
            对所有 $0\le m<n$, 都有 $X_{0,n}\le X_{0,m}+X_{m,n}$. 事实上, 根据分支-选择步骤的定义,
        }
        \begin{equation*}
            \min W_{0,n}
            \ge
            \min W_{0,m}+\min W_{m,n-m}
        \end{equation*}
        \bilingual{
            holds since it compares the minimum of two populations of particles at time $n$ obtained by performing the branching-selection procedure specified by \Cref{corollary:2} starting from the initial population distribution at time $m$ of $W_{0,m}$ and $N\delta_{\min W_{0,m}}$, respectively.
        }{
            成立, 此因不等号两侧比较了两种按照\Cref{corollary:2} 具体说明的分支-选择步骤所得到的时刻 $n$ 的粒子总数的最小值, 它们在时刻 $m$ 初始粒子分布分别为 $W_{0,m}$ 和 $N\delta_{\min W_{0,m}}$.
        }
        \item (\bilingual{Stationarity}{平稳性}).
        \bilingual{
            $(X_{nk,(n+1)k})_{n\ge 1}$ is a stationary sequence for each $k\ge 1$. Indeed, $X_{nk,(n+1)k}=-\min W_{nk,k}$. By definition of $W$, for each $k\ge 1$, the sequence $(W_{nk,k})_{n\ge 1}$ is i.i.d., since, for any $1\le n_1<n_2$, we have the fact that $n_1k+k-1<n_2k$ and that $W_{n_1k,k}$ is constructed from the point processes
        }{
            对任一 $k\ge 1$, $(X_{nk,(n+1)k})_{n\ge 1}$ 是一个平稳序列. 事实上, 回顾 $X$ 的定义有 $X_{nk,(n+1)k}=-\min W_{nk,k}$. 根据 $W$ 的定义, 对每个 $k\ge 1$, 序列 $(W_{nk,k})_{n\ge 1}$ 是独立同分布的, 此因对于任意 $1\le n_1<n_2$, 都有 $n_1k+k-1<n_2k$ 且 $W_{n_1k,k}$ 是由如下点过程构建的
        }
        $$
            \left\{
                \left(\Xi_{n_1k+l,i,j}\right)_{j\ge 1}:
                l\in[\![0,k-1]\!],\,
                i\in[\![0,N]\!]
            \right\},
        $$
        \bilingual{
            whereas $W_{n_2k,k}$ is constructed in the same way from the point processes
        }{
            然而 $W_{n_2k,k}$ 是以同样的方式由如下点过程构建的
        }
        $$
            \left\{
                \left(\Xi_{n_2k+l,i,j}\right)_{j\ge 1}:
                l\in[\![0,k-1]\!],\,
                i\in[\![0,N]\!]
            \right\},
        $$
        \bilingual{
            and these two sets of point processes are independently identically distributed. Hence, the sequence $(X_{nk,(n+1)k})_{n\ge 1}$ is also i.i.d. The result then follows from the fact that i.i.d.~sequences are stationary.
        }{
            且这两组点过程是独立同分布的. 从而, 序列 $(X_{nk,(n+1)k})_{n\ge 1}$ 也是独立同分布的. 结果由独立同分布序列是平稳的这一事实得出.
        }
        \item (\bilingual{Ergodicity}{遍历性}).
        \bilingual{
            Moreover, $(X_{nk,(n+1)k})_{n\ge 1}$ is an ergodic sequence for each $k\ge 1$. Indeed, since the sequence is i.i.d. as stated in the last item, the result follows from the fact that i.i.d.~sequences are ergodic.
        }{
            进一步, 对任一 $k\ge 1$, $(X_{nk,(n+1)k})_{n\ge 1}$ 是一个遍历序列. 事实上, 因为该序列是独立同分布的 (如上一项所述), 结果由独立同分布序列是遍历的这一事实得出.
        }
        \item (\bilingual{Distribution invariance}{分布不变性}).
        \bilingual{
            The distribution of $\{X_{m,m+k}:k\ge 1\}$ does not depend on $m$. Indeed, $X_{m,m+k}=-\min W_{m,k}$, and by definition of $W$, for any $0\le m_1<m_2$, $W_{m_1,k}$ and $W_{m_2,k}$ are identically distributed, since $W_{m_1,k}$ is constructed from the point processes
        }{
            $\{X_{m,m+k}:k\ge 1\}$ 的分布不依赖于 $m$. 事实上, 回顾 $X$ 的定义有 $X_{m,m+k}=-\min W_{m,k}$, 根据 $W$ 的定义, 对任意 $0\le m_1<m_2$, $W_{m_1,k}$ 和 $W_{m_2,k}$ 是同分布的, 此因 $W_{m_1,k}$ 是由如下点过程构建的
        }
        $$
            \left\{
                \left(\Xi_{m_1+l,i,j}\right)_{j\ge 1}:
                l\in[\![0,k-1]\!],\,
                i\in[\![0,N]\!]
            \right\},
        $$
        \bilingual{
            while $W_{m_2,k}$ is constructed in the same way from the point processes
        }{
            同时 $W_{m_2,k}$ 是以同样的方式由如下点过程构建的
        }
        $$
            \left\{
                \left(\Xi_{m_2+l,i,j}\right)_{j\ge 1}:
                l\in[\![0,k-1]\!],\,
                i\in[\![0,N]\!]
            \right\},
        $$
        \bilingual{
            and these two sets of point processes are identically distributed. Hence, the distribution of $\{X_{m,m+k}:k\ge 1\}$ does not depend on $m$.
        }{
            且这两组点过程是同分布的. 从而, $\{X_{m,m+k}:k\ge 1\}$ 的分布不依赖于 $m$.
        }
        \item (\bilingual{Finite expectation, Part 1}{有限期望, 部分 1}).
        \bilingual{
            For each $n\ge 1$, $\mathbb E[X_{0,n}]\ge cn$ for some
            $c> -\infty$. Follow recursively one currently selected particle
            and, before the next selection, its leftmost child. If $R_k$
            denotes the displacement used at step $k$, then
            $(R_k)_{k\ge0}$ are independent copies of $\Xi_1$ and selection
            can only move the population minimum further to the left.
            Consequently,
        }{
            对每个 $n\ge 1$, $\mathbb E[X_{0,n}]\ge cn$ 对某个
            $c>-\infty$ 成立. 递归地选取当前一个被保留的粒子, 并在下一次选择前追踪其最左子代. 若第 $k$ 步采用的位移记为 $R_k$, 则
            $(R_k)_{k\ge0}$ 是 $\Xi_1$ 的独立同分布副本, 且选择操作只可能使种群最小位置进一步向左. 因此
        }
        \begin{equation*}
            \min W_{0,n}\le\sum_{k=0}^{n-1}R_k,
            \qquad
            X_{0,n}\ge-\sum_{k=0}^{n-1}R_k.
        \end{equation*}
        \bilingual{
            Under \Cref{assumption:psi-one},
            $\mathbb E[(\Xi_1)_-]<\infty$ by
            \Cref{lemma:moment-1:neg}; together with
            \Cref{assumption:moment-1}, $\Xi_1$ is integrable. Hence
            $\mathbb E[X_{0,n}]\ge-n\mathbb E[\Xi_1]$.
        }{
            在\Cref{assumption:psi-one} 下, 由
            \Cref{lemma:moment-1:neg} 有
            $\mathbb E[(\Xi_1)_-]<\infty$; 再结合
            \Cref{assumption:moment-1}, $\Xi_1$ 可积. 因而
            $\mathbb E[X_{0,n}]\ge-n\mathbb E[\Xi_1]$.
        }
        \item (\bilingual{Finite expectation, Part 2}{有限期望, 部分 2}).
        \bilingual{
            For each $n\ge1$, $\mathbb E[|X_{0,n}|]<\infty$.
            Subadditivity gives
        }{
            对每个 $n\ge1$, $\mathbb E[|X_{0,n}|]<\infty$. 次可加性给出
        }
        \begin{equation*}
            (X_{0,n})_+
            \le
            \sum_{k=1}^n(X_{k-1,k})_+.
        \end{equation*}
        \bilingual{
            The designated-line inequality from Part 1 gives the other side,
        }{
            部分 1 中的指定谱系不等式给出另一侧,
        }
        \begin{equation*}
            (X_{0,n})_-
            \le
            \left(\sum_{k=0}^{n-1}R_k\right)_+
            \le
            \sum_{k=0}^{n-1}(R_k)_+.
        \end{equation*}
        \bilingual{
            For the one-step positive part,
        }{
            对于一步正部,
        }
        \begin{align*}
            (X_{0,1})_+
            &=\left(\min_{i\in[\![1,N]\!]}\Xi_{0,i,1}\right)_-
            =\max_{i\in[\![1,N]\!]}(\Xi_{0,i,1})_-
            \le\sum_{i=1}^N(\Xi_{0,i,1})_-.
        \end{align*}
        \bilingual{
            Thus $\mathbb E[(X_{0,1})_+]\le
            N\mathbb E[(\Xi_1)_-]<\infty$ by
            \Cref{lemma:moment-1:neg}, while
            $\mathbb E[(X_{0,n})_-]\le
            n\mathbb E[(\Xi_1)_+]<\infty$ by
            \Cref{assumption:moment-1}. Stationarity and the positive-part
            bound now prove $\mathbb E[|X_{0,n}|]<\infty$. This argument does
            not use an expectation identity before integrability is known.
        }{
            因此由\Cref{lemma:moment-1:neg} 有
            $\mathbb E[(X_{0,1})_+]\le
            N\mathbb E[(\Xi_1)_-]<\infty$; 同时由
            \Cref{assumption:moment-1} 有
            $\mathbb E[(X_{0,n})_-]\le
            n\mathbb E[(\Xi_1)_+]<\infty$. 再结合平稳性与正部上界, 即得
            $\mathbb E[|X_{0,n}|]<\infty$. 此论证没有在可积性得到证明之前使用期望恒等式.
        }
    \end{itemize}
    \bilingual{
        We conclude that the hypotheses of Kingman's subadditive ergodic theorem are satisfied, and thus
    }{
        定理的假设已被验证, 因此 Kingman 次可加遍历定理适用, 从而
    }
    \begin{equation*}
        -v_N
        :=
        \lim_{n\to\infty}
        \frac 1n
        X_{0,n}
        =
        -\lim_{n\to\infty}
        \frac 1n
        \min W_{0,n}
    \end{equation*}
    \bilingual{
        exists both almost surely and in $\mathrm L^1(\mathbb P)$, and is constant. By the definition of $W$, $\mathcal M^N_n$ has the same distribution as $W_{0,n}$ for all $n\ge 0$, and thus the desired result follows.
    }{
        在几乎必然的意义下与在 $\mathrm L^1(\mathbb P)$ 意义下都存在, 且是常数. 根据 $W$ 的定义, 所有 $n\ge 0$, $\mathcal M^N_n$ 与 $W_{0,n}$ 同分布, 从而命题得证.
    }
\end{proof}

\bilingual{
    We extend the last result to more general initial population distributions by using the coupling results in the last subsection.
}{
    利用上一小节的耦合结论, 上述结果可被推广到更一般的初始粒子分布.
}
\begin{corollary}
    \label{corollary:leftmost-speed:N}
    \bilingual{
        Under \Cref{assumption:psi-one,assumption:moment-1} (i.e. $\psi(1)=0$ and $\mathbb E[(\Xi_1)_+]<\infty$), there exists $v_N=v_N(\Xi)\in\mathbb R$ such that, with probability one, and in $\mathrm L^1(\mathbb P)$, for the $N$-BRW $\mathcal V^N$ with initial population $\mathcal Z^N_0\in \mathcal C(N)$,
    }{
        在\Cref{assumption:psi-one,assumption:moment-1} (即 $\psi(1)=0$ 且 $\mathbb E[(\Xi_1)_+]<\infty$) 的条件下, 存在 $v_N=v_N(\Xi)\in\mathbb R$ 使得对于以任一 $\mathcal Z^N_0\in \mathcal C(N)$ 为初始粒子分布的 $N$-分支随机游走 $\mathcal V^N$, 几乎必然成立, 且在 $\mathrm L^1(\mathbb P)$ 意义下成立,
    }
    {
        \setlength{\belowdisplayskip}{-10pt}
        \begin{equation*}
            \lim_{n\to\infty}
            \frac 1n
            \mathcal M^N_n
            =
            v_N.
        \end{equation*}
    }
\end{corollary}

\begin{proof}
    \quad

    \bilingual{
        Note that
    }{
        注意到
    }
    \begin{equation*}
        \mathcal Z'_0
        :=
        N\delta_{\max\mathcal Z^N_0}
        \succcurlyeq
        \mathcal Z^N_0
        \succcurlyeq
        N\delta_{\min\mathcal Z^N_0}
        =:
        \mathcal Z''_0.
    \end{equation*}
    \bilingual{
        By \Cref{corollary:2} and \Cref{corollary:translation}, for another $N$-BRW $\mathcal V$ with initial population distribution $N\delta_0$, there exists a coupling of all these four $N$-BRWs so that
    }{
        根据\Cref{corollary:2} 与\Cref{corollary:translation}, 对另一个以 $N\delta_0$ 为初始粒子分布的 $N$-分支随机游走 $\mathcal V$, 存在一个耦合使得这四个 $N$-分支随机游走满足如下关系:
    }
    \begin{equation*}
        \max\mathcal Z_0^N
        +
        \min\mathcal V_n
        =
        \min \mathcal Z'_n
        \ge \min\mathcal Z^N_n
        \ge \min\mathcal Z''_n
        =
        \min\mathcal Z_0^N
        +
        \min\mathcal V_n,
    \end{equation*}
    \bilingual{
        almost surely for all $n\ge 0$. By \Cref{proposition:2},
        there exists $v_N=v_N(\Xi)\in\mathbb R$ such that, with probability one, and in $\mathrm L^1(\mathbb P)$, for the $N$-BRW $\mathcal V$ with initial population $N\delta_0$, $\lim_{n\to\infty}\frac 1n\min\mathcal V_n=v_N$. It immediately follows that
    }{
        对所有 $n\ge 0$ 几乎必然成立. 根据\Cref{proposition:2}, 存在 $v_N=v_N(\Xi)\in\mathbb R$ 使得对于以 $N\delta_0$ 为初始粒子分布的 $N$-分支随机游走 $\mathcal V$, 几乎必然成立, 且在 $\mathrm L^1(\mathbb P)$ 意义下成立 $\lim_{n\to\infty}\frac 1n\min\mathcal V_n=v_N$. 由此立即得到
    }
    \begin{equation*}
        \lim_{n\to\infty}
        \frac 1n\min\mathcal Z'_n
        =
        \lim_{n\to\infty}
        \frac 1n\min\mathcal Z''_n
        =
        \lim_{n\to\infty}
        \frac 1n\min\mathcal Z^N_n
        =
        v_N
    \end{equation*}
    \bilingual{
        almost surely and in $\mathrm L^1(\mathbb P)$, which completes the proof.
    }{
        几乎必然成立, 且在 $\mathrm L^1(\mathbb P)$ 意义下成立, 由此命题得证.
    }
\end{proof}

\bilingual{
    We further extend the last result to initial population distributions with total mass less than $N$.
}{
    我们进一步将上述结果推广到总质量小于 $N$ 的初始粒子分布.
}

\begin{corollary}
    \label{corollary:leftmost-speed:N0}
    \bilingual{
        Consider \Cref{assumption:psi-one,assumption:moment-1} (i.e. $\psi(1)=0$ and $\mathbb E[(\Xi_1)_+]<\infty$) and an additional \Cref{assumption:psi-zero} (i.e., $\psi(0)>0$). There exists $v_N=v_N(\Xi)\in\mathbb R$ such that, with probability one, and in $\mathrm L^1(\mathbb P)$, for the $N$-BRW $\mathcal V^N$ with initial population $\mathcal Z^N_0\in \mathcal C(N_0)$ and $N_0\le N$,
    }{
        考虑\Cref{assumption:psi-one,assumption:moment-1} (即 $\psi(1)=0$ 且 $\mathbb E[(\Xi_1)_+]<\infty$) 和额外的\Cref{assumption:psi-zero} (即 $\psi(0)>0$). 存在 $v_N=v_N(\Xi)\in\mathbb R$ 使得对于具初始粒子分布 $\mathcal Z^N_0\in \mathcal C(N_0)$ 且 $N_0\le N$ 的 $N$-分支随机游走 $\mathcal V^N$,
    }
    {
        \setlength{\belowdisplayskip}{-10pt}
        \begin{equation*}
            \lim_{n\to\infty}
            \frac 1n
            \mathcal M^N_n
            =
            v_N.
        \end{equation*}
    }
\end{corollary}

\begin{proof}
    \bilingual{
        Under \Cref{assumption:psi-zero}, i.e., $\psi(0)>0$ or $\mathbb E[\#\Xi]>0$, we have $\mathbb P(\#\Xi\ge 2)>0$. Then there exists some time $n_0$ such that $\#\mathcal Z^N_{n_0}=N$. Then the result follows from \Cref{corollary:leftmost-speed:N}.
    }{
        在\Cref{assumption:psi-zero} 下, 即 $\psi(0)>0$ 或 $\mathbb E[\#\Xi]>0$, 我们有 $\mathbb P(\#\Xi\ge 2)>0$. 于是存在某个时间 $n_0$ 使得 $\#\mathcal Z^N_{n_0}=N$. 由此结果可直接从\Cref{corollary:leftmost-speed:N} 得出.
    }
\end{proof}

% !TeX root = ../../main.tex
\subsection{\bilingual{Coupling of \texorpdfstring{$N$}{N}-BRW and Killed BRW}{\texorpdfstring{$N$}{N}-分支随机游走与经灭杀的分支随机游走的耦合}}
\begin{lemma}
    \label{lemma:4.8}
    \bilingual{
        Let $\mathcal V^{\mathrm k}$ be the branching random walk $\mathcal V$ killed according to a given rule, such that the generation size process $\#\mathcal Z^{\mathrm k}$ is adapted to the natural filtration generated by $\mathcal V$. Assume that at time $0$, the population of $\mathcal V^{\mathrm k}$ is included in (or equal to) that of $\mathcal V^N$. Then there exists a coupling of $\mathcal V^N$ and $\mathcal V^{\mathrm k}$ such that
    }{
        设 $\mathcal V^{\mathrm k}$ 是按给定灭杀规则经灭杀的分支随机游走,~ 使得其粒子总数过程 $\#\mathcal Z^{\mathrm k}$ 适应于 $\mathcal V$ 的自然流.~ 假设在时刻 $0$,~ $\mathcal V^{\mathrm k}$ 的粒子分布包含于 (或等于) $\mathcal V^N$ 的粒子分布.~ 则存在 $\mathcal V^N$ 和 $\mathcal V^{\mathrm k}$ 的一个耦合使得
    }
    {
        \setlength{\belowdisplayskip}{-10pt}
        \begin{equation*}
            \mathcal M^N_k
            \le
            \mathcal M^{\mathrm k}_k,
            \quad
            k\le\tau:=\inf\left\{
                k\ge 0:
                \#\mathcal Z^{\mathrm k}_k
                >
                N
            \right\}.
        \end{equation*}
    }
\end{lemma}

\begin{proof}
    \quad

    \bilingual{
        For two finite counting measures on $\mathbb R$, i.e.,
    }{
        对于 $\mathbb R$ 上的两个有限计数测度,~ 即
    }
    $
        \mu=\sum_{i=1}^{a(\mu)}\delta_{x_i}$,
        % \quad
        $\nu=\sum_{i=1}^{a(\nu)}\delta_{y_i}$,
    % \end{equation*}
    \bilingual{
        where $x_1\le x_2\le\cdots\le x_{a(\mu)}$ and $y_1\le y_2\le\cdots\le y_{a(\nu)}$. It suffices to construct a coupling of $\mathcal V^N$ and $\mathcal V^{\mathrm k}$ such that
    }{
        其中 $x_1\le x_2\le\cdots\le x_{a(\mu)}$ 且 $y_1\le y_2\le\cdots\le y_{a(\nu)}$.~ 只需构建 $\mathcal V^N$ 和 $\mathcal V^{\mathrm k}$ 的一个耦合使得
    }
    \begin{equation}
        \mathcal V^{\mathrm k}_k
        \succcurlyeq
        \mathcal V^N_k,\quad
        k\le\tau.
        \label{eq:4.12}
    \end{equation}
    \bilingual{
        We proceed by induction on $k$. By assumption, \eqref{eq:4.12} holds for $k=0$. Assume that the coupling \eqref{eq:4.12} has been built up to some time $k<\tau$. We may write $\mathcal V^{\mathrm k}_k=\sum_{i=1}^{a_k}\delta_{x_i}$ and $\mathcal V^N_k=\sum_{i=1}^{b_k}\delta_{y_i}$ such that $a_k\le b_k$ and $x_i\ge y_i$ for all $i\in[\![1,a_k]\!]$. If $a_k=0$, $\tau=\infty$, and \eqref{eq:4.12} holds for all times after $k$ by definition. Consider the non-degenerate case where $a_k\ge 1$. We construct the coupling at time $k+1$ as follows.
    }{
        对 $k$ 归纳证明.~ 根据假设,~ \eqref{eq:4.12} 在 $k=0$ 时成立.~ 假设耦合 \eqref{eq:4.12} 已经构建到某个时刻 $k<\tau$.~ 此时可以写 $\mathcal V^{\mathrm k}_k=\sum_{i=1}^{a_k}\delta_{x_i}$ 和 $\mathcal V^N_k=\sum_{i=1}^{b_k}\delta_{y_i}$ 使得 $a_k\le b_k$ 且对于所有 $i\in[\![1,a_k]\!]$,~ 都有 $x_i\ge y_i$.~ 如果 $a_k=0$,~ 则 $\tau=\infty$,~ 根据定义 \eqref{eq:4.12} 在 $k$ 之后的所有时刻都成立.~ 考虑非退化情况 $a_k\ge 1$.~ 我们如下构建时刻 $k+1$ 的耦合.
    }
    \begin{itemize}
        \item
        \bilingual{
            For $1\le i\le a_k$, each $x_i$ particle of $\mathcal V^{\mathrm k}_k$ produces a collection of offspring $x_i+\Xi_{i,k}$ at time $k+1$ in $\mathcal V$, where $(\Xi_{i,k})_{i\ge 1,k\ge 0}$ are i.i.d.~copies of $\Xi$. Those offspring that are non-killed according to the given rule (if the set is not empty) constitute $\mathcal V^{\mathrm k}_{k+1}$. Note that if the killing cannot be performed here due to the use of future information in the killing rule, and the number of offspring is not well defined, then the proof cannot proceed, which is exactly the situation excluded by this lemma.
        }{
            对 $1\le i\le a_k$,~ $\mathcal V^{\mathrm k}_k$ 的每个粒子 $x_i$ 在时刻 $k+1$ 产生一组后代 $x_i+\Xi_{i,k}$,~ 其中 $(\Xi_{i,k})_{i\ge 1,k\ge 0}$ 是 $\Xi$ 的独立同分布的拷贝. 根据给定的灭杀规则,~ 那些未被灭杀的后代 (如果该集合非空) 构成 $\mathcal V^{\mathrm k}_{k+1}$. ~注意,~ 如果这里由于灭杀规则用到了未来的信息而无法执行灭杀,~ 从而无法确定余下粒子的数量,~ 这里的证明就无法继续下去,~ 这也就是本引理所排除的情形.
        }
        \item
        \bilingual{
            For $\mathcal V^N$, each $y_i$, $i\in[\![1,b_k]\!]$, produces a collection of offspring $y_i+\Xi_{i,k}$ at time $k+1$, and we select the $N$ leftmost particles to form $\mathcal V^N_{k+1}$. In particular, $\#\mathcal V^N_{k+1}=N$.
        }{
            对 $\mathcal V^N$,~ 每个 $y_i$,~ $i\in[\![1,b_k]\!]$,~ 在时刻 $k+1$ 产生一组后代 $y_i+\Xi_{i,k}$,~ 选取最左侧的 $N$ 个粒子组成 $\mathcal V^N_{k+1}$.~ 特别地,~ $\#\mathcal V^N_{k+1}=N$.
        }
        \item
        \bilingual{
            If the population of $\mathcal V^{\mathrm k}_{k+1}$ does not exceed $N$, then necessarily $\mathcal V^{\mathrm k}_{k+1}\succcurlyeq\mathcal V^N_{k+1}$, and \eqref{eq:4.12} holds at time $k+1$. Otherwise, we have $\#\mathcal V^{\mathrm k}_{k+1}>N$, which implies that $\tau=k+1$ and \eqref{eq:4.12} holds as well.
        }{
            若 $\mathcal V^{\mathrm k}_{k+1}$ 的粒子总数不超过 $N$,~ 则必然 $\mathcal V^{\mathrm k}_{k+1}\succcurlyeq\mathcal V^N_{k+1}$,~ 于是 \eqref{eq:4.12} 在时刻 $k+1$ 成立.~ 否则,~ 有 $\#\mathcal V^{\mathrm k}_{k+1}>N$,~ 此时 $\tau=k+1$ 且 \eqref{eq:4.12} 同样成立.
        }
        \qedhere
    \end{itemize}
\end{proof}

\begin{figure}[H]
    \centering
    \begin{tikzpicture}[
        xscale=3.3,
        yscale=1.1,
        >=Stealth,
        % 样式定义
        node-both/.style={circle, fill=blue!70!black, draw=green!60!black, line width=1.2pt, inner sep=1.2pt},
        node-N/.style={circle, fill=blue!70!black, inner sep=1.8pt},
        node-k/.style={circle, fill=green!60!black, inner sep=1.8pt},
        node-N-dead/.style={circle, draw=blue!70!black, fill=white, inner sep=1.4pt, cross out, draw=blue!80!black, line width=1.2pt},
        node-k-dead/.style={circle, draw=green!60!black, fill=white, inner sep=1.4pt, cross out, draw=green!60!black, line width=1.2pt},
        node-dead/.style={circle, draw=black!40, fill=white, inner sep=1.4pt},
        % 连线样式
        edge-N/.style={thick, blue!70!black},
        edge-k/.style={thick, green!60!black},
        edge-N-dashed/.style={dashed, blue!70!black},
        edge-k-dashed/.style={dashed, green!60!black},
        edge-both/.style={
            thick,
            blue!70!black, % 底色：蓝色
            postaction={
                draw,
                green!60!black,
                % on 2pt 代表绿色段长度，off 2pt 代表空隙（露出底下的蓝色）
                % 若要更频繁，将 2pt 改为 1pt
                dash pattern=on 5pt off 5pt
            }
        },
        edge-dashed/.style={dashed, black!30},
        % 灭杀边界
        kill-line/.style={red!80!black, dotted, ultra thick}
    ]

        % --- 坐标轴与背景 ---
        % \draw[->, gray!50, thick] (-0.2, -4.5) -- (-0.2, 2.5);
        \foreach \x in {0,1,2,3} \draw[-, black!20, dashed] (\x, -4.5) -- (\x, 2.5);
        \node[anchor=south] at (0, 2.7) {$k=0$};
        \node[anchor=south] at (1, 2.7) {$k=1$};
        \node[anchor=south] at (2, 2.7) {$k=2$};
        \node[anchor=south] at (3, 2.7) {$\tau$};

        % 灭杀屏障 (Vk 专用)
        % \draw[kill-line] (0, 0.2) -- (1, 0.2);

        % ================= n=0 =================
        \node[node-both] (n0) at (0, 0) {};
        \node[left] at (n0) {{\color{black}$\Xi_{1,0}$}};

        % ================= n=1 (N=2) =================
        % 产生4个后代。VN选前2个，Vk选1,3名（第2名过界被杀）
        \node[node-both] (n1-1) at (1, -1.2) {};
        \node[below] at (n1-1) {{\color{black}$\Xi_{1,1}$}};
        \node[node-N]    (n1-2) at (1, -0.2) {}; % VN保留，但Vk在此处被灭杀（过界或规则）
        \node[below] at (n1-2) {{\color{blue!70!black}$\Xi_{2,1}$}};
        \node[node-k]    (n1-3) at (1, 0.4)  {}; % VN淘汰（排名第3），但Vk保留
        \node[above] at (n1-3) {{\color{green!60!black}$\Xi_{2,1}$}};
        \node[node-N-dead] (n1-4) at (1, 1.5)  {};

        \draw[edge-both] (n0) -- (n1-1);
        \draw[edge-N] (n0) -- (n1-2); \draw[edge-N-dashed] (n0) -- (n1-2);
        \draw[edge-N-dashed] (n0) -- (n1-3); \draw[edge-k] (n0) -- (n1-3);
        \draw[edge-N-dashed] (n0) -- (n1-4);

        % ================= n=2 (N=2) =================
        % 耦合发生：VN从 y1, y2 繁殖；Vk从 x1, x3 繁殖

        % 1. VN 的前沿从 n1-2 爆发（产生极低值）
        \node[node-N] (n2-N-min) at (2, -3.0) {};
        \draw[edge-N] (n1-2) -- (n2-N-min);
        \draw[edge-N] (n1-2) -- (2, -2.2);
        \node[node-N] at (2, -2.2) {};
        \draw[edge-N-dashed] (n1-2) -- (2, -0.4);
        \node[node-N-dead] at (2, -0.4) {};

        % 2. Vk 的前沿从 n1-1 延续
        \node[node-k] (n2-k-min) at (2, -0.6) {};
        \node[below] at (n2-k-min) {{\color{green!60!black}$\Xi_{1,2}$}};
        \node[node-N-dead] at (2, -0.6) {};
        \draw[edge-both] (n1-1) -- (n2-k-min);

        % 3. 展现绿色来自被蓝色去掉的祖先
        \node[node-k] (n2-k-from-dead) at (2, 0.2) {};
        \node[above] at (n2-k-from-dead) {{\color{green!60!black}$\Xi_{2,2}$}};
        \draw[edge-k] (n1-3) -- (n2-k-from-dead);
        \node[node-k-dead] at (2, -2.4) {};
        \draw[edge-k-dashed] (n1-3) -- (2, -2.4);
        \node[node-k-dead] at (2, -1.6) {};
        \draw[edge-k-dashed] (n1-3) -- (2, -1.6);

        % ================= n=3 (核心耦合：左侧同步) =================
        % 蓝色最小值在 -3.0，绿色最小值在 -1.8
        % 定义相同的位移向量 Xi: {-1.2, -0.2, 0.5, 1.2}

        % --- 绿色最左侧的四分叉 (必须与上方蓝色四分叉完全平行) ---
        \node[node-k] (n3-k-1) at (3, -1.8) {}; % -0.6 + (-1.2)
        \node[node-k] (n3-k-2) at (3, -0.8) {}; % -0.6 + (-0.2)
        \node[node-k-dead] (n3-k-3) at (3, -0.1) {};
        \node[node-k-dead] (n3-k-4) at (3, 0.4) {};
        \draw[edge-k] (n2-k-min) -- (n3-k-1);
        \draw[edge-k] (n2-k-min) -- (n3-k-2);
        \draw[edge-k-dashed] (n2-k-min) -- (n3-k-3);
        \draw[edge-k-dashed] (n2-k-min) -- (n3-k-4);

        % --- 蓝色最左侧的四分叉 ---
        \node[node-N] (n3-N-1) at (3, -4.2) {}; % -3.0 + (-1.2)
        \node[node-N] (n3-N-2) at (3, -3.2) {}; % -3.0 + (-0.2)
        \node[node-N-dead] (n3-N-3) at (3, -2.5) {};
        \node[node-N-dead] (n3-N-4) at (3, -2.0) {};
        \draw[edge-N] (n2-N-min) -- (n3-N-1);
        \draw[edge-N] (n2-N-min) -- (n3-N-2);
        \draw[edge-N-dashed] (n2-N-min) -- (n3-N-3);
        \draw[edge-N-dashed] (n2-N-min) -- (n3-N-4);
        \node[below] at (n2-N-min) {{\color{blue!70!black}$\Xi_{1,2}$}};

        % 绿色其它支系
        \node[node-k] (n3-k-other) at (3, 1.3) {};
        \draw[edge-k] (n2-k-from-dead) -- (n3-k-other);
        \draw[edge-N-dashed] (2, -2.2) -- (3, -1.1);
        \node[node-N-dead] at (3, -1.1) {};
        \node[right] at (2, -2.2) {{\color{blue!70!black}$\Xi_{2,2}$}};

        % ================= 标注 =================
        % 最小值对比
        \draw[blue!70!black, <-, thick] (3.05, -4.2) -- (3.15, -4.2) node[right] {$\mathcal{M}^2_3$};
        \draw[green!60!black, <-, thick] (3.05, -1.8) -- (3.15, -1.8) node[right] {$\mathcal{M}^{\mathrm{k}}_3$};

        % % 耦合箭头：显示位移的一致性
        % \draw[<->, gray, thin, bend left=20] (2.1, -3.0) to node[midway, right, font=\tiny] {$\Xi_{1,n}$} (2.1, -1.8);
        % \draw[stealth-stealth, orange, dashed] (2.9, -4.2) -- (2.9, -3.0) node[midway, left, font=\tiny] {Parallel $\Xi$};

    \end{tikzpicture}
    \caption{\bilingual{Illustration of a \texorpdfstring{$2$}{2}-BRW (blue) and a killed BRW (green) coupling}{一个\texorpdfstring{$2$-BRW}{2-BRW} (蓝色) 与一个经灭杀的\texorpdfstring{BRW}{BRW} (绿色) 的耦合示例图}}
\end{figure}


% !TEX root = ../main.tex

\section{Main Theorems}

\begin{theorem}
    \label{proposition:3.1}
    Fix $a\in[0,1]$ and $t>0$. If $\mathcal V^N=\mathcal V^{N|a}$ starts with $\lfloor N^a\rfloor$ particles located at $0$, then for any sequence $\ell_N$ satisfying $\lim_{N\to\infty}\frac{\ell_N}{\log^3N}=t$, 
    {
        \setlength{\belowdisplayskip}{0pt}
        \begin{equation}
            \frac 1{\log N}
            \min_{|u|=\ell_N}
            \mathcal V^N(u)
            \xrightarrow{\mathrm L^2}
            -a+\frac{\pi^2}2\sigma^2t.
        \end{equation}
    }
\end{theorem}

\begin{proof}[Proof of \Cref{proposition:3.1}]
    \quad

    The result would follow if we can show that
    \begin{align}
        \lim_{N\to\infty}
        \frac 1{\log N}
        \mathbb E\left[
            \min_{|u|=\ell_N}\mathcal V^N(u)
        \right]
        &
        =
        -a+\frac{\pi^2}2\sigma^2t,
        \label{eq:4.13}
        \\
        \lim_{N\to\infty}
        \frac 1{\log^2 N}
        \mathrm{Var}\!\left[
            \min_{|u|=\ell_N}\mathcal V^N(u)
        \right]
        &
        =
        0.
        \label{eq:4.14}
    \end{align}
    Indeed, for
    \begin{equation*}
        \mathcal M^N_n:=\min_{|u|=n}\mathcal V^N(u), \quad
        \mathcal X_N:=\frac 1{\log N}\mathcal M^N_{\ell_N}, \quad
        \mu=\mu(t):=-a+\frac{\pi^2}2\sigma^2t,
    \end{equation*}
    we have
    \begin{equation*}
        \mathbb E\left[
            \left(
                \mathcal X_N-\mu
            \right)^2
        \right]
        =
        \mathrm{Var}\!\left[
            \mathcal X_N
        \right]
        +
        \left(
            \mathbb E\left[
                \mathcal X_N
            \right]
            -\mu
        \right)^2
        \xrightarrow{}
        0,
        \quad
        N\to\infty.
    \end{equation*}

    \quad

    Further, it suffices to show that for any $\zeta<\mu<\zeta'$, there exists some $\varepsilon=\varepsilon(\zeta,\zeta')>0$ such that for some $N_0=N_0(\varepsilon,\zeta,\zeta')>0$ and all $N\ge N_0$,
    \begin{align}
        \mathbb P(\mathcal X_N\ge\zeta')
        =
        \mathbb P\left(
            \mathcal M^N_{\ell_N}
            \ge
            \zeta'\log N
        \right)
        &
        \le
        N^{-\varepsilon},
        \label{eq:4.15}
        \\
        \mathbb P(\mathcal X_N\le\zeta)
        =
        \mathbb P\left(
            \mathcal M^N_{\ell_N}
            \le
            \zeta\log N
        \right)
        &
        \le
        N^{-\varepsilon}.
        \label{eq:4.16}
    \end{align}
    Indeed, observe that $\mathcal M^N_n$ is (stochastically) smaller than $\mathcal V^1$ at time $n$, where $\mathcal V^1$ is the branching random walk in which only the leftmost particle is kept at each step. Therefore, $\mathcal V^1$ is a random walk with step distribution $\Xi_1$. Under \Cref{assumption:psi-one,assumption:moment-4} (i.e. $\psi(1)=0$ and $\mathbb E[(\Xi_1)_+^4]<\infty$), we can show that
    \begin{equation*}
        \mathbb E[(\mathcal M^N_{n})_-^4]=O_n(1),
        \quad
        \mathbb E[(\mathcal M^N_{n})_+^4]=O_n(n^4),
    \end{equation*}
    as $n$ goes to infinity (see \Cref{lemma:M:moment-4:neg,lemma:M:moment-4:pos} for details),
    and thus
    \begin{equation*}
        \mathbb E[(\mathcal M^N_{n})^4]
        \le
        \mathbb E[(\mathcal M^N_{n})_-^4]
        +
        \mathbb E[(\mathcal M^N_{n})_+^4]
        =
        O_n(n^4).
    \end{equation*}
    Now, let for $N\ge 1$,
    \begin{equation*}
        H_N
        :=\left\{
            \zeta\log N\le \mathcal M^N_{\ell_N}\le \zeta'\log N
        \right\}
        =
        \left\{
            \zeta\le \mathcal X_N\le \zeta'
        \right\}.
    \end{equation*}
    According to \eqref{eq:4.15} and \eqref{eq:4.16}, for all $N\ge N_0$,
    $
        \mathbb P(
            H_N^{\mathrm c}
        )
        \le
        2N^{-\varepsilon}
    $, and thus
    \begin{equation*}
        \mathbb P(
            H_N^{\mathrm c}
        )
        =
        O_N(N^{-\varepsilon}),
        \quad
        \lim_{N\to\infty}\mathbb P(H_N)=1.
    \end{equation*}
    By Cauchy--Schwarz inequality and the fact that $\ell_N=O_N(\log^3 N)$, we have
    \begin{align*}
        \mathbb E\left[
            \mathcal X_N^2
            {\bf 1}_{H_N^{\mathrm c}}
        \right]
        &
        \le
        \frac{
            \mathbb E\left[
                \left(
                    \mathcal M^N_{\ell_N}
                \right)^4
            \right]^{\frac 12}
            \!\!\!
            \cdot
            \mathbb P\!\left(
                H_N^{\mathrm c}
            \right)^{\frac 12}
        }{\log^2 N}
        =
        \frac{
            O_N(\ell_N^2)
            \cdot
            O_N(N^{-\varepsilon/2})
        }{\log^2 N}
        =
        o_N(1),
    \end{align*}
    and thus
    \begin{equation*}
        \left|
            \mathbb E\left[
                \mathcal X_N
                {\bf 1}_{H_N^{\mathrm c}}
            \right]
        \right|
        \le
        \mathbb E\left[
            |\mathcal X_N|
            {\bf 1}_{H_N^{\mathrm c}}
        \right]
        \le
        \mathbb E\left[
            \mathcal X_N^2
            {\bf 1}_{H_N^{\mathrm c}}
        \right]^{\frac 12}
        =
        o_N(1).
    \end{equation*}
    On the other hand, by definition of $H_N$,
    \begin{align*}
        \zeta\mathbb P(H_N)
        \le
        \mathbb E\left[
            \mathcal X_N
            {\bf 1}_{H_N}
        \right]
        \le
        \zeta'\mathbb P(H_N),
        \\
        \zeta^2\mathbb P(H_N)
        \le
        \mathbb E\left[
            \mathcal X_N^2
            {\bf 1}_{H_N}
        \right]
        \le
        \zeta'^2\mathbb P(H_N),
    \end{align*}
    Let $N\to\infty$, then $\zeta,\zeta'\to \mu$, we get
    \begin{align*}
        \lim_{N\to\infty}
        \mathbb E[ \mathcal X_N ]
        =
        \lim_{N\to\infty}
        \mathbb E[ \mathcal X_N {\bf 1}_{H_N} ]
        &
        =
        \mu,\\
        \lim_{N\to\infty}
        \mathbb E[ \mathcal X_N^2 ]
        =
        \lim_{N\to\infty}
        \mathbb E[ \mathcal X_N^2 {\bf 1}_{H_N} ]
        &
        =
        \mu^2,
    \end{align*}
    yielding \eqref{eq:4.13} and \eqref{eq:4.14}.


    \quad

    We will first show \eqref{eq:4.15} for the case of $a\in(0,1)$ in the first sub-section, then for the case of $a=1$ and $a=0$ in the second sub-section, and finally show \eqref{eq:4.16} for all $a\in[0,1]$ in the third sub-section.
\end{proof}

\subsection{Proof of \eqref{eq:4.15} for \texorpdfstring{$a\in(0,1)$}{a∈(0,1)}}

\begin{proof}[Proof of \eqref{eq:4.15} for the case of \texorpdfstring{$a\in(0,1)$}{a∈(0,1)}]
    \quad

    We now proceed to show \eqref{eq:4.15} for the case of $a\in(0,1)$. 
    
    In light of \Cref{lemma:4.8}, we construct a killed branching random walk $\mathcal V^{\mathrm{k}}$ (killed according to a given rule specified later). Denote $\mathcal M^{\mathrm{k}}_n:=\min_{|u|=n}\mathcal V^{\mathrm{k}}(u)$ for $n\ge 0$. So long as $\mathcal V^{\mathrm{k}}$ satisfies the conditions listed in \Cref{lemma:4.8}, we have
    \begin{equation*}
        \mathcal M^N_{\ell_N}
        \le
        \mathcal M^{\mathrm{k}}_{\ell_N}
        \quad
        \text{on the event $\{\tau>\ell_N\}$,}
    \end{equation*}
    where $\tau$ is the first time that the generation size of $\mathcal V^{\mathrm{k}}$ exceeds $N$. Note that the left-hand side in \eqref{eq:4.15} 
    \begin{align}
        &\quad
        \mathbb P\left(
            \mathcal M^N_{\ell_N}
            \ge
            \zeta'\log N
        \right)
        \nonumber
        \\
        &
        =
        \mathbb P\left(
            \mathcal M^N_{\ell_N}
            \ge
            \zeta'\log N,\,
            \tau\le\ell_N
        \right)
        +
        \mathbb P\left(
            \mathcal M^N_{\ell_N}
            \ge
            \zeta'\log N,\,
            \tau>\ell_N
        \right)
        \nonumber
        \\ &
        \le
        \mathbb P\left(
            \tau\le\ell_N
        \right)
        +
        \mathbb P\left(
            \mathcal M^{\mathrm{k}}_{\ell_N}
            \ge
            \zeta'\log N
        \right).
        \label{eq:prop:3.1:1}
    \end{align}
    Therefore, we need to construct a $\mathcal V^{\mathrm{k}}$ such that, with high probability,
    \begin{itemize}
        \item $\tau>\ell_N$, i.e., the generation size of $\mathcal V^{\mathrm{k}}$ does not exeed $N$ till time $\ell_N$;
        \item $\mathcal M^{\mathrm{k}}_{\ell_N}<\zeta'\log N$, i.e., the minimum position of $\mathcal V^{\mathrm{k}}$ is below $\zeta'\log N$ at time $\ell_N$;
    \end{itemize}
    and, the probability of deviating from these two normal situations, i.e., the right-hand side of \eqref{eq:prop:3.1:1}, decays at least polynomially with respect to $N$. Note that, the left-hand side of \eqref{eq:prop:3.1:1} is naturally hard to compute, since the selection mechanism underlying the $N$-BRW requires that we keep track of the whole population at each step, meaning that when we are determining whether a particular particle should be preserved, we have to look at the positions of all other particles, including the descendants of different ancesters at time $0$, which is hard to analyze. On the other hand, the right-hand side of \eqref{eq:prop:3.1:1} is easier to compute, since we can construct $\mathcal V^{\mathrm{k}}$ with a simpler selection mechanism, so that when we are determining whether a particular particle should be preserved, we only need to look at the position of itself, and thus we can analyze $\mathcal V^{\mathrm{k}}$ by using the many-to-one formula and Mogulskii's estimate for random walks.

    We first state our construction of the killed branching random walk $\mathcal V^{\mathrm{k}}$. Choose and fix $c\in(0,t^{-1/3})$, $b\in(0,a)$, and $v>0$ such that
    \begin{equation*}
        ct^{\frac 13}+(a-b)+v<1,\quad
        -b+\frac{\pi^2}2\sigma^2\left(c^{-2}t^{\frac 13}\right)+2v\left(1+t^{\frac 13}\right)<\zeta'.
    \end{equation*}
    Note that these inequalites are possible as $c\uparrow t^{-1/3}$, $b\uparrow a$, and $v\downarrow 0$. Also note that they are dependant on $a$, $t$, and $\zeta'$.
    \begin{itemize}
        \item At time $0$, start with $\lfloor N^a\rfloor$ particles located at $0$; and they branch accoding to $\Xi$; note that these make $\mathcal V^{\mathrm{k}}$ satisfy the condition of \Cref{lemma:4.8};
        \item From time $0$ to $\ell:=\lfloor\log^{5/2}N\rfloor$, kill particles outside the tube $[-b\log N,0]$;
        \item At time $\ell$, further kill particles outside the interval $[-b\log N, (-b+v)\log N]$;
        \item From time $\ell+1$ to $\ell_N$, for any alive particle in the last step at position $y$, kill particles outside a tilted tube with a small drift, i.e., $y+[r_j,s_j]$ for $j\in[\![\ell,\ell_N]\!]$, where $r_j:=\frac{\pi^2}{2}\sigma^2c^{-2}L^{-2/3}(j-\ell)$, $s_j:=r_j+cL^{1/3}$ and $L:=\ell_N-\ell$.
        \item At time $\ell_N$, kill all particles outside the interval $y+[r_{\ell_N},r_{\ell_N}+vL^{1/3}]$ for any particle alive after the last step with an ancester of time $\ell$ at position $y$.
    \end{itemize}
    For the record, we would denote the set of particles alive after each operation respectively as
    \begin{equation*}
        \mathcal Z^{\mathrm{k}}_n,\,
        n\in[\![0,\ell]\!],\quad
        \mathcal Z^{\mathrm{k}}_{\ell+},\quad
        \mathcal Z^{\mathrm{k}}_n,\,
        n\in[\![\ell+1,\ell_N]\!],\quad
        \mathcal Z^{\mathrm{k}}_{\ell_N+}.
    \end{equation*}
    Note that
    \begin{equation*}
        \mathcal Z^{\mathrm{k}}_{\ell+}
        \subset
        \mathcal Z^{\mathrm{k}}_{\ell},\quad
        \mathcal Z^{\mathrm{k}}_{\ell_N+}
        \subset
        \mathcal Z^{\mathrm{k}}_{\ell_N}.
    \end{equation*}

    Next, we calculate an upper bound of the first term $\mathbb P(\tau\le\ell_N)$ in \eqref{eq:prop:3.1:1} for the constructed killed branching random walk $\mathcal V^{\mathrm{k}}$. Note that, by union bound and Markov's inequality,
    \begin{align*}
        \mathbb P(\tau\le\ell_N)
        &
        =
        \mathbb P\left(
            \begin{aligned}
                &\;\;\;\;\;
                (
                    \exists n\in[\![1,\ell-1]\!]:
                    \#\mathcal Z^{\mathrm{k}}_n>N
                )\\
                &\vee\;\;
                (
                    \#\mathcal Z^{\mathrm{k}}_{\ell+}>N
                )\\
                &\vee\;\;
                (
                    \exists n\in[\![\ell+1,\ell_N-1]\!]:
                    \#\mathcal Z^{\mathrm{k}}_n>N
                )\\
                &\vee\;\;
                (
                    \#\mathcal Z^{\mathrm{k}}_{\ell_N+}>N
                )
            \end{aligned}
        \right)
        \\ &
        \le
        \mathbb P\left(
            \exists n\in[\![1,\ell_N]\!]:
            \#\mathcal Z^{\mathrm{k}}_n>N
        \right)
        \\ &
        \le
        \sum_{n=1}^{\ell_N}
        \mathbb P\left(
            \#\mathcal Z^{\mathrm{k}}_n>N
        \right)
        % \\ &
        \le
        \sum_{n=1}^{\ell_N}
        \frac{
            \mathbb E\left[
                \#\mathcal Z^{\mathrm{k}}_n
            \right]
        }N.
    \end{align*}
    Let $\eta:=\max\{a,a-b+v+ct^{1/3}\}<1$. First we consider the operation from time $0$ to $\ell$, and the expected number $\mathbb E[\#\mathcal Z^{\mathrm{k}}_n]$ of particles alive till time $n\in[\![1,\ell]\!]$ is smaller than $N^{a+o_N(1)}$. Indeed, for any $n\in[\![1,\ell]\!]$, by \eqref{eq:4.10} in \Cref{lemma:4.7}, we have
    \begin{align*}
        \mathbb E\left[
            \#\mathcal Z^{\mathrm{k}}_n
        \right]
        &
        =
        \mathbb E\left[
            \sum_{|u|=n}
            {\bf 1}_{
                \{
                    -b\log N
                    \le
                    \mathcal V^{\mathrm{k}}(u_j)
                    \le
                    0,\,
                    \forall j\in[\![1,n]\!]
                \}
            }
        \right]
        \\
        &
        =
        \mathbb E\left[
            \sum_{i=1}^{\lfloor N^a\rfloor}
            \sum_{|u|=n}
            {\bf 1}_{
                \{
                    -b\log N
                    \le
                    \mathcal V^{\mathrm{k}|(i)}(u_j)
                    \le
                    0,\,
                    \forall j\in[\![1,n]\!]
                \}
            }
        \right]
        \\
        &
        =
        \sum_{i=1}^{\lfloor N^a\rfloor}
        \mathbb E\left[
            \sum_{|u|=n}
            {\bf 1}_{
                \{
                    -b\log N
                    \le
                    \mathcal V^{\mathrm{k}|(i)}(u_j)
                    \le
                    0,\,
                    \forall j\in[\![1,n]\!]
                \}
            }
        \right]
        \\
        &
        \le
        \sum_{i=1}^{\lfloor N^a\rfloor}
        \mathbb E\left[
            \sum_{n=1}^\infty
            \sum_{|u|=n}
            {\bf 1}_{
                \{
                    -b\log N
                    \le
                    \mathcal V^{\mathrm{k}|(i)}(u_j)
                    \le
                    0,\,
                    \forall j\in[\![1,n]\!]
                \}
            }
        \right]
        \\
        &
        =
        \lfloor N^a\rfloor
        O_N(\log N)
        =
        N^a
        O_N(\log N)
        =
        N^{a+o_N(1)}
        \le
        N^{\eta+o_N(1)},
    \end{align*}
    where, for the first equality, as defined, $\mathcal V^{\mathrm{k}}(u)=\mathcal V^{\mathrm{k}}(u,i)$ with $i$ the index of tree $\boldsymbol{\mathcal T}(i)$ which $u\in\mathbb V$ is on, and the summation is over all $i\in\mathbb I$ and $u\in\boldsymbol{\mathcal T}(i)$ with $|u|=n$. Then we consider the operation at time $\ell$, and the expected number $\mathbb E[\#\mathcal Z^{\mathrm{k}}_{\ell+}]$ of particles alive in the interval $[-b\log N, (-b+v)\log N]$ is of order $N^{a-b+v+o_N(1)}$. Indeed, by \eqref{eq:4.11} in \Cref{lemma:4.7}, we have
    \begin{align*}
        \mathbb E\left[\#\mathcal Z^{\mathrm{k}}_{\ell+}\right]
        &
        =
        \mathbb E\left[
            \sum_{|u|=\ell}
            {\bf 1}_{
                \{
                    -b\log N
                    \le
                    \mathcal V^{\mathrm{k}}(u_j)
                    \le
                    0,\,
                    \forall j\in[\![1,\ell]\!]
                \}
            }
            {\bf 1}_{
                \{
                    \mathcal V^{\mathrm{k}}(u)
                    \le
                    (-b+v)\log N
                \}
            }
        \right]
        \\
        &
        =
        \mathbb E\left[
            \sum_{i=1}^{\lfloor N^a\rfloor}
            \sum_{|u|=\ell}
            {\bf 1}_{
                \{
                    -b\log N
                    \le
                    \mathcal V^{\mathrm{k}|(i)}(u_j)
                    \le
                    0,\,
                    \forall j\in[\![1,\ell]\!]
                \}
            }
            {\bf 1}_{
                \{
                    \mathcal V^{\mathrm{k}|(i)}(u)
                    \le
                    (-b+v)\log N
                \}
            }
        \right]
        \\
        &
        =
        \sum_{i=1}^{\lfloor N^a\rfloor}
        \mathbb E\left[
            \sum_{|u|=\ell}
            {\bf 1}_{
                \{
                    -b\log N
                    \le
                    \mathcal V^{\mathrm{k}|(i)}(u_j)
                    \le
                    0,\,
                    \forall j\in[\![1,\ell]\!]
                \}
            }
            {\bf 1}_{
                \{
                    \mathcal V^{\mathrm{k}|(i)}(u)
                    \le
                    (-b+v)\log N
                \}
            }
        \right]
        \\
        &
        =
        \lfloor N^a\rfloor
        N^{-b+v+o_N(1)}
        =
        N^{a-b+v+o_N(1)}.
    \end{align*}
    Now we consider the operation from time $\ell+1$ to $\ell_N$, and the expected number of particles alive till time $n\in[\![\ell+1,\ell_N]\!]$ is smaller than $N^{a-b+v+ct^{1/3}+o_N(1)}$. Indeed, by \eqref{eq:4.11} in \Cref{lemma:4.7}, \eqref{eq:4.8} in \Cref{lemma:4.6}, and the branching property, we have, for any $n\in[\![\ell+1,\ell_N]\!]$,
    \begin{align*}
        \mathbb E\left[\#\mathcal Z^{\mathrm{k}}_n\right]
        &
        =
        \mathbb E\left[
            \sum_{|u|=n}
            \begin{aligned}
                {\bf 1}_{
                    \{
                        -b\log N
                        \le
                        \mathcal V^{\mathrm{k}}(u_j)
                        \le
                        0,\,
                        \forall j\in[\![1,\ell]\!]
                    \}
                }
                {\bf 1}_{
                    \{
                        \mathcal V^{\mathrm{k}}(u_\ell)
                        \le
                        (-b+v)\log N
                    \}
                }
                \\
                {\bf 1}_{
                    \{
                        r_j
                        \le
                        \mathcal V^{\mathrm{k}}(u_j)
                        -
                        \mathcal V^{\mathrm{k}}(u_\ell)
                        \le s_j,\,
                        \forall j\in[\![\ell+1,n]\!]
                    \}
                }
            \end{aligned}
        \right]
        \\
        &
        =
        \sum_{i=1}^{\lfloor N^a\rfloor}
        \mathbb E\left[
            \sum_{|u|=n}
            \begin{aligned}
                {\bf 1}_{
                    \{
                        -b\log N
                        \le
                        \mathcal V^{\mathrm{k}|(i)}(u_j)
                        \le
                        0,\,
                        \forall j\in[\![1,\ell]\!]
                    \}
                }
                {\bf 1}_{
                    \{
                        \mathcal V^{\mathrm{k}|(i)}(u_\ell)
                        \le
                        (-b+v)\log N
                    \}
                }
                \\
                {\bf 1}_{
                    \{
                        r_j
                        \le\mathcal V^{\mathrm{k}|(i)}(u_j)-
                        \mathcal V^{\mathrm{k}|(i)}(u_\ell)
                        \le s_j,\,
                        \forall j\in[\![\ell+1,n]\!]
                    \}
                }
            \end{aligned}
        \right]
        \\
        &
        =
        \sum_{i=1}^{\lfloor N^a\rfloor}
        \mathbb E\left[
            \sum_{|\nu|=\ell}
            \sum_{
                |u|=n,\,
                u_\ell=\nu
            }
            \!\!\!\!\!\!\!\!\!
            \begin{aligned}
                {\bf 1}_{
                    \{
                        -b\log N
                        \le
                        \mathcal V^{\mathrm{k}|(i)}(\nu_j)
                        \le
                        0,\,
                        \forall j\in[\![1,\ell]\!]
                    \}
                }
                {\bf 1}_{
                    \{
                        \mathcal V^{\mathrm{k}|(i)}(\nu)
                        \le
                        (-b+v)\log N
                    \}
                }
                \\
                {\bf 1}_{
                    \{
                        r_j
                        \le\mathcal V^{\mathrm{k}|(i)}(u_j)-
                        \mathcal V^{\mathrm{k}|(i)}(u_\ell)
                        \le s_j,\,
                        \forall j\in[\![\ell+1,n]\!]
                    \}
                }
            \end{aligned}
        \right]
        \\ &
        =
        \sum_{i=1}^{\lfloor N^a\rfloor}
        \mathbb E\left[
            \begin{aligned}
                \sum_{|\nu|=\ell}
                {\bf 1}_{
                    \{
                        -b\log N
                        \le
                        \mathcal V^{\mathrm{k}|(i)}(\nu_j)
                        \le
                        0,\,
                        \forall j\in[\![1,\ell]\!]
                    \}
                }
                {\bf 1}_{
                    \{
                        \mathcal V^{\mathrm{k}|(i)}(\nu)
                        \le
                        (-b+v)\log N
                    \}
                }
                \\
                \times\;
                \mathbb E\left(
                    \sum_{
                        |u|=n,\,
                        u_\ell=\nu
                    }
                    \!\!\!\!\!
                    {\bf 1}_{
                        \{
                            r_j
                            \le\mathcal V^{\mathrm{k}|(i)}(u_j)-
                            \mathcal V^{\mathrm{k}|(i)}(u_\ell)
                            \le s_j,\,
                            \forall j\in[\![\ell+1,n]\!]
                        \}
                    }
                    \middle|
                    \mathcal F_\ell
                \right)
            \end{aligned}
        \right]
        \\ &
        =
        \sum_{i=1}^{\lfloor N^a\rfloor}
        \mathbb E\left[
            \begin{aligned}
                \sum_{|\nu|=\ell}
                {\bf 1}_{
                    \{
                        -b\log N
                        \le
                        \mathcal V^{\mathrm{k}|(i)}(\nu_j)
                        \le
                        0,\,
                        \forall j\in[\![1,\ell]\!]
                    \}
                }
                {\bf 1}_{
                    \{
                        \mathcal V^{\mathrm{k}|(i)}(\nu)
                        \le
                        (-b+v)\log N
                    \}
                }
                \\
                \times\;
                \mathbb E\left[
                    \sum_{|x|=n-\ell}
                    {\bf 1}_{
                        \{
                            r_{j+\ell}
                            \le\mathcal V(x_j)
                            \le s_{j+\ell},\,
                            \forall j\in[\![1,n-\ell]\!]
                        \}
                    }
                \right]
            \end{aligned}
        \right]
        \\ &
        \le
        \mathbb E\left[
            \sum_{|x|\le L}
            {\bf 1}_{
                \{
                    r_{j+\ell}
                    \le\mathcal V(x_j)
                    \le s_{j+\ell},\,
                    \forall j\in[\![1,L]\!]
                \}
            }
        \right]
        \\ &
        \quad\times\;
        \mathbb E\left[
            \sum_{|\nu|=\ell}
            {\bf 1}_{
                \{
                    -b\log N
                    \le
                    \mathcal V^{\mathrm{k}}(\nu_j)
                    \le
                    0,\,
                    \forall j\in[\![1,\ell]\!]
                \}
            }
            {\bf 1}_{
                \{
                    \mathcal V^{\mathrm{k}}(\nu)
                    \le
                    (-b+v)\log N
                \}
            }
        \right]
        \\ &
        =
        e^{(c+o_L(1))L^{1/3}}
        \mathbb E\left[
            \#\mathcal Z^{\mathrm{k}}_{\ell+}
        \right]
        =
        e^{(c+o_L(1))L^{1/3}}
        N^{a-b+v+o_N(1)}
        \\ &
        =
        e^{(c+o_N(1))(t^{1/3}\log N+o_N(\log N))}
        N^{a-b+v+o_N(1)}
        \\ &
        =
        N^{a-b+v+ct^{1/3}+o_N(1)}
        \le
        N^{\eta+o_N(1)},
    \end{align*}
    where we need to calculate the order of $L^{1/3}$ as
    \begin{align*}
        L^{1/3}
        &
        =
        \left(
            \ell_N-\lfloor\log^{5/2}N\rfloor
        \right)^{1/3}
        \\ &
        =
        \left(
            (t+o_L(1))\log^3N+o_N(\log^3 N)
        \right)^{1/3}
        \\ &
        =
        t^{1/3}\log N+o_N(\log N),
    \end{align*}
    and $\eta=\max\{a,a-b+v+ct^{1/3}\}<1$. We have shown that
    \begin{equation*}
        \mathbb E\left[
            \#\mathcal Z^{\mathrm{k}}_n
        \right]
        \le
        N^{\eta+o_N(1)},\quad
        n\in[\![1,\ell_N]\!].
    \end{equation*}
    By union bound and Markov's inequality, we conclude that
    \begin{equation}
        \mathbb P\left(
            \tau\le \ell_N
        \right)
        \le
        \sum_{n=1}^{\ell_N}
        \frac{
            \mathbb E\left[
                \#\mathcal Z^{\mathrm{k}}_n
            \right]
        }{
            N
        }
        \le
        \ell_N
        N^{-(1-\eta)+o_N(1)}
        =
        N^{-(1-\eta)+o_N(1)}.
        \label{eq:4.19}
    \end{equation}
    We have calculated the small upper bound for the first term in \eqref{eq:prop:3.1:1}, the probability of the irregular event $\tau\le\ell_N$, and shown that this probability decays at least polynominally with respect to $N$.
    
    Then we calculate a small upper bound of the second and last term in \eqref{eq:prop:3.1:1}, i.e., $\mathbb P(\mathcal M^{\mathrm{k}}_{\ell_N}\ge\zeta'\log N)$, for the constructed killed branching random walk $\mathcal V^{\mathrm{k}}$. Note that, for any $u\in\mathcal Z^{\mathrm{k}}_{\ell_N+}$,
    \begin{align*}
        \mathcal V^{\mathrm{k}}(u)
        &
        \le
        \mathcal V^{\mathrm{k}}(u_\ell)
        +r_{\ell_N}
        +vL^{1/3}
        \\ &
        \le
        (-b+v)\log N
        +\frac{\pi^2}2\sigma^2c^{-2}\ell_N^{1/3}
        +v\ell_N^{1/3}
        \\ &
        \le
        (-b+v)\log N
        +\frac{\pi^2}2\sigma^2c^{-2}(8t\log^3N)^{1/3}
        +v(8t\log^3N)^{1/3}
        \\ &
        \le
        \left(
            -b+\frac{\pi^2}2\sigma^2c^{-2}t^{1/3}+2v\left(1+t^{1/3}\right)
        \right)
        \log N
        % \\ &
        <\zeta'\log N,
    \end{align*}
    where $\ell_N\le 8t\log^3N$ holds for some $N_1=N_1(t)>0$, and all $N\ge N_1$. It means that any alive particle of the killed branching random walk $\mathcal V^{\mathrm{k}}$ at time $\ell_N$ is below the threshold $\zeta'\log N$. Hence,
    \begin{align*}
        \mathbb P\left(
            \mathcal M^{\mathrm{k}}_{\ell_N}
            \ge
            \zeta'\log N
        \right)
        \le
        \mathbb P\left(
            \#\mathcal Z^{\mathrm{k}}_{\ell_N+}
            =
            0
        \right).
    \end{align*}
    We proceed to calculate a small upper bound for the probability of the irregular event $\mathcal Z^{\mathrm{k}}_{\ell_N+}=0$. Recall that $\mathbb E[\#\mathcal Z^{\mathrm{k}}_{\ell_+}]=N^{a-b+v+o_N(1)}$. Fix $\delta\in(0,\frac 12(a-b+v))$. Denote $\lfloor N^\delta\rfloor$ pairwise different elements in $\mathcal Z^{\mathrm{k}}_{\ell_+}$ as $\nu_i$, $i\in[\![1,\lfloor N^\delta\rfloor]\!]$. We have
    \begin{align*}
        &\quad
        \mathbb P\left(
            \#\mathcal Z^{\mathrm{k}}_{\ell_N+}
            =
            0
        \right)
        \\
        &
        =
        \mathbb E\left[
            {\bf 1}_{
                \left\{
                    \#\mathcal Z^{\mathrm{k}}_{\ell_N+}
                    =
                    0
                \right\}
            }
            \left(
                {\bf 1}_{
                    \left\{
                        \#\mathcal Z^{\mathrm{k}}_{\ell_+}
                        \ge
                        N^\delta
                    \right\}
                }
                +
                {\bf 1}_{
                    \left\{
                        \#\mathcal Z^{\mathrm{k}}_{\ell_+}
                        <
                        N^\delta
                    \right\}
                }
            \right)
        \right]
        \\ &
        \le
        \mathbb E\left[
            {\bf 1}_{
                \left\{
                    \#\mathcal Z^{\mathrm{k}}_{\ell_N+}
                    =
                    0
                \right\}
            }
            {\bf 1}_{
                \left\{
                    \#\mathcal Z^{\mathrm{k}}_{\ell_+}
                    \ge
                    N^\delta
                \right\}
            }
        \right]
        +
        \mathbb P\left(
            \#\mathcal Z^{\mathrm{k}}_{\ell+}
            <
            N^\delta
        \right)
        \\ &
        =
        \mathbb E\left[
            \mathbb E\left(
                \prod_{
                    \nu\in
                    \mathcal Z^{\mathrm{k}}_{\ell+}
                }
                \!\!\!
                {\bf 1}_{
                    \left\{
                        \forall|u|=\ell_N,\,
                        u_\ell=\nu:\,
                        u\notin
                        \mathcal Z^{\mathrm{k}}_{\ell_N+}
                    \right\}
                }
                {\bf 1}_{
                    \left\{
                        \#\mathcal Z^{\mathrm{k}}_{\ell+}
                        \ge
                        N^\delta
                    \right\}
                }
                \middle|
                \mathscr F_\ell
            \right)
        \right]
        +
        \mathbb P\left(
            \#\mathcal Z^{\mathrm{k}}_{\ell+}
            <
            N^\delta
        \right)
        \\ &
        \le
        \mathbb E\left[
            \mathbb E\left(
                \prod_{i=1}^{
                    \lfloor N^\delta\rfloor
                }
                {\bf 1}_{
                    \left\{
                        \forall|u|=\ell_N,\,
                        u_\ell=\nu_i:\,
                        u\notin
                        \mathcal Z^{\mathrm{k}}_{\ell_N+}
                    \right\}
                }
                \middle|
                \mathscr F_\ell
            \right)
        \right]
        +
        \mathbb P\left(
            \#\mathcal Z^{\mathrm{k}}_{\ell+}
            <
            N^\delta
        \right)
        \\ &
        =
        \mathbb E\left[
            \prod_{i=1}^{
                \lfloor N^\delta\rfloor
            }
            \mathbb E\left(
                {\bf 1}_{
                    \left\{
                        \forall|u|=\ell_N,\,
                        u_\ell=\nu_i:\,
                        u\notin
                        \mathcal Z^{\mathrm{k}}_{\ell_N+}
                    \right\}
                }
                \middle|
                \mathscr F_\ell
            \right)
        \right]
        +
        \mathbb P\left(
            \#\mathcal Z^{\mathrm{k}}_{\ell+}
            <
            N^\delta
        \right),
    \end{align*}
    where we would proceed to use the branching property and \eqref{eq:4.7} in \Cref{lemma:4.6} to compute, for any $i\in[\![1,\lfloor N^\delta\rfloor]\!]$,
    \begin{align*}
        &\quad
        \mathbb E\left(
            {\bf 1}_{
                \left\{
                    \forall|u|=\ell_N,\,
                    u_\ell=\nu_i:\,
                    u\notin
                    \mathcal Z^{\mathrm{k}}_{\ell_N+}
                \right\}
            }
            \middle|
            \mathscr F_\ell
        \right)
        \\
        &
        =
        1-
        \mathbb E\left(
            {\bf 1}_{
                \left\{
                    \exists|u|=\ell_N:\,
                    u_\ell=\nu_i,\,
                    u\in
                    \mathcal Z^{\mathrm{k}}_{\ell_N+}
                \right\}
            }
            \middle|
            \mathscr F_\ell
        \right)
        \\ &
        =
        1-
        \mathbb E\left(
            {\bf 1}_{
                \left\{
                    \exists|u|=L:\,
                    u_\ell=\nu_i,\,
                    r_j
                    \le
                    \mathcal V^{\mathrm{k}}(u_j)
                    -
                    \mathcal V^{\mathrm{k}}(u_\ell)
                    \le
                    s_j,\,
                    \forall j\in[\![1,L]\!],\,
                    \mathcal V^{\mathrm{k}}(u)
                    -
                    \mathcal V^{\mathrm{k}}(u_\ell)
                    \le
                    r_{\ell_N}+vL^{1/3}
                \right\}
            }
            \middle|
            \mathscr F_\ell
        \right)
        \\ &
        =
        1-
        \mathbb E\left[
            {\bf 1}_{
                \left\{
                    \exists|x|=L:\,
                    r_{j+\ell}
                    \le
                    \mathcal V(x_j)
                    \le
                    s_{j+\ell},\,
                    \forall j\in[\![1,L]\!],\,
                    \mathcal V(x)
                    \le
                    r_{\ell_N}+vL^{1/3}
                \right\}
            }
        \right]
        \\ &
        =
        1-e^{-L^{1/3}o_L(1)}
        \le
        e^{
            -e^{
                -L^{\text{\fontsize{6pt}{5pt}$1/3$}}
                o_{\text{\fontsize{5pt}{5pt}$L$}}{\text{\fontsize{6pt}{5pt}$(1)$}}
            }
        }
        =
        e^{
            -N^{
                o_{\text{\fontsize{5pt}{5pt}$N$}}{\text{\fontsize{6pt}{5pt}$(1)$}}
            }
        },
    \end{align*}
    where the last inequality holds since $1-x\le e^{-x}$ for all $x\in[0,1]$, and the last equality holds since $L^{1/3}=t^{1/3}\log N+o_N(\log N)$ as computed above. Hence,
    \begin{align*}
        \mathbb P\left(
            \#\mathcal Z^{\mathrm{k}}_{\ell_N+}
            =
            0
        \right)
        &
        \le
        \mathbb E\left[
            \prod_{i=1}^{
                \lfloor N^\delta\rfloor
            }
            e^{
                -N^{
                    o_{\text{\fontsize{5pt}{5pt}$N$}}{\text{\fontsize{6pt}{5pt}$(1)$}}
                }
            }
        \right]
        +
        \mathbb P\left(
            \#\mathcal Z^{\mathrm{k}}_{\ell+}
            <
            N^\delta
        \right)
        \\ &
        =
        e^{
            -N^{
                \delta+
                o_{\text{\fontsize{5pt}{5pt}$N$}}{\text{\fontsize{6pt}{5pt}$(1)$}}
            }
        }
        +
        \mathbb P\left(
            \#\mathcal Z^{\mathrm{k}}_{\ell+}
            <
            N^\delta
        \right).
    \end{align*}
    To calculate an upper bound for the second term in the last equation, we would use the fact that $\#\mathcal Z^{\mathrm{k}}_{\ell+}$ is stochastically larger than a binomial random variable $X$ and then apply Chernoff's bound for binomial distribution. Indeed, according to \eqref{eq:4.9} in \Cref{lemma:4.7}, we know that the probability for a branching random walk starting with one particle at $0$ to have at least one particle staying in the tube $[-b\log N,0]$ till time $\ell$ and ending up in the interval $[-b\log N, (-b+v)\log N]$ at time $\ell$ is of order $N^{-b+v+o_N(1)}$. Therefore, the number $\#\mathcal Z^{\mathrm{k}}_{\ell+}$ of particles alive after the same operations for $\mathcal V^{\mathrm{k}}$ starting with $\lfloor N^a\rfloor$ particles is (stochastically) larger than the number of particles at time $0$ to have at least one descendant alive after the same operations, which is a binomial random variable $X$ with parameters $\lfloor N^a\rfloor$ as the trial number and the aforementioned probability of order $N^{-b+v+o_N(1)}$ as the success probability, whose mean $\mathbb E[X]$ is of order $N^{a-b+v+o_N(1)}$. By Chernoff's bound for binomial distribution, 
    \begin{align*}
        \mathbb P\left(
            \#\mathcal Z^{\mathrm{k}}_{\ell+}
            <
            N^\delta
        \right)
        &
        \le
        \mathbb P\left(
            \#\mathcal Z^{\mathrm{k}}_{\ell+}
            \le
            \lfloor N^\delta\rfloor
        \right)
        \\ &
        \le
        \mathbb P\left(
            X
            \le
            \lfloor N^\delta\rfloor
        \right)
        =
        \mathbb P\left(
            X
            \le
            \left(
                1
                -
                \lambda
            \right)
            \mathbb E[X]
        \right)
        \\ &
        \le
        e^{
            -\frac 12
            \lambda^2
            \mathbb E[X]
        }
        \le
        e^{
            -\frac 18
            \mathbb E[X]
        }
        \le
        e^{
            -N^{
                \delta+
                o_{\text{\fontsize{5pt}{5pt}$N$}}{\text{\fontsize{6pt}{5pt}$(1)$}}
            }
        },
    \end{align*}
    where
    \begin{equation*}
        \lambda
        :=
        1
        -
        \frac{
            \lfloor N^\delta\rfloor
        }{
            \mathbb E[X]
        }
        \ge
        \frac 12,
        \quad
        N^\delta
        \le
        \frac 18
        N^{a-b+v}
    \end{equation*}
    holds for some $N_2=N_2(a,t,\zeta')>N_1$ and all $N\ge N_2$. We have shown that, for all $N\ge N_2$,
    \begin{equation*}
        \mathbb P\left(
            \mathcal M^{\mathrm{k}}_{\ell_N}
            \ge
            \zeta'\log N
        \right)
        \le
        \mathbb P\left(
            \#\mathcal Z^{\mathrm{k}}_{\ell_N+}
            =
            0
        \right)
        \le
        2
        e^{
            -N^{
                \delta+
                o_{\text{\fontsize{5pt}{5pt}$N$}}{\text{\fontsize{6pt}{5pt}$(1)$}}
            }
        }
        =
        e^{
            -N^{
                \delta+
                o_{\text{\fontsize{5pt}{5pt}$N$}}{\text{\fontsize{6pt}{5pt}$(1)$}}
            }
        }.
    \end{equation*}
    Together with \eqref{eq:4.19} and \eqref{eq:prop:3.1:1}, we have, for some $\varepsilon=\varepsilon(a,t,\zeta')\in(0,\frac 12(1-\eta))$, some $N_0=N_0(a,t,\zeta')>N_2$, and all $N\ge N_0$,
    \begin{align*}
        \mathbb P\left(
            \mathcal M^N_{\ell_N}
            \ge
            \zeta'\log N
        \right)
        &
        \le
        \mathbb P\left(
            \tau\le\ell_N
        \right)
        +
        \mathbb P\left(
            \mathcal M^{\mathrm{k}}_{\ell_N}
            \ge
            \zeta'\log N
        \right)
        \\ &
        \le
        N^{-(1-\eta)+o_N(1)}
        +
        e^{
            -N^{
                \delta+
                o_{\text{\fontsize{5pt}{5pt}$N$}}{\text{\fontsize{6pt}{5pt}$(1)$}}
            }
        }
        \le
        N^{-\varepsilon},
    \end{align*}
    which completes the proof of \eqref{eq:4.15} for the case $a\in(0,1)$.
\end{proof}

\subsection{Proof of \eqref{eq:4.15} for the case of \texorpdfstring{$a\in\{0,1\}$}{a∈{0,1}}}

\begin{proof}[Proof of \eqref{eq:4.15} for the case of \texorpdfstring{$a\in\{0,1\}$}{a∈{0,1}}]
    \quad

    Case $a=1$ follows the proven case immediately since starting with fewer particles (stochastically) increases the minimum position, i.e., for the same $\varepsilon$, $N_0$ and all $N\ge N_0$,
    \begin{equation*}
        \mathbb P\left(
            \mathcal M^{N|1}_{\ell_N}
            \ge
            \zeta'\log N
        \right)
        \le
        \mathbb P\left(
            \mathcal M^{N|1/2}_{\ell_N}
            \ge
            \zeta'\log N
        \right)
        \le
        N^{-\varepsilon}.
    \end{equation*}
    
    \quad

    Now we proceed to show \eqref{eq:4.15} for the case of $a=0$. Here, we only prove under the additional assumptions that $\#\Xi\ge 2$ a.s. (\Cref{assumption:at-least-two-children}) and there exists some $\lambda>0$ such that
    $\mathbb E[e^{\lambda\Xi_2}]<\infty$ (\Cref{assumption:moment-exp-2}). Under \Cref{assumption:at-least-two-children}, we consider the binary branching random walk $\mathcal V^{\mathrm{b}}$ on the same probability space starting with one particle at $0$ with offspring law given by $(\Xi_1,\Xi_2)$, i.e., the two left-most children in $\Xi$. Let $\mathcal L^{\mathrm{b}}_k$ be the maximum of such binary branching random walk at time $k\ge 0$. By union bound and standard large deviation estimates (where we use \Cref{assumption:moment-exp-2}) for the sum of i.i.d.~random variables, for some fixed
    \begin{equation*}
        c
        >
        \max\left\{
            0,\;
            \frac{
                2+\log 2
                +
                \log\mathbb E\left[
                    e^{\lambda\Xi_2}
                \right]
            }\lambda
        \right\},
    \end{equation*}
    we have
    \begin{align*}
        \mathbb P\left(
            \mathcal L^{\mathrm{b}}_k
            \ge
            ck
        \right)
        &
        =
        \mathbb P\left(
            \exists|u|=k:\,
            \mathcal V^{\mathrm{b}}(u)\ge ck
        \right)
        % \\ &
        \le
        \sum_{|u|=k}
        \mathbb P\left(
            V''(u)\ge ck
        \right)
        \\[-10pt] &
        =
        \sum_{|u|=k}
        \mathbb P\left(
            \sum_{j=1}^k
            \left(
                \mathcal V^{\mathrm{b}}(u_j)
                -
                V''(u_{j-1})
            \right)
            \ge
            ck
        \right)
        \\ &
        \le
        \sum_{|u|=k}
        \mathbb P\left(
            \sum_{j=1}^k
            \left(
                \Xi_{u_{j-1}}
            \right)_2
            \ge
            ck
        \right)
        \\ &
        =
        \sum_{|u|=k}
        \mathbb P\left(
            e^{
                \lambda
                \sum_{j=1}^k
                \left(
                    \Xi_{u_{j-1}}
                \right)_2
            }
            \ge
            e^{\lambda ck}
        \right)
        \\ &
        \le
        \sum_{|u|=k}
        \frac{
            \mathbb E\left[
                e^{
                    \lambda
                    \sum_{j=1}^k
                    \left(
                        \Xi_{u_{j-1}}
                    \right)_2
                }
            \right]
        }{
            e^{\lambda ck}
        }
        % \\ &
        =
        2^k
        \frac{
            \left(
                \mathbb E\left[
                    e^{\lambda\Xi_2}
                \right]
            \right)^k
        }{
            e^{\lambda ck}
        }
        % \\ &
        =
        2^k
        e^{
            -k\left(
                c\lambda-
                \log\mathbb E\left[
                    e^{\lambda\Xi_2}
                \right]
            \right)
        }
        \\ &
        \le
        2^ke^{-k(2+\log 2)}
        =
        e^{-2k}.
    \end{align*}
    Now choose and fix some $\delta,\delta'\in(0,1)$ such that
    \begin{equation*}
        \zeta'-c\delta
        >
        \frac{\pi^2}2\sigma^2t,
        \quad
        \delta'<\frac 12\delta\log 2,
    \end{equation*}
    and we take $k=k_N:=\lfloor\delta\log N\rfloor$ in the above calculation. We have the relation
    \begin{equation*}
        N^{\delta'}
        <
        2^k
        =
        \#\mathcal Z^{\mathrm{b}}_k
        <
        N^\delta
        <
        N
    \end{equation*}
    holds for some $N_1=N_1(\delta,\delta')=N_1(t,\zeta')>0$ and all $N\ge N_1$. Indeed,
    \begin{equation*}
        N^{\delta'}
        <
        2^{\frac 12\delta\log N}
        <
        2^{\lfloor\delta\log N\rfloor}
        =
        2^k
        \le
        2^{\delta\log N}
        <
        e^{\delta\log N}
        =
        N^\delta
        <
        N.
    \end{equation*}
    We can consider an underlying $\mathcal V$ on the same probability space. We can thus view the original $N$-BRW $\mathcal V^{N|0}$ and the binary branching random walk $\mathcal V^{\mathrm{b}}$ defined above as subtrees of $\mathcal V$. Furthermore, we construct a translated and killed $N$-BRW $\mathcal V^{\mathrm{tk}}$ whose translation and selection mechanism happened at time $k$ are specified as follows: 
    \begin{itemize}
        \item We translate the left-most $\lfloor N^{\delta'}\rfloor$ particles at time $k$ (and their descendants) in the original $N$-BRW $\mathcal V^{N|0}$ to the position of $\mathcal L^{\mathrm{b}}_k$, i.e., the right-most position of the binary branching random walk $\mathcal V^{\mathrm{b}}$ at time $k$; Note that since $N^{\delta'}<2^k$ holds for all $N\ge N_1$, all these particles are translated to the right when $N\ge N_1$.
        \item We kill all other particles at time $k$ in the original $N$-BRW $\mathcal V^{N|0}$.
    \end{itemize}
    It is then clear that the translated and killed $N$-BRW $\mathcal V^{\mathrm{tk}}$ has its minimum position at time $\ell_N$ point-wise larger than the minimum position in question. Hence, by the triangular inequality,
    \begin{align*}
        % &\quad
        \mathbb P\left(
            \mathcal M^{N|0}_{\ell_N}
            \ge
            \zeta'\log N
        \right)
        % \\
        &
        \le
        \mathbb P\left(
            \mathcal M^{\mathrm{tk}}_{\ell_N}
            \ge
            \zeta'\log N
        \right)
        \\ &
        =
        \mathbb P\left(
            \mathcal L^{\mathrm{b}}_k
            +
            \left(
                \mathcal M^{\mathrm{tk}}_{\ell_N}
                -
                \mathcal L^{\mathrm{b}}_k
            \right)
            \ge
            c\delta\log N
            +
            (\zeta'-c\delta)\log N
        \right)
        \\ &
        \le
        \mathbb P\left(
            \mathcal L^{\mathrm{b}}_k
            \ge
            c\delta\log N
        \right)
        +
        \mathbb P\left(
            \mathcal M^{\mathrm{tk}}_{\ell_N}
            -
            \mathcal L^{\mathrm{b}}_k
            \ge
            (\zeta'-c\delta)\log N
        \right).
    \end{align*}
    For the first term, for all $N\ge N_1$,
    \begin{equation*}
        \mathbb P\left(
            \mathcal L^{\mathrm{b}}_k
            \ge
            c\delta\log N
        \right)
        \le
        \mathbb P\left(
            \mathcal L^{\mathrm{b}}_k
            \ge
            c\lfloor\delta\log N\rfloor
        \right)
        \le 
        e^{-2\lfloor\delta\log N\rfloor}
        <
        e^{-\delta\log N}
        =
        N^{-\delta}.
    \end{equation*}
    For the second term, by the branching property and the construction of $\mathcal V^{\mathrm{tk}}$,
    \begin{align*}
        \mathbb P\left(
            \mathcal M^{\mathrm{tk}}_{\ell_N}
            -
            \mathcal L^{\mathrm{b}}_k
            \ge
            (\zeta'-c\delta)\log N
        \right)
        &
        =
        \mathbb E\left[
            \mathbb E\left(
                {\bf 1}_{
                    \left\{
                        \mathcal M^{\mathrm{tk}}_{\ell_N}
                        -
                        \mathcal L^{\mathrm{b}}_k
                        \ge
                        (\zeta'-c\delta)\log N
                    \right\}
                }
                \middle|
                \mathcal F_k
            \right)
        \right]
        \\ &
        =
        \mathbb P\left(
            \mathcal M^{N|\delta'}_{\ell_N-k}
            \ge
            (\zeta'-c\delta)\log N
        \right)
        \le
        N^{-\varepsilon'},
    \end{align*}
    where $\mathcal V^{N|\delta'}$ is on the same probability space starting with $\lfloor N^{\delta'}\rfloor$ particles at time $0$ and position $0$, and the last inequality is by the proven case of $\delta'\in(0,1)$ and the choice of $\delta$ such that $\zeta'-c\delta>\frac{\pi^2}2\sigma^2t$, and it holds for some $\varepsilon'=\varepsilon'(\delta,\delta',t,\zeta')=\varepsilon'(t,\zeta')>0$, some $N_2=N_2(\delta,\delta',t,\zeta')=N_2(t,\zeta')>N_1$, and all $N\ge N_2$, and we must also verify
    \begin{equation*}
        \lim_{N\to\infty}\frac{\ell_N-k_N}{\log^3N}
        =
        \lim_{N\to\infty}\frac{(t+o(1))\log^3N+o(\log^3N)}{\log^3N}
        =
        t.
    \end{equation*}
    Hence,
    \begin{align*}
        \mathbb P\left(
            \mathcal M^{N|0}_{\ell_N}
            \ge
            \zeta'\log N
        \right)
        &
        \le
        \mathbb P\left(
            \mathcal L^{\mathrm{b}}_k
            \ge
            c\delta\log N
        \right)
        +
        \mathbb P\left(
            \mathcal M^{\mathrm{tk}}_{\ell_N}
            -
            \mathcal L^{\mathrm{b}}_k
            \ge
            (\zeta'-c\delta)\log N
        \right)
        \\ &
        \le
        N^{-\delta}
        +
        N^{-\varepsilon'}
        \le
        N^{-\varepsilon},        
    \end{align*}
    where the last inequality is by the choice of $\varepsilon\in(0,\frac 12\min\{\delta,\varepsilon'\})$ such that $\varepsilon<\delta$ and $\varepsilon<\varepsilon'$, and it holds for some $N_0=N_0(t,\zeta')>N_2$ and all $N\ge N_0$. We are now done with the proof of \eqref{eq:4.15} for all $a\in[0,1]$.
\end{proof}

\subsection{Proof of \eqref{eq:4.16}}

\begin{proof}[Proof of \eqref{eq:4.16}]
    \quad

    Now we proceed to show \eqref{eq:4.16}. Let $b>1$ in this part. Let
    \begin{align*}
        F_N
        &
        :=
        \left\{
            \exists j,n\in[\![0,\ell_N]\!],\,
            \exists |u|=n:\,
            j\le n,\,
            \mathcal V^{N|a}(u_j)-\mathcal V^{N|a}(u)
            >
            b\log N
        \right\}
        \\ &
        =
        \left\{
            \begin{aligned}
                &\exists j\in[\![0,\ell_N-1]\!],\,
                \exists k\in[\![1,\#\mathcal Z^{N|a}_j]\!],
                \\ &\exists n\in[\![j+1,\ell_N]\!],\,
                \exists |u|=n:
                \\ &\quad
                    u_j=\nu^{(k)}_j,
                \\ &\quad
                    \mathcal V^{N|a}(u_j)
                    -
                    \mathcal V^{N|a}(u_i)
                    \le
                    b\log N,\,
                    \forall i\in[\![j+1,n-1]\!],
                \\ &\quad
                    \mathcal V^{N|a}(u_j)
                    -
                    \mathcal V^{N|a}(u)
                    >
                    b\log N
            \end{aligned}
        \right\}
    \end{align*}
    be the event that there exists a position drop of over $b\log N$ from an ancester to its descendant, where we denote
    \begin{equation*}
        \mathcal Z^{N|a}_j
        =:
        \left\{
            \nu^{(k)}_j:\,
            k\in\left[\!\!\left[1,\#\mathcal Z^{N|a}_j\right]\!\!\right]
        \right\}.
    \end{equation*}
    Hence, by the union bound for $j$, the branching property for $\mathcal V$, and the many-to-one formula for $\mathcal V$,
    \begin{align*}
        \mathbb P(F_N)
        &
        \le
        \sum_{j=0}^{\ell_N-1}
        \mathbb E\left[
            \mathbb E\left(
                \sum_{k=1}^{\#\mathcal Z^{N|a}_j}
                {\bf 1}{
                    \left\{
                        \!\text{\scriptsize $
                        \begin{aligned}
                            &\exists n\in[\![j+1,\ell_N]\!],\,
                            \exists |u|=n:\,
                            u_j=\nu^{(k)}_j,
                            \\[-6pt] &
                                \mathcal V^{N|a}(u_j)
                                -
                                \mathcal V^{N|a}(u_i)
                                \le
                                b\log N,\,
                                \forall i\in[\![j+1,n-1]\!],
                            \\[-6pt] &
                                \mathcal V^{N|a}(u_j)
                                -
                                \mathcal V^{N|a}(u)
                                >
                                b\log N
                        \end{aligned}
                        $}
                    \right\}
                }
                \middle|
                \mathscr F_j
            \right)
        \right]
        \\ &
        =
        \sum_{j=0}^{\ell_N-1}
        \!\!
        \mathbb E\!\left[
            \!
            \mathbb E\!\left(
                \sum_{m=1}^N
                \!\!
                {\bf 1}_{
                    \left\{
                        \#\mathcal Z^{N|a}_j
                        =
                        m
                    \right\}
                }
                \sum_{k=1}^m
                \!
                {\bf 1}{
                    \left\{
                        \!\text{\scriptsize $
                        \begin{aligned}
                            &\exists n\in[\![j+1,\ell_N]\!],\,
                            \exists |u|=n:\,
                            u_j=\nu^{(k)}_j,
                            \\[-6pt] &
                                \mathcal V^{N|a}(u_j)
                                -
                                \mathcal V^{N|a}(u_i)
                                \le
                                b\log N,\,
                                \forall i\in[\![j+1,n-1]\!],
                            \\[-6pt] &
                                \mathcal V^{N|a}(u_j)
                                -
                                \mathcal V^{N|a}(u)
                                >
                                b\log N
                        \end{aligned}
                        $}
                    \right\}
                }
                \middle|
                \mathscr F_j
            \right)
        \right]
        \\ &
        =
        \sum_{j=0}^{\ell_N-1}
        \!\!
        \mathbb E\!\left[
            \!
            \sum_{m=1}^N
            \!\!
            {\bf 1}_{
                \left\{
                    \#\mathcal Z^{N|a}_j
                    =
                    m
                \right\}
            }
            \sum_{k=1}^m
            \!
            \mathbb E\!\left(
                {\bf 1}{
                    \left\{
                        \!\text{\scriptsize $
                        \begin{aligned}
                            &\exists n\in[\![j+1,\ell_N]\!],\,
                            \exists |u|=n:\,
                            u_j=\nu^{(k)}_j,
                            \\[-6pt] &
                                \mathcal V^{N|a}(u_j)
                                -
                                \mathcal V^{N|a}(u_i)
                                \le
                                b\log N,\,
                                \forall i\in[\![j+1,n-1]\!],
                            \\[-6pt] &
                                \mathcal V^{N|a}(u_j)
                                -
                                \mathcal V^{N|a}(u)
                                >
                                b\log N
                        \end{aligned}
                        $}
                    \right\}
                }
                \middle|
                \mathscr F_j
            \right)
        \right]
        \\ &
        \le
        \sum_{j=0}^{\ell_N-1}
        \!\!
        \mathbb E\!\left[
            \!
            \sum_{m=1}^N
            \!\!
            {\bf 1}_{
                \left\{
                    \#\mathcal Z^{N|a}_j
                    =
                    m
                \right\}
            }
            \sum_{k=1}^m
            \!
            \mathbb E\!\left(
                {\bf 1}{
                    \left\{
                        \!\text{\scriptsize $
                        \begin{aligned}
                            &\exists n\in[\![j+1,\ell_N]\!],\,
                            \exists |u|=n:\,
                            u_j=\nu^{(k)}_j,
                            \\[-6pt] &
                                \mathcal V(u_j)
                                -
                                \mathcal V(u_i)
                                \le
                                b\log N,\,
                                \forall i\in[\![j+1,n-1]\!],
                            \\[-6pt] &
                                \mathcal V(u_j)
                                -
                                \mathcal V(u)
                                >
                                b\log N
                        \end{aligned}
                        $}
                    \right\}
                }
                \middle|
                \mathscr F_j
            \right)
        \right]
        \\ &
        =
        \sum_{j=0}^{\ell_N-1}
        \!\!
        \mathbb E\!\left[
            \!
            \sum_{m=1}^N
            \!\!
            {\bf 1}_{
                \left\{
                    \#\mathcal Z^{N|a}_j
                    =
                    m
                \right\}
            }
            \sum_{k=1}^m
            \!
            \mathbb P\left(
                \text{\scriptsize $
                \begin{aligned}
                    &\exists n\in[\![1,\ell_N-j]\!],\,
                    \exists |u|=n:
                    \\[-6pt] &
                        \mathcal V(u_i)
                        \ge
                        -b\log N,\,
                        \forall i\in[\![1,n-1]\!],
                    \\[-6pt] &
                        \mathcal V(u)
                        <
                        -b\log N
                \end{aligned}
                $}
            \right)
        \right]
        \\ &
        \le
        N
        \sum_{j=0}^{\ell_N-1}
        \mathbb P\left(
            \text{\scriptsize $
            \begin{aligned}
                &\exists n\in[\![1,\ell_N-j]\!],\,
                \exists |u|=n:
                \\[-6pt] &
                    \mathcal V(u_i)
                    \ge
                    -b\log N,\,
                    \forall i\in[\![1,n-1]\!],
                \\[-6pt] &
                    \mathcal V(u)
                    <
                    -b\log N
            \end{aligned}
            $}
        \right)
        \\ &
        \le
        N
        \sum_{j=0}^{\ell_N-1}
        \sum_{n=1}^{\ell_N-j}
        \mathbb E\left[
            \sum_{|u|=n}
            {\bf 1}_{
                \left\{
                    \mathcal V(u_i)
                    \ge
                    -b\log N,\,
                    \forall i\in[\![1,n-1]\!],\,
                    \mathcal V(u)
                    <
                    -b\log N
                \right\}
            }
        \right]
        \\ &
        =
        N
        \sum_{j=0}^{\ell_N-1}
        \sum_{n=1}^{\ell_N-j}
        \mathbb E\left[
            e^{Y_n}
            {\bf 1}_{
                \left\{
                    Y_i
                    \ge
                    -b\log N,\,
                    \forall i\in[\![1,n-1]\!],\,
                    Y_n
                    <
                    -b\log N
                \right\}
            }
        \right]
        \\ &
        \le
        N
        \ell_N
        N^{-b}
        \mathbb E\left[
            {\bf 1}_{
                \left\{
                    \exists n\in[\![1,\ell_N-j]\!]:\,
                    Y_i
                    \ge
                    -b\log N,\,
                    \forall i\in[\![1,n-1]\!],\,
                    Y_n
                    <
                    -b\log N
                \right\}
            }
        \right]
        \\[5pt] &
        \le N^{1-b}(t+o(1))\log^3N
        = N^{-(b-1)+o_N(1)}.
    \end{align*}
    For $a\in[0,1]$ and $\zeta<-a+\frac{\pi^2}2\sigma^2t$, we have
    \begin{equation*}
        \mathbb P\left(
            \mathcal M^{N|a}_{\ell_N}
            \ge
            \zeta\log N
        \right)
        \le
        \mathbb P\left(
            F_N
        \right)
        +
        \mathbb P\left(
            \mathcal M^{N|a}_{\ell_N}
            \ge
            \zeta\log N,\,
            F_N^{\mathrm c}
        \right),
    \end{equation*}
    and we proceed to calculate, by union bound and the many-to-one formula for $\mathcal V$,
    \begin{align*}
        &\quad
        \mathbb P\left(
            \mathcal M^{N|a}_{\ell_N}
            \ge
            \zeta\log N,\,
            F_N^{\mathrm c}
        \right)
        \\
        &
        =
        \mathbb E\!\left[
            \begin{aligned}
                &{\bf 1}\left\{
                    \exists k\in[\![1,\lfloor N^a\rfloor]\!],\,
                    \exists|u|=\ell_N,\,
                    \mathcal V^{N|a|(k)}(u)<\zeta\log N
                \right\}\\
                &{\bf 1}\left\{
                    \forall j,n\in[\![0,\ell_N]\!],\,
                    \forall|u|=n,\,
                    j\le n,\,
                    \mathcal V^{N|a}(u_j)-\mathcal V^{N|a}(u)
                    \le
                    b\log N
                \right\}
            \end{aligned}
        \right]
        \\ &
        =
        \mathbb E\!\left[
            \begin{aligned}
                &{\bf 1}\left\{
                    \begin{aligned}
                        &\exists k\in[\![1,\lfloor N^a\rfloor]\!],\,
                        \exists|u|=\ell_N,\,
                        \mathcal V^{N|a|(k)}(u)<\zeta\log N,
                        \\
                        &\forall i,j\in[\![0,\ell_N]\!],\,
                        i\le j\,\Rightarrow\,
                        \mathcal V^{N|a|(k)}(u_j)
                        -\mathcal V^{N|a|(k)}(u_i)
                        \ge
                        -b\log N
                    \end{aligned}
                \right\}\\
                &{\bf 1}\left\{
                    \forall j,n\in[\![0,\ell_N]\!],\,
                    \forall|u|=n,\,
                    j\le n\,\Rightarrow\,
                    \mathcal V^{N|a}(u_j)-\mathcal V^{N|a}(u)
                    \le
                    b\log N
                \right\}
            \end{aligned}
        \right]
        \\ &
        =
        \mathbb E\!\left[
            \begin{aligned}
                &{\bf 1}\left\{
                    \begin{aligned}
                        &\exists k\in[\![1,\lfloor N^a\rfloor]\!],\,
                        \exists|u|=\ell_N,\,
                        \mathcal V^{N|a|(k)}(u)<\zeta\log N,
                        \\
                        &\min_{1\le i\le j\le\ell_N}\left(
                            \mathcal V^{N|a|(k)}(u_j)
                            -\mathcal V^{N|a|(k)}(u_i)
                        \right)
                        \ge
                        -b\log N
                    \end{aligned}
                \right\}\\
                &{\bf 1}\left\{
                    \forall j,n\in[\![0,\ell_N]\!],\,
                    \forall|u|=n,\,
                    j\le n\,\Rightarrow\,
                    \mathcal V^{N|a}(u_j)-\mathcal V^{N|a}(u)
                    \le
                    b\log N
                \right\}
            \end{aligned}
        \right]
        \\ &
        \le
        \sum_{k=1}^{\lfloor N^a\rfloor}
        \mathbb E\!\left[
            \sum_{|u|=\ell_N}
            {\bf 1}\left\{
                \begin{aligned}
                    \mathcal V^{N|a|(k)}(u)<\zeta\log N,
                    \\
                    \min_{1\le i\le j\le\ell_N}\left(
                        \mathcal V^{N|a|(k)}(u_j)
                        -\mathcal V^{N|a|(k)}(u_i)
                    \right)
                    \ge
                    -b\log N
                \end{aligned}
            \right\}
        \right]
        \\ &
        =
        \sum_{k=1}^{\lfloor N^a\rfloor}
        \mathbb E\!\left[
            e^{Y_{\ell_N}}
            {\bf 1}\left\{
                \begin{aligned}
                    Y_{\ell_N}<\zeta\log N,
                    \\
                    \min_{1\le i\le j\le\ell_N}\left(
                        Y_j
                        -Y_i
                    \right)
                    \ge
                    -b\log N
                \end{aligned}
            \right\}
        \right]
        \\ &
        \le
        N^{a+\zeta}
        \mathbb P\left(
            Y_{\ell_N}<\zeta\log N,\,
            \min_{1\le i\le j\le\ell_N}\left(
                Y_j
                -Y_i
            \right)
            \ge
            -b\log N
        \right)
        \\ &
        \le
        N^{a+\zeta}
        \mathbb P\left(
            Y_{\ell_N}<(b+\zeta+a)\log N,\,
            \min_{1\le i\le j\le\ell_N}\left(
                Y_j
                -Y_i
            \right)
            \ge
            -b\log N
        \right)
        \\ &
        \le
        N^{a+\zeta}
        e^{-d\frac{\ell_N}{(b\log N)^2}}
        =
        N^{a+\zeta-dt b^{-2}+o_N(1)},
    \end{align*}
    where the last inequality holds for some $N_1=N_1(a,t,\zeta)>0$ and all $N\ge N_1$ by applying \Cref{lemma:2.3} with $\Delta=b\log N$, $L=\ell_N$, $0<c=\frac b{b+\zeta+a}<1$ and any $d\in(0,\frac{\pi^2}2\sigma^2)$. We may need to make sure that $L\ge(b\log N)^{2+c}$ when $N\ge N_1$. We can choose $b$ sufficiently close to $1$ and $d$ sufficiently close to $\frac{\pi^2}2\sigma^2$ such that $a+\zeta-dtb^{-2}<0$. Hence,
    \begin{equation*}
        \mathbb P\left(
            \mathcal M^{N|a}_{\ell_N}
            \ge
            \zeta\log N
        \right)
        \le
        N^{-(b-1)+o_N(1)}
        +
        N^{a+\zeta-dt b^{-2}+o_N(1)}
        \le
        N^{-\varepsilon},
    \end{equation*}
    where we chooce some $\varepsilon\in(0,\frac 12\min\{b-1,dtb^{-2}-a-\zeta\})$ such that $\varepsilon<b-1$ and $\varepsilon<dtb^{-2}-a-\zeta$, and it holds for some $N_0=N_0(a,t,\zeta)>N_1$ and all $N\ge N_0$. We are now done with the proof of \eqref{eq:4.16} for all $a\in[0,1]$.
\end{proof}

\subsection{Corollary on the asymptotic speed of the left-most particle}

\begin{theorem}
    \label{theorem:speed}
    Consider a branching random walk $\mathcal V^{N|a}$ that starts with $\lfloor N^a\rfloor$ particle at position $0$. Under all assumptions, the speed of the left-most particle in $\mathcal V^{N|a}$ exists and is given by
    \begin{equation*}
        v_N
        =
        \lim_{n\to\infty}
        \frac{1}{n}\mathcal M^{N|a}_n
        =
        \frac{\pi^2}2\sigma^2\frac{1}{\log^2 N}
        +
        o_N\left(\frac{1}{\log^2 N}\right).
    \end{equation*}
\end{theorem}

\begin{proof}[Proof of \Cref{theorem:speed}]
    \quad

    The existence of the speed is guaranteed by \Cref{corollary:leftmost-speed:N0} under all assumptions. We proceed to calculate the speed by showing that there exists an upper bound and a lower bound of the same order. For both bounds, we are to construct, with respect to $\mathcal V^{N|a}$, a branching random walk, $\mathcal V^{\mathrm{upper}}$ and $\mathcal V^{\mathrm{lower}}$ respectively, such that the minimum position in $\mathcal V^{\mathrm{upper}}$ is point-wise (or stochastically) larger than the minimum position in question, and the minimum position in $\mathcal V^{\mathrm{lower}}$ is point-wise (or stochastically) smaller than the minimum position in question. We then calculate the speed of the left-most particle in $\mathcal V^{\mathrm{upper}}$ and $\mathcal V^{\mathrm{lower}}$ using the proved result, \Cref{proposition:3.1}, to get the desired upper and lower bounds respectively.

    For any $t>0$, choose and fix some sequence $\ell_N=\ell_N(t)$ such that
    \begin{equation*}
        \lim_{N\to\infty}\frac{\ell_N}{\log^3N}=t.
    \end{equation*}

    For the upper bound, we construct, with respect to $\mathcal V^{N|a}$, an upper branching random walk $\mathcal V^{\mathrm{upper}}$ as follows:
    \begin{itemize}
        \item At time $0$, we kill all particles but one in the original $N$-BRW $\mathcal V^{N|a}$, so that there is only one particle at time $0$ and position $0$ in $\mathcal V^{\mathrm{upper}}$;
        \item From time $0$ till $\ell_N$, we let the process evolve as an $N$-BRW, i.e., at each step, we keep the left-most $N$ particles and kill all other particles;
        \item At time $\ell_N$, we kill all particles but one left-most particle in $\mathcal V^{\mathrm{upper}}$;
        \item From time $\ell_N$ onward, we let the process evolve as an $N$-BRW and repeat the killing procedure at all times $k\ell_N$ for $k\in\mathbb N$.
    \end{itemize}
    By definition of $\mathcal V^{\mathrm{upper}}$, we have, by \Cref{corollary:2} and by induction on $k$, on the same probability space,
    \begin{equation*}
        \mathcal V^{\mathrm{upper}}_{k\ell_N}
        \succcurlyeq
        \mathcal V^{N|a}_{k\ell_N},
        \quad\forall k\in\mathbb N,
    \end{equation*}
    and, by some coupling of $\mathcal V^{\mathrm{upper}}$ and $\mathcal V^{N|a}$,
    \begin{equation*}
        \mathcal V^{\mathrm{upper}}_n
        \succcurlyeq
        \mathcal V^{N|a}_n,
        \quad\forall n\in[\![k\ell_N,(k+1)\ell_N]\!],\,\forall k\in\mathbb N.
    \end{equation*}
    Hence, there is some coupling of $\mathcal V^{\mathrm{upper}}$ and $\mathcal V^{N|a}$ such that
    \begin{equation*}
        \mathcal M^{\mathrm{upper}}_n
        \ge
        \mathcal M^{N|a}_n,
        \quad\forall n\in\mathbb N.
    \end{equation*}
    Furthermore, by definition of $\mathcal V^{\mathrm{upper}}$, we have
    \begin{equation*}
        \mathcal V^{\mathrm{upper}}_{(k+1)\ell_N+n}
        -
        \mathcal M^{\mathrm{upper}}_{k\ell_N}
        \stackrel{\mathrm{d}}{=}
        \mathcal V^{N|0}_{\ell_N},
        \quad k\in\mathbb N,
    \end{equation*}
    and they are independent for different $k$'s, where, both sides of the equality is regarded as random finite counting measures, and the right-hand side is the population measure obtained from an $N$-BRW $\mathcal V^{N|0}$ starting with $N^0=1$ particle at position $0$ on the same probability space. Hence, 
    \begin{equation*}
        \left(
            \min\left(
                \mathcal V^{\mathrm{upper}}_{(k+1)\ell_N}
                -
                \mathcal M^{\mathrm{upper}}_{k\ell_N}
            \right):
            k\in\mathbb N
        \right)
    \end{equation*}
    is an i.i.d.~sequence of random variables with the same distribution as $\min\mathcal V^{N|0}_{\ell_N}$. To use the Strong Law of Large Numbers (SLLN), we need to check that, given $N$,
    \begin{equation*}
        \mathbb E\left[
            \min\mathcal V^{N|0}_{\ell_N}
        \right]
        =
        \mathbb E\left[
            \mathcal M^{N|0}_{\ell_N}
        \right]
        <
        \infty.
    \end{equation*}
    Indeed, under \Cref{assumption:psi-one,assumption:moment-1}, i.e., $\psi(1)=0$ and $\mathbb E[(\Xi_1)_+]<\infty$, we have
    \begin{equation*}
        \left|
            \mathbb E\left[
                    \mathcal M^{N|0}_n
            \right]
        \right|
        \le
        \mathbb E\left[
            \left(
                \mathcal M^{N|0}_n
            \right)_+
        \right]
        +
        \mathbb E\left[
            \left(
                \mathcal M^{N|0}_n
            \right)_-
        \right]
        <
        \infty,
    \end{equation*}
   see \Cref{lemma:M:moment-1:neg,lemma:M:moment-1:pos} under the assumptions for details. Hence, by SLLN and the proven result of \Cref{proposition:3.1} for $a=0$,
    \begin{align*}
        v_N
        &=
        \lim_{n\to\infty}
        \frac{1}{n}\mathcal M^{N|a}_n
        =
        \lim_{k\to\infty}
        \frac 1{k\ell_N}
        \mathcal M^{N|a}_{k\ell_N}
        \\ &
        \le
        \limsup_{k\to\infty}
        \frac 1{k\ell_N}
        \mathcal M^{\mathrm{upper}}_{k\ell_N}
        % \\ &
        =
        \frac 1{\ell_N}
        \limsup_{k\to\infty}
        \frac 1{k}
        \left(
            \sum_{i=1}^k
            \left(
                \mathcal M^{\mathrm{upper}}_{i\ell_N}
                -
                \mathcal M^{\mathrm{upper}}_{(i-1)\ell_N}
            \right)
        \right)
        \\ &
        =
        \frac 1{\ell_N}
        \limsup_{k\to\infty}
        \frac 1{k}
        \left(
            \sum_{i=1}^k
            \min
            \left(
                \mathcal V^{\mathrm{upper}}_{i\ell_N}
                -
                \mathcal M^{\mathrm{upper}}_{(i-1)\ell_N}
            \right)
        \right)
        \\ &
        =
        \frac 1{\ell_N}
        \mathbb E\left[
            \min\mathcal V^{N|0}_{\ell_N}
        \right]
        % \\ &
        =
        \frac 1{\ell_N}
        \mathbb E\left[
            \mathcal M^{N|0}_{\ell_N}
        \right]
        \\ &
        =
        \frac 1{t\log^3N+o(\log^3N)}
        \left(
            \frac{\pi^2}2\sigma^2t\log N
            +
            o_N(\log N)
        \right)
        \\ &
        =
        \frac{\pi^2}2\sigma^2\frac{1}{\log^2 N}
        +
        o_N\left(\frac{1}{\log^2 N}\right).
    \end{align*}

    Similarly, for the lower bound, we construct, by modifying $\mathcal V^{N|1}$ (i.e. an $N$-BRW starting with $N^1=N$ particles at position $0$ on the same probability space such that $\mathcal V^{N|a}$ is a subtree of $\mathcal V^{N|1}$), a lower branching random walk $\mathcal V^{\mathrm{lower}}$ for $\mathcal V^{N|a}$ as follows:
    \begin{itemize}
        \item At time $0$, we keep all particles in the original $N$-BRW $\mathcal V^{N|1}$ and let the process evolve as $\mathcal V^{N|1}$ from $0$ till $\ell_N$;
        \item At time $\ell_N$, even if the population size $\#\mathcal Z^{N|a}_{\ell_N}$ of $\mathcal V^{N|a}$ does not reach $N$, we translate all the $N$ particles in $\mathcal V^{\mathrm{lower}}_{\ell_N}$ to the left-most position $\min\mathcal V^{\mathrm{lower}}_{\ell_N}$;
        \item From time $\ell_N$ onward, we let the process evolve as $\mathcal V^{N|1}$ and repeat the translation procedure at all times $k\ell_N$ for $k\in\mathbb N$.
    \end{itemize}
    By definition of $\mathcal V^{\mathrm{lower}}$, we have, by \Cref{corollary:2} and by induction on $k$, on the same probability space,
    \begin{equation*}
        \mathcal V^{N|a}_{k\ell_N}
        \succcurlyeq
        \mathcal V^{\mathrm{lower}}_{k\ell_N},
        \quad\forall k\in\mathbb N,
    \end{equation*}
    and, by some coupling of $\mathcal V^{\mathrm{lower}}$ and $\mathcal V^{N|a}$,
    \begin{equation*}
        \mathcal V^{N|a}_n
        \succcurlyeq
        \mathcal V^{\mathrm{lower}}_n,
        \quad\forall n\in[\![k\ell_N,(k+1)\ell_N]\!],\,\forall k\in\mathbb N.
    \end{equation*}
    Hence, there is some coupling of $\mathcal V^{\mathrm{lower}}$ and $\mathcal V^{N|a}$ such that
    \begin{equation*}
        \mathcal M^{N|a}_n
        \ge
        \mathcal M^{\mathrm{lower}}_n,
        \quad\forall n\in\mathbb N.
    \end{equation*}
    Furthermore, by definition of $\mathcal V^{\mathrm{lower}}$, we have
    \begin{equation*}
        \mathcal V^{\mathrm{lower}}_{(k+1)\ell_N}
        -
        \mathcal M^{\mathrm{lower}}_{k\ell_N}
        \stackrel{\mathrm{d}}{=}
        \mathcal V^{N|1}_{\ell_N},
        \quad k\in\mathbb N,
    \end{equation*}
    and they are independent for different $k$'s. Hence,
    \begin{equation*}
        \left(
            \min\left(
                \mathcal V^{\mathrm{lower}}_{(k+1)\ell_N}
                -
                \mathcal M^{\mathrm{lower}}_{k\ell_N}
            \right):
            k\in\mathbb N
        \right)
    \end{equation*}
    is an i.i.d.~sequence of random variables with the same distribution as $\min\mathcal V^{N|1}_{\ell_N}$. To use SLLN, we need to check that, given $N$,
    \begin{equation*}
        \mathbb E\left[
            \min\mathcal V^{N|1}_{\ell_N}
        \right]
        =
        \mathbb E\left[
            \mathcal M^{N|1}_{\ell_N}
        \right]
        <
        \infty.
    \end{equation*}
    Indeed, since $\mathcal M^{N|1}_n\le\mathcal M^{N|0}_n$ holds (stocatistically) for all $n$, we have
    \begin{equation*}
        \mathbb E\left[
            \left(
                \mathcal M^{N|1}_n
            \right)_+
        \right]
        \le
        \mathbb E\left[
            \left(
                \mathcal M^{N|0}_n
            \right)_+
        \right]
        <
        \infty,
    \end{equation*}
    and
    \begin{align*}
        \mathbb E\left[
            \left(
                \mathcal M^{N|1}_n
            \right)_-
        \right]
        &
        \le
        \mathbb E\left[
            e^{
                \left(
                    \mathcal M^{N|1}_n
                \right)_-
            }
        \right]
        % \\ &
        \le
        \mathbb E\left[
            e^{
                -\mathcal M^{N|1}_n
            }
            +
            1
        \right]
        \\ &
        \le
        1
        +
        \sum_{|u|=n}
        \mathbb E\left[
            e^{
                -\mathcal V^{N|1}(u)
            }
        \right]
        % \\ &
        =
        1
        +
        N\mathbb E[W_n]
        =
        1+N<\infty,
    \end{align*}
    where $W_n$ is the additive martingale, and $\mathbb E[W_n]=1$ is proven in \Cref{lemma:M:exp-neg-moment} under \Cref{assumption:psi-one}, i.e., $\psi(1)=0$.
    Hence, by SLLN and the proven result of \Cref{proposition:3.1} for $a=1$,
    \begin{align*}
        v_N
        &=
        \lim_{n\to\infty}
        \frac{1}{n}\mathcal M^{N|a}_n
        =
        \lim_{k\to\infty}
        \frac 1{k\ell_N}
        \mathcal M^{N|a}_{k\ell_N}
        \\ &
        \le
        \liminf_{k\to\infty}
        \frac 1{k\ell_N}
        \mathcal M^{\mathrm{lower}}_{k\ell_N}
        % \\ &
        =
        \frac 1{\ell_N}
        \liminf_{k\to\infty}
        \frac 1{k}
        \left(
            \sum_{i=1}^k
            \left(
                \mathcal M^{\mathrm{lower}}_{i\ell_N}
                -
                \mathcal M^{\mathrm{lower}}_{(i-1)\ell_N}
            \right)
        \right)
        \\ &
        =
        \frac 1{\ell_N}
        \liminf_{k\to\infty}
        \frac 1{k}
        \left(
            \sum_{i=1}^k
            \min
            \left(
                \mathcal V^{\mathrm{lower}}_{i\ell_N}
                -
                \mathcal M^{\mathrm{lower}}_{(i-1)\ell_N}
            \right)
        \right)
        \\ &
        =
        \frac 1{\ell_N}
        \mathbb E\left[
            \min\mathcal V^{N|1}_{\ell_N}
        \right]
        % \\ &
        =
        \frac 1{\ell_N}
        \mathbb E\left[
            \mathcal M^{N|1}_{\ell_N}
        \right]
        \\ &
        =
        \frac 1{t\log^3N+o(\log^3N)}
        \left(
            \left(
                -1+
                \frac{\pi^2}2\sigma^2t
            \right)\log N
            +
            o_N(\log N)
        \right)
        \\ &
        =
        \left(
            -\frac 1t+
            \frac{\pi^2}2\sigma^2
        \right)
        \frac{1}{\log^2 N}
        +
        o_N\left(\frac{1}{\log^2 N}\right),
    \end{align*}
    and let $t\to\infty$ to get the desired lower bound. We are now done with the proof of \Cref{theorem:speed}.
\end{proof}
\section{Other Methods}

\chapter{Conclusions and Future Work}
\section{Conclusions}
\section{Future Work}


% 附录
\appendix
% !TeX root = ../main.tex
\chapter{Miscellaneous Lemmas}

\section{Many-to-one Formula}

\section{Lemmas on Log-Laplace Transforms}

    Consider the point process $\Xi$ governing the branching random walk. Let $\psi$ be the log-Laplace transform defined by
    \begin{equation}
        \psi(\lambda)
        :=
        \log\mathbb E\left[
            \sum_{x\in\Xi}e^{-\lambda x}
        \right]
        \in(-\infty,\infty],\quad
        \lambda\in\mathbb R.
    \end{equation}
    And it can be viewed as the log-Laplace transform 
    \begin{equation}
        \psi(\lambda)=\log\int_{\mathbb R}e^{-\lambda x}\,\mu(\mathrm dx)
    \end{equation}
    of the measure
    \begin{equation}
        \mu(A)
        :=
        \mathbb E\left[
            \sum_{x\in\Xi}\mathbb I[x\in A]
        \right],
        \quad
        A\in\mathscr B(\mathbb R).
    \end{equation}
    \begin{lemma} \it
        (i) The domain $\{\lambda:\psi(\lambda)<\infty\}$ is an interval and $\psi$ is convex.
        (ii) $\psi$ is $C^\infty$ on the interior of its domain.
        (iii）$\psi$ is analytic on the interior of its domain.
    \end{lemma}
    Proof is omitted. Note that (i) follows from H\"older's inequality of $\mu$; (ii) follows from Mean Value Theorem and Dominated Convergence Theorem and the fact that for any $n\ge 0$, $|x|^ne^{-\lambda x}\le C_ne^{-\lambda'x}$ for some $\lambda'>\lambda$ and constant $C_n$; and (iii) follows from Fubini-Tonelli Theorem.

    Now consider the random walk with step distribution $\xi$ obtained from the many-to-one lemma. Let $\kappa$ be the log-Laplace transform defined by
    \begin{equation}
        \kappa(\theta):=\log\mathbb Ee^{\theta\xi},\quad
        \theta\in\mathbb R.
    \end{equation}
    By (???), We can get
    \begin{equation}
        \kappa(\theta)
        =\mathbb E\left[
            \sum_{x\in\Xi}e^{\theta x}e^{-x}
        \right]
        =\psi(-\theta+1).
    \end{equation}
    We immediately obtain
    \begin{lemma}\it
        The domain of $\kappa$ is an interval, $\kappa$ is convex, and is $C^\infty$ and analytic on the interior of its domain.
    \end{lemma}

\section{Lemmas on Applying Mogul'skii's Theorem}

    Consider some independent identically distributed random variables $(\xi_i)_{i\ge 1}$ and some sequence $(x(n))$ indexed by $n$. Let
    \begin{equation}
        s_n(t):=\frac{1}{x(n)}\sum_{i=0}^{\lfloor nt\rfloor}\xi_i,\quad
        t\in[0,1],\,\xi_0=0.
        \label{eq:mogulskii-s-n-t}
    \end{equation}
    $s_n$ is thus defined as a random real function taking values in $D(0,1)$ (i.e., real functions on $[0,1]$ such that càdlàg on $[0,1)$ and left-continuous at $1$).
    
    Mogul'skii defined some functional $H_\alpha$ for $\alpha\in(0,2]$ on some set $\mathfrak N$ of $D(0,1)$. We state here only the definition for $G\in\mathfrak N$ of the form
    \begin{equation}
        G=\{f\in D(0,1):f(0)=0,\,S_2(t)<f(t)<S_1(t),\,\forall t\in[0,1]\},
        \label{eq:mogulskii-G}
    \end{equation}
    where $S_1,S_2$ are right-continuous and piecewise constant such that
    \begin{equation}
        S_2(t-0)\wedge S_2(t+0)
        <
        S_1(t-0)\wedge S_1(t+0),\quad
        t\in[0,1].
    \end{equation}
    and we have the definition
    \begin{equation}
        H_\alpha(G):=\int_0^1(S_1(t)-S_2(t))^{-\alpha}\,\mathrm dt.
        \label{eq:mogulskii-h-alpha}
    \end{equation}
    
    Recall the Mogul'skii's Theorem for $\alpha=2$. 
    
    \begin{theorem*}[Mogul'skii's Theorem]\it 
        \label{theorem:mogulskii}
        Suppose that $\mathbb E\xi_i=0$, $\mathbb E\xi_i^2=1$ and that the sequence $(x(n))$ is such that
        \begin{equation*}
            \lim_{n\to\infty}x(n)=\infty,\quad
            \lim_{n\to\infty}x(n)n^{-1/2}=0.
        \end{equation*}
        Then the relation
        \begin{equation*}
            \log\mathbb P(s_n\in G)\sim-\frac{\pi^2}{2}H_2(G)nx^{-2}(n)
        \end{equation*}
        holds for set $G\in\mathfrak N$.
    \end{theorem*}
    
    Now consider a sequence $(L_\Delta)_{\Delta\ge 0}$ indexed by $\Delta$ with $\lim_{\Delta\to\infty}L_\Delta=\infty$. Basically, when applying Mogul'skii's Theorem, we are letting $n$ there to be $L_\Delta$ and $x(n)$ there to be $\Delta$, to form an inverse sequence. We are in this section to make rigorous this statement.
    
    Let
    \begin{align}
        x(n)&:=\inf\{\Delta\ge 0: L_\Delta+1\ge n\},&n\ge 0,
        \\ y(n)&:= x(n+1)-1,&n\ge 0.
    \end{align}
    
    \begin{lemma}\it\label{lemma:mogulskii-x-y}
    The two sequences $x$ and $y$ satisfy
    \begin{equation}
        \lim_{n\to\infty}x(n)=\lim_{n\to\infty}y(n)=\infty.
    \end{equation}
    If we further require that $\lim_{\Delta\to\infty}\frac{L_\Delta}{\Delta^2}=\infty$, we have
    \begin{equation}
        \lim_{n\to\infty}x(n)n^{-1/2}
        =\lim_{n\to\infty}y(n)n^{-1/2}
        =0.
    \end{equation}
    \end{lemma}
    The proof for Lemma \ref{lemma:mogulskii-x-y} is omitted. We shall use the fact $L_{x(n)}\ge n-1$ and thus $0<\frac{x^2(n)}{n}<\frac{x^2(n)}{L_{x(n)}}$.
    
    Consider the random walk $(Y_n)_{n\ge 0}$ with zero mean and finite variance $0<\sigma^2<\infty$ obtained from the many-to-one lemma. For any $a\in(0,1)$, we denote in this section
    \begin{align}
        q(\Delta,n)&:=\log\mathbb P\left(
            -a\Delta\le Y_j\le (1-a)\Delta,\,
            \forall j\in[\![1,n]\!]
        \right),\\
        p(\Delta,n)&:=\frac{\Delta^2}nq(\Delta,n).
    \end{align}
    For any $\Delta\ge 0$, there exists $n=n(\Delta)=L_\Delta+1$ such that $x(n)\le\Delta<x(n+1)$,  and we can get by monotonicity of probability that
    \begin{lemma} \it\label{lemma:mogulskii-mono}
        For $\Delta_1\le\Delta_2$ such that $L:=L_{\Delta_1}=L_{\Delta_2}$, we have $q(\Delta_1,L)\le q(\Delta_2,L)$.
    \end{lemma}
    
    We are now able to get to the main lemma with an additional condition for $(L_\Delta)$.
    \begin{lemma}
        \it
        \label{lemma:mogulskii-cal}
        \quad
        \begin{enumerate}[(i)]
            \item By Mogul'skii's Theorem, we have, indexed by $n$,
                \begin{equation}
                    \lim_{n\to\infty}p(x(n),n-1)
                    =\lim_{n\to\infty}p(y(n),n-1)
                    =-\frac{\pi^2}2\sigma^2.
                \end{equation}
            \item For $n=n(\Delta)$, and indexed by $\Delta$,
                \begin{equation}
                    \lim_{\Delta\to\infty}p(x(n),n-1)
                    =\lim_{\Delta\to\infty}p(y(n),n-1)
                    =-\frac{\pi^2}2\sigma^2.
                \end{equation}
            \item Under the additional condition for $(L_\Delta)$ that, when $\Delta\to\infty$,
                \begin{equation}
                    N_\Delta:=\#\left\{\hat\Delta\ge 0: L_\Delta=L_{\hat\Delta}\right\}
                    =O_\Delta(1),
                \end{equation}
                we have
                \begin{equation}
                    \lim_{\Delta\to\infty}p(\Delta,L_\Delta)=-\frac{\pi^2}2\sigma^2.
                \end{equation}
            \item Under the same condition as (iii), for $v>0$,
                \begin{align}
                    &\quad
                    \lim_{\Delta\to\infty}
                    \frac{\Delta^2}{L_\Delta}
                    \log\mathbb P\left(
                        -a\Delta\le Y_j\le (1-a)\Delta,\,
                        \forall j\in[\![1,L_\Delta]\!],\,
                        Y_{L_\Delta}\in[-a\Delta,(v-a)\Delta]
                    \right)
                    \nonumber \\ &
                    =
                    -\frac{\pi^2}{2}\sigma^2.
                \end{align}
        \end{enumerate}
    \end{lemma}
    \begin{proof}[引理 \ref{lemma:mogulskii-cal}]
        For (i) and (ii), we prove here only the result for $x$. The result for $y$ is the same since the conditions for Mogul'skii's Theorem (Theorem \ref{theorem:mogulskii}) are satisfied by both $x$ and $y$ as stated in Lemma \ref{lemma:mogulskii-x-y}.
        \begin{enumerate}[(i)]
            \item  We denote here the step random variables of $(Y_i)$ as $(\sigma\xi_i)$. Define $s_n$ as in \eqref{eq:mogulskii-s-n-t}. We can express $q(x(n),n-1)$ in $s_n$ as
            \begin{align}
                q(x(n),n-1)
                &=
                \log\mathbb P
                \left(
                    \frac 1\sigma(-a)\le s_n(t)\le\frac 1\sigma(1-a),\,
                    t\in[0,1]
                \right) \nonumber
                \\ &=:\log\mathbb P(s_n\in G_0),
                \quad n\ge 1.
            \end{align}
            For each $m\ge 1$, we apply Mogul'skii's Theorem for $G=G(m)$ of the form \eqref{eq:mogulskii-G} with
            \begin{equation}
                S_1=\frac{1}\sigma\left(1-a-\frac 1{2m}\right),\,
                S_2=\frac{1}\sigma\left(-a+\frac 1{2m}\right),
            \end{equation}
            and we can use \eqref{eq:mogulskii-h-alpha} to calculate the corresponding $H_2(G)$ as
            \begin{equation}
                H_2(G)=\sigma^2\left(1-\frac 1m\right)^{-2}.
            \end{equation}
            Hence, by Mogul'skii's Theorem, we have
            \begin{align}
                \lim_{n\to\infty}
                \frac{x^2(n)}{n}
                \log\mathbb P(s_n\in G(m))
                =
                -\frac{\pi^2}2\sigma^2\left(1-\frac 1m\right)^{-2}.
            \end{align}
            Since $G(m)\subset G_0$, by monotonicity of probability, 
            \begin{align}
                \liminf_{n\to\infty}p(x(n),n-1)
                &\ge\liminf_{n\to\infty}\frac{x^2(n)}n\log\mathbb P(s_n\in G(m))
                \nonumber\\ &
                =-\frac{\pi^2}2\sigma^2\left(1-\frac 1m\right)^{-2}.
            \end{align}
            Let $m\to\infty$ and we get
            \begin{equation}
                \liminf_{n\to\infty}p(x(n),n-1)\ge-\frac{\pi^2}{2}\sigma^2.
            \end{equation}
            Similarly, we apply Mogul'skii's Theorem for $G=\widetilde G(m)$ of the form \eqref{eq:mogulskii-G} with
            \begin{equation}
                S_1=\frac{1}\sigma\left(1-a+\frac 1{2m}\right),\,
                S_2=\frac{1}\sigma\left(-a-\frac 1{2m}\right),
            \end{equation}
            and we can get
            \begin{equation}
                \limsup_{n\to\infty}p(x(n),n-1)\le-\frac{\pi^2}2\sigma^2\left(1-\frac 1m\right)^{-2}.
            \end{equation}
            Let $m\to\infty$ and the result would follow.
            \item When $\Delta\to\infty$, $n=n(\Delta)$ goes to infity because of the original condition on $(L_\Delta)$, and the result follows from (i) for subsequences.
            \item By Lemma \ref{lemma:mogulskii-mono} and since $x(n)\le\Delta\le x(n+1)-1=y(n)$, we have the lower and upper bound for $p(\Delta,L_\Delta)$ as
            \begin{align}
                p(\Delta,L_\Delta)
                &=\frac{\Delta^2}{L_\Delta}q(\Delta,L_\Delta)
                \ge\frac{\Delta^2}{L_\Delta}q(x(n),L_\Delta)
                =\frac{\Delta^2}{x^2(n)}p(x(n),n-1)
                \\
                p(\Delta,L_\Delta)
                &=\frac{\Delta^2}{L_\Delta}q(\Delta,L_\Delta)
                \le\frac{\Delta^2}{L_\Delta}q(y(n),L_\Delta)
                =\frac{\Delta^2}{y^2(n)}p(y(n),n-1)
            \end{align}
            Let $\Delta\to\infty$ in both equations and it suffices to show that
            \begin{equation}
                \lim_{\Delta\to\infty}\frac{\Delta}{x(n)}
                =\lim_{\Delta\to\infty}\frac{\Delta}{y(n)}
                =1.
            \end{equation}
            Indeed, we have $x(n)\le\Delta\le x(n)+N_\Delta$ and $y(n)-N_\Delta\le\Delta\le y(n)$, and the result follows from the additional condition.
            \item The proof is analagous to (i) to (iii). We only need to calculate the value of $H_2(G)$ for $G$ of the form \eqref{eq:mogulskii-G} with
            \begin{equation}
                S_1(t)=\begin{cases}
                    \frac 1\sigma\left(1-a\mp\frac{v\wedge 1}{2m}\right),
                    &t\in[0,\frac{n-1}n),\\
                    \frac 1\sigma\left(v-a\mp\frac{v\wedge 1}{2m}\right),
                    &t\in[\frac{n-1}n,1],
                \end{cases}
                \quad
                S_2=\frac 1\sigma\left(-a\pm\frac{v\wedge 1}{2m}\right).
            \end{equation}
            We use \eqref{eq:mogulskii-h-alpha} to calculate
            \begin{align}
                H_2(G)
                &=
                \sigma^2\left(1-\frac{v\wedge 1}m\right)^{-2}\frac{n-1}n
                +
                \sigma^2\left(v-\frac{v\wedge 1}m\right)^{-2}\frac 1n
                \\ &\sim
                \sigma^2\left(1-\frac{v\wedge 1}m\right)^{-2},
            \end{align}
            when $n\to\infty$.
        \end{enumerate}
    \end{proof}

    \begin{example}
        \label{example:mogulskii-L-1}
        For $v>0$, $u>0$ and $L_\Delta=\lfloor u\Delta^{2+v}\rfloor$, we have all the conditions satisfied, i.e.,
        \begin{equation*}
            \lim_{\Delta\to\infty}\frac{L_\Delta}{\Delta^2}=\infty,\,
            N_\Delta = O_\Delta(1).
        \end{equation*}
        Indeed, $L_\Delta>u\Delta^{2+v}-1$, and $N_{\Delta}=1$ for $\Delta$ large enough.
    \end{example}

    \begin{example}
        \label{example:mogulskii-L-2}
        For $v>0$, and $L_\Delta=\lfloor\Delta^{2+v}\rfloor-\lfloor\Delta^2\rfloor$, we have all the conditions satisfied. Indeed, $L_\Delta>\Delta^{2+v}-\Delta^2-2$, and $N_\Delta=1$ for $\Delta$ large enough.
    \end{example}
\chapter{Codes for Simulation}

% 结尾部分
\backmatter

% 用于盲审的论文需隐去致谢、发表论文、科研成果、



% 打印参考文献列表
\printbibliography
\renewcommand{\bibname}{参考文献}

\begin{acknowledgements}
  致谢
\end{acknowledgements}


% 发表论文及科研成果
% \chapter*{攻读学位期间研究成果}
% \addcontentsline{toc}{chapter}{\normalfont \sffamily 攻读学位期间研究成果}
% \begin{enumerate}
% \item Authors. Title. \textit{Journal Name}, xx(x),2022: xxx–xxx,
% \item Authors. Title. Submitted to \textit{Journal Name}, arXiv: xxxxxx, 2022,
% \end{enumerate}

\end{document}





